% !TeX program = pdflatex
% !TeX encoding = UTF-8
\documentclass[11pt,aspectratio=169,t]{beamer}
% Все презентации курса компилируются только pdflatex.
\usepackage{iftex}
\RequirePDFTeX
\usepackage{cmap}
\usepackage[T2A]{fontenc}
\usepackage[utf8]{inputenc}
\usepackage[english,russian]{babel}
\usepackage{PTSerif}
\usepackage{amsmath,amsfonts,amssymb}
\usepackage{graphicx,microtype,anyfontsize,tikz}
\usepackage{ragged2e}
\usetikzlibrary{arrows.meta,calc}

\usetheme{default}
\usefonttheme{serif}
\DeclareMathOperator{\sen}{sen}
\bibliographystyle{apalike}

% Общие сведения; номер лекции, её название и дата задаются в N.tex.
\newcommand{\coursename}{Введение в маневрирование КА}
\author{Чиняев Владимир Николаевич}
\institute{Московский физико-технический институт}

% Книжное оформление с голубым акцентом.
\definecolor{coursepaper}{HTML}{FBF9F3}
\definecolor{courseink}{HTML}{302B28}
\definecolor{courseaccent}{HTML}{187AA5}
\definecolor{coursemuted}{HTML}{6D6560}
\definecolor{courserule}{HTML}{A9CADA}
\definecolor{coursetint}{HTML}{EDF4F6}
\tikzset{
  course diagram/.style={font=\fontsize{9.5}{11.5}\selectfont,
    line cap=round,line join=round,>=Latex},
  orbit line/.style={draw=courseaccent,line width=1pt},
  aux line/.style={draw=courserule,line width=.5pt},
  vector line/.style={draw=courseink,line width=.65pt,->},
  angle line/.style={draw=courseaccent,line width=.8pt,->},
  orbit point/.style={circle,fill=courseaccent,inner sep=1.7pt}
}
\setbeamercolor{normal text}{fg=courseink,bg=coursepaper}
\setbeamercolor{background canvas}{bg=coursepaper}
\setbeamercolor{structure}{fg=courseaccent}
\setbeamercolor{alerted text}{fg=courseaccent}
\setbeamercolor{frametitle}{fg=courseink,bg=coursepaper}
\setbeamercolor{footline}{fg=coursemuted,bg=coursepaper}
\setbeamercolor{caption name}{fg=courseaccent}

% Основной текст — 11/14 pt; подписи к рисункам — 9.5/11.5 pt.
\let\coursenormalsize\normalsize
\renewcommand{\normalsize}{\coursenormalsize\fontsize{11}{14}\selectfont}
\setbeamerfont{normal text}{size*={11pt}{14pt}}
\setbeamerfont{frametitle}{size*={20pt}{23pt},series=\mdseries}
\setbeamerfont{title}{size*={27pt}{31pt},series=\mdseries}
\setbeamerfont{subtitle}{size*={12.8pt}{15pt},series=\mdseries}
\setbeamerfont{author}{size*={10.8pt}{13pt}}
\setbeamerfont{institute}{size*={9.2pt}{11pt}}
\setbeamerfont{date}{size*={9pt}{11pt}}
\setbeamerfont{footline}{size*={6.8pt}{8pt}}
\setbeamerfont{caption}{size*={9.5pt}{11.5pt}}
\setbeamerfont{itemize/enumerate body}{size*={11pt}{14pt}}
\setbeamerfont{itemize/enumerate subbody}{size*={11pt}{14pt}}
\setbeamerfont{itemize/enumerate subsubbody}{size*={11pt}{14pt}}
\AtBeginDocument{\normalsize}

\setbeamersize{text margin left=6mm,text margin right=6mm}
\setbeamertemplate{navigation symbols}{}
\setbeamertemplate{headline}{}
\setbeamertemplate{caption}[numbered]
\setbeamertemplate{itemize items}[circle]

% Стандартные блоки Beamer позволяют группировать текст и формулы.
\setbeamertemplate{blocks}[default]
\setbeamerfont{block title}{size*={11pt}{14pt},series=\bfseries}
\setbeamerfont{block body}{size*={11pt}{14pt}}
\setbeamercolor{block title}{fg=courseaccent,bg=coursepaper}
\setbeamercolor{block body}{fg=courseink,bg=coursetint}
\setbeamercolor{block title alerted}{parent=block title}
\setbeamercolor{block body alerted}{parent=block body}
\setbeamercolor{block title example}{parent=block title}
\setbeamercolor{block body example}{parent=block body}

% Компактный заголовок; его высота учитывается при размещении содержания.
\setbeamertemplate{frametitle}{%
  \vskip4.5mm%
  \begin{beamercolorbox}[wd=\textwidth,sep=0pt]{frametitle}
    \usebeamerfont{frametitle}\strut\insertframetitle\strut\par
    \vskip2mm%
    {\color{courseaccent}\hrule height .55pt}
  \end{beamercolorbox}%
}

% Нижний колонтитул: курс слева, лекция и номер слайда справа.
\setbeamertemplate{footline}{%
  \begin{beamercolorbox}[wd=\paperwidth,ht=5.5mm,dp=0pt]{footline}
    \begin{tikzpicture}[remember picture,overlay,x=1mm,y=1mm]
      \begin{scope}[shift={(current page.south west)}]
        \draw[courserule,line width=.3pt] (6,5.5)--(154,5.5);
        \node[text=coursemuted,anchor=west,inner sep=0pt,
          font=\usebeamerfont{footline}] at (6,2.8) {\coursename};
        \node[text=coursemuted,anchor=east,inner sep=0pt,
          font=\usebeamerfont{footline}] at (154,2.8)
          {Лекция \arabic{lecture}\quad\textbullet\quad\insertframenumber};
      \end{scope}
    \end{tikzpicture}%
  \end{beamercolorbox}%
}

% Титульник рассчитан на 16:9 (160 × 90 мм), без нижнего колонтитула.
\setbeamertemplate{title page}{%
  \begin{tikzpicture}[remember picture,overlay,x=1mm,y=1mm]
    \begin{scope}[shift={(current page.south west)}]
      \node[text=coursemuted,anchor=north,font=\usebeamerfont{institute}]
        at (80,82) {\insertinstitute};
      \draw[courseaccent,line width=.55pt] (71,71)--(89,71);
      \node[text=courseink,align=center,text width=125mm,
        font=\usebeamerfont{title}] at (80,53)
        {\hyphenpenalty=10000\coursename\par};
      \node[text=courseaccent,align=center,text width=145mm,font=\usebeamerfont{subtitle}]
        at (80,32) {Лекция \arabic{lecture}. \inserttitle};
      \node[text=courseink,font=\usebeamerfont{author}]
        at (80,19) {\insertauthor};
      \node[text=coursemuted,font=\usebeamerfont{date}]
        at (80,9) {\insertdate};
    \end{scope}
  \end{tikzpicture}%
}

\setlength{\abovedisplayskip}{5pt}
\setlength{\belowdisplayskip}{5pt}
\setlength{\abovedisplayshortskip}{3pt}
\setlength{\belowdisplayshortskip}{4pt}
\newcommand{\eqcompact}{\fontsize{10}{12}\selectfont}
\newcommand{\figcaption}[1]{{\usebeamerfont{caption}\color{coursemuted}#1\par}}
\newcommand{\derivationmathspacing}{%
  \setlength{\abovedisplayskip}{2pt}%
  \setlength{\belowdisplayskip}{3pt}%
  \setlength{\abovedisplayshortskip}{2pt}%
  \setlength{\belowdisplayshortskip}{2pt}%
  \setlength{\jot}{2pt}}
\setcounter{lecture}{6}
\title{Околокруговое приближение.\\Задача перехода.}
\date{27 октября 2026 г.}

\begin{document}
\begin{frame}[plain]
\relax
\titlepage
\end{frame}

% Перенесено из лекции 5, слайд 22.
\begin{frame}{Базис-вектор Лоудена}
\relax
\label{primerfull}
\relax
\derivationmathspacing
Перейдём к константам Баранова:
{\derivationmathspacing
\[
\boxed{\begin{aligned}
 \lambda_1&=l_1,&\lambda_2&=l_3+2l_6,&\lambda_3&=l_2,\\
 \lambda_4&=-l_4,&\lambda_5&=l_5,&\lambda_6&=l_6.
\end{aligned}}
\]
}{\derivationmathspacing
\[
\boxed{\begin{aligned}
 s_r(\varphi)&=\lambda_3\sin\varphi-\lambda_2\cos\varphi+2\lambda_6,\\
 s_t(\varphi)&=2\lambda_1+2\lambda_3\cos\varphi
                         +2\lambda_2\sin\varphi-3\lambda_6\varphi,\\
 s_z(\varphi)&=\lambda_4\sin\varphi+\lambda_5\cos\varphi.
\end{aligned}}
\]
}
\medskip
\justifying
Вектор $\boldsymbol s=(s_r,s_t,s_z)$ называется базис-вектором Лоудена.
У Баранова те же компоненты обозначаются также $\lambda,\mu,\nu$;
им соответствуют $s_r,s_t,s_z$ в таком же порядке.
\end{frame}


% Перенесено из лекции 5, слайд 28.
\begin{frame}{Достаточные условия глобального минимума}
\relax
\derivationmathspacing
\justifying
В околокруговом приближении манёвр с произвольным числом импульсов
обеспечивает глобальный минимум $J=\sum_j\Delta\widehat V_j$, если выполнены следующие условия.
\begin{enumerate}
\setlength{\itemsep}{0pt}
 \item Импульсы выполняют все шесть граничных
 условий, а их фазы принадлежат
 допустимому интервалу $[\varphi_i,\varphi_f]$.
 \item Существует один набор постоянных
 $\lambda_1,\ldots,\lambda_6$, задающий базис-вектор
 по формулам слайда~\ref{primerfull}, для которого в точке
 каждого ненулевого импульса выполнено:
 {\eqcompact\derivationmathspacing
\[
\boldsymbol s(\varphi_j)
     =\frac{\Delta\widehat{\boldsymbol V}_j}{\Delta\widehat V_j},
  \qquad \|\boldsymbol s(\varphi_j)\|=1.
\]
}
 \item Для того же базис-вектора выполнено достаточное
 условие:
 {\eqcompact\derivationmathspacing
\[
\|\boldsymbol s(\varphi)\|\leq1
  \qquad\text{при всех }\varphi\in[\varphi_i,\varphi_f].
\]
}
\end{enumerate}
\end{frame}

% Перенесено из лекции 5, слайд 29.
\begin{frame}{Достаточные условия глобального минимума}
\relax
\derivationmathspacing
\justifying
Внутри интервала в точках ненулевых импульсов также выполнено
условие стационарности:
{\eqcompact\derivationmathspacing
\[
\left.\frac{\mathrm d}{\mathrm d\varphi}\|\boldsymbol s(\varphi)\|^2
          \right|_{\varphi_j}=0.
\]
}
Оно следует из второго и третьего пунктов: в точке импульса
квадрат модуля достигает своего максимального значения $1$.
На фиксированных концах интервала равенство производной нулю
не требуется.

Перечисленная совокупность \emph{достаточна} для глобального
минимума без предположения регулярности. При независимости дифференциалов граничных ограничений
она также \emph{необходима}: для оптимального
манёвра существуют множители, удовлетворяющие этим условиям.
Таким образом, для регулярного случая получены необходимые
и достаточные условия оптимальности.
\end{frame}

% Методичка, § 3.1: постановка задачи перехода.
\begin{frame}{Постановка задачи перехода}
\relax
\label{fiveboundaries}
\derivationmathspacing
Заданы начальная и конечная орбиты. Начало маневрирования,
конечное положение аппарата на орбите и время перехода свободны.
\begin{columns}[T,onlytextwidth]
\begin{column}{.53\textwidth}
\setlength{\jot}{0pt}
\[
\boxed{\begin{aligned}
 2\sum\nolimits_j\Delta\widehat V_{tj}&=\Delta a,\\
 \sum\nolimits_j(\Delta\widehat V_{rj}\sin\varphi_j
       +2\Delta\widehat V_{tj}\cos\varphi_j)&=\Delta e_x,\\
 \sum\nolimits_j(-\Delta\widehat V_{rj}\cos\varphi_j
       +2\Delta\widehat V_{tj}\sin\varphi_j)&=\Delta e_y,\\
 \sum\nolimits_j\Delta\widehat V_{zj}\sin\varphi_j&=0,\\
 \sum\nolimits_j\Delta\widehat V_{zj}\cos\varphi_j&=\Delta i.
\end{aligned}}
\]
\smallskip
$\Delta i$ --- взаимный угол плоскостей; $V_0$ --- круговая скорость.
\end{column}
\begin{column}{.43\textwidth}
\begin{block}{Критерий качества}
\setlength{\jot}{0pt}
\[
\begin{gathered}
 J=\sum\nolimits_{j=1}^N\Delta\widehat V_j\longrightarrow\min,\\
 \Delta V_\Sigma=V_0J,\\
 \Delta\widehat V_j=\lVert\Delta\widehat{\boldsymbol V}_j\rVert\\
 =\sqrt{(\Delta\widehat V_{rj})^2+(\Delta\widehat V_{tj})^2
                         +(\Delta\widehat V_{zj})^2},\\
 \Delta\widehat{\boldsymbol V}_j=\Delta\boldsymbol V_j^{\text{\normalfont имп}}/V_0.
\end{gathered}
\]
\end{block}
\end{column}
\end{columns}
\end{frame}

\begin{frame}{Базис-вектор в общем случае}
\relax
\derivationmathspacing
При $\lambda_2^2+\lambda_3^2>0$ вводим вспомогательный угол $\alpha$:
{\eqcompact\derivationmathspacing
\setlength{\jot}{0pt}
\[
\cos\alpha=\frac{\lambda_3}{\sqrt{\lambda_2^2+\lambda_3^2}},
 \qquad
 \sin\alpha=\frac{\lambda_2}{\sqrt{\lambda_2^2+\lambda_3^2}}.
\]
}
Тогда полный базис-вектор принимает вид
{\eqcompact\derivationmathspacing
\setlength{\jot}{0pt}
\[
\boxed{\begin{aligned}
 s_r&=\sqrt{\lambda_2^2+\lambda_3^2}\,\sin(\varphi-\alpha)+2\lambda_6,\\
 s_t&=2\lambda_1+2\sqrt{\lambda_2^2+\lambda_3^2}\,
                       \cos(\varphi-\alpha)-3\lambda_6\varphi,\\
 s_z&=\frac{\lambda_3\lambda_4-\lambda_2\lambda_5}
                {\sqrt{\lambda_2^2+\lambda_3^2}}\,\sin(\varphi-\alpha)+\frac{\lambda_2\lambda_4+\lambda_3\lambda_5}
                {\sqrt{\lambda_2^2+\lambda_3^2}}\,\cos(\varphi-\alpha).
\end{aligned}}
\]
}
\end{frame}

% eq:transfer-zero-multiplier, eq:primer-components.
\begin{frame}{Базис-вектор в задаче перехода}
\relax
\derivationmathspacing
{\justifying
Конечное положение аппарата на орбите свободно.
Соответствующий множитель равен нулю: $l_6=\lambda_6=0$.
\par}
{\eqcompact\derivationmathspacing
\[
\boxed{\begin{aligned}
 s_r&=\sqrt{\lambda_2^2+\lambda_3^2}\,\sin(\varphi-\alpha),\\
 s_t&=2\lambda_1+2\sqrt{\lambda_2^2+\lambda_3^2}\,\cos(\varphi-\alpha),\\
 s_z&=\frac{\lambda_3\lambda_4-\lambda_2\lambda_5}
 {\sqrt{\lambda_2^2+\lambda_3^2}}\,\sin(\varphi-\alpha)
 +\frac{\lambda_2\lambda_4+\lambda_3\lambda_5}
 {\sqrt{\lambda_2^2+\lambda_3^2}}\,\cos(\varphi-\alpha).
\end{aligned}}
\]
}
\medskip
{\justifying
Ось $X$ направлена к восходящему узлу конечной плоскости относительно
начальной; угол $\varphi$ отсчитывается от оси $X$.
\par}
\end{frame}

% eq:planar-hodograph; eq:primer-phase-shift из главы 2.
\begin{frame}{Базис-вектор в компланарной задаче}
\relax
\derivationmathspacing
{\justifying
При $\Delta i=0$ ось $X$ выбирается произвольно.
Нормальные импульсы не меняют первые три условия и лишь увеличивают
затраты. Полагаем $\lambda_4=\lambda_5=0$.
\par}
\[
\boxed{\begin{aligned}
 s_r&=\sqrt{\lambda_2^2+\lambda_3^2}\sin(\varphi-\alpha),\\
 s_t&=2\lambda_1+2\sqrt{\lambda_2^2+\lambda_3^2}\cos(\varphi-\alpha),
 \qquad s_z=0.
\end{aligned}}
\]
\[
 \cos\alpha=\frac{\lambda_3}{\sqrt{\lambda_2^2+\lambda_3^2}},\qquad
 \sin\alpha=\frac{\lambda_2}{\sqrt{\lambda_2^2+\lambda_3^2}}.
\]
\smallskip
{\justifying
Годограф --- эллипс с центром на оси $s_t$ и отношением полуосей $2:1$.
При $\lambda_2=\lambda_3=0$ угол $\alpha$ не определён,
а годограф вырождается в точку.
\par}
\end{frame}

% Компланарный переход: manual/chapter-3-algorithms.tex, sec:planar.
% Этот фрагмент включается после кадра с видом компланарного базис-вектора.

\begin{frame}{Компланарный переход: типы годографа}
\relax
\centering
% Годографы из методички; сохранены координаты всех кривых и осей.
\begin{tikzpicture}[course diagram,scale=.85]
\begin{scope}[shift={(0.0,0)}]
\draw[coursemuted!60,dashed] (0,0) circle[radius=1.4cm];
\draw[->,coursemuted] (-1.7,0)--(1.85,0) node[right] {$s_t$};
\draw[->,coursemuted] (0,-1.65)--(0,1.75) node[above] {$s_r$};
\draw[courseaccent,line width=1.15pt] plot coordinates {(1.19000,-0.31500) (1.20087,-0.30762) (1.21147,-0.30015) (1.22182,-0.29259) (1.23190,-0.28493) (1.24171,-0.27720) (1.25125,-0.26937) (1.26052,-0.26147) (1.26950,-0.25348) (1.27821,-0.24542) (1.28664,-0.23728) (1.29478,-0.22907) (1.30263,-0.22079) (1.31019,-0.21245) (1.31746,-0.20404) (1.32443,-0.19556) (1.33111,-0.18703) (1.33749,-0.17844) (1.34357,-0.16980) (1.34934,-0.16110) (1.35481,-0.15235) (1.35998,-0.14356) (1.36484,-0.13473) (1.36938,-0.12585) (1.37362,-0.11694) (1.37755,-0.10799) (1.38116,-0.09900) (1.38446,-0.08999) (1.38744,-0.08095) (1.39011,-0.07189) (1.39246,-0.06280) (1.39449,-0.05369) (1.39621,-0.04457) (1.39761,-0.03543) (1.39868,-0.02629) (1.39944,-0.01713) (1.39988,-0.00797) (1.40000,0.00119) (1.39980,0.01035) (1.39927,0.01951) (1.39843,0.02867) (1.39727,0.03781) (1.39579,0.04694) (1.39400,0.05606) (1.39188,0.06516) (1.38945,0.07425) (1.38670,0.08331) (1.38363,0.09234) (1.38025,0.10134) (1.37656,0.11032) (1.37255,0.11926) (1.36823,0.12817) (1.36360,0.13703) (1.35866,0.14586) (1.35342,0.15463) (1.34787,0.16337) (1.34202,0.17205) (1.33586,0.18068) (1.32940,0.18926) (1.32265,0.19777) (1.31560,0.20623) (1.30825,0.21463) (1.30061,0.22295) (1.29269,0.23122) (1.28447,0.23941) (1.27597,0.24752) (1.26719,0.25557) (1.25813,0.26353) (1.24880,0.27142) (1.23918,0.27922) (1.22930,0.28693) (1.21915,0.29456) (1.20874,0.30210) (1.19806,0.30955) (1.18713,0.31690) (1.17594,0.32416) (1.16450,0.33132) (1.15281,0.33838) (1.14088,0.34533) (1.12870,0.35218) (1.11629,0.35892) (1.10365,0.36555) (1.09077,0.37207) (1.07767,0.37848) (1.06435,0.38477) (1.05081,0.39095) (1.03706,0.39700) (1.02310,0.40294) (1.00893,0.40875) (0.99456,0.41444) (0.98000,0.42000) (0.96524,0.42543) (0.95030,0.43074) (0.93517,0.43591) (0.91987,0.44095) (0.90439,0.44586) (0.88874,0.45063) (0.87293,0.45526) (0.85696,0.45975) (0.84084,0.46411) (0.82456,0.46832) (0.80815,0.47239) (0.79159,0.47631) (0.77489,0.48010) (0.75807,0.48373) (0.74113,0.48722) (0.72406,0.49056) (0.70688,0.49375) (0.68959,0.49678) (0.67220,0.49967) (0.65471,0.50241) (0.63713,0.50499) (0.61946,0.50742) (0.60170,0.50969) (0.58387,0.51181) (0.56597,0.51377) (0.54801,0.51558) (0.52998,0.51723) (0.51190,0.51872) (0.49377,0.52006) (0.47560,0.52123) (0.45738,0.52225) (0.43914,0.52310) (0.42087,0.52380) (0.40257,0.52434) (0.38426,0.52472) (0.36594,0.52494) (0.34762,0.52500) (0.32929,0.52490) (0.31097,0.52464) (0.29267,0.52422) (0.27438,0.52364) (0.25611,0.52290) (0.23787,0.52200) (0.21967,0.52094) (0.20151,0.51972) (0.18339,0.51835) (0.16532,0.51682) (0.14731,0.51513) (0.12936,0.51328) (0.11148,0.51127) (0.09367,0.50912) (0.07594,0.50680) (0.05829,0.50433) (0.04073,0.50171) (0.02327,0.49893) (0.00590,0.49601) (-0.01136,0.49293) (-0.02851,0.48970) (-0.04555,0.48632) (-0.06246,0.48280) (-0.07925,0.47913) (-0.09591,0.47531) (-0.11243,0.47134) (-0.12881,0.46724) (-0.14505,0.46299) (-0.16113,0.45860) (-0.17706,0.45407) (-0.19283,0.44940) (-0.20844,0.44459) (-0.22387,0.43965) (-0.23913,0.43458) (-0.25421,0.42937) (-0.26910,0.42403) (-0.28381,0.41857) (-0.29832,0.41297) (-0.31264,0.40725) (-0.32675,0.40141) (-0.34066,0.39544) (-0.35436,0.38935) (-0.36784,0.38315) (-0.38110,0.37682) (-0.39415,0.37039) (-0.40696,0.36384) (-0.41955,0.35718) (-0.43190,0.35041) (-0.44401,0.34353) (-0.45588,0.33655) (-0.46750,0.32947) (-0.47888,0.32228) (-0.49000,0.31500) (-0.50087,0.30762) (-0.51147,0.30015) (-0.52182,0.29259) (-0.53190,0.28493) (-0.54171,0.27720) (-0.55125,0.26937) (-0.56052,0.26147) (-0.56950,0.25348) (-0.57821,0.24542) (-0.58664,0.23728) (-0.59478,0.22907) (-0.60263,0.22079) (-0.61019,0.21245) (-0.61746,0.20404) (-0.62443,0.19556) (-0.63111,0.18703) (-0.63749,0.17844) (-0.64357,0.16980) (-0.64934,0.16110) (-0.65481,0.15235) (-0.65998,0.14356) (-0.66484,0.13473) (-0.66938,0.12585) (-0.67362,0.11694) (-0.67755,0.10799) (-0.68116,0.09900) (-0.68446,0.08999) (-0.68744,0.08095) (-0.69011,0.07189) (-0.69246,0.06280) (-0.69449,0.05369) (-0.69621,0.04457) (-0.69761,0.03543) (-0.69868,0.02629) (-0.69944,0.01713) (-0.69988,0.00797) (-0.70000,-0.00119) (-0.69980,-0.01035) (-0.69927,-0.01951) (-0.69843,-0.02867) (-0.69727,-0.03781) (-0.69579,-0.04694) (-0.69400,-0.05606) (-0.69188,-0.06516) (-0.68945,-0.07425) (-0.68670,-0.08331) (-0.68363,-0.09234) (-0.68025,-0.10134) (-0.67656,-0.11032) (-0.67255,-0.11926) (-0.66823,-0.12817) (-0.66360,-0.13703) (-0.65866,-0.14586) (-0.65342,-0.15463) (-0.64787,-0.16337) (-0.64202,-0.17205) (-0.63586,-0.18068) (-0.62940,-0.18926) (-0.62265,-0.19777) (-0.61560,-0.20623) (-0.60825,-0.21463) (-0.60061,-0.22295) (-0.59269,-0.23122) (-0.58447,-0.23941) (-0.57597,-0.24752) (-0.56719,-0.25557) (-0.55813,-0.26353) (-0.54880,-0.27142) (-0.53918,-0.27922) (-0.52930,-0.28693) (-0.51915,-0.29456) (-0.50874,-0.30210) (-0.49806,-0.30955) (-0.48713,-0.31690) (-0.47594,-0.32416) (-0.46450,-0.33132) (-0.45281,-0.33838) (-0.44088,-0.34533) (-0.42870,-0.35218) (-0.41629,-0.35892) (-0.40365,-0.36555) (-0.39077,-0.37207) (-0.37767,-0.37848) (-0.36435,-0.38477) (-0.35081,-0.39095) (-0.33706,-0.39700) (-0.32310,-0.40294) (-0.30893,-0.40875) (-0.29456,-0.41444) (-0.28000,-0.42000) (-0.26524,-0.42543) (-0.25030,-0.43074) (-0.23517,-0.43591) (-0.21987,-0.44095) (-0.20439,-0.44586) (-0.18874,-0.45063) (-0.17293,-0.45526) (-0.15696,-0.45975) (-0.14084,-0.46411) (-0.12456,-0.46832) (-0.10815,-0.47239) (-0.09159,-0.47631) (-0.07489,-0.48010) (-0.05807,-0.48373) (-0.04113,-0.48722) (-0.02406,-0.49056) (-0.00688,-0.49375) (0.01041,-0.49678) (0.02780,-0.49967) (0.04529,-0.50241) (0.06287,-0.50499) (0.08054,-0.50742) (0.09830,-0.50969) (0.11613,-0.51181) (0.13403,-0.51377) (0.15199,-0.51558) (0.17002,-0.51723) (0.18810,-0.51872) (0.20623,-0.52006) (0.22440,-0.52123) (0.24262,-0.52225) (0.26086,-0.52310) (0.27913,-0.52380) (0.29743,-0.52434) (0.31574,-0.52472) (0.33406,-0.52494) (0.35238,-0.52500) (0.37071,-0.52490) (0.38903,-0.52464) (0.40733,-0.52422) (0.42562,-0.52364) (0.44389,-0.52290) (0.46213,-0.52200) (0.48033,-0.52094) (0.49849,-0.51972) (0.51661,-0.51835) (0.53468,-0.51682) (0.55269,-0.51513) (0.57064,-0.51328) (0.58852,-0.51127) (0.60633,-0.50912) (0.62406,-0.50680) (0.64171,-0.50433) (0.65927,-0.50171) (0.67673,-0.49893) (0.69410,-0.49601) (0.71136,-0.49293) (0.72851,-0.48970) (0.74555,-0.48632) (0.76246,-0.48280) (0.77925,-0.47913) (0.79591,-0.47531) (0.81243,-0.47134) (0.82881,-0.46724) (0.84505,-0.46299) (0.86113,-0.45860) (0.87706,-0.45407) (0.89283,-0.44940) (0.90844,-0.44459) (0.92387,-0.43965) (0.93913,-0.43458) (0.95421,-0.42937) (0.96910,-0.42403) (0.98381,-0.41857) (0.99832,-0.41297) (1.01264,-0.40725) (1.02675,-0.40141) (1.04066,-0.39544) (1.05436,-0.38935) (1.06784,-0.38315) (1.08110,-0.37682) (1.09415,-0.37039) (1.10696,-0.36384) (1.11955,-0.35718) (1.13190,-0.35041) (1.14401,-0.34353) (1.15588,-0.33655) (1.16750,-0.32947) (1.17888,-0.32228) (1.19000,-0.31500)};
\filldraw[fill=courseink,draw=coursepaper,line width=.4pt] (1.40000,0.00000) circle[radius=2.3pt];
\node[below left,inner sep=2pt] at (0,0) {$0$};
\node[align=center,text width=4.3cm,anchor=north] at (0,-1.9) {а) Смещённый эллипс\\Касающиеся орбиты};
\end{scope}
\begin{scope}[shift={(5.5,0)}]
\draw[coursemuted!60,dashed] (0,0) circle[radius=1.4cm];
\draw[->,coursemuted] (-1.7,0)--(1.85,0) node[right] {$s_t$};
\draw[->,coursemuted] (0,-1.65)--(0,1.75) node[above] {$s_r$};
\filldraw[fill=courseink,draw=coursepaper,line width=.4pt] (1.40000,0.00000) circle[radius=2.3pt];
\filldraw[fill=courseink,draw=coursepaper,line width=.4pt] (1.40000,0.00000) circle[radius=2.3pt];
\node[below left,inner sep=2pt] at (0,0) {$0$};
\node[align=center,text width=4.3cm,anchor=north] at (0,-1.9) {б) Точка\\Непересекающиеся орбиты};
\end{scope}
\begin{scope}[shift={(11.0,0)}]
\draw[coursemuted!60,dashed] (0,0) circle[radius=1.4cm];
\draw[->,coursemuted] (-1.7,0)--(1.85,0) node[right] {$s_t$};
\draw[->,coursemuted] (0,-1.65)--(0,1.75) node[above] {$s_r$};
\draw[courseaccent,line width=1.15pt] plot coordinates {(1.25220,-0.31305) (1.26293,-0.30207) (1.27329,-0.29101) (1.28325,-0.27985) (1.29282,-0.26861) (1.30200,-0.25729) (1.31078,-0.24589) (1.31917,-0.23441) (1.32715,-0.22287) (1.33472,-0.21125) (1.34190,-0.19957) (1.34866,-0.18783) (1.35501,-0.17604) (1.36095,-0.16418) (1.36647,-0.15228) (1.37158,-0.14034) (1.37627,-0.12835) (1.38054,-0.11632) (1.38439,-0.10425) (1.38781,-0.09216) (1.39082,-0.08003) (1.39340,-0.06788) (1.39556,-0.05571) (1.39729,-0.04353) (1.39860,-0.03133) (1.39948,-0.01912) (1.39993,-0.00690) (1.39996,0.00531) (1.39956,0.01753) (1.39874,0.02974) (1.39748,0.04194) (1.39581,0.05413) (1.39371,0.06630) (1.39118,0.07845) (1.38823,0.09058) (1.38486,0.10268) (1.38106,0.11475) (1.37685,0.12678) (1.37221,0.13878) (1.36716,0.15073) (1.36169,0.16264) (1.35580,0.17450) (1.34951,0.18630) (1.34280,0.19805) (1.33568,0.20974) (1.32816,0.22136) (1.32023,0.23292) (1.31190,0.24440) (1.30317,0.25581) (1.29404,0.26714) (1.28452,0.27840) (1.27460,0.28956) (1.26430,0.30064) (1.25362,0.31163) (1.24255,0.32252) (1.23110,0.33331) (1.21928,0.34400) (1.20709,0.35459) (1.19453,0.36507) (1.18160,0.37544) (1.16832,0.38569) (1.15468,0.39583) (1.14068,0.40584) (1.12634,0.41574) (1.11166,0.42550) (1.09664,0.43514) (1.08128,0.44464) (1.06560,0.45401) (1.04959,0.46324) (1.03326,0.47233) (1.01662,0.48127) (0.99966,0.49007) (0.98241,0.49872) (0.96485,0.50721) (0.94700,0.51556) (0.92886,0.52374) (0.91044,0.53177) (0.89174,0.53963) (0.87276,0.54733) (0.85353,0.55486) (0.83403,0.56223) (0.81428,0.56942) (0.79428,0.57644) (0.77404,0.58328) (0.75356,0.58995) (0.73285,0.59643) (0.71192,0.60274) (0.69078,0.60886) (0.66942,0.61479) (0.64786,0.62054) (0.62610,0.62610) (0.60415,0.63147) (0.58202,0.63664) (0.55971,0.64162) (0.53722,0.64641) (0.51458,0.65100) (0.49178,0.65539) (0.46883,0.65958) (0.44573,0.66357) (0.42250,0.66736) (0.39915,0.67095) (0.37567,0.67433) (0.35207,0.67750) (0.32837,0.68047) (0.30457,0.68323) (0.28067,0.68579) (0.25669,0.68813) (0.23263,0.69027) (0.20851,0.69219) (0.18431,0.69391) (0.16006,0.69541) (0.13577,0.69670) (0.11143,0.69778) (0.08705,0.69865) (0.06266,0.69930) (0.03824,0.69974) (0.01381,0.69997) (-0.01063,0.69998) (-0.03506,0.69978) (-0.05948,0.69937) (-0.08388,0.69874) (-0.10826,0.69790) (-0.13260,0.69685) (-0.15691,0.69559) (-0.18116,0.69411) (-0.20536,0.69243) (-0.22950,0.69053) (-0.25357,0.68842) (-0.27756,0.68611) (-0.30146,0.68358) (-0.32528,0.68084) (-0.34899,0.67790) (-0.37260,0.67475) (-0.39610,0.67140) (-0.41947,0.66784) (-0.44272,0.66408) (-0.46583,0.66011) (-0.48880,0.65595) (-0.51162,0.65158) (-0.53429,0.64702) (-0.55679,0.64226) (-0.57912,0.63730) (-0.60128,0.63215) (-0.62325,0.62681) (-0.64504,0.62127) (-0.66662,0.61555) (-0.68801,0.60964) (-0.70918,0.60354) (-0.73014,0.59726) (-0.75088,0.59080) (-0.77139,0.58416) (-0.79166,0.57734) (-0.81169,0.57034) (-0.83147,0.56317) (-0.85100,0.55583) (-0.87028,0.54832) (-0.88928,0.54064) (-0.90802,0.53280) (-0.92648,0.52480) (-0.94465,0.51663) (-0.96254,0.50831) (-0.98014,0.49983) (-0.99744,0.49120) (-1.01443,0.48242) (-1.03111,0.47350) (-1.04748,0.46443) (-1.06354,0.45522) (-1.07926,0.44587) (-1.09466,0.43638) (-1.10973,0.42676) (-1.12445,0.41701) (-1.13884,0.40714) (-1.15288,0.39714) (-1.16656,0.38702) (-1.17989,0.37678) (-1.19286,0.36643) (-1.20547,0.35596) (-1.21771,0.34539) (-1.22958,0.33471) (-1.24108,0.32393) (-1.25220,0.31305) (-1.26293,0.30207) (-1.27329,0.29101) (-1.28325,0.27985) (-1.29282,0.26861) (-1.30200,0.25729) (-1.31078,0.24589) (-1.31917,0.23441) (-1.32715,0.22287) (-1.33472,0.21125) (-1.34190,0.19957) (-1.34866,0.18783) (-1.35501,0.17604) (-1.36095,0.16418) (-1.36647,0.15228) (-1.37158,0.14034) (-1.37627,0.12835) (-1.38054,0.11632) (-1.38439,0.10425) (-1.38781,0.09216) (-1.39082,0.08003) (-1.39340,0.06788) (-1.39556,0.05571) (-1.39729,0.04353) (-1.39860,0.03133) (-1.39948,0.01912) (-1.39993,0.00690) (-1.39996,-0.00531) (-1.39956,-0.01753) (-1.39874,-0.02974) (-1.39748,-0.04194) (-1.39581,-0.05413) (-1.39371,-0.06630) (-1.39118,-0.07845) (-1.38823,-0.09058) (-1.38486,-0.10268) (-1.38106,-0.11475) (-1.37685,-0.12678) (-1.37221,-0.13878) (-1.36716,-0.15073) (-1.36169,-0.16264) (-1.35580,-0.17450) (-1.34951,-0.18630) (-1.34280,-0.19805) (-1.33568,-0.20974) (-1.32816,-0.22136) (-1.32023,-0.23292) (-1.31190,-0.24440) (-1.30317,-0.25581) (-1.29404,-0.26714) (-1.28452,-0.27840) (-1.27460,-0.28956) (-1.26430,-0.30064) (-1.25362,-0.31163) (-1.24255,-0.32252) (-1.23110,-0.33331) (-1.21928,-0.34400) (-1.20709,-0.35459) (-1.19453,-0.36507) (-1.18160,-0.37544) (-1.16832,-0.38569) (-1.15468,-0.39583) (-1.14068,-0.40584) (-1.12634,-0.41574) (-1.11166,-0.42550) (-1.09664,-0.43514) (-1.08128,-0.44464) (-1.06560,-0.45401) (-1.04959,-0.46324) (-1.03326,-0.47233) (-1.01662,-0.48127) (-0.99966,-0.49007) (-0.98241,-0.49872) (-0.96485,-0.50721) (-0.94700,-0.51556) (-0.92886,-0.52374) (-0.91044,-0.53177) (-0.89174,-0.53963) (-0.87276,-0.54733) (-0.85353,-0.55486) (-0.83403,-0.56223) (-0.81428,-0.56942) (-0.79428,-0.57644) (-0.77404,-0.58328) (-0.75356,-0.58995) (-0.73285,-0.59643) (-0.71192,-0.60274) (-0.69078,-0.60886) (-0.66942,-0.61479) (-0.64786,-0.62054) (-0.62610,-0.62610) (-0.60415,-0.63147) (-0.58202,-0.63664) (-0.55971,-0.64162) (-0.53722,-0.64641) (-0.51458,-0.65100) (-0.49178,-0.65539) (-0.46883,-0.65958) (-0.44573,-0.66357) (-0.42250,-0.66736) (-0.39915,-0.67095) (-0.37567,-0.67433) (-0.35207,-0.67750) (-0.32837,-0.68047) (-0.30457,-0.68323) (-0.28067,-0.68579) (-0.25669,-0.68813) (-0.23263,-0.69027) (-0.20851,-0.69219) (-0.18431,-0.69391) (-0.16006,-0.69541) (-0.13577,-0.69670) (-0.11143,-0.69778) (-0.08705,-0.69865) (-0.06266,-0.69930) (-0.03824,-0.69974) (-0.01381,-0.69997) (0.01063,-0.69998) (0.03506,-0.69978) (0.05948,-0.69937) (0.08388,-0.69874) (0.10826,-0.69790) (0.13260,-0.69685) (0.15691,-0.69559) (0.18116,-0.69411) (0.20536,-0.69243) (0.22950,-0.69053) (0.25357,-0.68842) (0.27756,-0.68611) (0.30146,-0.68358) (0.32528,-0.68084) (0.34899,-0.67790) (0.37260,-0.67475) (0.39610,-0.67140) (0.41947,-0.66784) (0.44272,-0.66408) (0.46583,-0.66011) (0.48880,-0.65595) (0.51162,-0.65158) (0.53429,-0.64702) (0.55679,-0.64226) (0.57912,-0.63730) (0.60128,-0.63215) (0.62325,-0.62681) (0.64504,-0.62127) (0.66662,-0.61555) (0.68801,-0.60964) (0.70918,-0.60354) (0.73014,-0.59726) (0.75088,-0.59080) (0.77139,-0.58416) (0.79166,-0.57734) (0.81169,-0.57034) (0.83147,-0.56317) (0.85100,-0.55583) (0.87028,-0.54832) (0.88928,-0.54064) (0.90802,-0.53280) (0.92648,-0.52480) (0.94465,-0.51663) (0.96254,-0.50831) (0.98014,-0.49983) (0.99744,-0.49120) (1.01443,-0.48242) (1.03111,-0.47350) (1.04748,-0.46443) (1.06354,-0.45522) (1.07926,-0.44587) (1.09466,-0.43638) (1.10973,-0.42676) (1.12445,-0.41701) (1.13884,-0.40714) (1.15288,-0.39714) (1.16656,-0.38702) (1.17989,-0.37678) (1.19286,-0.36643) (1.20547,-0.35596) (1.21771,-0.34539) (1.22958,-0.33471) (1.24108,-0.32393) (1.25220,-0.31305)};
\filldraw[fill=courseink,draw=coursepaper,line width=.4pt] (1.40000,0.00000) circle[radius=2.3pt];
\filldraw[fill=courseink,draw=coursepaper,line width=.4pt] (-1.40000,0.00000) circle[radius=2.3pt];
\node[below left,inner sep=2pt] at (0,0) {$0$};
\node[align=center,text width=4.3cm,anchor=north] at (0,-1.9) {в) Центральный эллипс\\Пересекающиеся орбиты};
\end{scope}
\end{tikzpicture}

\par\medskip
\figcaption{Пунктир --- единичная окружность; отмечены точки касания.
Для смещённого эллипса и точки показан разгонный вариант.}
\end{frame}

% eq:planar-hodograph; eq:planar-boundaries.
\begin{frame}{Оптимальные компланарные импульсы}
\relax
{\justifying
Годограф должен лежать внутри единичной окружности.
Импульсы возможны в точках касания с ней.
Наиболее удалены от начала координат вершины эллипса на оси $s_t$;
в них $s_r=0$. Поэтому оптимальные импульсы трансверсальны.
\par}
\begin{block}{Граничные условия при нулевых радиальных компонентах}
\derivationmathspacing
\[
\begin{aligned}
 2\sum\nolimits_j\Delta\widehat V_{tj}&=\Delta a,\\
 2\sum\nolimits_j\Delta\widehat V_{tj}\cos\varphi_j&=\Delta e_x,\\
 2\sum\nolimits_j\Delta\widehat V_{tj}\sin\varphi_j&=\Delta e_y.
\end{aligned}
\]
\end{block}
\end{frame}

% eq:phi-e; eq:planar-intersection.
\begin{frame}{Взаимное расположение орбит}
\relax
При $\Delta e=\sqrt{(\Delta e_x)^2+(\Delta e_y)^2}>0$ вводим угол
оптимальной коррекции эксцентриситета:
\derivationmathspacing
\[
 \cos\varphi_e=\frac{\Delta e_x}{\Delta e},\qquad
 \sin\varphi_e=\frac{\Delta e_y}{\Delta e},\qquad
 \varphi_e=\operatorname{atan2}(\Delta e_y,\Delta e_x).
\]
Вычитая радиусы заданных орбит при одном и том же угле, получаем
\[
 \frac{r_f(\varphi)-r_i(\varphi)}{r_0}
 =\Delta a-\Delta e_x\cos\varphi-\Delta e_y\sin\varphi
 =\Delta a-\Delta e\cos(\varphi-\varphi_e).
\]
\begin{block}{Классификация в околокруговом приближении}
$|\Delta a|=\Delta e>0$ --- касающиеся орбиты.\par
$|\Delta a|>\Delta e$ --- непересекающиеся орбиты.\par
$|\Delta a|<\Delta e$ --- пересекающиеся орбиты.
\end{block}
\end{frame}

% eq:planar-one-contact; sec:planar-tangent.
\begin{frame}{Касающиеся орбиты: геометрия}
\relax
Пусть $|\Delta a|=\Delta e>0$.
{\justifying
Разность радиусов достигает нуля, не меняя знака:
при $\Delta a>0$ --- в точке $\varphi_e$,
при $\Delta a<0$ --- в точке $\varphi_e+\pi$.
В общей точке возможен переход одним импульсом.
\par}
\begin{block}{Смещённый эллипс с одним касанием}
\derivationmathspacing
\[
 2|\lambda_1|+2\sqrt{\lambda_2^2+\lambda_3^2}=1,\qquad
 \lambda_1\ne0,\quad\lambda_2^2+\lambda_3^2>0.
\]
При $\lambda_1>0$ касание в $(s_r,s_t)=(0,1)$ --- разгонный импульс.\par
При $\lambda_1<0$ касание в $(s_r,s_t)=(0,-1)$ --- тормозной импульс.
\end{block}
{\justifying
Другая вершина эллипса остаётся внутри единичной окружности:
годограф выделяет одно место импульса за оборот.
\par}
\end{frame}

\begin{frame}{Касающиеся орбиты: годограф}
\relax
\centering
% Рисунок из методички; координаты кривых, осей и отметок сохранены.
\begin{tikzpicture}[course diagram,scale=.72]
\node at (2.65,6.85) {Годограф};
\node at (11.6,6.85) {Модули импульсов};
\node at (11.6,4.1) {Модуль базис-вектора};
\draw[coursemuted!60,dashed] (2.65,3.6) circle[radius=1.7cm];
\draw[->,coursemuted] (.6,3.6)--(4.85,3.6) node[right] {$s_t$};
\draw[->,coursemuted] (2.65,1.55)--(2.65,5.8) node[above] {$s_r$};
\draw[courseaccent,line width=1.15pt] plot coordinates {(4.09500,3.21750) (4.10820,3.22646) (4.12108,3.23553) (4.13364,3.24472) (4.14588,3.25401) (4.15779,3.26340) (4.16938,3.27290) (4.18063,3.28250) (4.19154,3.29220) (4.20211,3.30199) (4.21234,3.31187) (4.22223,3.32184) (4.23176,3.33189) (4.24095,3.34203) (4.24977,3.35224) (4.25824,3.36253) (4.26635,3.37289) (4.27410,3.38332) (4.28148,3.39382) (4.28849,3.40438) (4.29513,3.41500) (4.30140,3.42567) (4.30730,3.43640) (4.31282,3.44718) (4.31797,3.45800) (4.32274,3.46887) (4.32712,3.47978) (4.33113,3.49073) (4.33475,3.50170) (4.33799,3.51271) (4.34085,3.52375) (4.34331,3.53480) (4.34540,3.54588) (4.34709,3.55697) (4.34840,3.56808) (4.34932,3.57920) (4.34985,3.59032) (4.35000,3.60145) (4.34975,3.61257) (4.34912,3.62370) (4.34810,3.63481) (4.34669,3.64591) (4.34489,3.65700) (4.34271,3.66808) (4.34014,3.67913) (4.33719,3.69016) (4.33385,3.70116) (4.33012,3.71213) (4.32602,3.72306) (4.32153,3.73396) (4.31667,3.74482) (4.31142,3.75563) (4.30580,3.76639) (4.29981,3.77711) (4.29344,3.78777) (4.28670,3.79837) (4.27959,3.80892) (4.27211,3.81940) (4.26427,3.82981) (4.25607,3.84015) (4.24751,3.85042) (4.23859,3.86062) (4.22932,3.87073) (4.21969,3.88076) (4.20972,3.89071) (4.19940,3.90057) (4.18873,3.91033) (4.17773,3.92000) (4.16639,3.92958) (4.15472,3.93905) (4.14273,3.94842) (4.13040,3.95768) (4.11776,3.96684) (4.10479,3.97588) (4.09152,3.98481) (4.07793,3.99362) (4.06404,4.00232) (4.04984,4.01088) (4.03535,4.01933) (4.02057,4.02765) (4.00550,4.03583) (3.99014,4.04388) (3.97451,4.05180) (3.95860,4.05958) (3.94243,4.06722) (3.92599,4.07472) (3.90929,4.08208) (3.89233,4.08928) (3.87513,4.09634) (3.85768,4.10325) (3.84000,4.11000) (3.82208,4.11660) (3.80394,4.12304) (3.78557,4.12932) (3.76698,4.13544) (3.74819,4.14140) (3.72919,4.14719) (3.70999,4.15281) (3.69060,4.15827) (3.67102,4.16356) (3.65126,4.16867) (3.63132,4.17361) (3.61121,4.17838) (3.59094,4.18297) (3.57052,4.18739) (3.54994,4.19162) (3.52922,4.19567) (3.50835,4.19955) (3.48736,4.20324) (3.46624,4.20674) (3.44500,4.21007) (3.42365,4.21320) (3.40220,4.21615) (3.38064,4.21891) (3.35899,4.22148) (3.33725,4.22387) (3.31544,4.22606) (3.29355,4.22806) (3.27159,4.22988) (3.24958,4.23150) (3.22751,4.23292) (3.20539,4.23416) (3.18324,4.23520) (3.16105,4.23605) (3.13884,4.23670) (3.11660,4.23716) (3.09436,4.23743) (3.07210,4.23750) (3.04985,4.23738) (3.02761,4.23706) (3.00538,4.23655) (2.98317,4.23584) (2.96099,4.23495) (2.93885,4.23385) (2.91674,4.23257) (2.89469,4.23109) (2.87269,4.22942) (2.85075,4.22756) (2.82888,4.22551) (2.80708,4.22327) (2.78537,4.22083) (2.76374,4.21821) (2.74221,4.21540) (2.72078,4.21240) (2.69946,4.20922) (2.67825,4.20585) (2.65716,4.20230) (2.63620,4.19856) (2.61538,4.19464) (2.59469,4.19054) (2.57415,4.18625) (2.55377,4.18179) (2.53354,4.17716) (2.51348,4.17235) (2.49358,4.16736) (2.47387,4.16220) (2.45434,4.15687) (2.43499,4.15137) (2.41585,4.14570) (2.39690,4.13986) (2.37816,4.13386) (2.35963,4.12770) (2.34132,4.12138) (2.32323,4.11490) (2.30538,4.10826) (2.28775,4.10146) (2.27037,4.09452) (2.25323,4.08742) (2.23634,4.08018) (2.21971,4.07278) (2.20334,4.06525) (2.18723,4.05757) (2.17139,4.04976) (2.15583,4.04180) (2.14055,4.03371) (2.12556,4.02549) (2.11085,4.01714) (2.09643,4.00867) (2.08232,4.00007) (2.06851,3.99134) (2.05500,3.98250) (2.04180,3.97354) (2.02892,3.96447) (2.01636,3.95528) (2.00412,3.94599) (1.99221,3.93660) (1.98062,3.92710) (1.96937,3.91750) (1.95846,3.90780) (1.94789,3.89801) (1.93766,3.88813) (1.92777,3.87816) (1.91824,3.86811) (1.90905,3.85797) (1.90023,3.84776) (1.89176,3.83747) (1.88365,3.82711) (1.87590,3.81668) (1.86852,3.80618) (1.86151,3.79562) (1.85487,3.78500) (1.84860,3.77433) (1.84270,3.76360) (1.83718,3.75282) (1.83203,3.74200) (1.82726,3.73113) (1.82288,3.72022) (1.81887,3.70927) (1.81525,3.69830) (1.81201,3.68729) (1.80915,3.67625) (1.80669,3.66520) (1.80460,3.65412) (1.80291,3.64303) (1.80160,3.63192) (1.80068,3.62080) (1.80015,3.60968) (1.80000,3.59855) (1.80025,3.58743) (1.80088,3.57630) (1.80190,3.56519) (1.80331,3.55409) (1.80511,3.54300) (1.80729,3.53192) (1.80986,3.52087) (1.81281,3.50984) (1.81615,3.49884) (1.81988,3.48787) (1.82398,3.47694) (1.82847,3.46604) (1.83333,3.45518) (1.83858,3.44437) (1.84420,3.43361) (1.85019,3.42289) (1.85656,3.41223) (1.86330,3.40163) (1.87041,3.39108) (1.87789,3.38060) (1.88573,3.37019) (1.89393,3.35985) (1.90249,3.34958) (1.91141,3.33938) (1.92068,3.32927) (1.93031,3.31924) (1.94028,3.30929) (1.95060,3.29943) (1.96127,3.28967) (1.97227,3.28000) (1.98361,3.27042) (1.99528,3.26095) (2.00727,3.25158) (2.01960,3.24232) (2.03224,3.23316) (2.04521,3.22412) (2.05848,3.21519) (2.07207,3.20638) (2.08596,3.19768) (2.10016,3.18912) (2.11465,3.18067) (2.12943,3.17235) (2.14450,3.16417) (2.15986,3.15612) (2.17549,3.14820) (2.19140,3.14042) (2.20757,3.13278) (2.22401,3.12528) (2.24071,3.11792) (2.25767,3.11072) (2.27487,3.10366) (2.29232,3.09675) (2.31000,3.09000) (2.32792,3.08340) (2.34606,3.07696) (2.36443,3.07068) (2.38302,3.06456) (2.40181,3.05860) (2.42081,3.05281) (2.44001,3.04719) (2.45940,3.04173) (2.47898,3.03644) (2.49874,3.03133) (2.51868,3.02639) (2.53879,3.02162) (2.55906,3.01703) (2.57948,3.01261) (2.60006,3.00838) (2.62078,3.00433) (2.64165,3.00045) (2.66264,2.99676) (2.68376,2.99326) (2.70500,2.98993) (2.72635,2.98680) (2.74780,2.98385) (2.76936,2.98109) (2.79101,2.97852) (2.81275,2.97613) (2.83456,2.97394) (2.85645,2.97194) (2.87841,2.97012) (2.90042,2.96850) (2.92249,2.96708) (2.94461,2.96584) (2.96676,2.96480) (2.98895,2.96395) (3.01116,2.96330) (3.03340,2.96284) (3.05564,2.96257) (3.07790,2.96250) (3.10015,2.96262) (3.12239,2.96294) (3.14462,2.96345) (3.16683,2.96416) (3.18901,2.96505) (3.21115,2.96615) (3.23326,2.96743) (3.25531,2.96891) (3.27731,2.97058) (3.29925,2.97244) (3.32112,2.97449) (3.34292,2.97673) (3.36463,2.97917) (3.38626,2.98179) (3.40779,2.98460) (3.42922,2.98760) (3.45054,2.99078) (3.47175,2.99415) (3.49284,2.99770) (3.51380,3.00144) (3.53462,3.00536) (3.55531,3.00946) (3.57585,3.01375) (3.59623,3.01821) (3.61646,3.02284) (3.63652,3.02765) (3.65642,3.03264) (3.67613,3.03780) (3.69566,3.04313) (3.71501,3.04863) (3.73415,3.05430) (3.75310,3.06014) (3.77184,3.06614) (3.79037,3.07230) (3.80868,3.07862) (3.82677,3.08510) (3.84462,3.09174) (3.86225,3.09854) (3.87963,3.10548) (3.89677,3.11258) (3.91366,3.11982) (3.93029,3.12722) (3.94666,3.13475) (3.96277,3.14243) (3.97861,3.15024) (3.99417,3.15820) (4.00945,3.16629) (4.02444,3.17451) (4.03915,3.18286) (4.05357,3.19133) (4.06768,3.19993) (4.08149,3.20866) (4.09500,3.21750)};
\filldraw[fill=courseink,draw=coursepaper,line width=.4pt] (4.35000,3.60000) circle[radius=2pt];
\node[above,inner sep=3pt] at (4.35000,3.60000) {1};
\node[below left,inner sep=2pt] at (2.65,3.6) {$0$};
\draw[coursemuted!25,thin] plot coordinates {(7.10000,0.85000) (16.50000,0.85000)};
\node[left,inner sep=3pt] at (7.1,0.85) {$0$};
\draw[coursemuted!25,thin] plot coordinates {(7.10000,2.15000) (16.50000,2.15000)};
\node[left,inner sep=3pt] at (7.1,2.15) {$0{,}5$};
\draw[courseink!65,dashed] plot coordinates {(7.10000,3.45000) (16.50000,3.45000)};
\node[left,inner sep=3pt] at (7.1,3.45) {$1$};
\draw[->,coursemuted] plot coordinates {(7.10000,0.85000) (16.60000,0.85000)};
\draw[->,coursemuted] plot coordinates {(7.10000,0.85000) (7.10000,3.84000)};
\node[rotate=90] at (5.9,2.2) {$\|\boldsymbol s(\varphi)\|$};
\draw[courseaccent,line width=1.15pt] plot coordinates {(7.10000,3.13612) (7.11306,3.14420) (7.12611,3.15219) (7.13917,3.16010) (7.15222,3.16791) (7.16528,3.17562) (7.17833,3.18325) (7.19139,3.19078) (7.20444,3.19821) (7.21750,3.20555) (7.23056,3.21280) (7.24361,3.21994) (7.25667,3.22699) (7.26972,3.23394) (7.28278,3.24079) (7.29583,3.24754) (7.30889,3.25419) (7.32194,3.26074) (7.33500,3.26718) (7.34806,3.27352) (7.36111,3.27976) (7.37417,3.28590) (7.38722,3.29193) (7.40028,3.29785) (7.41333,3.30367) (7.42639,3.30938) (7.43944,3.31499) (7.45250,3.32049) (7.46556,3.32588) (7.47861,3.33116) (7.49167,3.33633) (7.50472,3.34139) (7.51778,3.34635) (7.53083,3.35119) (7.54389,3.35592) (7.55694,3.36054) (7.57000,3.36504) (7.58306,3.36944) (7.59611,3.37372) (7.60917,3.37789) (7.62222,3.38194) (7.63528,3.38588) (7.64833,3.38971) (7.66139,3.39342) (7.67444,3.39702) (7.68750,3.40050) (7.70056,3.40386) (7.71361,3.40711) (7.72667,3.41024) (7.73972,3.41325) (7.75278,3.41615) (7.76583,3.41893) (7.77889,3.42159) (7.79194,3.42414) (7.80500,3.42657) (7.81806,3.42887) (7.83111,3.43106) (7.84417,3.43313) (7.85722,3.43508) (7.87028,3.43692) (7.88333,3.43863) (7.89639,3.44022) (7.90944,3.44169) (7.92250,3.44305) (7.93556,3.44428) (7.94861,3.44539) (7.96167,3.44639) (7.97472,3.44726) (7.98778,3.44801) (8.00083,3.44864) (8.01389,3.44916) (8.02694,3.44955) (8.04000,3.44982) (8.05306,3.44997) (8.06611,3.45000) (8.07917,3.44990) (8.09222,3.44969) (8.10528,3.44936) (8.11833,3.44891) (8.13139,3.44833) (8.14444,3.44764) (8.15750,3.44682) (8.17056,3.44589) (8.18361,3.44483) (8.19667,3.44365) (8.20972,3.44236) (8.22278,3.44094) (8.23583,3.43941) (8.24889,3.43775) (8.26194,3.43598) (8.27500,3.43408) (8.28806,3.43207) (8.30111,3.42994) (8.31417,3.42769) (8.32722,3.42532) (8.34028,3.42283) (8.35333,3.42022) (8.36639,3.41750) (8.37944,3.41466) (8.39250,3.41170) (8.40556,3.40862) (8.41861,3.40543) (8.43167,3.40212) (8.44472,3.39870) (8.45778,3.39516) (8.47083,3.39150) (8.48389,3.38773) (8.49694,3.38385) (8.51000,3.37985) (8.52306,3.37573) (8.53611,3.37151) (8.54917,3.36717) (8.56222,3.36271) (8.57528,3.35815) (8.58833,3.35347) (8.60139,3.34868) (8.61444,3.34378) (8.62750,3.33877) (8.64056,3.33365) (8.65361,3.32842) (8.66667,3.32309) (8.67972,3.31764) (8.69278,3.31209) (8.70583,3.30642) (8.71889,3.30066) (8.73194,3.29478) (8.74500,3.28880) (8.75806,3.28272) (8.77111,3.27653) (8.78417,3.27024) (8.79722,3.26384) (8.81028,3.25734) (8.82333,3.25074) (8.83639,3.24404) (8.84944,3.23724) (8.86250,3.23034) (8.87556,3.22334) (8.88861,3.21624) (8.90167,3.20904) (8.91472,3.20175) (8.92778,3.19436) (8.94083,3.18687) (8.95389,3.17929) (8.96694,3.17162) (8.98000,3.16385) (8.99306,3.15599) (9.00611,3.14804) (9.01917,3.14000) (9.03222,3.13187) (9.04528,3.12365) (9.05833,3.11535) (9.07139,3.10695) (9.08444,3.09847) (9.09750,3.08991) (9.11056,3.08126) (9.12361,3.07252) (9.13667,3.06371) (9.14972,3.05481) (9.16278,3.04583) (9.17583,3.03678) (9.18889,3.02764) (9.20194,3.01843) (9.21500,3.00914) (9.22806,2.99977) (9.24111,2.99033) (9.25417,2.98082) (9.26722,2.97124) (9.28028,2.96158) (9.29333,2.95185) (9.30639,2.94206) (9.31944,2.93220) (9.33250,2.92227) (9.34556,2.91227) (9.35861,2.90221) (9.37167,2.89209) (9.38472,2.88191) (9.39778,2.87166) (9.41083,2.86136) (9.42389,2.85099) (9.43694,2.84058) (9.45000,2.83010) (9.46306,2.81957) (9.47611,2.80899) (9.48917,2.79836) (9.50222,2.78767) (9.51528,2.77694) (9.52833,2.76616) (9.54139,2.75534) (9.55444,2.74447) (9.56750,2.73355) (9.58056,2.72260) (9.59361,2.71160) (9.60667,2.70057) (9.61972,2.68950) (9.63278,2.67839) (9.64583,2.66725) (9.65889,2.65607) (9.67194,2.64487) (9.68500,2.63363) (9.69806,2.62237) (9.71111,2.61109) (9.72417,2.59977) (9.73722,2.58844) (9.75028,2.57708) (9.76333,2.56571) (9.77639,2.55432) (9.78944,2.54291) (9.80250,2.53149) (9.81556,2.52005) (9.82861,2.50861) (9.84167,2.49716) (9.85472,2.48570) (9.86778,2.47424) (9.88083,2.46277) (9.89389,2.45131) (9.90694,2.43985) (9.92000,2.42839) (9.93306,2.41693) (9.94611,2.40549) (9.95917,2.39405) (9.97222,2.38263) (9.98528,2.37122) (9.99833,2.35983) (10.01139,2.34846) (10.02444,2.33711) (10.03750,2.32578) (10.05056,2.31448) (10.06361,2.30320) (10.07667,2.29196) (10.08972,2.28075) (10.10278,2.26958) (10.11583,2.25845) (10.12889,2.24735) (10.14194,2.23630) (10.15500,2.22530) (10.16806,2.21434) (10.18111,2.20344) (10.19417,2.19259) (10.20722,2.18180) (10.22028,2.17107) (10.23333,2.16040) (10.24639,2.14980) (10.25944,2.13927) (10.27250,2.12880) (10.28556,2.11842) (10.29861,2.10811) (10.31167,2.09788) (10.32472,2.08774) (10.33778,2.07768) (10.35083,2.06771) (10.36389,2.05784) (10.37694,2.04806) (10.39000,2.03839) (10.40306,2.02882) (10.41611,2.01935) (10.42917,2.01000) (10.44222,2.00076) (10.45528,1.99163) (10.46833,1.98263) (10.48139,1.97375) (10.49444,1.96499) (10.50750,1.95637) (10.52056,1.94788) (10.53361,1.93953) (10.54667,1.93132) (10.55972,1.92325) (10.57278,1.91534) (10.58583,1.90757) (10.59889,1.89996) (10.61194,1.89250) (10.62500,1.88521) (10.63806,1.87808) (10.65111,1.87112) (10.66417,1.86433) (10.67722,1.85771) (10.69028,1.85127) (10.70333,1.84501) (10.71639,1.83893) (10.72944,1.83304) (10.74250,1.82733) (10.75556,1.82182) (10.76861,1.81650) (10.78167,1.81137) (10.79472,1.80644) (10.80778,1.80171) (10.82083,1.79718) (10.83389,1.79285) (10.84694,1.78873) (10.86000,1.78481) (10.87306,1.78110) (10.88611,1.77760) (10.89917,1.77431) (10.91222,1.77122) (10.92528,1.76835) (10.93833,1.76568) (10.95139,1.76323) (10.96444,1.76099) (10.97750,1.75895) (10.99056,1.75712) (11.00361,1.75551) (11.01667,1.75410) (11.02972,1.75289) (11.04278,1.75189) (11.05583,1.75110) (11.06889,1.75050) (11.08194,1.75011) (11.09500,1.74991) (11.10806,1.74991) (11.12111,1.75010) (11.13417,1.75048) (11.14722,1.75104) (11.16028,1.75179) (11.17333,1.75272) (11.18639,1.75383) (11.19944,1.75511) (11.21250,1.75656) (11.22556,1.75818) (11.23861,1.75997) (11.25167,1.76191) (11.26472,1.76400) (11.27778,1.76625) (11.29083,1.76865) (11.30389,1.77119) (11.31694,1.77386) (11.33000,1.77667) (11.34306,1.77962) (11.35611,1.78268) (11.36917,1.78587) (11.38222,1.78918) (11.39528,1.79260) (11.40833,1.79612) (11.42139,1.79975) (11.43444,1.80349) (11.44750,1.80731) (11.46056,1.81123) (11.47361,1.81523) (11.48667,1.81932) (11.49972,1.82348) (11.51278,1.82772) (11.52583,1.83203) (11.53889,1.83640) (11.55194,1.84083) (11.56500,1.84533) (11.57806,1.84987) (11.59111,1.85447) (11.60417,1.85911) (11.61722,1.86379) (11.63028,1.86852) (11.64333,1.87327) (11.65639,1.87806) (11.66944,1.88288) (11.68250,1.88772) (11.69556,1.89258) (11.70861,1.89745) (11.72167,1.90235) (11.73472,1.90725) (11.74778,1.91216) (11.76083,1.91707) (11.77389,1.92199) (11.78694,1.92690) (11.80000,1.93182) (11.81306,1.93672) (11.82611,1.94161) (11.83917,1.94649) (11.85222,1.95136) (11.86528,1.95621) (11.87833,1.96104) (11.89139,1.96584) (11.90444,1.97062) (11.91750,1.97537) (11.93056,1.98010) (11.94361,1.98479) (11.95667,1.98944) (11.96972,1.99406) (11.98278,1.99864) (11.99583,2.00319) (12.00889,2.00768) (12.02194,2.01214) (12.03500,2.01655) (12.04806,2.02090) (12.06111,2.02521) (12.07417,2.02947) (12.08722,2.03368) (12.10028,2.03783) (12.11333,2.04192) (12.12639,2.04595) (12.13944,2.04993) (12.15250,2.05384) (12.16556,2.05769) (12.17861,2.06148) (12.19167,2.06520) (12.20472,2.06886) (12.21778,2.07245) (12.23083,2.07597) (12.24389,2.07942) (12.25694,2.08279) (12.27000,2.08610) (12.28306,2.08933) (12.29611,2.09248) (12.30917,2.09556) (12.32222,2.09857) (12.33528,2.10149) (12.34833,2.10434) (12.36139,2.10711) (12.37444,2.10979) (12.38750,2.11240) (12.40056,2.11492) (12.41361,2.11736) (12.42667,2.11972) (12.43972,2.12200) (12.45278,2.12418) (12.46583,2.12629) (12.47889,2.12831) (12.49194,2.13024) (12.50500,2.13208) (12.51806,2.13383) (12.53111,2.13550) (12.54417,2.13708) (12.55722,2.13857) (12.57028,2.13997) (12.58333,2.14128) (12.59639,2.14250) (12.60944,2.14362) (12.62250,2.14466) (12.63556,2.14561) (12.64861,2.14646) (12.66167,2.14722) (12.67472,2.14789) (12.68778,2.14847) (12.70083,2.14896) (12.71389,2.14935) (12.72694,2.14965) (12.74000,2.14986) (12.75306,2.14997) (12.76611,2.15000) (12.77917,2.14993) (12.79222,2.14976) (12.80528,2.14951) (12.81833,2.14916) (12.83139,2.14872) (12.84444,2.14818) (12.85750,2.14756) (12.87056,2.14684) (12.88361,2.14603) (12.89667,2.14513) (12.90972,2.14413) (12.92278,2.14305) (12.93583,2.14187) (12.94889,2.14061) (12.96194,2.13925) (12.97500,2.13780) (12.98806,2.13627) (13.00111,2.13464) (13.01417,2.13293) (13.02722,2.13113) (13.04028,2.12924) (13.05333,2.12727) (13.06639,2.12520) (13.07944,2.12306) (13.09250,2.12082) (13.10556,2.11851) (13.11861,2.11610) (13.13167,2.11362) (13.14472,2.11105) (13.15778,2.10841) (13.17083,2.10568) (13.18389,2.10287) (13.19694,2.09998) (13.21000,2.09701) (13.22306,2.09397) (13.23611,2.09085) (13.24917,2.08766) (13.26222,2.08439) (13.27528,2.08104) (13.28833,2.07763) (13.30139,2.07414) (13.31444,2.07059) (13.32750,2.06697) (13.34056,2.06328) (13.35361,2.05952) (13.36667,2.05570) (13.37972,2.05181) (13.39278,2.04787) (13.40583,2.04386) (13.41889,2.03980) (13.43194,2.03567) (13.44500,2.03150) (13.45806,2.02726) (13.47111,2.02298) (13.48417,2.01864) (13.49722,2.01426) (13.51028,2.00983) (13.52333,2.00535) (13.53639,2.00083) (13.54944,1.99627) (13.56250,1.99166) (13.57556,1.98703) (13.58861,1.98235) (13.60167,1.97764) (13.61472,1.97290) (13.62778,1.96814) (13.64083,1.96334) (13.65389,1.95853) (13.66694,1.95369) (13.68000,1.94883) (13.69306,1.94396) (13.70611,1.93907) (13.71917,1.93417) (13.73222,1.92926) (13.74528,1.92435) (13.75833,1.91943) (13.77139,1.91452) (13.78444,1.90960) (13.79750,1.90470) (13.81056,1.89980) (13.82361,1.89491) (13.83667,1.89004) (13.84972,1.88519) (13.86278,1.88037) (13.87583,1.87557) (13.88889,1.87079) (13.90194,1.86605) (13.91500,1.86135) (13.92806,1.85669) (13.94111,1.85207) (13.95417,1.84750) (13.96722,1.84298) (13.98028,1.83852) (13.99333,1.83412) (14.00639,1.82978) (14.01944,1.82550) (14.03250,1.82130) (14.04556,1.81718) (14.05861,1.81314) (14.07167,1.80918) (14.08472,1.80531) (14.09778,1.80153) (14.11083,1.79785) (14.12389,1.79428) (14.13694,1.79080) (14.15000,1.78744) (14.16306,1.78420) (14.17611,1.78107) (14.18917,1.77807) (14.20222,1.77519) (14.21528,1.77245) (14.22833,1.76985) (14.24139,1.76738) (14.25444,1.76506) (14.26750,1.76289) (14.28056,1.76088) (14.29361,1.75902) (14.30667,1.75732) (14.31972,1.75579) (14.33278,1.75442) (14.34583,1.75323) (14.35889,1.75222) (14.37194,1.75138) (14.38500,1.75072) (14.39806,1.75025) (14.41111,1.74997) (14.42417,1.74988) (14.43722,1.74999) (14.45028,1.75029) (14.46333,1.75079) (14.47639,1.75149) (14.48944,1.75239) (14.50250,1.75349) (14.51556,1.75481) (14.52861,1.75632) (14.54167,1.75805) (14.55472,1.75998) (14.56778,1.76213) (14.58083,1.76448) (14.59389,1.76704) (14.60694,1.76982) (14.62000,1.77280) (14.63306,1.77599) (14.64611,1.77940) (14.65917,1.78301) (14.67222,1.78683) (14.68528,1.79085) (14.69833,1.79508) (14.71139,1.79951) (14.72444,1.80415) (14.73750,1.80898) (14.75056,1.81401) (14.76361,1.81924) (14.77667,1.82467) (14.78972,1.83028) (14.80278,1.83608) (14.81583,1.84207) (14.82889,1.84825) (14.84194,1.85460) (14.85500,1.86113) (14.86806,1.86784) (14.88111,1.87472) (14.89417,1.88177) (14.90722,1.88899) (14.92028,1.89636) (14.93333,1.90390) (14.94639,1.91159) (14.95944,1.91944) (14.97250,1.92743) (14.98556,1.93558) (14.99861,1.94386) (15.01167,1.95228) (15.02472,1.96084) (15.03778,1.96953) (15.05083,1.97835) (15.06389,1.98730) (15.07694,1.99636) (15.09000,2.00555) (15.10306,2.01485) (15.11611,2.02426) (15.12917,2.03379) (15.14222,2.04341) (15.15528,2.05314) (15.16833,2.06297) (15.18139,2.07289) (15.19444,2.08290) (15.20750,2.09300) (15.22056,2.10319) (15.23361,2.11346) (15.24667,2.12381) (15.25972,2.13424) (15.27278,2.14474) (15.28583,2.15531) (15.29889,2.16595) (15.31194,2.17665) (15.32500,2.18741) (15.33806,2.19823) (15.35111,2.20911) (15.36417,2.22004) (15.37722,2.23102) (15.39028,2.24205) (15.40333,2.25312) (15.41639,2.26424) (15.42944,2.27539) (15.44250,2.28658) (15.45556,2.29781) (15.46861,2.30907) (15.48167,2.32035) (15.49472,2.33167) (15.50778,2.34301) (15.52083,2.35437) (15.53389,2.36576) (15.54694,2.37716) (15.56000,2.38857) (15.57306,2.40000) (15.58611,2.41144) (15.59917,2.42289) (15.61222,2.43435) (15.62528,2.44581) (15.63833,2.45727) (15.65139,2.46874) (15.66444,2.48020) (15.67750,2.49166) (15.69056,2.50312) (15.70361,2.51457) (15.71667,2.52600) (15.72972,2.53743) (15.74278,2.54885) (15.75583,2.56025) (15.76889,2.57163) (15.78194,2.58299) (15.79500,2.59434) (15.80806,2.60566) (15.82111,2.61696) (15.83417,2.62824) (15.84722,2.63948) (15.86028,2.65070) (15.87333,2.66189) (15.88639,2.67305) (15.89944,2.68417) (15.91250,2.69526) (15.92556,2.70631) (15.93861,2.71733) (15.95167,2.72830) (15.96472,2.73924) (15.97778,2.75013) (15.99083,2.76098) (16.00389,2.77178) (16.01694,2.78253) (16.03000,2.79324) (16.04306,2.80390) (16.05611,2.81450) (16.06917,2.82506) (16.08222,2.83556) (16.09528,2.84600) (16.10833,2.85639) (16.12139,2.86673) (16.13444,2.87700) (16.14750,2.88721) (16.16056,2.89736) (16.17361,2.90745) (16.18667,2.91748) (16.19972,2.92744) (16.21278,2.93734) (16.22583,2.94716) (16.23889,2.95692) (16.25194,2.96661) (16.26500,2.97623) (16.27806,2.98578) (16.29111,2.99526) (16.30417,3.00466) (16.31722,3.01398) (16.33028,3.02323) (16.34333,3.03241) (16.35639,3.04150) (16.36944,3.05052) (16.38250,3.05945) (16.39556,3.06831) (16.40861,3.07708) (16.42167,3.08577) (16.43472,3.09438) (16.44778,3.10290) (16.46083,3.11133) (16.47389,3.11968) (16.48694,3.12794) (16.50000,3.13612)};
\draw[->,coursemuted] plot coordinates {(7.10000,4.65000) (16.60000,4.65000)};
\draw[->,coursemuted] plot coordinates {(7.10000,4.65000) (7.10000,6.50000)};
\draw[coursemuted!25,thin] plot coordinates {(7.10000,4.65000) (16.50000,4.65000)};
\node[left,inner sep=3pt] at (7.1,4.65000) {$0$};
\draw[coursemuted!25,thin] plot coordinates {(7.10000,5.45000) (16.50000,5.45000)};
\node[left,inner sep=3pt] at (7.1,5.45000) {$5$};
\draw[coursemuted!25,thin] plot coordinates {(7.10000,6.25000) (16.50000,6.25000)};
\node[left,inner sep=3pt] at (7.1,6.25000) {$10$};
\node[rotate=90] at (5.9,5.55) {$\Delta V_j$, м/с};
\draw[coursemuted!16,thin] plot coordinates {(7.10000,0.80000) (7.10000,6.25000)};
\node[below,inner sep=3pt] at (7.10000,0.85) {$0$};
\draw[coursemuted!16,thin] plot coordinates {(9.45000,0.80000) (9.45000,6.25000)};
\node[below,inner sep=3pt] at (9.45000,0.85) {$\pi/2$};
\draw[coursemuted!16,thin] plot coordinates {(11.80000,0.80000) (11.80000,6.25000)};
\node[below,inner sep=3pt] at (11.80000,0.85) {$\pi$};
\draw[coursemuted!16,thin] plot coordinates {(14.15000,0.80000) (14.15000,6.25000)};
\node[below,inner sep=3pt] at (14.15000,0.85) {$3\pi/2$};
\draw[coursemuted!16,thin] plot coordinates {(16.50000,0.80000) (16.50000,6.25000)};
\node[below,inner sep=3pt] at (16.50000,0.85) {$2\pi$};
\node at (11.8,.15) {Угол $\varphi$, рад};
\draw[courseink!70,densely dashed,thin] plot coordinates {(8.06271,0.85000) (8.06271,3.45000)};
\draw[courseink,line width=1pt] plot coordinates {(8.06271,4.65000) (8.06271,5.85737)};
\fill[courseink] (8.06271,5.85737) circle[radius=1.8pt];
\fill[courseink] (8.06271,3.45000) circle[radius=1.8pt];
\node[above,inner sep=2pt] at (8.06271,5.85737) {1};
\end{tikzpicture}

\par
\figcaption{Во всех примерах: $a_i=r_0=7000$ км, $V_0=7546{,}053$ м/с.\\
$\Delta a=0{,}002$, $\Delta e_x=0{,}0016$, $\Delta e_y=0{,}0012$, $\Delta i=0$.}
\end{frame}

% eq:planar-tangent-impulse; вывод из eq:planar-boundaries.
% После кадра «Касающиеся орбиты: годограф».
% manual/chapter-3-algorithms.tex, sec:planar-tangent.
\begin{frame}{Касающиеся орбиты}
\relax
\derivationmathspacing
Для одного трансверсального импульса три граничных условия имеют вид
\[
\begin{aligned}
 2\Delta\widehat V_t&=\Delta a,\\
 2\Delta\widehat V_t\cos\varphi_1&=\Delta e_x,\\
 2\Delta\widehat V_t\sin\varphi_1&=\Delta e_y.
\end{aligned}
\]
Первое равенство сразу определяет компоненту импульса:
\[
 \Delta\widehat V_t=\frac{\Delta a}{2}.
\]
\justifying
Чтобы найти место импульса и проверить совместность условий,
возведём второе и третье равенства в квадрат и сложим.
\end{frame}

\begin{frame}{Касающиеся орбиты}
\relax
\derivationmathspacing
Сумма квадратов двух условий на эксцентриситет даёт
\[
\begin{aligned}
 4(\Delta\widehat V_t)^2(\cos^2\varphi_1+\sin^2\varphi_1)
   &=(\Delta e_x)^2+(\Delta e_y)^2,\\
 4(\Delta\widehat V_t)^2&=\Delta e^2.
\end{aligned}
\]
С учётом $2\Delta\widehat V_t=\Delta a$ получаем
$(\Delta a)^2=\Delta e^2$, то есть $|\Delta a|=\Delta e$.
При $\Delta e>0$ делим условия на $\Delta a$:
\[
\begin{aligned}
 \cos\varphi_1&=\frac{\Delta e_x}{\Delta a}
   =\operatorname{sgn}(\Delta a)\cos\varphi_e,\\
 \sin\varphi_1&=\frac{\Delta e_y}{\Delta a}
   =\operatorname{sgn}(\Delta a)\sin\varphi_e.
\end{aligned}
\]
При $\Delta a<0$ оба знака меняются: угол сдвигается на $\pi$.
\end{frame}

\begin{frame}{Касающиеся орбиты}
\relax
\derivationmathspacing
\begin{block}{Один трансверсальный импульс при $|\Delta a|=\Delta e>0$}
\[
\begin{gathered}
 \Delta\widehat V_r=0,\qquad
 \Delta\widehat V_t=\frac{\Delta a}{2},\qquad
 J=|\Delta\widehat V_t|=\frac{|\Delta a|}{2},\\[2pt]
 \varphi_1=\begin{cases}
   \varphi_e,&\Delta a>0,\\
   \varphi_e+\pi,&\Delta a<0.
 \end{cases}
\end{gathered}
\]
\end{block}
\justifying
Все три граничных условия выполнены. Импульс сонаправлен
единичному базис-вектору в единственной точке касания годографа.
На остальном обороте его модуль меньше единицы;
достаточные условия оптимальности выполнены.
\par\medskip
При $\Delta a=\Delta e=0$ заданные орбиты совпадают,
поэтому манёвр не требуется.
\end{frame}


% eq:planar-point-primer; sec:planar-disjoint.
\begin{frame}{Непересекающиеся орбиты}
\relax
Пусть $|\Delta a|>\Delta e$. Одна орбита целиком снаружи другой.
\begin{block}{Годограф вырождается в точку}
\derivationmathspacing
\[
\begin{gathered}
 \boldsymbol s=(0,\operatorname{sgn}\Delta a,0),\\
 \lambda_1=\tfrac12\operatorname{sgn}\Delta a,\qquad
 \lambda_2=\lambda_3=\lambda_4=\lambda_5=0.
\end{gathered}
\]
Оба импульса разгонные при $\Delta a>0$ и тормозные при $\Delta a<0$.
\end{block}
{\justifying
Модуль базис-вектора равен единице на всём обороте,
поэтому он не выделяет места импульсов. Локальные компоненты
постоянны, трансверсальное направление поворачивается вместе с аппаратом.
\par}
\end{frame}

\begin{frame}{Непересекающиеся орбиты: годограф}
\relax
\centering
% Рисунок из методички; координаты кривых, осей и отметок сохранены.
\begin{tikzpicture}[course diagram,scale=.72]
\node at (2.65,6.85) {Годограф};
\node at (11.6,6.85) {Модули импульсов};
\node at (11.6,4.1) {Модуль базис-вектора};
\draw[coursemuted!60,dashed] (2.65,3.6) circle[radius=1.7cm];
\draw[->,coursemuted] (.6,3.6)--(4.85,3.6) node[right] {$s_t$};
\draw[->,coursemuted] (2.65,1.55)--(2.65,5.8) node[above] {$s_r$};
\draw[courseaccent,line width=1.15pt] plot coordinates {(4.35000,3.60000) (4.35000,3.60000) (4.35000,3.60000) (4.35000,3.60000) (4.35000,3.60000) (4.35000,3.60000) (4.35000,3.60000) (4.35000,3.60000) (4.35000,3.60000) (4.35000,3.60000) (4.35000,3.60000) (4.35000,3.60000) (4.35000,3.60000) (4.35000,3.60000) (4.35000,3.60000) (4.35000,3.60000) (4.35000,3.60000) (4.35000,3.60000) (4.35000,3.60000) (4.35000,3.60000) (4.35000,3.60000) (4.35000,3.60000) (4.35000,3.60000) (4.35000,3.60000) (4.35000,3.60000) (4.35000,3.60000) (4.35000,3.60000) (4.35000,3.60000) (4.35000,3.60000) (4.35000,3.60000) (4.35000,3.60000) (4.35000,3.60000) (4.35000,3.60000) (4.35000,3.60000) (4.35000,3.60000) (4.35000,3.60000) (4.35000,3.60000) (4.35000,3.60000) (4.35000,3.60000) (4.35000,3.60000) (4.35000,3.60000) (4.35000,3.60000) (4.35000,3.60000) (4.35000,3.60000) (4.35000,3.60000) (4.35000,3.60000) (4.35000,3.60000) (4.35000,3.60000) (4.35000,3.60000) (4.35000,3.60000) (4.35000,3.60000) (4.35000,3.60000) (4.35000,3.60000) (4.35000,3.60000) (4.35000,3.60000) (4.35000,3.60000) (4.35000,3.60000) (4.35000,3.60000) (4.35000,3.60000) (4.35000,3.60000) (4.35000,3.60000) (4.35000,3.60000) (4.35000,3.60000) (4.35000,3.60000) (4.35000,3.60000) (4.35000,3.60000) (4.35000,3.60000) (4.35000,3.60000) (4.35000,3.60000) (4.35000,3.60000) (4.35000,3.60000) (4.35000,3.60000) (4.35000,3.60000) (4.35000,3.60000) (4.35000,3.60000) (4.35000,3.60000) (4.35000,3.60000) (4.35000,3.60000) (4.35000,3.60000) (4.35000,3.60000) (4.35000,3.60000) (4.35000,3.60000) (4.35000,3.60000) (4.35000,3.60000) (4.35000,3.60000) (4.35000,3.60000) (4.35000,3.60000) (4.35000,3.60000) (4.35000,3.60000) (4.35000,3.60000) (4.35000,3.60000) (4.35000,3.60000) (4.35000,3.60000) (4.35000,3.60000) (4.35000,3.60000) (4.35000,3.60000) (4.35000,3.60000) (4.35000,3.60000) (4.35000,3.60000) (4.35000,3.60000) (4.35000,3.60000) (4.35000,3.60000) (4.35000,3.60000) (4.35000,3.60000) (4.35000,3.60000) (4.35000,3.60000) (4.35000,3.60000) (4.35000,3.60000) (4.35000,3.60000) (4.35000,3.60000) (4.35000,3.60000) (4.35000,3.60000) (4.35000,3.60000) (4.35000,3.60000) (4.35000,3.60000) (4.35000,3.60000) (4.35000,3.60000) (4.35000,3.60000) (4.35000,3.60000) (4.35000,3.60000) (4.35000,3.60000) (4.35000,3.60000) (4.35000,3.60000) (4.35000,3.60000) (4.35000,3.60000) (4.35000,3.60000) (4.35000,3.60000) (4.35000,3.60000) (4.35000,3.60000) (4.35000,3.60000) (4.35000,3.60000) (4.35000,3.60000) (4.35000,3.60000) (4.35000,3.60000) (4.35000,3.60000) (4.35000,3.60000) (4.35000,3.60000) (4.35000,3.60000) (4.35000,3.60000) (4.35000,3.60000) (4.35000,3.60000) (4.35000,3.60000) (4.35000,3.60000) (4.35000,3.60000) (4.35000,3.60000) (4.35000,3.60000) (4.35000,3.60000) (4.35000,3.60000) (4.35000,3.60000) (4.35000,3.60000) (4.35000,3.60000) (4.35000,3.60000) (4.35000,3.60000) (4.35000,3.60000) (4.35000,3.60000) (4.35000,3.60000) (4.35000,3.60000) (4.35000,3.60000) (4.35000,3.60000) (4.35000,3.60000) (4.35000,3.60000) (4.35000,3.60000) (4.35000,3.60000) (4.35000,3.60000) (4.35000,3.60000) (4.35000,3.60000) (4.35000,3.60000) (4.35000,3.60000) (4.35000,3.60000) (4.35000,3.60000) (4.35000,3.60000) (4.35000,3.60000) (4.35000,3.60000) (4.35000,3.60000) (4.35000,3.60000) (4.35000,3.60000) (4.35000,3.60000) (4.35000,3.60000) (4.35000,3.60000) (4.35000,3.60000) (4.35000,3.60000) (4.35000,3.60000) (4.35000,3.60000) (4.35000,3.60000) (4.35000,3.60000) (4.35000,3.60000) (4.35000,3.60000) (4.35000,3.60000) (4.35000,3.60000) (4.35000,3.60000) (4.35000,3.60000) (4.35000,3.60000) (4.35000,3.60000) (4.35000,3.60000) (4.35000,3.60000) (4.35000,3.60000) (4.35000,3.60000) (4.35000,3.60000) (4.35000,3.60000) (4.35000,3.60000) (4.35000,3.60000) (4.35000,3.60000) (4.35000,3.60000) (4.35000,3.60000) (4.35000,3.60000) (4.35000,3.60000) (4.35000,3.60000) (4.35000,3.60000) (4.35000,3.60000) (4.35000,3.60000) (4.35000,3.60000) (4.35000,3.60000) (4.35000,3.60000) (4.35000,3.60000) (4.35000,3.60000) (4.35000,3.60000) (4.35000,3.60000) (4.35000,3.60000) (4.35000,3.60000) (4.35000,3.60000) (4.35000,3.60000) (4.35000,3.60000) (4.35000,3.60000) (4.35000,3.60000) (4.35000,3.60000) (4.35000,3.60000) (4.35000,3.60000) (4.35000,3.60000) (4.35000,3.60000) (4.35000,3.60000) (4.35000,3.60000) (4.35000,3.60000) (4.35000,3.60000) (4.35000,3.60000) (4.35000,3.60000) (4.35000,3.60000) (4.35000,3.60000) (4.35000,3.60000) (4.35000,3.60000) (4.35000,3.60000) (4.35000,3.60000) (4.35000,3.60000) (4.35000,3.60000) (4.35000,3.60000) (4.35000,3.60000) (4.35000,3.60000) (4.35000,3.60000) (4.35000,3.60000) (4.35000,3.60000) (4.35000,3.60000) (4.35000,3.60000) (4.35000,3.60000) (4.35000,3.60000) (4.35000,3.60000) (4.35000,3.60000) (4.35000,3.60000) (4.35000,3.60000) (4.35000,3.60000) (4.35000,3.60000) (4.35000,3.60000) (4.35000,3.60000) (4.35000,3.60000) (4.35000,3.60000) (4.35000,3.60000) (4.35000,3.60000) (4.35000,3.60000) (4.35000,3.60000) (4.35000,3.60000) (4.35000,3.60000) (4.35000,3.60000) (4.35000,3.60000) (4.35000,3.60000) (4.35000,3.60000) (4.35000,3.60000) (4.35000,3.60000) (4.35000,3.60000) (4.35000,3.60000) (4.35000,3.60000) (4.35000,3.60000) (4.35000,3.60000) (4.35000,3.60000) (4.35000,3.60000) (4.35000,3.60000) (4.35000,3.60000) (4.35000,3.60000) (4.35000,3.60000) (4.35000,3.60000) (4.35000,3.60000) (4.35000,3.60000) (4.35000,3.60000) (4.35000,3.60000) (4.35000,3.60000) (4.35000,3.60000) (4.35000,3.60000) (4.35000,3.60000) (4.35000,3.60000) (4.35000,3.60000) (4.35000,3.60000) (4.35000,3.60000) (4.35000,3.60000) (4.35000,3.60000) (4.35000,3.60000) (4.35000,3.60000) (4.35000,3.60000) (4.35000,3.60000) (4.35000,3.60000) (4.35000,3.60000) (4.35000,3.60000) (4.35000,3.60000) (4.35000,3.60000) (4.35000,3.60000) (4.35000,3.60000) (4.35000,3.60000) (4.35000,3.60000) (4.35000,3.60000) (4.35000,3.60000) (4.35000,3.60000) (4.35000,3.60000) (4.35000,3.60000) (4.35000,3.60000) (4.35000,3.60000) (4.35000,3.60000) (4.35000,3.60000) (4.35000,3.60000) (4.35000,3.60000) (4.35000,3.60000) (4.35000,3.60000) (4.35000,3.60000) (4.35000,3.60000) (4.35000,3.60000) (4.35000,3.60000) (4.35000,3.60000) (4.35000,3.60000) (4.35000,3.60000) (4.35000,3.60000) (4.35000,3.60000) (4.35000,3.60000) (4.35000,3.60000) (4.35000,3.60000) (4.35000,3.60000) (4.35000,3.60000) (4.35000,3.60000) (4.35000,3.60000) (4.35000,3.60000) (4.35000,3.60000) (4.35000,3.60000) (4.35000,3.60000) (4.35000,3.60000) (4.35000,3.60000) (4.35000,3.60000) (4.35000,3.60000) (4.35000,3.60000) (4.35000,3.60000) (4.35000,3.60000) (4.35000,3.60000) (4.35000,3.60000) (4.35000,3.60000) (4.35000,3.60000) (4.35000,3.60000) (4.35000,3.60000) (4.35000,3.60000)};
\filldraw[fill=courseink,draw=coursepaper,line width=.4pt] (4.35000,3.60000) circle[radius=2pt];
\filldraw[fill=courseink,draw=coursepaper,line width=.4pt] (4.35000,3.60000) circle[radius=2pt];
\node[above,align=center] at (4.35,3.6) {1, 2};
\node[below left,inner sep=2pt] at (2.65,3.6) {$0$};
\draw[coursemuted!25,thin] plot coordinates {(7.10000,0.85000) (16.50000,0.85000)};
\node[left,inner sep=3pt] at (7.1,0.85) {$0$};
\draw[coursemuted!25,thin] plot coordinates {(7.10000,2.15000) (16.50000,2.15000)};
\node[left,inner sep=3pt] at (7.1,2.15) {$0{,}5$};
\draw[courseink!65,dashed] plot coordinates {(7.10000,3.45000) (16.50000,3.45000)};
\node[left,inner sep=3pt] at (7.1,3.45) {$1$};
\draw[->,coursemuted] plot coordinates {(7.10000,0.85000) (16.60000,0.85000)};
\draw[->,coursemuted] plot coordinates {(7.10000,0.85000) (7.10000,3.84000)};
\node[rotate=90] at (5.9,2.2) {$\|\boldsymbol s(\varphi)\|$};
\draw[courseaccent,line width=1.15pt] plot coordinates {(7.10000,3.45000) (7.11306,3.45000) (7.12611,3.45000) (7.13917,3.45000) (7.15222,3.45000) (7.16528,3.45000) (7.17833,3.45000) (7.19139,3.45000) (7.20444,3.45000) (7.21750,3.45000) (7.23056,3.45000) (7.24361,3.45000) (7.25667,3.45000) (7.26972,3.45000) (7.28278,3.45000) (7.29583,3.45000) (7.30889,3.45000) (7.32194,3.45000) (7.33500,3.45000) (7.34806,3.45000) (7.36111,3.45000) (7.37417,3.45000) (7.38722,3.45000) (7.40028,3.45000) (7.41333,3.45000) (7.42639,3.45000) (7.43944,3.45000) (7.45250,3.45000) (7.46556,3.45000) (7.47861,3.45000) (7.49167,3.45000) (7.50472,3.45000) (7.51778,3.45000) (7.53083,3.45000) (7.54389,3.45000) (7.55694,3.45000) (7.57000,3.45000) (7.58306,3.45000) (7.59611,3.45000) (7.60917,3.45000) (7.62222,3.45000) (7.63528,3.45000) (7.64833,3.45000) (7.66139,3.45000) (7.67444,3.45000) (7.68750,3.45000) (7.70056,3.45000) (7.71361,3.45000) (7.72667,3.45000) (7.73972,3.45000) (7.75278,3.45000) (7.76583,3.45000) (7.77889,3.45000) (7.79194,3.45000) (7.80500,3.45000) (7.81806,3.45000) (7.83111,3.45000) (7.84417,3.45000) (7.85722,3.45000) (7.87028,3.45000) (7.88333,3.45000) (7.89639,3.45000) (7.90944,3.45000) (7.92250,3.45000) (7.93556,3.45000) (7.94861,3.45000) (7.96167,3.45000) (7.97472,3.45000) (7.98778,3.45000) (8.00083,3.45000) (8.01389,3.45000) (8.02694,3.45000) (8.04000,3.45000) (8.05306,3.45000) (8.06611,3.45000) (8.07917,3.45000) (8.09222,3.45000) (8.10528,3.45000) (8.11833,3.45000) (8.13139,3.45000) (8.14444,3.45000) (8.15750,3.45000) (8.17056,3.45000) (8.18361,3.45000) (8.19667,3.45000) (8.20972,3.45000) (8.22278,3.45000) (8.23583,3.45000) (8.24889,3.45000) (8.26194,3.45000) (8.27500,3.45000) (8.28806,3.45000) (8.30111,3.45000) (8.31417,3.45000) (8.32722,3.45000) (8.34028,3.45000) (8.35333,3.45000) (8.36639,3.45000) (8.37944,3.45000) (8.39250,3.45000) (8.40556,3.45000) (8.41861,3.45000) (8.43167,3.45000) (8.44472,3.45000) (8.45778,3.45000) (8.47083,3.45000) (8.48389,3.45000) (8.49694,3.45000) (8.51000,3.45000) (8.52306,3.45000) (8.53611,3.45000) (8.54917,3.45000) (8.56222,3.45000) (8.57528,3.45000) (8.58833,3.45000) (8.60139,3.45000) (8.61444,3.45000) (8.62750,3.45000) (8.64056,3.45000) (8.65361,3.45000) (8.66667,3.45000) (8.67972,3.45000) (8.69278,3.45000) (8.70583,3.45000) (8.71889,3.45000) (8.73194,3.45000) (8.74500,3.45000) (8.75806,3.45000) (8.77111,3.45000) (8.78417,3.45000) (8.79722,3.45000) (8.81028,3.45000) (8.82333,3.45000) (8.83639,3.45000) (8.84944,3.45000) (8.86250,3.45000) (8.87556,3.45000) (8.88861,3.45000) (8.90167,3.45000) (8.91472,3.45000) (8.92778,3.45000) (8.94083,3.45000) (8.95389,3.45000) (8.96694,3.45000) (8.98000,3.45000) (8.99306,3.45000) (9.00611,3.45000) (9.01917,3.45000) (9.03222,3.45000) (9.04528,3.45000) (9.05833,3.45000) (9.07139,3.45000) (9.08444,3.45000) (9.09750,3.45000) (9.11056,3.45000) (9.12361,3.45000) (9.13667,3.45000) (9.14972,3.45000) (9.16278,3.45000) (9.17583,3.45000) (9.18889,3.45000) (9.20194,3.45000) (9.21500,3.45000) (9.22806,3.45000) (9.24111,3.45000) (9.25417,3.45000) (9.26722,3.45000) (9.28028,3.45000) (9.29333,3.45000) (9.30639,3.45000) (9.31944,3.45000) (9.33250,3.45000) (9.34556,3.45000) (9.35861,3.45000) (9.37167,3.45000) (9.38472,3.45000) (9.39778,3.45000) (9.41083,3.45000) (9.42389,3.45000) (9.43694,3.45000) (9.45000,3.45000) (9.46306,3.45000) (9.47611,3.45000) (9.48917,3.45000) (9.50222,3.45000) (9.51528,3.45000) (9.52833,3.45000) (9.54139,3.45000) (9.55444,3.45000) (9.56750,3.45000) (9.58056,3.45000) (9.59361,3.45000) (9.60667,3.45000) (9.61972,3.45000) (9.63278,3.45000) (9.64583,3.45000) (9.65889,3.45000) (9.67194,3.45000) (9.68500,3.45000) (9.69806,3.45000) (9.71111,3.45000) (9.72417,3.45000) (9.73722,3.45000) (9.75028,3.45000) (9.76333,3.45000) (9.77639,3.45000) (9.78944,3.45000) (9.80250,3.45000) (9.81556,3.45000) (9.82861,3.45000) (9.84167,3.45000) (9.85472,3.45000) (9.86778,3.45000) (9.88083,3.45000) (9.89389,3.45000) (9.90694,3.45000) (9.92000,3.45000) (9.93306,3.45000) (9.94611,3.45000) (9.95917,3.45000) (9.97222,3.45000) (9.98528,3.45000) (9.99833,3.45000) (10.01139,3.45000) (10.02444,3.45000) (10.03750,3.45000) (10.05056,3.45000) (10.06361,3.45000) (10.07667,3.45000) (10.08972,3.45000) (10.10278,3.45000) (10.11583,3.45000) (10.12889,3.45000) (10.14194,3.45000) (10.15500,3.45000) (10.16806,3.45000) (10.18111,3.45000) (10.19417,3.45000) (10.20722,3.45000) (10.22028,3.45000) (10.23333,3.45000) (10.24639,3.45000) (10.25944,3.45000) (10.27250,3.45000) (10.28556,3.45000) (10.29861,3.45000) (10.31167,3.45000) (10.32472,3.45000) (10.33778,3.45000) (10.35083,3.45000) (10.36389,3.45000) (10.37694,3.45000) (10.39000,3.45000) (10.40306,3.45000) (10.41611,3.45000) (10.42917,3.45000) (10.44222,3.45000) (10.45528,3.45000) (10.46833,3.45000) (10.48139,3.45000) (10.49444,3.45000) (10.50750,3.45000) (10.52056,3.45000) (10.53361,3.45000) (10.54667,3.45000) (10.55972,3.45000) (10.57278,3.45000) (10.58583,3.45000) (10.59889,3.45000) (10.61194,3.45000) (10.62500,3.45000) (10.63806,3.45000) (10.65111,3.45000) (10.66417,3.45000) (10.67722,3.45000) (10.69028,3.45000) (10.70333,3.45000) (10.71639,3.45000) (10.72944,3.45000) (10.74250,3.45000) (10.75556,3.45000) (10.76861,3.45000) (10.78167,3.45000) (10.79472,3.45000) (10.80778,3.45000) (10.82083,3.45000) (10.83389,3.45000) (10.84694,3.45000) (10.86000,3.45000) (10.87306,3.45000) (10.88611,3.45000) (10.89917,3.45000) (10.91222,3.45000) (10.92528,3.45000) (10.93833,3.45000) (10.95139,3.45000) (10.96444,3.45000) (10.97750,3.45000) (10.99056,3.45000) (11.00361,3.45000) (11.01667,3.45000) (11.02972,3.45000) (11.04278,3.45000) (11.05583,3.45000) (11.06889,3.45000) (11.08194,3.45000) (11.09500,3.45000) (11.10806,3.45000) (11.12111,3.45000) (11.13417,3.45000) (11.14722,3.45000) (11.16028,3.45000) (11.17333,3.45000) (11.18639,3.45000) (11.19944,3.45000) (11.21250,3.45000) (11.22556,3.45000) (11.23861,3.45000) (11.25167,3.45000) (11.26472,3.45000) (11.27778,3.45000) (11.29083,3.45000) (11.30389,3.45000) (11.31694,3.45000) (11.33000,3.45000) (11.34306,3.45000) (11.35611,3.45000) (11.36917,3.45000) (11.38222,3.45000) (11.39528,3.45000) (11.40833,3.45000) (11.42139,3.45000) (11.43444,3.45000) (11.44750,3.45000) (11.46056,3.45000) (11.47361,3.45000) (11.48667,3.45000) (11.49972,3.45000) (11.51278,3.45000) (11.52583,3.45000) (11.53889,3.45000) (11.55194,3.45000) (11.56500,3.45000) (11.57806,3.45000) (11.59111,3.45000) (11.60417,3.45000) (11.61722,3.45000) (11.63028,3.45000) (11.64333,3.45000) (11.65639,3.45000) (11.66944,3.45000) (11.68250,3.45000) (11.69556,3.45000) (11.70861,3.45000) (11.72167,3.45000) (11.73472,3.45000) (11.74778,3.45000) (11.76083,3.45000) (11.77389,3.45000) (11.78694,3.45000) (11.80000,3.45000) (11.81306,3.45000) (11.82611,3.45000) (11.83917,3.45000) (11.85222,3.45000) (11.86528,3.45000) (11.87833,3.45000) (11.89139,3.45000) (11.90444,3.45000) (11.91750,3.45000) (11.93056,3.45000) (11.94361,3.45000) (11.95667,3.45000) (11.96972,3.45000) (11.98278,3.45000) (11.99583,3.45000) (12.00889,3.45000) (12.02194,3.45000) (12.03500,3.45000) (12.04806,3.45000) (12.06111,3.45000) (12.07417,3.45000) (12.08722,3.45000) (12.10028,3.45000) (12.11333,3.45000) (12.12639,3.45000) (12.13944,3.45000) (12.15250,3.45000) (12.16556,3.45000) (12.17861,3.45000) (12.19167,3.45000) (12.20472,3.45000) (12.21778,3.45000) (12.23083,3.45000) (12.24389,3.45000) (12.25694,3.45000) (12.27000,3.45000) (12.28306,3.45000) (12.29611,3.45000) (12.30917,3.45000) (12.32222,3.45000) (12.33528,3.45000) (12.34833,3.45000) (12.36139,3.45000) (12.37444,3.45000) (12.38750,3.45000) (12.40056,3.45000) (12.41361,3.45000) (12.42667,3.45000) (12.43972,3.45000) (12.45278,3.45000) (12.46583,3.45000) (12.47889,3.45000) (12.49194,3.45000) (12.50500,3.45000) (12.51806,3.45000) (12.53111,3.45000) (12.54417,3.45000) (12.55722,3.45000) (12.57028,3.45000) (12.58333,3.45000) (12.59639,3.45000) (12.60944,3.45000) (12.62250,3.45000) (12.63556,3.45000) (12.64861,3.45000) (12.66167,3.45000) (12.67472,3.45000) (12.68778,3.45000) (12.70083,3.45000) (12.71389,3.45000) (12.72694,3.45000) (12.74000,3.45000) (12.75306,3.45000) (12.76611,3.45000) (12.77917,3.45000) (12.79222,3.45000) (12.80528,3.45000) (12.81833,3.45000) (12.83139,3.45000) (12.84444,3.45000) (12.85750,3.45000) (12.87056,3.45000) (12.88361,3.45000) (12.89667,3.45000) (12.90972,3.45000) (12.92278,3.45000) (12.93583,3.45000) (12.94889,3.45000) (12.96194,3.45000) (12.97500,3.45000) (12.98806,3.45000) (13.00111,3.45000) (13.01417,3.45000) (13.02722,3.45000) (13.04028,3.45000) (13.05333,3.45000) (13.06639,3.45000) (13.07944,3.45000) (13.09250,3.45000) (13.10556,3.45000) (13.11861,3.45000) (13.13167,3.45000) (13.14472,3.45000) (13.15778,3.45000) (13.17083,3.45000) (13.18389,3.45000) (13.19694,3.45000) (13.21000,3.45000) (13.22306,3.45000) (13.23611,3.45000) (13.24917,3.45000) (13.26222,3.45000) (13.27528,3.45000) (13.28833,3.45000) (13.30139,3.45000) (13.31444,3.45000) (13.32750,3.45000) (13.34056,3.45000) (13.35361,3.45000) (13.36667,3.45000) (13.37972,3.45000) (13.39278,3.45000) (13.40583,3.45000) (13.41889,3.45000) (13.43194,3.45000) (13.44500,3.45000) (13.45806,3.45000) (13.47111,3.45000) (13.48417,3.45000) (13.49722,3.45000) (13.51028,3.45000) (13.52333,3.45000) (13.53639,3.45000) (13.54944,3.45000) (13.56250,3.45000) (13.57556,3.45000) (13.58861,3.45000) (13.60167,3.45000) (13.61472,3.45000) (13.62778,3.45000) (13.64083,3.45000) (13.65389,3.45000) (13.66694,3.45000) (13.68000,3.45000) (13.69306,3.45000) (13.70611,3.45000) (13.71917,3.45000) (13.73222,3.45000) (13.74528,3.45000) (13.75833,3.45000) (13.77139,3.45000) (13.78444,3.45000) (13.79750,3.45000) (13.81056,3.45000) (13.82361,3.45000) (13.83667,3.45000) (13.84972,3.45000) (13.86278,3.45000) (13.87583,3.45000) (13.88889,3.45000) (13.90194,3.45000) (13.91500,3.45000) (13.92806,3.45000) (13.94111,3.45000) (13.95417,3.45000) (13.96722,3.45000) (13.98028,3.45000) (13.99333,3.45000) (14.00639,3.45000) (14.01944,3.45000) (14.03250,3.45000) (14.04556,3.45000) (14.05861,3.45000) (14.07167,3.45000) (14.08472,3.45000) (14.09778,3.45000) (14.11083,3.45000) (14.12389,3.45000) (14.13694,3.45000) (14.15000,3.45000) (14.16306,3.45000) (14.17611,3.45000) (14.18917,3.45000) (14.20222,3.45000) (14.21528,3.45000) (14.22833,3.45000) (14.24139,3.45000) (14.25444,3.45000) (14.26750,3.45000) (14.28056,3.45000) (14.29361,3.45000) (14.30667,3.45000) (14.31972,3.45000) (14.33278,3.45000) (14.34583,3.45000) (14.35889,3.45000) (14.37194,3.45000) (14.38500,3.45000) (14.39806,3.45000) (14.41111,3.45000) (14.42417,3.45000) (14.43722,3.45000) (14.45028,3.45000) (14.46333,3.45000) (14.47639,3.45000) (14.48944,3.45000) (14.50250,3.45000) (14.51556,3.45000) (14.52861,3.45000) (14.54167,3.45000) (14.55472,3.45000) (14.56778,3.45000) (14.58083,3.45000) (14.59389,3.45000) (14.60694,3.45000) (14.62000,3.45000) (14.63306,3.45000) (14.64611,3.45000) (14.65917,3.45000) (14.67222,3.45000) (14.68528,3.45000) (14.69833,3.45000) (14.71139,3.45000) (14.72444,3.45000) (14.73750,3.45000) (14.75056,3.45000) (14.76361,3.45000) (14.77667,3.45000) (14.78972,3.45000) (14.80278,3.45000) (14.81583,3.45000) (14.82889,3.45000) (14.84194,3.45000) (14.85500,3.45000) (14.86806,3.45000) (14.88111,3.45000) (14.89417,3.45000) (14.90722,3.45000) (14.92028,3.45000) (14.93333,3.45000) (14.94639,3.45000) (14.95944,3.45000) (14.97250,3.45000) (14.98556,3.45000) (14.99861,3.45000) (15.01167,3.45000) (15.02472,3.45000) (15.03778,3.45000) (15.05083,3.45000) (15.06389,3.45000) (15.07694,3.45000) (15.09000,3.45000) (15.10306,3.45000) (15.11611,3.45000) (15.12917,3.45000) (15.14222,3.45000) (15.15528,3.45000) (15.16833,3.45000) (15.18139,3.45000) (15.19444,3.45000) (15.20750,3.45000) (15.22056,3.45000) (15.23361,3.45000) (15.24667,3.45000) (15.25972,3.45000) (15.27278,3.45000) (15.28583,3.45000) (15.29889,3.45000) (15.31194,3.45000) (15.32500,3.45000) (15.33806,3.45000) (15.35111,3.45000) (15.36417,3.45000) (15.37722,3.45000) (15.39028,3.45000) (15.40333,3.45000) (15.41639,3.45000) (15.42944,3.45000) (15.44250,3.45000) (15.45556,3.45000) (15.46861,3.45000) (15.48167,3.45000) (15.49472,3.45000) (15.50778,3.45000) (15.52083,3.45000) (15.53389,3.45000) (15.54694,3.45000) (15.56000,3.45000) (15.57306,3.45000) (15.58611,3.45000) (15.59917,3.45000) (15.61222,3.45000) (15.62528,3.45000) (15.63833,3.45000) (15.65139,3.45000) (15.66444,3.45000) (15.67750,3.45000) (15.69056,3.45000) (15.70361,3.45000) (15.71667,3.45000) (15.72972,3.45000) (15.74278,3.45000) (15.75583,3.45000) (15.76889,3.45000) (15.78194,3.45000) (15.79500,3.45000) (15.80806,3.45000) (15.82111,3.45000) (15.83417,3.45000) (15.84722,3.45000) (15.86028,3.45000) (15.87333,3.45000) (15.88639,3.45000) (15.89944,3.45000) (15.91250,3.45000) (15.92556,3.45000) (15.93861,3.45000) (15.95167,3.45000) (15.96472,3.45000) (15.97778,3.45000) (15.99083,3.45000) (16.00389,3.45000) (16.01694,3.45000) (16.03000,3.45000) (16.04306,3.45000) (16.05611,3.45000) (16.06917,3.45000) (16.08222,3.45000) (16.09528,3.45000) (16.10833,3.45000) (16.12139,3.45000) (16.13444,3.45000) (16.14750,3.45000) (16.16056,3.45000) (16.17361,3.45000) (16.18667,3.45000) (16.19972,3.45000) (16.21278,3.45000) (16.22583,3.45000) (16.23889,3.45000) (16.25194,3.45000) (16.26500,3.45000) (16.27806,3.45000) (16.29111,3.45000) (16.30417,3.45000) (16.31722,3.45000) (16.33028,3.45000) (16.34333,3.45000) (16.35639,3.45000) (16.36944,3.45000) (16.38250,3.45000) (16.39556,3.45000) (16.40861,3.45000) (16.42167,3.45000) (16.43472,3.45000) (16.44778,3.45000) (16.46083,3.45000) (16.47389,3.45000) (16.48694,3.45000) (16.50000,3.45000)};
\draw[->,coursemuted] plot coordinates {(7.10000,4.65000) (16.60000,4.65000)};
\draw[->,coursemuted] plot coordinates {(7.10000,4.65000) (7.10000,6.50000)};
\draw[coursemuted!25,thin] plot coordinates {(7.10000,4.65000) (16.50000,4.65000)};
\node[left,inner sep=3pt] at (7.1,4.65000) {$0$};
\draw[coursemuted!25,thin] plot coordinates {(7.10000,5.45000) (16.50000,5.45000)};
\node[left,inner sep=3pt] at (7.1,5.45000) {$6{,}5$};
\draw[coursemuted!25,thin] plot coordinates {(7.10000,6.25000) (16.50000,6.25000)};
\node[left,inner sep=3pt] at (7.1,6.25000) {$13$};
\node[rotate=90] at (5.9,5.55) {$\Delta V_j$, м/с};
\draw[coursemuted!16,thin] plot coordinates {(7.10000,0.80000) (7.10000,6.25000)};
\node[below,inner sep=3pt] at (7.10000,0.85) {$0$};
\draw[coursemuted!16,thin] plot coordinates {(9.45000,0.80000) (9.45000,6.25000)};
\node[below,inner sep=3pt] at (9.45000,0.85) {$\pi/2$};
\draw[coursemuted!16,thin] plot coordinates {(11.80000,0.80000) (11.80000,6.25000)};
\node[below,inner sep=3pt] at (11.80000,0.85) {$\pi$};
\draw[coursemuted!16,thin] plot coordinates {(14.15000,0.80000) (14.15000,6.25000)};
\node[below,inner sep=3pt] at (14.15000,0.85) {$3\pi/2$};
\draw[coursemuted!16,thin] plot coordinates {(16.50000,0.80000) (16.50000,6.25000)};
\node[below,inner sep=3pt] at (16.50000,0.85) {$2\pi$};
\node at (11.8,.15) {Угол $\varphi$, рад};
\draw[courseink!70,densely dashed,thin] plot coordinates {(7.88333,0.85000) (7.88333,3.45000)};
\draw[courseink,line width=1pt] plot coordinates {(7.88333,4.65000) (7.88333,5.88411)};
\fill[courseink] (7.88333,5.88411) circle[radius=1.8pt];
\fill[courseink] (7.88333,3.45000) circle[radius=1.8pt];
\node[above,inner sep=2pt] at (7.88333,5.88411) {1};
\draw[courseink!70,densely dashed,thin] plot coordinates {(12.17019,0.85000) (12.17019,3.45000)};
\draw[courseink,line width=1pt] plot coordinates {(12.17019,4.65000) (12.17019,5.27338)};
\fill[courseink] (12.17019,5.27338) circle[radius=1.8pt];
\fill[courseink] (12.17019,3.45000) circle[radius=1.8pt];
\node[above,inner sep=2pt] at (12.17019,5.27338) {2};
\end{tikzpicture}

\par
\figcaption{$\Delta a=0{,}004$, $\Delta e_x=0{,}001$, $\Delta e_y=0{,}001$, $\Delta i=0$.\\
Для примера выбран угол первого импульса $\varphi_1=30^\circ$.}
\end{frame}

% eq:planar-free-first; частный случай круговых орбит: строки 361--364.
% После кадра «Непересекающиеся орбиты: годограф».
% manual/chapter-3-algorithms.tex, sec:planar-disjoint.
\begin{frame}{Непересекающиеся орбиты}
\relax
\derivationmathspacing
Для двух трансверсальных импульсов граничные условия равны
\[
\begin{aligned}
 2(\Delta\widehat V_{t1}+\Delta\widehat V_{t2})&=\Delta a,\\
 2\Delta\widehat V_{t1}\cos\varphi_1
  +2\Delta\widehat V_{t2}\cos\varphi_2&=\Delta e_x,\\
 2\Delta\widehat V_{t1}\sin\varphi_1
  +2\Delta\widehat V_{t2}\sin\varphi_2&=\Delta e_y.
\end{aligned}
\]
\justifying
Обе компоненты имеют знак $\Delta a$, поэтому первое равенство даёт
\par
\[
 J=|\Delta\widehat V_{t1}|+|\Delta\widehat V_{t2}|
   =\frac{|\Delta a|}{2}.
\]
\justifying
Угол $\varphi_1$ выбираем произвольно. После этого неизвестны
две трансверсальные компоненты и угол $\varphi_2$.
\end{frame}

\begin{frame}{Непересекающиеся орбиты}
\relax
\derivationmathspacing
Из условия на большую полуось выражаем второй импульс:
\[
 2\Delta\widehat V_{t2}=\Delta a-2\Delta\widehat V_{t1}.
\]
В двух остальных условиях переносим вклад первого импульса вправо:
\[
\begin{aligned}
 2\Delta\widehat V_{t2}\cos\varphi_2
   &=\Delta e_x-2\Delta\widehat V_{t1}\cos\varphi_1,\\
 2\Delta\widehat V_{t2}\sin\varphi_2
   &=\Delta e_y-2\Delta\widehat V_{t1}\sin\varphi_1.
\end{aligned}
\]
Возводим их в квадрат и складываем, исключая угол $\varphi_2$:
\[
\begin{aligned}
 (\Delta a-2\Delta\widehat V_{t1})^2
   &=(\Delta e_x-2\Delta\widehat V_{t1}\cos\varphi_1)^2\\
   &\quad +(\Delta e_y-2\Delta\widehat V_{t1}\sin\varphi_1)^2.
\end{aligned}
\]
\end{frame}

\begin{frame}{Непересекающиеся орбиты}
\relax
\derivationmathspacing
Раскрываем квадраты в полученном равенстве:
\[
\begin{aligned}
 (\Delta a)^2-4\Delta a\,\Delta\widehat V_{t1}
       +4(\Delta\widehat V_{t1})^2
   &=(\Delta e_x)^2+(\Delta e_y)^2\\
   &\quad-4\Delta\widehat V_{t1}
       (\Delta e_x\cos\varphi_1+\Delta e_y\sin\varphi_1)\\
   &\quad+4(\Delta\widehat V_{t1})^2
       (\cos^2\varphi_1+\sin^2\varphi_1).
\end{aligned}
\]
\justifying
Так как $\cos^2\varphi_1+\sin^2\varphi_1=1$, с обеих сторон
стоит одно и то же слагаемое $4(\Delta\widehat V_{t1})^2$.
Сокращаем его и учитываем $(\Delta e_x)^2+(\Delta e_y)^2=\Delta e^2$:
\par
\[
\begin{aligned}
 (\Delta a)^2-4\Delta a\,\Delta\widehat V_{t1}
   &=\Delta e^2-4\Delta\widehat V_{t1}
        (\Delta e_x\cos\varphi_1+\Delta e_y\sin\varphi_1).
\end{aligned}
\]
Получено линейное уравнение относительно $\Delta\widehat V_{t1}$.
\end{frame}

\begin{frame}{Непересекающиеся орбиты}
\relax
\derivationmathspacing
Переносим слагаемые с $\Delta\widehat V_{t1}$ в одну сторону:
\[
\begin{aligned}
 4\Delta\widehat V_{t1}
  (\Delta e_x\cos\varphi_1+\Delta e_y\sin\varphi_1-\Delta a)
   &=\Delta e^2-(\Delta a)^2.
\end{aligned}
\]
\begin{block}{Трансверсальные компоненты импульсов}
\[
\begin{aligned}
 \Delta\widehat V_{t1}
   &=\frac{\Delta e^2-(\Delta a)^2}
       {4(\Delta e_x\cos\varphi_1+\Delta e_y\sin\varphi_1-\Delta a)},\\[3pt]
 \Delta\widehat V_{t2}
   &=\frac{\Delta a}{2}-\Delta\widehat V_{t1}.
\end{aligned}
\]
\end{block}
\justifying
Условие на большую полуось выполнено. Определим второй угол.
\end{frame}

\begin{frame}{Непересекающиеся орбиты}
\relax
\derivationmathspacing
Делим преобразованные условия на эксцентриситет на
$2\Delta\widehat V_{t2}$:
\[
\boxed{\begin{aligned}
 \cos\varphi_2
   &=\frac{\Delta e_x-2\Delta\widehat V_{t1}\cos\varphi_1}
       {2\Delta\widehat V_{t2}},\\[3pt]
 \sin\varphi_2
   &=\frac{\Delta e_y-2\Delta\widehat V_{t1}\sin\varphi_1}
       {2\Delta\widehat V_{t2}}.
\end{aligned}}
\]
\justifying
Сумма квадратов обеспечивает $\cos^2\varphi_2+\sin^2\varphi_2=1$.
Четверть угла определяем по обеим компонентам:
\par
\[
 \varphi_2=\operatorname{atan2}(\sin\varphi_2,\cos\varphi_2).
\]
Радиальные компоненты обоих импульсов равны нулю.
\end{frame}

\begin{frame}{Непересекающиеся орбиты}
\relax
\derivationmathspacing
При $|\Delta a|>\Delta e$ для любого выбранного $\varphi_1$ имеем
\[
 \bigl|\Delta e_x\cos\varphi_1+\Delta e_y\sin\varphi_1\bigr|
 \leq\Delta e<|\Delta a|.
\]
\justifying
Поэтому знаменатель формулы для $\Delta\widehat V_{t1}$
не обращается в нуль; обе найденные компоненты имеют знак $\Delta a$.
В частности, $\Delta\widehat V_{t2}\ne0$, и второй угол определён.
\par\medskip
\[
\begin{gathered}
 \boldsymbol s=(0,\operatorname{sgn}\Delta a,0),\qquad
 \|\boldsymbol s\|=1,\\[2pt]
 J=|\Delta\widehat V_{t1}|+|\Delta\widehat V_{t2}|
   =\left|\Delta\widehat V_{t1}+\Delta\widehat V_{t2}\right|
   =\frac{|\Delta a|}{2}.
\end{gathered}
\]
\justifying
Импульсы сонаправлены базис-вектору, все граничные условия выполнены.
Разным $\varphi_1$ соответствуют разные переходные орбиты
при одинаковых оптимальных затратах.
\end{frame}

\begin{frame}{Обобщённый переход Гомана}
\relax
\derivationmathspacing
Перицентры компланарных околокруговых орбит сонаправлены.\par
Ось $X$ --- к перицентрам: $\Delta e_y=0$, $\Delta e_x=e_f-e_i$,
$|\Delta a|>|\Delta e_x|$.
\begin{columns}[T,onlytextwidth]
\begin{column}{.59\textwidth}
\smallskip
В околокруговом приближении при $\varphi_1=0$:
\[
\begin{aligned}
 \Delta\widehat V_{t1}
   &=\frac{(\Delta e_x)^2-(\Delta a)^2}{4(\Delta e_x-\Delta a)}
     =\frac{\Delta a+\Delta e_x}{4},\\
 \Delta\widehat V_{t2}
   &=\frac{\Delta a}{2}-\Delta\widehat V_{t1}
     =\frac{\Delta a-\Delta e_x}{4}.
\end{aligned}
\]
Подстановка в условия для второго угла:
\[
\begin{gathered}
 \cos\varphi_2=-1,\quad \sin\varphi_2=0,\\
 \varphi_2=\pi,\qquad J=\frac{|\Delta a|}{2}.
\end{gathered}
\]
\smallskip
При $e_i=e_f$ оба импульса равны $\Delta a/4$.\par
\end{column}
\begin{column}{.39\textwidth}
\centering
% Схема с усиленной эллиптичностью: a_i=12, e_i=.65; a_f=20, e_f=.60.
% Это геометрическая иллюстрация, не численный пример околокругового приближения.
% Общий фокус O=(0,0); центры (-a_i e_i,0), (-a_f e_f,0).
% r_pi=4.2, r_af=32; переходная дуга имеет a_tr=18.1,
% c_tr=13.9, b_tr=sqrt(r_pi*r_af), касания в противоположных апсидах.
\begin{tikzpicture}[course diagram,x=1mm,y=1mm,scale=.84]
  \path[use as bounding box] (-37,-32) rectangle (16,19);
  \pgfmathsetmacro{\htAi}{12}
  \pgfmathsetmacro{\htEi}{.65}
  \pgfmathsetmacro{\htAf}{20}
  \pgfmathsetmacro{\htEf}{.60}
  \pgfmathsetmacro{\htBi}{\htAi*sqrt(1-\htEi*\htEi)}
  \pgfmathsetmacro{\htBf}{\htAf*sqrt(1-\htEf*\htEf)}
  \pgfmathsetmacro{\htRp}{\htAi*(1-\htEi)}
  \pgfmathsetmacro{\htRa}{\htAf*(1+\htEf)}
  \pgfmathsetmacro{\htAt}{(\htRp+\htRa)/2}
  \pgfmathsetmacro{\htCt}{(\htRa-\htRp)/2}
  \pgfmathsetmacro{\htBt}{sqrt(\htRp*\htRa)}
  \coordinate (htO) at (0,0);
  \coordinate (htP) at (\htRp,0);
  \coordinate (htA) at (-\htRa,0);

  \draw[aux line,dashed] (-34,0)--(10,0);
  \draw[draw=courseink,line width=.75pt]
    ({-\htAi*\htEi},0) ellipse[x radius=\htAi,y radius=\htBi];
  \draw[draw=courseink,line width=.75pt,densely dashed]
    ({-\htAf*\htEf},0) ellipse[x radius=\htAf,y radius=\htBf];
  \draw[orbit line]
    plot[domain=0:180,samples=121,variable=\u]
      ({-\htCt+\htAt*cos(\u)},{\htBt*sin(\u)});
  \draw[angle line,line width=1pt]
    plot[domain=75:96,samples=25,variable=\u]
      ({-\htCt+\htAt*cos(\u)},{\htBt*sin(\u)});

  % Длинные подписи вынесены в легенду; в поле схемы только O и номера.
  \node[circle,fill=courseink,inner sep=1.5pt] at (htO) {};
  \node[anchor=north east,inner sep=0pt] at (-.8,-1.2) {$O$};
  \node[orbit point] at (htP) {};
  \node[orbit point] at (htA) {};
  \draw[vector line,line width=.85pt] (htP)--++(0,7);
  \draw[vector line,line width=.85pt] (htA)--++(0,-7);
  \node[inner sep=0pt] at (11,7.3) {$1$};
  \node[inner sep=0pt] at (-34.5,-6.5) {$2$};

  \draw[draw=courseink,line width=.75pt] (-29,-20.5)--(-22,-20.5);
  \node[anchor=west,inner sep=0pt] at (-19.5,-20.5) {Начальная $i$};
  \draw[draw=courseink,line width=.75pt,densely dashed]
    (-29,-25)--(-22,-25);
  \node[anchor=west,inner sep=0pt] at (-19.5,-25) {Конечная $f$};
  \draw[orbit line] (-29,-29.5)--(-22,-29.5);
  \node[anchor=west,inner sep=0pt] at (-19.5,-29.5) {Переходная};
\end{tikzpicture}

\par
\figcaption{$O$ --- общий фокус орбит.\\Эллиптичность усилена.}
\end{column}
\end{columns}
\end{frame}


% eq:planar-central-primer; вывод углов: строки 392--419.
\begin{frame}{Пересекающиеся орбиты: геометрия}
\relax
\derivationmathspacing
Пусть $|\Delta a|<\Delta e$.
Годограф --- центральный эллипс с двумя касаниями:
$\lambda_1=0$, $\sqrt{\lambda_2^2+\lambda_3^2}=1/2$.
\[
 s_r=\tfrac12\sin(\varphi-\alpha),\qquad
 s_t=\cos(\varphi-\alpha),\qquad s_z=0.
\]
{\justifying
В одной точке касания $s_t=1$, в другой $s_t=-1$.
Первый импульс разгонный, второй --- тормозной;
места импульсов отличаются на $\pi$:
\par}
\[
 \varphi_1=\alpha,\qquad \varphi_2=\alpha+\pi.
\]
\medskip
{\justifying
Изменения большой полуоси частично компенсируются,
а изменения вектора эксцентриситета складываются в одном направлении.
Угол $\alpha$ определим из граничных условий.
\par}
\end{frame}

\begin{frame}{Пересекающиеся орбиты: годограф}
\relax
\centering
% Рисунок из методички; координаты кривых, осей и отметок сохранены.
\begin{tikzpicture}[course diagram,scale=.72]
\node at (2.65,6.85) {Годограф};
\node at (11.6,6.85) {Модули импульсов};
\node at (11.6,4.1) {Модуль базис-вектора};
\draw[coursemuted!60,dashed] (2.65,3.6) circle[radius=1.7cm];
\draw[->,coursemuted] (.6,3.6)--(4.85,3.6) node[right] {$s_t$};
\draw[->,coursemuted] (2.65,1.55)--(2.65,5.8) node[above] {$s_r$};
\draw[courseaccent,line width=1.15pt] plot coordinates {(4.17053,3.21987) (4.18356,3.23319) (4.19613,3.24663) (4.20823,3.26018) (4.21986,3.27383) (4.23100,3.28758) (4.24167,3.30142) (4.25185,3.31535) (4.26154,3.32938) (4.27074,3.34348) (4.27944,3.35766) (4.28765,3.37192) (4.29537,3.38624) (4.30258,3.40063) (4.30928,3.41508) (4.31549,3.42959) (4.32118,3.44415) (4.32637,3.45876) (4.33104,3.47341) (4.33520,3.48810) (4.33885,3.50282) (4.34199,3.51757) (4.34461,3.53235) (4.34671,3.54715) (4.34830,3.56196) (4.34937,3.57678) (4.34992,3.59162) (4.34995,3.60645) (4.34947,3.62129) (4.34847,3.63611) (4.34695,3.65093) (4.34491,3.66573) (4.34236,3.68051) (4.33929,3.69526) (4.33571,3.70999) (4.33161,3.72468) (4.32700,3.73934) (4.32188,3.75395) (4.31626,3.76852) (4.31012,3.78303) (4.30348,3.79749) (4.29633,3.81189) (4.28869,3.82622) (4.28054,3.84049) (4.27190,3.85468) (4.26276,3.86879) (4.25313,3.88283) (4.24302,3.89677) (4.23242,3.91063) (4.22133,3.92439) (4.20977,3.93805) (4.19773,3.95161) (4.18522,3.96506) (4.17225,3.97840) (4.15881,3.99163) (4.14491,4.00474) (4.13055,4.01772) (4.11575,4.03058) (4.10050,4.04330) (4.08480,4.05589) (4.06867,4.06834) (4.05211,4.08065) (4.03512,4.09281) (4.01770,4.10482) (3.99987,4.11668) (3.98163,4.12838) (3.96299,4.13992) (3.94394,4.15130) (3.92450,4.16250) (3.90467,4.17354) (3.88446,4.18440) (3.86388,4.19508) (3.84292,4.20559) (3.82160,4.21590) (3.79993,4.22603) (3.77790,4.23597) (3.75553,4.24572) (3.73282,4.25527) (3.70978,4.26462) (3.68643,4.27376) (3.66275,4.28270) (3.63877,4.29144) (3.61448,4.29996) (3.58990,4.30827) (3.56504,4.31636) (3.53989,4.32424) (3.51448,4.33189) (3.48880,4.33933) (3.46287,4.34653) (3.43668,4.35351) (3.41026,4.36026) (3.38361,4.36678) (3.35673,4.37307) (3.32964,4.37912) (3.30234,4.38493) (3.27485,4.39050) (3.24716,4.39583) (3.21929,4.40092) (3.19125,4.40577) (3.16304,4.41037) (3.13468,4.41472) (3.10616,4.41883) (3.07751,4.42268) (3.04873,4.42629) (3.01983,4.42964) (2.99082,4.43274) (2.96170,4.43559) (2.93248,4.43818) (2.90318,4.44052) (2.87381,4.44260) (2.84436,4.44443) (2.81486,4.44599) (2.78530,4.44730) (2.75571,4.44836) (2.72608,4.44915) (2.69643,4.44968) (2.66677,4.44996) (2.63709,4.44998) (2.60743,4.44973) (2.57778,4.44923) (2.54814,4.44847) (2.51854,4.44745) (2.48898,4.44618) (2.45947,4.44464) (2.43002,4.44285) (2.40063,4.44081) (2.37132,4.43850) (2.34210,4.43594) (2.31297,4.43313) (2.28394,4.43006) (2.25502,4.42674) (2.22622,4.42317) (2.19755,4.41934) (2.16902,4.41527) (2.14064,4.41095) (2.11241,4.40638) (2.08435,4.40157) (2.05646,4.39651) (2.02874,4.39121) (2.00122,4.38567) (1.97390,4.37989) (1.94678,4.37387) (1.91987,4.36761) (1.89319,4.36112) (1.86674,4.35440) (1.84053,4.34745) (1.81456,4.34028) (1.78885,4.33287) (1.76340,4.32525) (1.73822,4.31740) (1.71332,4.30934) (1.68870,4.30105) (1.66438,4.29256) (1.64035,4.28385) (1.61664,4.27494) (1.59324,4.26582) (1.57016,4.25649) (1.54741,4.24697) (1.52499,4.23725) (1.50292,4.22734) (1.48120,4.21723) (1.45983,4.20694) (1.43883,4.19646) (1.41819,4.18580) (1.39793,4.17496) (1.37806,4.16395) (1.35856,4.15276) (1.33947,4.14141) (1.32077,4.12989) (1.30248,4.11821) (1.28459,4.10637) (1.26713,4.09438) (1.25008,4.08224) (1.23346,4.06995) (1.21727,4.05752) (1.20152,4.04495) (1.18621,4.03224) (1.17135,4.01940) (1.15693,4.00643) (1.14297,3.99334) (1.12947,3.98013) (1.11644,3.96681) (1.10387,3.95337) (1.09177,3.93982) (1.08014,3.92617) (1.06900,3.91242) (1.05833,3.89858) (1.04815,3.88465) (1.03846,3.87062) (1.02926,3.85652) (1.02056,3.84234) (1.01235,3.82808) (1.00463,3.81376) (0.99742,3.79937) (0.99072,3.78492) (0.98451,3.77041) (0.97882,3.75585) (0.97363,3.74124) (0.96896,3.72659) (0.96480,3.71190) (0.96115,3.69718) (0.95801,3.68243) (0.95539,3.66765) (0.95329,3.65285) (0.95170,3.63804) (0.95063,3.62322) (0.95008,3.60838) (0.95005,3.59355) (0.95053,3.57871) (0.95153,3.56389) (0.95305,3.54907) (0.95509,3.53427) (0.95764,3.51949) (0.96071,3.50474) (0.96429,3.49001) (0.96839,3.47532) (0.97300,3.46066) (0.97812,3.44605) (0.98374,3.43148) (0.98988,3.41697) (0.99652,3.40251) (1.00367,3.38811) (1.01131,3.37378) (1.01946,3.35951) (1.02810,3.34532) (1.03724,3.33121) (1.04687,3.31717) (1.05698,3.30323) (1.06758,3.28937) (1.07867,3.27561) (1.09023,3.26195) (1.10227,3.24839) (1.11478,3.23494) (1.12775,3.22160) (1.14119,3.20837) (1.15509,3.19526) (1.16945,3.18228) (1.18425,3.16942) (1.19950,3.15670) (1.21520,3.14411) (1.23133,3.13166) (1.24789,3.11935) (1.26488,3.10719) (1.28230,3.09518) (1.30013,3.08332) (1.31837,3.07162) (1.33701,3.06008) (1.35606,3.04870) (1.37550,3.03750) (1.39533,3.02646) (1.41554,3.01560) (1.43612,3.00492) (1.45708,2.99441) (1.47840,2.98410) (1.50007,2.97397) (1.52210,2.96403) (1.54447,2.95428) (1.56718,2.94473) (1.59022,2.93538) (1.61357,2.92624) (1.63725,2.91730) (1.66123,2.90856) (1.68552,2.90004) (1.71010,2.89173) (1.73496,2.88364) (1.76011,2.87576) (1.78552,2.86811) (1.81120,2.86067) (1.83713,2.85347) (1.86332,2.84649) (1.88974,2.83974) (1.91639,2.83322) (1.94327,2.82693) (1.97036,2.82088) (1.99766,2.81507) (2.02515,2.80950) (2.05284,2.80417) (2.08071,2.79908) (2.10875,2.79423) (2.13696,2.78963) (2.16532,2.78528) (2.19384,2.78117) (2.22249,2.77732) (2.25127,2.77371) (2.28017,2.77036) (2.30918,2.76726) (2.33830,2.76441) (2.36752,2.76182) (2.39682,2.75948) (2.42619,2.75740) (2.45564,2.75557) (2.48514,2.75401) (2.51470,2.75270) (2.54429,2.75164) (2.57392,2.75085) (2.60357,2.75032) (2.63323,2.75004) (2.66291,2.75002) (2.69257,2.75027) (2.72222,2.75077) (2.75186,2.75153) (2.78146,2.75255) (2.81102,2.75382) (2.84053,2.75536) (2.86998,2.75715) (2.89937,2.75919) (2.92868,2.76150) (2.95790,2.76406) (2.98703,2.76687) (3.01606,2.76994) (3.04498,2.77326) (3.07378,2.77683) (3.10245,2.78066) (3.13098,2.78473) (3.15936,2.78905) (3.18759,2.79362) (3.21565,2.79843) (3.24354,2.80349) (3.27126,2.80879) (3.29878,2.81433) (3.32610,2.82011) (3.35322,2.82613) (3.38013,2.83239) (3.40681,2.83888) (3.43326,2.84560) (3.45947,2.85255) (3.48544,2.85972) (3.51115,2.86713) (3.53660,2.87475) (3.56178,2.88260) (3.58668,2.89066) (3.61130,2.89895) (3.63562,2.90744) (3.65965,2.91615) (3.68336,2.92506) (3.70676,2.93418) (3.72984,2.94351) (3.75259,2.95303) (3.77501,2.96275) (3.79708,2.97266) (3.81880,2.98277) (3.84017,2.99306) (3.86117,3.00354) (3.88181,3.01420) (3.90207,3.02504) (3.92194,3.03605) (3.94144,3.04724) (3.96053,3.05859) (3.97923,3.07011) (3.99752,3.08179) (4.01541,3.09363) (4.03287,3.10562) (4.04992,3.11776) (4.06654,3.13005) (4.08273,3.14248) (4.09848,3.15505) (4.11379,3.16776) (4.12865,3.18060) (4.14307,3.19357) (4.15703,3.20666) (4.17053,3.21987)};
\filldraw[fill=courseink,draw=coursepaper,line width=.4pt] (4.35000,3.60000) circle[radius=2pt];
\node[above,inner sep=3pt] at (4.35000,3.60000) {1};
\filldraw[fill=courseink,draw=coursepaper,line width=.4pt] (0.95000,3.60000) circle[radius=2pt];
\node[above,inner sep=3pt] at (0.95000,3.60000) {2};
\node[below left,inner sep=2pt] at (2.65,3.6) {$0$};
\draw[coursemuted!25,thin] plot coordinates {(7.10000,0.85000) (16.50000,0.85000)};
\node[left,inner sep=3pt] at (7.1,0.85) {$0$};
\draw[coursemuted!25,thin] plot coordinates {(7.10000,2.15000) (16.50000,2.15000)};
\node[left,inner sep=3pt] at (7.1,2.15) {$0{,}5$};
\draw[courseink!65,dashed] plot coordinates {(7.10000,3.45000) (16.50000,3.45000)};
\node[left,inner sep=3pt] at (7.1,3.45) {$1$};
\draw[->,coursemuted] plot coordinates {(7.10000,0.85000) (16.60000,0.85000)};
\draw[->,coursemuted] plot coordinates {(7.10000,0.85000) (7.10000,3.84000)};
\node[rotate=90] at (5.9,2.2) {$\|\boldsymbol s(\varphi)\|$};
\draw[courseaccent,line width=1.15pt] plot coordinates {(7.10000,3.24708) (7.11306,3.25440) (7.12611,3.26161) (7.13917,3.26869) (7.15222,3.27565) (7.16528,3.28248) (7.17833,3.28919) (7.19139,3.29577) (7.20444,3.30223) (7.21750,3.30856) (7.23056,3.31476) (7.24361,3.32083) (7.25667,3.32677) (7.26972,3.33258) (7.28278,3.33826) (7.29583,3.34380) (7.30889,3.34921) (7.32194,3.35449) (7.33500,3.35963) (7.34806,3.36463) (7.36111,3.36950) (7.37417,3.37423) (7.38722,3.37882) (7.40028,3.38328) (7.41333,3.38759) (7.42639,3.39176) (7.43944,3.39580) (7.45250,3.39969) (7.46556,3.40344) (7.47861,3.40705) (7.49167,3.41051) (7.50472,3.41383) (7.51778,3.41701) (7.53083,3.42005) (7.54389,3.42294) (7.55694,3.42568) (7.57000,3.42828) (7.58306,3.43074) (7.59611,3.43305) (7.60917,3.43521) (7.62222,3.43722) (7.63528,3.43909) (7.64833,3.44081) (7.66139,3.44239) (7.67444,3.44382) (7.68750,3.44510) (7.70056,3.44623) (7.71361,3.44721) (7.72667,3.44805) (7.73972,3.44873) (7.75278,3.44927) (7.76583,3.44966) (7.77889,3.44991) (7.79194,3.45000) (7.80500,3.44994) (7.81806,3.44974) (7.83111,3.44939) (7.84417,3.44889) (7.85722,3.44824) (7.87028,3.44744) (7.88333,3.44650) (7.89639,3.44540) (7.90944,3.44416) (7.92250,3.44277) (7.93556,3.44124) (7.94861,3.43955) (7.96167,3.43772) (7.97472,3.43575) (7.98778,3.43362) (8.00083,3.43135) (8.01389,3.42894) (8.02694,3.42637) (8.04000,3.42367) (8.05306,3.42081) (8.06611,3.41782) (8.07917,3.41468) (8.09222,3.41139) (8.10528,3.40796) (8.11833,3.40439) (8.13139,3.40068) (8.14444,3.39682) (8.15750,3.39283) (8.17056,3.38869) (8.18361,3.38441) (8.19667,3.38000) (8.20972,3.37544) (8.22278,3.37075) (8.23583,3.36591) (8.24889,3.36095) (8.26194,3.35584) (8.27500,3.35060) (8.28806,3.34522) (8.30111,3.33972) (8.31417,3.33407) (8.32722,3.32830) (8.34028,3.32239) (8.35333,3.31635) (8.36639,3.31019) (8.37944,3.30389) (8.39250,3.29747) (8.40556,3.29092) (8.41861,3.28424) (8.43167,3.27744) (8.44472,3.27051) (8.45778,3.26346) (8.47083,3.25629) (8.48389,3.24900) (8.49694,3.24159) (8.51000,3.23406) (8.52306,3.22641) (8.53611,3.21865) (8.54917,3.21077) (8.56222,3.20278) (8.57528,3.19467) (8.58833,3.18646) (8.60139,3.17813) (8.61444,3.16970) (8.62750,3.16116) (8.64056,3.15251) (8.65361,3.14376) (8.66667,3.13491) (8.67972,3.12595) (8.69278,3.11690) (8.70583,3.10775) (8.71889,3.09850) (8.73194,3.08916) (8.74500,3.07972) (8.75806,3.07020) (8.77111,3.06058) (8.78417,3.05088) (8.79722,3.04109) (8.81028,3.03121) (8.82333,3.02126) (8.83639,3.01122) (8.84944,3.00110) (8.86250,2.99091) (8.87556,2.98065) (8.88861,2.97031) (8.90167,2.95990) (8.91472,2.94942) (8.92778,2.93888) (8.94083,2.92827) (8.95389,2.91761) (8.96694,2.90688) (8.98000,2.89610) (8.99306,2.88526) (9.00611,2.87437) (9.01917,2.86343) (9.03222,2.85245) (9.04528,2.84142) (9.05833,2.83035) (9.07139,2.81924) (9.08444,2.80809) (9.09750,2.79691) (9.11056,2.78570) (9.12361,2.77447) (9.13667,2.76321) (9.14972,2.75192) (9.16278,2.74062) (9.17583,2.72931) (9.18889,2.71798) (9.20194,2.70664) (9.21500,2.69530) (9.22806,2.68395) (9.24111,2.67261) (9.25417,2.66127) (9.26722,2.64994) (9.28028,2.63863) (9.29333,2.62732) (9.30639,2.61604) (9.31944,2.60478) (9.33250,2.59355) (9.34556,2.58235) (9.35861,2.57119) (9.37167,2.56006) (9.38472,2.54898) (9.39778,2.53795) (9.41083,2.52697) (9.42389,2.51604) (9.43694,2.50518) (9.45000,2.49438) (9.46306,2.48366) (9.47611,2.47301) (9.48917,2.46244) (9.50222,2.45195) (9.51528,2.44155) (9.52833,2.43125) (9.54139,2.42105) (9.55444,2.41095) (9.56750,2.40096) (9.58056,2.39109) (9.59361,2.38133) (9.60667,2.37171) (9.61972,2.36221) (9.63278,2.35285) (9.64583,2.34363) (9.65889,2.33456) (9.67194,2.32565) (9.68500,2.31689) (9.69806,2.30829) (9.71111,2.29987) (9.72417,2.29162) (9.73722,2.28354) (9.75028,2.27566) (9.76333,2.26797) (9.77639,2.26047) (9.78944,2.25318) (9.80250,2.24610) (9.81556,2.23923) (9.82861,2.23257) (9.84167,2.22615) (9.85472,2.21995) (9.86778,2.21398) (9.88083,2.20825) (9.89389,2.20277) (9.90694,2.19754) (9.92000,2.19256) (9.93306,2.18783) (9.94611,2.18337) (9.95917,2.17917) (9.97222,2.17524) (9.98528,2.17159) (9.99833,2.16821) (10.01139,2.16511) (10.02444,2.16229) (10.03750,2.15976) (10.05056,2.15752) (10.06361,2.15556) (10.07667,2.15390) (10.08972,2.15253) (10.10278,2.15145) (10.11583,2.15067) (10.12889,2.15019) (10.14194,2.15000) (10.15500,2.15011) (10.16806,2.15052) (10.18111,2.15122) (10.19417,2.15222) (10.20722,2.15351) (10.22028,2.15510) (10.23333,2.15698) (10.24639,2.15915) (10.25944,2.16161) (10.27250,2.16435) (10.28556,2.16738) (10.29861,2.17068) (10.31167,2.17427) (10.32472,2.17812) (10.33778,2.18225) (10.35083,2.18665) (10.36389,2.19130) (10.37694,2.19622) (10.39000,2.20139) (10.40306,2.20680) (10.41611,2.21247) (10.42917,2.21837) (10.44222,2.22451) (10.45528,2.23088) (10.46833,2.23747) (10.48139,2.24429) (10.49444,2.25132) (10.50750,2.25856) (10.52056,2.26600) (10.53361,2.27364) (10.54667,2.28148) (10.55972,2.28950) (10.57278,2.29770) (10.58583,2.30608) (10.59889,2.31463) (10.61194,2.32335) (10.62500,2.33223) (10.63806,2.34126) (10.65111,2.35044) (10.66417,2.35976) (10.67722,2.36922) (10.69028,2.37882) (10.70333,2.38854) (10.71639,2.39838) (10.72944,2.40834) (10.74250,2.41841) (10.75556,2.42858) (10.76861,2.43886) (10.78167,2.44923) (10.79472,2.45970) (10.80778,2.47025) (10.82083,2.48088) (10.83389,2.49159) (10.84694,2.50237) (10.86000,2.51321) (10.87306,2.52412) (10.88611,2.53509) (10.89917,2.54611) (10.91222,2.55718) (10.92528,2.56829) (10.93833,2.57944) (10.95139,2.59063) (10.96444,2.60186) (10.97750,2.61311) (10.99056,2.62439) (11.00361,2.63568) (11.01667,2.64700) (11.02972,2.65832) (11.04278,2.66966) (11.05583,2.68100) (11.06889,2.69235) (11.08194,2.70369) (11.09500,2.71503) (11.10806,2.72636) (11.12111,2.73768) (11.13417,2.74899) (11.14722,2.76027) (11.16028,2.77154) (11.17333,2.78278) (11.18639,2.79400) (11.19944,2.80519) (11.21250,2.81634) (11.22556,2.82746) (11.23861,2.83854) (11.25167,2.84958) (11.26472,2.86058) (11.27778,2.87153) (11.29083,2.88243) (11.30389,2.89328) (11.31694,2.90408) (11.33000,2.91482) (11.34306,2.92550) (11.35611,2.93613) (11.36917,2.94669) (11.38222,2.95718) (11.39528,2.96761) (11.40833,2.97796) (11.42139,2.98825) (11.43444,2.99846) (11.44750,3.00859) (11.46056,3.01865) (11.47361,3.02863) (11.48667,3.03852) (11.49972,3.04834) (11.51278,3.05806) (11.52583,3.06770) (11.53889,3.07725) (11.55194,3.08671) (11.56500,3.09608) (11.57806,3.10535) (11.59111,3.11453) (11.60417,3.12361) (11.61722,3.13259) (11.63028,3.14147) (11.64333,3.15024) (11.65639,3.15892) (11.66944,3.16748) (11.68250,3.17595) (11.69556,3.18430) (11.70861,3.19254) (11.72167,3.20068) (11.73472,3.20870) (11.74778,3.21661) (11.76083,3.22440) (11.77389,3.23208) (11.78694,3.23964) (11.80000,3.24708) (11.81306,3.25440) (11.82611,3.26161) (11.83917,3.26869) (11.85222,3.27565) (11.86528,3.28248) (11.87833,3.28919) (11.89139,3.29577) (11.90444,3.30223) (11.91750,3.30856) (11.93056,3.31476) (11.94361,3.32083) (11.95667,3.32677) (11.96972,3.33258) (11.98278,3.33826) (11.99583,3.34380) (12.00889,3.34921) (12.02194,3.35449) (12.03500,3.35963) (12.04806,3.36463) (12.06111,3.36950) (12.07417,3.37423) (12.08722,3.37882) (12.10028,3.38328) (12.11333,3.38759) (12.12639,3.39176) (12.13944,3.39580) (12.15250,3.39969) (12.16556,3.40344) (12.17861,3.40705) (12.19167,3.41051) (12.20472,3.41383) (12.21778,3.41701) (12.23083,3.42005) (12.24389,3.42294) (12.25694,3.42568) (12.27000,3.42828) (12.28306,3.43074) (12.29611,3.43305) (12.30917,3.43521) (12.32222,3.43722) (12.33528,3.43909) (12.34833,3.44081) (12.36139,3.44239) (12.37444,3.44382) (12.38750,3.44510) (12.40056,3.44623) (12.41361,3.44721) (12.42667,3.44805) (12.43972,3.44873) (12.45278,3.44927) (12.46583,3.44966) (12.47889,3.44991) (12.49194,3.45000) (12.50500,3.44994) (12.51806,3.44974) (12.53111,3.44939) (12.54417,3.44889) (12.55722,3.44824) (12.57028,3.44744) (12.58333,3.44650) (12.59639,3.44540) (12.60944,3.44416) (12.62250,3.44277) (12.63556,3.44124) (12.64861,3.43955) (12.66167,3.43772) (12.67472,3.43575) (12.68778,3.43362) (12.70083,3.43135) (12.71389,3.42894) (12.72694,3.42637) (12.74000,3.42367) (12.75306,3.42081) (12.76611,3.41782) (12.77917,3.41468) (12.79222,3.41139) (12.80528,3.40796) (12.81833,3.40439) (12.83139,3.40068) (12.84444,3.39682) (12.85750,3.39283) (12.87056,3.38869) (12.88361,3.38441) (12.89667,3.38000) (12.90972,3.37544) (12.92278,3.37075) (12.93583,3.36591) (12.94889,3.36095) (12.96194,3.35584) (12.97500,3.35060) (12.98806,3.34522) (13.00111,3.33972) (13.01417,3.33407) (13.02722,3.32830) (13.04028,3.32239) (13.05333,3.31635) (13.06639,3.31019) (13.07944,3.30389) (13.09250,3.29747) (13.10556,3.29092) (13.11861,3.28424) (13.13167,3.27744) (13.14472,3.27051) (13.15778,3.26346) (13.17083,3.25629) (13.18389,3.24900) (13.19694,3.24159) (13.21000,3.23406) (13.22306,3.22641) (13.23611,3.21865) (13.24917,3.21077) (13.26222,3.20278) (13.27528,3.19467) (13.28833,3.18646) (13.30139,3.17813) (13.31444,3.16970) (13.32750,3.16116) (13.34056,3.15251) (13.35361,3.14376) (13.36667,3.13491) (13.37972,3.12595) (13.39278,3.11690) (13.40583,3.10775) (13.41889,3.09850) (13.43194,3.08916) (13.44500,3.07972) (13.45806,3.07020) (13.47111,3.06058) (13.48417,3.05088) (13.49722,3.04109) (13.51028,3.03121) (13.52333,3.02126) (13.53639,3.01122) (13.54944,3.00110) (13.56250,2.99091) (13.57556,2.98065) (13.58861,2.97031) (13.60167,2.95990) (13.61472,2.94942) (13.62778,2.93888) (13.64083,2.92827) (13.65389,2.91761) (13.66694,2.90688) (13.68000,2.89610) (13.69306,2.88526) (13.70611,2.87437) (13.71917,2.86343) (13.73222,2.85245) (13.74528,2.84142) (13.75833,2.83035) (13.77139,2.81924) (13.78444,2.80809) (13.79750,2.79691) (13.81056,2.78570) (13.82361,2.77447) (13.83667,2.76321) (13.84972,2.75192) (13.86278,2.74062) (13.87583,2.72931) (13.88889,2.71798) (13.90194,2.70664) (13.91500,2.69530) (13.92806,2.68395) (13.94111,2.67261) (13.95417,2.66127) (13.96722,2.64994) (13.98028,2.63863) (13.99333,2.62732) (14.00639,2.61604) (14.01944,2.60478) (14.03250,2.59355) (14.04556,2.58235) (14.05861,2.57119) (14.07167,2.56006) (14.08472,2.54898) (14.09778,2.53795) (14.11083,2.52697) (14.12389,2.51604) (14.13694,2.50518) (14.15000,2.49438) (14.16306,2.48366) (14.17611,2.47301) (14.18917,2.46244) (14.20222,2.45195) (14.21528,2.44155) (14.22833,2.43125) (14.24139,2.42105) (14.25444,2.41095) (14.26750,2.40096) (14.28056,2.39109) (14.29361,2.38133) (14.30667,2.37171) (14.31972,2.36221) (14.33278,2.35285) (14.34583,2.34363) (14.35889,2.33456) (14.37194,2.32565) (14.38500,2.31689) (14.39806,2.30829) (14.41111,2.29987) (14.42417,2.29162) (14.43722,2.28354) (14.45028,2.27566) (14.46333,2.26797) (14.47639,2.26047) (14.48944,2.25318) (14.50250,2.24610) (14.51556,2.23923) (14.52861,2.23257) (14.54167,2.22615) (14.55472,2.21995) (14.56778,2.21398) (14.58083,2.20825) (14.59389,2.20277) (14.60694,2.19754) (14.62000,2.19256) (14.63306,2.18783) (14.64611,2.18337) (14.65917,2.17917) (14.67222,2.17524) (14.68528,2.17159) (14.69833,2.16821) (14.71139,2.16511) (14.72444,2.16229) (14.73750,2.15976) (14.75056,2.15752) (14.76361,2.15556) (14.77667,2.15390) (14.78972,2.15253) (14.80278,2.15145) (14.81583,2.15067) (14.82889,2.15019) (14.84194,2.15000) (14.85500,2.15011) (14.86806,2.15052) (14.88111,2.15122) (14.89417,2.15222) (14.90722,2.15351) (14.92028,2.15510) (14.93333,2.15698) (14.94639,2.15915) (14.95944,2.16161) (14.97250,2.16435) (14.98556,2.16738) (14.99861,2.17068) (15.01167,2.17427) (15.02472,2.17812) (15.03778,2.18225) (15.05083,2.18665) (15.06389,2.19130) (15.07694,2.19622) (15.09000,2.20139) (15.10306,2.20680) (15.11611,2.21247) (15.12917,2.21837) (15.14222,2.22451) (15.15528,2.23088) (15.16833,2.23747) (15.18139,2.24429) (15.19444,2.25132) (15.20750,2.25856) (15.22056,2.26600) (15.23361,2.27364) (15.24667,2.28148) (15.25972,2.28950) (15.27278,2.29770) (15.28583,2.30608) (15.29889,2.31463) (15.31194,2.32335) (15.32500,2.33223) (15.33806,2.34126) (15.35111,2.35044) (15.36417,2.35976) (15.37722,2.36922) (15.39028,2.37882) (15.40333,2.38854) (15.41639,2.39838) (15.42944,2.40834) (15.44250,2.41841) (15.45556,2.42858) (15.46861,2.43886) (15.48167,2.44923) (15.49472,2.45970) (15.50778,2.47025) (15.52083,2.48088) (15.53389,2.49159) (15.54694,2.50237) (15.56000,2.51321) (15.57306,2.52412) (15.58611,2.53509) (15.59917,2.54611) (15.61222,2.55718) (15.62528,2.56829) (15.63833,2.57944) (15.65139,2.59063) (15.66444,2.60186) (15.67750,2.61311) (15.69056,2.62439) (15.70361,2.63568) (15.71667,2.64700) (15.72972,2.65832) (15.74278,2.66966) (15.75583,2.68100) (15.76889,2.69235) (15.78194,2.70369) (15.79500,2.71503) (15.80806,2.72636) (15.82111,2.73768) (15.83417,2.74899) (15.84722,2.76027) (15.86028,2.77154) (15.87333,2.78278) (15.88639,2.79400) (15.89944,2.80519) (15.91250,2.81634) (15.92556,2.82746) (15.93861,2.83854) (15.95167,2.84958) (15.96472,2.86058) (15.97778,2.87153) (15.99083,2.88243) (16.00389,2.89328) (16.01694,2.90408) (16.03000,2.91482) (16.04306,2.92550) (16.05611,2.93613) (16.06917,2.94669) (16.08222,2.95718) (16.09528,2.96761) (16.10833,2.97796) (16.12139,2.98825) (16.13444,2.99846) (16.14750,3.00859) (16.16056,3.01865) (16.17361,3.02863) (16.18667,3.03852) (16.19972,3.04834) (16.21278,3.05806) (16.22583,3.06770) (16.23889,3.07725) (16.25194,3.08671) (16.26500,3.09608) (16.27806,3.10535) (16.29111,3.11453) (16.30417,3.12361) (16.31722,3.13259) (16.33028,3.14147) (16.34333,3.15024) (16.35639,3.15892) (16.36944,3.16748) (16.38250,3.17595) (16.39556,3.18430) (16.40861,3.19254) (16.42167,3.20068) (16.43472,3.20870) (16.44778,3.21661) (16.46083,3.22440) (16.47389,3.23208) (16.48694,3.23964) (16.50000,3.24708)};
\draw[->,coursemuted] plot coordinates {(7.10000,4.65000) (16.60000,4.65000)};
\draw[->,coursemuted] plot coordinates {(7.10000,4.65000) (7.10000,6.50000)};
\draw[coursemuted!25,thin] plot coordinates {(7.10000,4.65000) (16.50000,4.65000)};
\node[left,inner sep=3pt] at (7.1,4.65000) {$0$};
\draw[coursemuted!25,thin] plot coordinates {(7.10000,5.45000) (16.50000,5.45000)};
\node[left,inner sep=3pt] at (7.1,5.45000) {$4$};
\draw[coursemuted!25,thin] plot coordinates {(7.10000,6.25000) (16.50000,6.25000)};
\node[left,inner sep=3pt] at (7.1,6.25000) {$8$};
\node[rotate=90] at (5.9,5.55) {$\Delta V_j$, м/с};
\draw[coursemuted!16,thin] plot coordinates {(7.10000,0.80000) (7.10000,6.25000)};
\node[below,inner sep=3pt] at (7.10000,0.85) {$0$};
\draw[coursemuted!16,thin] plot coordinates {(9.45000,0.80000) (9.45000,6.25000)};
\node[below,inner sep=3pt] at (9.45000,0.85) {$\pi/2$};
\draw[coursemuted!16,thin] plot coordinates {(11.80000,0.80000) (11.80000,6.25000)};
\node[below,inner sep=3pt] at (11.80000,0.85) {$\pi$};
\draw[coursemuted!16,thin] plot coordinates {(14.15000,0.80000) (14.15000,6.25000)};
\node[below,inner sep=3pt] at (14.15000,0.85) {$3\pi/2$};
\draw[coursemuted!16,thin] plot coordinates {(16.50000,0.80000) (16.50000,6.25000)};
\node[below,inner sep=3pt] at (16.50000,0.85) {$2\pi$};
\node at (11.8,.15) {Угол $\varphi$, рад};
\draw[courseink!70,densely dashed,thin] plot coordinates {(7.79364,0.85000) (7.79364,3.45000)};
\draw[courseink,line width=1pt] plot coordinates {(7.79364,4.65000) (7.79364,5.87098)};
\fill[courseink] (7.79364,5.87098) circle[radius=1.8pt];
\fill[courseink] (7.79364,3.45000) circle[radius=1.8pt];
\node[above,inner sep=2pt] at (7.79364,5.87098) {1};
\draw[courseink!70,densely dashed,thin] plot coordinates {(12.49364,0.85000) (12.49364,3.45000)};
\draw[courseink,line width=1pt] plot coordinates {(12.49364,4.65000) (12.49364,5.11637)};
\fill[courseink] (12.49364,5.11637) circle[radius=1.8pt];
\fill[courseink] (12.49364,3.45000) circle[radius=1.8pt];
\node[above,inner sep=2pt] at (12.49364,5.11637) {2};
\end{tikzpicture}

\par
\figcaption{$\Delta a=0{,}001$, $\Delta e_x=0{,}002$, $\Delta e_y=0{,}001$, $\Delta i=0$.\\
Пунктир отмечает уровень $\|\boldsymbol s\|=1$; точки пронумерованы по импульсам.}
\end{frame}

% eq:planar-impulses; стоимость и смысл мест импульсов: строки 421--442.
% После кадра «Пересекающиеся орбиты: годограф».
% В предшествующем кадре геометрии оставить только форму годографа,
% противоположные знаки и углы phi_1=alpha, phi_2=alpha+pi.
% manual/chapter-3-algorithms.tex, sec:planar-crossing.
\begin{frame}{Пересекающиеся орбиты}
\relax
\derivationmathspacing
Нумеруем разгонный импульс первым, тормозной --- вторым:
\[
 \Delta\widehat V_{t1}>0,\qquad\Delta\widehat V_{t2}<0,\qquad
 \varphi_1=\alpha,\qquad\varphi_2=\alpha+\pi.
\]
Подстановка в условия на эксцентриситет даёт
\[
\begin{aligned}
 2\Delta\widehat V_{t1}\cos\alpha
   +2\Delta\widehat V_{t2}\cos(\alpha+\pi)&=\Delta e_x,\\
 2\Delta\widehat V_{t1}\sin\alpha
   +2\Delta\widehat V_{t2}\sin(\alpha+\pi)&=\Delta e_y.
\end{aligned}
\]
Так как $\cos(\alpha+\pi)=-\cos\alpha$ и
$\sin(\alpha+\pi)=-\sin\alpha$, получаем
\[
\begin{aligned}
 2(\Delta\widehat V_{t1}-\Delta\widehat V_{t2})\cos\alpha&=\Delta e_x,\\
 2(\Delta\widehat V_{t1}-\Delta\widehat V_{t2})\sin\alpha&=\Delta e_y.
\end{aligned}
\]
\end{frame}

\begin{frame}{Пересекающиеся орбиты}
\relax
\derivationmathspacing
Возводим два полученных условия в квадрат и складываем:
\[
\begin{aligned}
 4(\Delta\widehat V_{t1}-\Delta\widehat V_{t2})^2
   (\cos^2\alpha+\sin^2\alpha)
   &=(\Delta e_x)^2+(\Delta e_y)^2,\\[2pt]
 4(\Delta\widehat V_{t1}-\Delta\widehat V_{t2})^2&=\Delta e^2.
\end{aligned}
\]
\justifying
Поскольку первый импульс разгонный, а второй тормозной,
$\Delta\widehat V_{t1}-\Delta\widehat V_{t2}>0$.
Поэтому при извлечении корня выбираем положительный знак:
\par
\[
 \boxed{\quad
 2(\Delta\widehat V_{t1}-\Delta\widehat V_{t2})=\Delta e.
 \quad}
\]
Так условия на эксцентриситет определяют разность компонент.
\end{frame}

\begin{frame}{Пересекающиеся орбиты}
\relax
\derivationmathspacing
Делим условия на эксцентриситет на найденную величину
$2(\Delta\widehat V_{t1}-\Delta\widehat V_{t2})=\Delta e$:
\[
 \cos\alpha=\frac{\Delta e_x}{\Delta e},\qquad
 \sin\alpha=\frac{\Delta e_y}{\Delta e}.
\]
Это определение угла оптимальной коррекции эксцентриситета:
\[
 \alpha=\varphi_e=\operatorname{atan2}(\Delta e_y,\Delta e_x)
 \quad\text{по модулю }2\pi.
\]
\begin{block}{Места двух импульсов}
\[
 \varphi_1=\varphi_e,\qquad \varphi_2=\varphi_e+\pi.
\]
\end{block}
\justifying
Импульсы лежат на линии направления разности векторов эксцентриситета.
Оба места найдены из граничных условий.
\end{frame}

\begin{frame}{Пересекающиеся орбиты}
\relax
\derivationmathspacing
Условие на большую полуось задаёт сумму компонент,
условия на эксцентриситет --- их разность:
\[
\begin{cases}
 2(\Delta\widehat V_{t1}+\Delta\widehat V_{t2})=\Delta a,\\
 2(\Delta\widehat V_{t1}-\Delta\widehat V_{t2})=\Delta e.
\end{cases}
\]
Складывая эти равенства и вычитая второе из первого, получаем
\[
 4\Delta\widehat V_{t1}=\Delta a+\Delta e,\qquad
 4\Delta\widehat V_{t2}=\Delta a-\Delta e.
\]
\begin{block}{Два трансверсальных импульса}
\[
 \Delta\widehat V_{t1}=\frac{\Delta a+\Delta e}{4},\qquad
 \Delta\widehat V_{t2}=\frac{\Delta a-\Delta e}{4},\qquad
 \Delta\widehat V_{r1}=\Delta\widehat V_{r2}=0.
\]
\end{block}
\end{frame}

\begin{frame}{Пересекающиеся орбиты}
\relax
\derivationmathspacing
При $|\Delta a|<\Delta e$ найденные компоненты имеют нужные знаки:
\[
 \Delta\widehat V_{t1}=\frac{\Delta a+\Delta e}{4}>0,\qquad
 \Delta\widehat V_{t2}=\frac{\Delta a-\Delta e}{4}<0.
\]
Следовательно, сумма модулей равна
\[
\begin{aligned}
 J&=|\Delta\widehat V_{t1}|+|\Delta\widehat V_{t2}|\\
  &=\Delta\widehat V_{t1}-\Delta\widehat V_{t2}
   =\frac{\Delta a+\Delta e}{4}-\frac{\Delta a-\Delta e}{4}
   =\frac{\Delta e}{2}.
\end{aligned}
\]
\justifying
Граничные условия выполнены, направления совпадают с базис-вектором
в двух точках касания; между ними его модуль меньше единицы.
Достаточные условия оптимальности выполнены.
\par\medskip
Переход выполняется через промежуточную орбиту.
Места импульсов не являются точками пересечения заданных орбит.
\end{frame}


% Источник: manual/chapter-3-algorithms.tex, sec:spatial;
% fig:spatial-hodographs-overview. Рисунок перенесён из методички.
\begin{frame}{Некомпланарный переход: типы годографов}
\relax
\derivationmathspacing
\centering
% Годографы из методички; сохранены координаты всех кривых и осей.
\begin{tikzpicture}[course diagram,scale=.60]
\begin{scope}[shift={(0.0,-0.0)}]
\fill[coursetint!45!white] (0,0) circle[radius=1.9cm];
\filldraw[fill=courseaccent,fill opacity=.09,draw=courseaccent!30] (1.34531,0.06275)--(0.45802,2.79341)--(-1.34531,-0.06275)--(-0.45802,-2.79341)--cycle;
\draw[coursemuted!40,thin,densely dashed] plot coordinates {(-1.45548,-0.51614) (-1.41198,-0.53729) (-1.36675,-0.55779) (-1.31985,-0.57761) (-1.27135,-0.59673) (-1.22130,-0.61511) (-1.16976,-0.63275) (-1.11679,-0.64962) (-1.06247,-0.66570) (-1.00685,-0.68096) (-0.95000,-0.69540) (-0.89200,-0.70898) (-0.83291,-0.72171) (-0.77280,-0.73355) (-0.71175,-0.74451) (-0.64984,-0.75455) (-0.58713,-0.76367) (-0.52371,-0.77187) (-0.45965,-0.77912) (-0.39503,-0.78543) (-0.32993,-0.79078) (-0.26443,-0.79516) (-0.19860,-0.79858) (-0.13254,-0.80102) (-0.06631,-0.80249) (-0.00000,-0.80297) (0.06631,-0.80249) (0.13254,-0.80102) (0.19860,-0.79858) (0.26443,-0.79516) (0.32993,-0.79078) (0.39503,-0.78543) (0.45965,-0.77912) (0.52371,-0.77187) (0.58713,-0.76367) (0.64984,-0.75455) (0.71175,-0.74451) (0.77280,-0.73355) (0.83291,-0.72171) (0.89200,-0.70898) (0.95000,-0.69540) (1.00685,-0.68096) (1.06247,-0.66570) (1.11679,-0.64962) (1.16976,-0.63275) (1.22130,-0.61511) (1.27135,-0.59673) (1.31985,-0.57761) (1.36675,-0.55779) (1.41198,-0.53729) (1.45548,-0.51614) (1.49722,-0.49436) (1.53713,-0.47198) (1.57517,-0.44902) (1.61129,-0.42551) (1.64545,-0.40149) (1.67760,-0.37697) (1.70771,-0.35200) (1.73574,-0.32660) (1.76165,-0.30080) (1.78542,-0.27463) (1.80701,-0.24813) (1.82640,-0.22133) (1.84356,-0.19426) (1.85848,-0.16695) (1.87113,-0.13944) (1.88151,-0.11175) (1.88959,-0.08393) (1.89537,-0.05601) (1.89884,-0.02802) (1.90000,-0.00000) (1.89884,0.02802) (1.89537,0.05601) (1.88959,0.08393) (1.88151,0.11175) (1.87113,0.13944) (1.85848,0.16695) (1.84356,0.19426) (1.82640,0.22133) (1.80701,0.24813) (1.78542,0.27463) (1.76165,0.30080) (1.73574,0.32660) (1.70771,0.35200) (1.67760,0.37697) (1.64545,0.40149) (1.61129,0.42551) (1.57517,0.44902) (1.53713,0.47198) (1.49722,0.49436) (1.45548,0.51614) (1.41198,0.53729) (1.36675,0.55779) (1.31985,0.57761) (1.27135,0.59673) (1.22130,0.61511) (1.16976,0.63275) (1.11679,0.64962) (1.06247,0.66570) (1.00685,0.68096) (0.95000,0.69540) (0.89200,0.70898) (0.83291,0.72171) (0.77280,0.73355) (0.71175,0.74451) (0.64984,0.75455) (0.58713,0.76367) (0.52371,0.77187) (0.45965,0.77912) (0.39503,0.78543) (0.32993,0.79078) (0.26443,0.79516) (0.19860,0.79858) (0.13254,0.80102) (0.06631,0.80249) (0.00000,0.80297) (-0.06631,0.80249) (-0.13254,0.80102) (-0.19860,0.79858) (-0.26443,0.79516) (-0.32993,0.79078) (-0.39503,0.78543) (-0.45965,0.77912) (-0.52371,0.77187) (-0.58713,0.76367) (-0.64984,0.75455) (-0.71175,0.74451) (-0.77280,0.73355) (-0.83291,0.72171) (-0.89200,0.70898) (-0.95000,0.69540) (-1.00685,0.68096) (-1.06247,0.66570) (-1.11679,0.64962) (-1.16976,0.63275) (-1.22130,0.61511) (-1.27135,0.59673) (-1.31985,0.57761) (-1.36675,0.55779) (-1.41198,0.53729) (-1.45548,0.51614) (-1.49722,0.49436) (-1.53713,0.47198) (-1.57517,0.44902) (-1.61129,0.42551) (-1.64545,0.40149) (-1.67760,0.37697) (-1.70771,0.35200) (-1.73574,0.32660) (-1.76165,0.30080) (-1.78542,0.27463) (-1.80701,0.24813) (-1.82640,0.22133) (-1.84356,0.19426) (-1.85848,0.16695) (-1.87113,0.13944) (-1.88151,0.11175) (-1.88959,0.08393) (-1.89537,0.05601) (-1.89884,0.02802) (-1.90000,0.00000) (-1.89884,-0.02802) (-1.89537,-0.05601) (-1.88959,-0.08393) (-1.88151,-0.11175) (-1.87113,-0.13944) (-1.85848,-0.16695) (-1.84356,-0.19426) (-1.82640,-0.22133) (-1.80701,-0.24813) (-1.78542,-0.27463) (-1.76165,-0.30080) (-1.73574,-0.32660) (-1.70771,-0.35200) (-1.67760,-0.37697) (-1.64545,-0.40149) (-1.61129,-0.42551) (-1.57517,-0.44902) (-1.53713,-0.47198) (-1.49722,-0.49436) (-1.45548,-0.51614)};
\draw[coursemuted!55,thin] plot coordinates {(-1.90000,-0.00000) (-1.89884,-0.02802) (-1.89537,-0.05601) (-1.88959,-0.08393) (-1.88151,-0.11175) (-1.87113,-0.13944) (-1.85848,-0.16695) (-1.84356,-0.19426) (-1.82640,-0.22133) (-1.80701,-0.24813) (-1.78542,-0.27463) (-1.76165,-0.30080) (-1.73574,-0.32660) (-1.70771,-0.35200) (-1.67760,-0.37697) (-1.64545,-0.40149) (-1.61129,-0.42551) (-1.57517,-0.44902) (-1.53713,-0.47198) (-1.49722,-0.49436) (-1.45548,-0.51614) (-1.41198,-0.53729) (-1.36675,-0.55779) (-1.31985,-0.57761) (-1.27135,-0.59673) (-1.22130,-0.61511) (-1.16976,-0.63275) (-1.11679,-0.64962) (-1.06247,-0.66570) (-1.00685,-0.68096) (-0.95000,-0.69540) (-0.89200,-0.70898) (-0.83291,-0.72171) (-0.77280,-0.73355) (-0.71175,-0.74451) (-0.64984,-0.75455) (-0.58713,-0.76367) (-0.52371,-0.77187) (-0.45965,-0.77912) (-0.39503,-0.78543) (-0.32993,-0.79078) (-0.26443,-0.79516) (-0.19860,-0.79858) (-0.13254,-0.80102) (-0.06631,-0.80249) (-0.00000,-0.80297) (0.06631,-0.80249) (0.13254,-0.80102) (0.19860,-0.79858) (0.26443,-0.79516) (0.32993,-0.79078) (0.39503,-0.78543) (0.45965,-0.77912) (0.52371,-0.77187) (0.58713,-0.76367) (0.64984,-0.75455) (0.71175,-0.74451) (0.77280,-0.73355) (0.83291,-0.72171) (0.89200,-0.70898) (0.95000,-0.69540) (1.00685,-0.68096) (1.06247,-0.66570) (1.11679,-0.64962) (1.16976,-0.63275) (1.22130,-0.61511) (1.27135,-0.59673) (1.31985,-0.57761) (1.36675,-0.55779) (1.41198,-0.53729) (1.45548,-0.51614) (1.49722,-0.49436) (1.53713,-0.47198) (1.57517,-0.44902) (1.61129,-0.42551) (1.64545,-0.40149) (1.67760,-0.37697) (1.70771,-0.35200) (1.73574,-0.32660) (1.76165,-0.30080) (1.78542,-0.27463) (1.80701,-0.24813) (1.82640,-0.22133) (1.84356,-0.19426) (1.85848,-0.16695) (1.87113,-0.13944) (1.88151,-0.11175) (1.88959,-0.08393) (1.89537,-0.05601) (1.89884,-0.02802) (1.90000,-0.00000)};
\draw[coursemuted!40,thin,densely dashed] plot coordinates {(-1.45548,-0.51614) (-1.45460,-0.45573) (-1.45194,-0.39477) (-1.44751,-0.33332) (-1.44132,-0.27147) (-1.43337,-0.20928) (-1.42368,-0.14684) (-1.41225,-0.08422) (-1.39910,-0.02150) (-1.38425,0.04124) (-1.36771,0.10394) (-1.34950,0.16651) (-1.32965,0.22887) (-1.30818,0.29096) (-1.28512,0.35270) (-1.26049,0.41400) (-1.23432,0.47480) (-1.20665,0.53502) (-1.17751,0.59459) (-1.14694,0.65343) (-1.11497,0.71148) (-1.08164,0.76866) (-1.04699,0.82491) (-1.01106,0.88015) (-0.97391,0.93432) (-0.93557,0.98735) (-0.89609,1.03917) (-0.85551,1.08973) (-0.81390,1.13897) (-0.77129,1.18681) (-0.72774,1.23321) (-0.68331,1.27811) (-0.63804,1.32145) (-0.59200,1.36318) (-0.54523,1.40325) (-0.49780,1.44161) (-0.44977,1.47821) (-0.40119,1.51301) (-0.35211,1.54597) (-0.30261,1.57704) (-0.25274,1.60620) (-0.20256,1.63339) (-0.15214,1.65860) (-0.10153,1.68179) (-0.05080,1.70292) (-0.00000,1.72198) (0.05080,1.73895) (0.10153,1.75379) (0.15214,1.76650) (0.20256,1.77706) (0.25274,1.78545) (0.30261,1.79167) (0.35211,1.79570) (0.40119,1.79755) (0.44977,1.79720) (0.49780,1.79467) (0.54523,1.78995) (0.59200,1.78305) (0.63804,1.77397) (0.68331,1.76274) (0.72774,1.74935) (0.77129,1.73384) (0.81390,1.71621) (0.85551,1.69650) (0.89609,1.67471) (0.93557,1.65089) (0.97391,1.62505) (1.01106,1.59723) (1.04699,1.56747) (1.08164,1.53580) (1.11497,1.50226) (1.14694,1.46689) (1.17751,1.42973) (1.20665,1.39082) (1.23432,1.35023) (1.26049,1.30798) (1.28512,1.26415) (1.30818,1.21877) (1.32965,1.17191) (1.34950,1.12363) (1.36771,1.07397) (1.38425,1.02300) (1.39910,0.97079) (1.41225,0.91740) (1.42368,0.86288) (1.43337,0.80732) (1.44132,0.75077) (1.44751,0.69331) (1.45194,0.63500) (1.45460,0.57592) (1.45548,0.51614) (1.45460,0.45573) (1.45194,0.39477) (1.44751,0.33332) (1.44132,0.27147) (1.43337,0.20928) (1.42368,0.14684) (1.41225,0.08422) (1.39910,0.02150) (1.38425,-0.04124) (1.36771,-0.10394) (1.34950,-0.16651) (1.32965,-0.22887) (1.30818,-0.29096) (1.28512,-0.35270) (1.26049,-0.41400) (1.23432,-0.47480) (1.20665,-0.53502) (1.17751,-0.59459) (1.14694,-0.65343) (1.11497,-0.71148) (1.08164,-0.76866) (1.04699,-0.82491) (1.01106,-0.88015) (0.97391,-0.93432) (0.93557,-0.98735) (0.89609,-1.03917) (0.85551,-1.08973) (0.81390,-1.13897) (0.77129,-1.18681) (0.72774,-1.23321) (0.68331,-1.27811) (0.63804,-1.32145) (0.59200,-1.36318) (0.54523,-1.40325) (0.49780,-1.44161) (0.44977,-1.47821) (0.40119,-1.51301) (0.35211,-1.54597) (0.30261,-1.57704) (0.25274,-1.60620) (0.20256,-1.63339) (0.15214,-1.65860) (0.10153,-1.68179) (0.05080,-1.70292) (0.00000,-1.72198) (-0.05080,-1.73895) (-0.10153,-1.75379) (-0.15214,-1.76650) (-0.20256,-1.77706) (-0.25274,-1.78545) (-0.30261,-1.79167) (-0.35211,-1.79570) (-0.40119,-1.79755) (-0.44977,-1.79720) (-0.49780,-1.79467) (-0.54523,-1.78995) (-0.59200,-1.78305) (-0.63804,-1.77397) (-0.68331,-1.76274) (-0.72774,-1.74935) (-0.77129,-1.73384) (-0.81390,-1.71621) (-0.85551,-1.69650) (-0.89609,-1.67471) (-0.93557,-1.65089) (-0.97391,-1.62505) (-1.01106,-1.59723) (-1.04699,-1.56747) (-1.08164,-1.53580) (-1.11497,-1.50226) (-1.14694,-1.46689) (-1.17751,-1.42973) (-1.20665,-1.39082) (-1.23432,-1.35023) (-1.26049,-1.30798) (-1.28512,-1.26415) (-1.30818,-1.21877) (-1.32965,-1.17191) (-1.34950,-1.12363) (-1.36771,-1.07397) (-1.38425,-1.02300) (-1.39910,-0.97079) (-1.41225,-0.91740) (-1.42368,-0.86288) (-1.43337,-0.80732) (-1.44132,-0.75077) (-1.44751,-0.69331) (-1.45194,-0.63500) (-1.45460,-0.57592) (-1.45548,-0.51614)};
\draw[coursemuted!55,thin] plot coordinates {(-0.85467,-1.69692) (-0.89526,-1.67518) (-0.93477,-1.65140) (-0.97313,-1.62560) (-1.01031,-1.59783) (-1.04626,-1.56810) (-1.08094,-1.53647) (-1.11429,-1.50297) (-1.14629,-1.46763) (-1.17690,-1.43051) (-1.20607,-1.39164) (-1.23377,-1.35108) (-1.25996,-1.30887) (-1.28463,-1.26507) (-1.30772,-1.21972) (-1.32923,-1.17289) (-1.34911,-1.12463) (-1.36735,-1.07500) (-1.38392,-1.02406) (-1.39881,-0.97188) (-1.41200,-0.91851) (-1.42346,-0.86402) (-1.43319,-0.80847) (-1.44117,-0.75195) (-1.44740,-0.69450) (-1.45187,-0.63621) (-1.45456,-0.57715) (-1.45548,-0.51738) (-1.45463,-0.45698) (-1.45201,-0.39602) (-1.44762,-0.33459) (-1.44146,-0.27274) (-1.43355,-0.21056) (-1.42390,-0.14813) (-1.41250,-0.08551) (-1.39939,-0.02280) (-1.38457,0.03995) (-1.36807,0.10265) (-1.34989,0.16522) (-1.33008,0.22759) (-1.30864,0.28969) (-1.28561,0.35143) (-1.26101,0.41274) (-1.23487,0.47355) (-1.20724,0.53379) (-1.17813,0.59337) (-1.14758,0.65223) (-1.11564,0.71030) (-1.08233,0.76750) (-1.04771,0.82376) (-1.01182,0.87902) (-0.97469,0.93321) (-0.93637,0.98627) (-0.89691,1.03812) (-0.85636,1.08871) (-0.81476,1.13797) (-0.77218,1.18584) (-0.72865,1.23227) (-0.68423,1.27720) (-0.63898,1.32057) (-0.59295,1.36234) (-0.54620,1.40244) (-0.49879,1.44083) (-0.45076,1.47747) (-0.40219,1.51231) (-0.35313,1.54531) (-0.30363,1.57642) (-0.25377,1.60562) (-0.20360,1.63285) (-0.15318,1.65810) (-0.10257,1.68133) (-0.05184,1.70251) (-0.00105,1.72161) (0.04975,1.73862) (0.10049,1.75351) (0.15110,1.76626) (0.20153,1.77686) (0.25171,1.78530) (0.30159,1.79156) (0.35110,1.79564) (0.40018,1.79753) (0.44878,1.79723) (0.49682,1.79474) (0.54426,1.79007) (0.59104,1.78321) (0.63710,1.77418) (0.68239,1.76299) (0.72684,1.74965) (0.77040,1.73418) (0.81303,1.71660) (0.85467,1.69692)};
\draw[coursemuted!40,thin,densely dashed] plot coordinates {(1.22130,-0.61511) (1.22055,-0.55464) (1.21832,-0.49350) (1.21461,-0.43175) (1.20941,-0.36947) (1.20274,-0.30675) (1.19461,-0.24365) (1.18502,-0.18026) (1.17399,-0.11664) (1.16152,-0.05289) (1.14764,0.01094) (1.13237,0.07474) (1.11571,0.13846) (1.09769,0.20201) (1.07834,0.26531) (1.05767,0.32829) (1.03572,0.39087) (1.01250,0.45297) (0.98805,0.51452) (0.96239,0.57544) (0.93557,0.63567) (0.90760,0.69511) (0.87853,0.75371) (0.84838,0.81140) (0.81721,0.86809) (0.78503,0.92373) (0.75191,0.97824) (0.71786,1.03156) (0.68294,1.08362) (0.64719,1.13437) (0.61065,1.18373) (0.57336,1.23164) (0.53538,1.27806) (0.49675,1.32292) (0.45751,1.36617) (0.41771,1.40775) (0.37740,1.44762) (0.33663,1.48573) (0.29546,1.52202) (0.25392,1.55647) (0.21208,1.58901) (0.16997,1.61962) (0.12766,1.64825) (0.08519,1.67488) (0.04262,1.69947) (0.00000,1.72198) (-0.04262,1.74240) (-0.08519,1.76070) (-0.12766,1.77685) (-0.16997,1.79083) (-0.21208,1.80264) (-0.25392,1.81224) (-0.29546,1.81964) (-0.33663,1.82483) (-0.37740,1.82779) (-0.41771,1.82852) (-0.45751,1.82702) (-0.49675,1.82330) (-0.53538,1.81736) (-0.57336,1.80920) (-0.61065,1.79884) (-0.64719,1.78629) (-0.68294,1.77156) (-0.71786,1.75467) (-0.75191,1.73564) (-0.78503,1.71450) (-0.81721,1.69128) (-0.84838,1.66599) (-0.87853,1.63867) (-0.90760,1.60935) (-0.93557,1.57808) (-0.96239,1.54488) (-0.98805,1.50980) (-1.01250,1.47287) (-1.03572,1.43416) (-1.05767,1.39370) (-1.07834,1.35154) (-1.09769,1.30773) (-1.11571,1.26233) (-1.13237,1.21539) (-1.14764,1.16697) (-1.16152,1.11713) (-1.17399,1.06593) (-1.18502,1.01343) (-1.19461,0.95969) (-1.20274,0.90479) (-1.20941,0.84878) (-1.21461,0.79174) (-1.21832,0.73374) (-1.22055,0.67484) (-1.22130,0.61511) (-1.22055,0.55464) (-1.21832,0.49350) (-1.21461,0.43175) (-1.20941,0.36947) (-1.20274,0.30675) (-1.19461,0.24365) (-1.18502,0.18026) (-1.17399,0.11664) (-1.16152,0.05289) (-1.14764,-0.01094) (-1.13237,-0.07474) (-1.11571,-0.13846) (-1.09769,-0.20201) (-1.07834,-0.26531) (-1.05767,-0.32829) (-1.03572,-0.39087) (-1.01250,-0.45297) (-0.98805,-0.51452) (-0.96239,-0.57544) (-0.93557,-0.63567) (-0.90760,-0.69511) (-0.87853,-0.75371) (-0.84838,-0.81140) (-0.81721,-0.86809) (-0.78503,-0.92373) (-0.75191,-0.97824) (-0.71786,-1.03156) (-0.68294,-1.08362) (-0.64719,-1.13437) (-0.61065,-1.18373) (-0.57336,-1.23164) (-0.53538,-1.27806) (-0.49675,-1.32292) (-0.45751,-1.36617) (-0.41771,-1.40775) (-0.37740,-1.44762) (-0.33663,-1.48573) (-0.29546,-1.52202) (-0.25392,-1.55647) (-0.21208,-1.58901) (-0.16997,-1.61962) (-0.12766,-1.64825) (-0.08519,-1.67488) (-0.04262,-1.69947) (-0.00000,-1.72198) (0.04262,-1.74240) (0.08519,-1.76070) (0.12766,-1.77685) (0.16997,-1.79083) (0.21208,-1.80264) (0.25392,-1.81224) (0.29546,-1.81964) (0.33663,-1.82483) (0.37740,-1.82779) (0.41771,-1.82852) (0.45751,-1.82702) (0.49675,-1.82330) (0.53538,-1.81736) (0.57336,-1.80920) (0.61065,-1.79884) (0.64719,-1.78629) (0.68294,-1.77156) (0.71786,-1.75467) (0.75191,-1.73564) (0.78503,-1.71450) (0.81721,-1.69128) (0.84838,-1.66599) (0.87853,-1.63867) (0.90760,-1.60935) (0.93557,-1.57808) (0.96239,-1.54488) (0.98805,-1.50980) (1.01250,-1.47287) (1.03572,-1.43416) (1.05767,-1.39370) (1.07834,-1.35154) (1.09769,-1.30773) (1.11571,-1.26233) (1.13237,-1.21539) (1.14764,-1.16697) (1.16152,-1.11713) (1.17399,-1.06593) (1.18502,-1.01343) (1.19461,-0.95969) (1.20274,-0.90479) (1.20941,-0.84878) (1.21461,-0.79174) (1.21832,-0.73374) (1.22055,-0.67484) (1.22130,-0.61511)};
\draw[coursemuted!55,thin] plot coordinates {(0.63503,-1.79074) (0.67105,-1.77674) (0.70625,-1.76057) (0.74060,-1.74226) (0.77404,-1.72182) (0.80653,-1.69929) (0.83805,-1.67469) (0.86854,-1.64805) (0.89798,-1.61940) (0.92632,-1.58877) (0.95353,-1.55621) (0.97958,-1.52176) (1.00444,-1.48545) (1.02808,-1.44733) (1.05046,-1.40745) (1.07156,-1.36585) (1.09136,-1.32259) (1.10982,-1.27772) (1.12694,-1.23129) (1.14268,-1.18336) (1.15703,-1.13399) (1.16997,-1.08324) (1.18148,-1.03116) (1.19156,-0.97783) (1.20018,-0.92331) (1.20734,-0.86767) (1.21303,-0.81097) (1.21724,-0.75328) (1.21997,-0.69467) (1.22121,-0.63522) (1.22097,-0.57499) (1.21923,-0.51406) (1.21602,-0.45250) (1.21132,-0.39040) (1.20514,-0.32782) (1.19750,-0.26483) (1.18839,-0.20153) (1.17784,-0.13798) (1.16586,-0.07426) (1.15245,-0.01045) (1.13764,0.05337) (1.12144,0.11712) (1.10388,0.18074) (1.08497,0.24413) (1.06474,0.30722) (1.04322,0.36994) (1.02042,0.43222) (0.99638,0.49396) (0.97112,0.55510) (0.94469,0.61557) (0.91710,0.67528) (0.88839,0.73418) (0.85860,0.79217) (0.82777,0.84921) (0.79592,0.90521) (0.76311,0.96010) (0.72937,1.01383) (0.69474,1.06632) (0.65926,1.11751) (0.62298,1.16734) (0.58594,1.21575) (0.54818,1.26268) (0.50976,1.30807) (0.47072,1.35186) (0.43110,1.39401) (0.39096,1.43446) (0.35034,1.47316) (0.30930,1.51007) (0.26788,1.54513) (0.22613,1.57832) (0.18411,1.60958) (0.14186,1.63888) (0.09944,1.66618) (0.05690,1.69146) (0.01429,1.71467) (-0.02834,1.73580) (-0.07094,1.75481) (-0.11344,1.77168) (-0.15581,1.78639) (-0.19799,1.79893) (-0.23993,1.80927) (-0.28158,1.81741) (-0.32288,1.82334) (-0.36379,1.82704) (-0.40425,1.82852) (-0.44423,1.82777) (-0.48366,1.82480) (-0.52250,1.81960) (-0.56071,1.81218) (-0.59823,1.80256) (-0.63503,1.79074)};
\draw[coursemuted!50,thin,densely dashed] plot coordinates {(1.67381,0.59356) (0.00000,0.00000)};
\draw[->,courseink!75,thin] plot coordinates {(0.00000,0.00000) (-1.96490,-0.69679)};
\node[above left,inner sep=2pt] at (-1.96490,-0.69679) {$s_r$};
\draw[coursemuted!50,thin,densely dashed] plot coordinates {(-1.40449,0.70738) (0.00000,0.00000)};
\draw[->,courseink!75,thin] plot coordinates {(0.00000,0.00000) (1.64875,-0.83040)};
\node[below right,inner sep=2pt] at (1.64875,-0.83040) {$s_t$};
\draw[coursemuted!50,thin,densely dashed] plot coordinates {(0.00000,-1.98028) (0.00000,0.00000)};
\draw[->,courseink!75,thin] plot coordinates {(0.00000,0.00000) (0.00000,2.32468)};
\node[above,inner sep=2pt] at (0.00000,2.32468) {$s_z$};
\draw[courseaccent,line width=1.2pt] plot coordinates {(0.85063,1.34725) (0.85426,1.33547) (0.85763,1.32329) (0.86074,1.31070) (0.86358,1.29772) (0.86617,1.28434) (0.86849,1.27057) (0.87054,1.25641) (0.87233,1.24187) (0.87385,1.22696) (0.87511,1.21166) (0.87610,1.19600) (0.87683,1.17998) (0.87728,1.16359) (0.87747,1.14685) (0.87740,1.12976) (0.87705,1.11233) (0.87644,1.09456) (0.87556,1.07645) (0.87442,1.05802) (0.87300,1.03927) (0.87133,1.02019) (0.86938,1.00081) (0.86717,0.98112) (0.86470,0.96114) (0.86197,0.94086) (0.85897,0.92029) (0.85571,0.89945) (0.85219,0.87833) (0.84841,0.85694) (0.84437,0.83529) (0.84007,0.81339) (0.83552,0.79124) (0.83072,0.76885) (0.82566,0.74622) (0.82035,0.72337) (0.81479,0.70029) (0.80898,0.67701) (0.80292,0.65351) (0.79662,0.62982) (0.79008,0.60594) (0.78330,0.58187) (0.77628,0.55762) (0.76902,0.53321) (0.76152,0.50863) (0.75380,0.48389) (0.74585,0.45901) (0.73767,0.43399) (0.72926,0.40884) (0.72063,0.38356) (0.71178,0.35817) (0.70272,0.33267) (0.69344,0.30706) (0.68395,0.28136) (0.67425,0.25558) (0.66435,0.22972) (0.65424,0.20379) (0.64394,0.17779) (0.63344,0.15175) (0.62274,0.12565) (0.61186,0.09952) (0.60079,0.07336) (0.58954,0.04717) (0.57811,0.02097) (0.56650,-0.00523) (0.55471,-0.03143) (0.54276,-0.05763) (0.53065,-0.08381) (0.51837,-0.10996) (0.50594,-0.13608) (0.49335,-0.16215) (0.48061,-0.18818) (0.46772,-0.21415) (0.45469,-0.24005) (0.44153,-0.26588) (0.42822,-0.29163) (0.41479,-0.31730) (0.40123,-0.34286) (0.38755,-0.36832) (0.37376,-0.39367) (0.35984,-0.41890) (0.34582,-0.44400) (0.33169,-0.46896) (0.31746,-0.49379) (0.30314,-0.51846) (0.28872,-0.54297) (0.27422,-0.56732) (0.25963,-0.59150) (0.24496,-0.61549) (0.23022,-0.63930) (0.21540,-0.66292) (0.20053,-0.68633) (0.18559,-0.70953) (0.17059,-0.73252) (0.15554,-0.75528) (0.14045,-0.77781) (0.12531,-0.80011) (0.11013,-0.82216) (0.09492,-0.84396) (0.07968,-0.86551) (0.06442,-0.88679) (0.04914,-0.90780) (0.03384,-0.92854) (0.01853,-0.94899) (0.00322,-0.96915) (-0.01209,-0.98902) (-0.02741,-1.00859) (-0.04271,-1.02785) (-0.05800,-1.04679) (-0.07327,-1.06542) (-0.08852,-1.08372) (-0.10374,-1.10169) (-0.11893,-1.11933) (-0.13409,-1.13663) (-0.14920,-1.15358) (-0.16427,-1.17018) (-0.17929,-1.18642) (-0.19425,-1.20230) (-0.20916,-1.21781) (-0.22400,-1.23296) (-0.23877,-1.24772) (-0.25347,-1.26211) (-0.26809,-1.27612) (-0.28263,-1.28973) (-0.29709,-1.30295) (-0.31145,-1.31578) (-0.32572,-1.32820) (-0.33989,-1.34022) (-0.35396,-1.35183) (-0.36792,-1.36303) (-0.38176,-1.37382) (-0.39550,-1.38418) (-0.40911,-1.39412) (-0.42259,-1.40364) (-0.43595,-1.41274) (-0.44917,-1.42140) (-0.46226,-1.42963) (-0.47521,-1.43742) (-0.48801,-1.44478) (-0.50066,-1.45169) (-0.51316,-1.45816) (-0.52551,-1.46419) (-0.53769,-1.46978) (-0.54971,-1.47491) (-0.56156,-1.47960) (-0.57324,-1.48383) (-0.58475,-1.48761) (-0.59608,-1.49094) (-0.60723,-1.49382) (-0.61819,-1.49624) (-0.62897,-1.49821) (-0.63955,-1.49971) (-0.64994,-1.50077) (-0.66013,-1.50136) (-0.67011,-1.50150) (-0.67990,-1.50118) (-0.68948,-1.50040) (-0.69884,-1.49917) (-0.70800,-1.49748) (-0.71694,-1.49533) (-0.72566,-1.49273) (-0.73416,-1.48967) (-0.74243,-1.48616) (-0.75048,-1.48219) (-0.75830,-1.47778) (-0.76589,-1.47291) (-0.77325,-1.46760) (-0.78037,-1.46184) (-0.78726,-1.45563) (-0.79390,-1.44898) (-0.80030,-1.44189) (-0.80646,-1.43436) (-0.81237,-1.42639) (-0.81804,-1.41799) (-0.82345,-1.40916) (-0.82862,-1.39990) (-0.83353,-1.39021) (-0.83819,-1.38009) (-0.84259,-1.36956) (-0.84674,-1.35861) (-0.85063,-1.34725) (-0.85426,-1.33547) (-0.85763,-1.32329) (-0.86074,-1.31070) (-0.86358,-1.29772) (-0.86617,-1.28434) (-0.86849,-1.27057) (-0.87054,-1.25641) (-0.87233,-1.24187) (-0.87385,-1.22696) (-0.87511,-1.21166) (-0.87610,-1.19600) (-0.87683,-1.17998) (-0.87728,-1.16359) (-0.87747,-1.14685) (-0.87740,-1.12976) (-0.87705,-1.11233) (-0.87644,-1.09456) (-0.87556,-1.07645) (-0.87442,-1.05802) (-0.87300,-1.03927) (-0.87133,-1.02019) (-0.86938,-1.00081) (-0.86717,-0.98112) (-0.86470,-0.96114) (-0.86197,-0.94086) (-0.85897,-0.92029) (-0.85571,-0.89945) (-0.85219,-0.87833) (-0.84841,-0.85694) (-0.84437,-0.83529) (-0.84007,-0.81339) (-0.83552,-0.79124) (-0.83072,-0.76885) (-0.82566,-0.74622) (-0.82035,-0.72337) (-0.81479,-0.70029) (-0.80898,-0.67701) (-0.80292,-0.65351) (-0.79662,-0.62982) (-0.79008,-0.60594) (-0.78330,-0.58187) (-0.77628,-0.55762) (-0.76902,-0.53321) (-0.76152,-0.50863) (-0.75380,-0.48389) (-0.74585,-0.45901) (-0.73767,-0.43399) (-0.72926,-0.40884) (-0.72063,-0.38356) (-0.71178,-0.35817) (-0.70272,-0.33267) (-0.69344,-0.30706) (-0.68395,-0.28136) (-0.67425,-0.25558) (-0.66435,-0.22972) (-0.65424,-0.20379) (-0.64394,-0.17779) (-0.63344,-0.15175) (-0.62274,-0.12565) (-0.61186,-0.09952) (-0.60079,-0.07336) (-0.58954,-0.04717) (-0.57811,-0.02097) (-0.56650,0.00523) (-0.55471,0.03143) (-0.54276,0.05763) (-0.53065,0.08381) (-0.51837,0.10996) (-0.50594,0.13608) (-0.49335,0.16215) (-0.48061,0.18818) (-0.46772,0.21415) (-0.45469,0.24005) (-0.44153,0.26588) (-0.42822,0.29163) (-0.41479,0.31730) (-0.40123,0.34286) (-0.38755,0.36832) (-0.37376,0.39367) (-0.35984,0.41890) (-0.34582,0.44400) (-0.33169,0.46896) (-0.31746,0.49379) (-0.30314,0.51846) (-0.28872,0.54297) (-0.27422,0.56732) (-0.25963,0.59150) (-0.24496,0.61549) (-0.23022,0.63930) (-0.21540,0.66292) (-0.20053,0.68633) (-0.18559,0.70953) (-0.17059,0.73252) (-0.15554,0.75528) (-0.14045,0.77781) (-0.12531,0.80011) (-0.11013,0.82216) (-0.09492,0.84396) (-0.07968,0.86551) (-0.06442,0.88679) (-0.04914,0.90780) (-0.03384,0.92854) (-0.01853,0.94899) (-0.00322,0.96915) (0.01209,0.98902) (0.02741,1.00859) (0.04271,1.02785) (0.05800,1.04679) (0.07327,1.06542) (0.08852,1.08372) (0.10374,1.10169) (0.11893,1.11933) (0.13409,1.13663) (0.14920,1.15358) (0.16427,1.17018) (0.17929,1.18642) (0.19425,1.20230) (0.20916,1.21781) (0.22400,1.23296) (0.23877,1.24772) (0.25347,1.26211) (0.26809,1.27612) (0.28263,1.28973) (0.29709,1.30295) (0.31145,1.31578) (0.32572,1.32820) (0.33989,1.34022) (0.35396,1.35183) (0.36792,1.36303) (0.38176,1.37382) (0.39550,1.38418) (0.40911,1.39412) (0.42259,1.40364) (0.43595,1.41274) (0.44917,1.42140) (0.46226,1.42963) (0.47521,1.43742) (0.48801,1.44478) (0.50066,1.45169) (0.51316,1.45816) (0.52551,1.46419) (0.53769,1.46978) (0.54971,1.47491) (0.56156,1.47960) (0.57324,1.48383) (0.58475,1.48761) (0.59608,1.49094) (0.60723,1.49382) (0.61819,1.49624) (0.62897,1.49821) (0.63955,1.49971) (0.64994,1.50077) (0.66013,1.50136) (0.67011,1.50150) (0.67990,1.50118) (0.68948,1.50040) (0.69884,1.49917) (0.70800,1.49748) (0.71694,1.49533) (0.72566,1.49273) (0.73416,1.48967) (0.74243,1.48616) (0.75048,1.48219) (0.75830,1.47778) (0.76589,1.47291) (0.77325,1.46760) (0.78037,1.46184) (0.78726,1.45563) (0.79390,1.44898) (0.80030,1.44189) (0.80646,1.43436) (0.81237,1.42639) (0.81804,1.41799) (0.82345,1.40916) (0.82862,1.39990) (0.83353,1.39021) (0.83819,1.38009) (0.84259,1.36956) (0.84674,1.35861) (0.85063,1.34725)};
\draw[->,courseaccent,line width=.55pt] plot coordinates {(0.00000,0.00000) (0.85063,1.34725)};
\filldraw[fill=courseink,draw=coursepaper,line width=.4pt] (0.85063,1.34725) circle[radius=2.3pt];
\draw[->,courseaccent,line width=.55pt] plot coordinates {(0.00000,0.00000) (-0.85063,-1.34725)};
\filldraw[fill=courseink,draw=coursepaper,line width=.4pt] (-0.85063,-1.34725) circle[radius=2.3pt];
\draw[coursemuted!65,thin] (0,0) circle[radius=1.9cm];
\node[below left,inner sep=2pt] at (0,0) {$0$};
\node[align=center,text width=3.45cm,anchor=north] at (0,-3.15) {а) Узловое (NC)\\Эллипс с центром\\в начале координат};
\end{scope}
\begin{scope}[shift={(6.05,0)}]
\filldraw[fill=coursetint!45!white,draw=coursemuted!65,line width=.65pt] (0.00000,0.00000) circle[radius=1.9cm];
\draw[->,courseink!75,thin] plot coordinates {(0.00000,0.00000) (-1.90041,-1.63661)};
\node[below left,inner sep=2pt] at (-1.90041,-1.63661) {$s_r$};
\draw[->,courseink!75,thin] plot coordinates {(0.00000,0.00000) (2.50800,-0.00000)};
\node[right,inner sep=2pt] at (2.50800,-0.00000) {$s_t$};
\draw[->,courseink!75,thin] plot coordinates {(0.00000,0.00000) (-0.18772,2.50096)};
\node[above,inner sep=2pt] at (-0.18772,2.50096) {$s_z$};
\draw[coursemuted!60,thin,densely dashed] plot coordinates {(1.32544,1.36133) (1.32544,-1.36133)};
\draw[courseaccent,line width=1.2pt] plot coordinates {(1.27810,1.38984) (1.28339,1.38836) (1.28867,1.38647) (1.29395,1.38415) (1.29921,1.38141) (1.30447,1.37825) (1.30971,1.37467) (1.31494,1.37067) (1.32015,1.36625) (1.32535,1.36141) (1.33052,1.35617) (1.33568,1.35051) (1.34081,1.34443) (1.34592,1.33795) (1.35100,1.33106) (1.35606,1.32376) (1.36109,1.31607) (1.36608,1.30797) (1.37105,1.29947) (1.37598,1.29058) (1.38088,1.28129) (1.38574,1.27161) (1.39057,1.26155) (1.39535,1.25110) (1.40010,1.24027) (1.40480,1.22906) (1.40946,1.21748) (1.41407,1.20553) (1.41864,1.19321) (1.42316,1.18052) (1.42764,1.16748) (1.43206,1.15408) (1.43643,1.14033) (1.44074,1.12623) (1.44501,1.11179) (1.44921,1.09701) (1.45336,1.08190) (1.45745,1.06645) (1.46149,1.05069) (1.46546,1.03460) (1.46937,1.01820) (1.47321,1.00148) (1.47700,0.98446) (1.48071,0.96715) (1.48436,0.94953) (1.48794,0.93163) (1.49146,0.91344) (1.49490,0.89498) (1.49828,0.87624) (1.50158,0.85724) (1.50481,0.83798) (1.50796,0.81846) (1.51104,0.79869) (1.51404,0.77867) (1.51697,0.75842) (1.51982,0.73794) (1.52259,0.71724) (1.52529,0.69631) (1.52790,0.67518) (1.53043,0.65383) (1.53288,0.63229) (1.53525,0.61056) (1.53753,0.58864) (1.53973,0.56654) (1.54184,0.54427) (1.54387,0.52183) (1.54582,0.49924) (1.54768,0.47649) (1.54945,0.45359) (1.55113,0.43056) (1.55273,0.40740) (1.55424,0.38411) (1.55565,0.36071) (1.55698,0.33720) (1.55822,0.31358) (1.55937,0.28987) (1.56043,0.26607) (1.56140,0.24219) (1.56228,0.21823) (1.56306,0.19421) (1.56376,0.17013) (1.56436,0.14600) (1.56487,0.12182) (1.56529,0.09760) (1.56561,0.07336) (1.56584,0.04909) (1.56599,0.02481) (1.56603,0.00052) (1.56599,-0.02377) (1.56585,-0.04805) (1.56562,-0.07231) (1.56530,-0.09656) (1.56489,-0.12077) (1.56438,-0.14495) (1.56378,-0.16909) (1.56309,-0.19317) (1.56231,-0.21720) (1.56144,-0.24115) (1.56047,-0.26504) (1.55942,-0.28884) (1.55827,-0.31256) (1.55704,-0.33618) (1.55571,-0.35970) (1.55430,-0.38311) (1.55280,-0.40640) (1.55120,-0.42957) (1.54952,-0.45260) (1.54776,-0.47550) (1.54590,-0.49826) (1.54396,-0.52086) (1.54193,-0.54331) (1.53982,-0.56558) (1.53763,-0.58769) (1.53535,-0.60962) (1.53298,-0.63136) (1.53054,-0.65291) (1.52801,-0.67426) (1.52540,-0.69541) (1.52271,-0.71634) (1.51994,-0.73705) (1.51710,-0.75754) (1.51417,-0.77780) (1.51117,-0.79783) (1.50810,-0.81761) (1.50494,-0.83714) (1.50172,-0.85641) (1.49842,-0.87543) (1.49505,-0.89418) (1.49161,-0.91265) (1.48810,-0.93085) (1.48452,-0.94877) (1.48087,-0.96639) (1.47716,-0.98372) (1.47338,-1.00075) (1.46953,-1.01748) (1.46563,-1.03390) (1.46166,-1.05000) (1.45763,-1.06578) (1.45354,-1.08124) (1.44939,-1.09637) (1.44519,-1.11116) (1.44093,-1.12562) (1.43661,-1.13973) (1.43225,-1.15349) (1.42783,-1.16691) (1.42336,-1.17997) (1.41884,-1.19267) (1.41427,-1.20500) (1.40966,-1.21697) (1.40500,-1.22857) (1.40030,-1.23979) (1.39556,-1.25064) (1.39077,-1.26110) (1.38595,-1.27118) (1.38109,-1.28088) (1.37620,-1.29018) (1.37126,-1.29909) (1.36630,-1.30761) (1.36130,-1.31573) (1.35628,-1.32344) (1.35122,-1.33075) (1.34614,-1.33766) (1.34103,-1.34416) (1.33590,-1.35025) (1.33075,-1.35593) (1.32557,-1.36120) (1.32038,-1.36605) (1.31517,-1.37048) (1.30994,-1.37450) (1.30470,-1.37810) (1.29944,-1.38128) (1.29417,-1.38404) (1.28890,-1.38638) (1.28361,-1.38829) (1.27832,-1.38978) (1.27303,-1.39085) (1.26773,-1.39150) (1.26243,-1.39172) (1.25713,-1.39152) (1.25183,-1.39089) (1.24653,-1.38984) (1.24124,-1.38836) (1.23596,-1.38647) (1.23068,-1.38415) (1.22541,-1.38141) (1.22016,-1.37825) (1.21491,-1.37467) (1.20969,-1.37067) (1.20447,-1.36625) (1.19928,-1.36141) (1.19410,-1.35617) (1.18895,-1.35051) (1.18382,-1.34443) (1.17871,-1.33795) (1.17362,-1.33106) (1.16857,-1.32376) (1.16354,-1.31607) (1.15854,-1.30797) (1.15358,-1.29947) (1.14864,-1.29058) (1.14375,-1.28129) (1.13888,-1.27161) (1.13406,-1.26155) (1.12927,-1.25110) (1.12453,-1.24027) (1.11983,-1.22906) (1.11517,-1.21748) (1.11055,-1.20553) (1.10598,-1.19321) (1.10146,-1.18052) (1.09699,-1.16748) (1.09257,-1.15408) (1.08820,-1.14033) (1.08388,-1.12623) (1.07962,-1.11179) (1.07542,-1.09701) (1.07127,-1.08190) (1.06717,-1.06645) (1.06314,-1.05069) (1.05917,-1.03460) (1.05526,-1.01820) (1.05141,-1.00148) (1.04763,-0.98446) (1.04392,-0.96715) (1.04027,-0.94953) (1.03668,-0.93163) (1.03317,-0.91344) (1.02972,-0.89498) (1.02635,-0.87624) (1.02305,-0.85724) (1.01982,-0.83798) (1.01667,-0.81846) (1.01359,-0.79869) (1.01058,-0.77867) (1.00766,-0.75842) (1.00481,-0.73794) (1.00203,-0.71724) (0.99934,-0.69631) (0.99673,-0.67518) (0.99420,-0.65383) (0.99175,-0.63229) (0.98938,-0.61056) (0.98710,-0.58864) (0.98490,-0.56654) (0.98278,-0.54427) (0.98075,-0.52183) (0.97881,-0.49924) (0.97695,-0.47649) (0.97518,-0.45359) (0.97349,-0.43056) (0.97190,-0.40740) (0.97039,-0.38411) (0.96897,-0.36071) (0.96764,-0.33720) (0.96640,-0.31358) (0.96525,-0.28987) (0.96420,-0.26607) (0.96323,-0.24219) (0.96235,-0.21823) (0.96157,-0.19421) (0.96087,-0.17013) (0.96027,-0.14600) (0.95976,-0.12182) (0.95934,-0.09760) (0.95902,-0.07336) (0.95878,-0.04909) (0.95864,-0.02481) (0.95859,-0.00052) (0.95864,0.02377) (0.95878,0.04805) (0.95900,0.07231) (0.95933,0.09656) (0.95974,0.12077) (0.96025,0.14495) (0.96084,0.16909) (0.96153,0.19317) (0.96232,0.21720) (0.96319,0.24115) (0.96415,0.26504) (0.96521,0.28884) (0.96635,0.31256) (0.96759,0.33618) (0.96891,0.35970) (0.97033,0.38311) (0.97183,0.40640) (0.97342,0.42957) (0.97510,0.45260) (0.97687,0.47550) (0.97873,0.49826) (0.98067,0.52086) (0.98269,0.54331) (0.98481,0.56558) (0.98700,0.58769) (0.98928,0.60962) (0.99165,0.63136) (0.99409,0.65291) (0.99662,0.67426) (0.99923,0.69541) (1.00192,0.71634) (1.00468,0.73705) (1.00753,0.75754) (1.01045,0.77780) (1.01346,0.79783) (1.01653,0.81761) (1.01968,0.83714) (1.02291,0.85641) (1.02621,0.87543) (1.02958,0.89418) (1.03302,0.91265) (1.03653,0.93085) (1.04011,0.94877) (1.04376,0.96639) (1.04747,0.98372) (1.05125,1.00075) (1.05509,1.01748) (1.05900,1.03390) (1.06297,1.05000) (1.06700,1.06578) (1.07109,1.08124) (1.07524,1.09637) (1.07944,1.11116) (1.08370,1.12562) (1.08801,1.13973) (1.09238,1.15349) (1.09680,1.16691) (1.10127,1.17997) (1.10579,1.19267) (1.11035,1.20500) (1.11497,1.21697) (1.11963,1.22857) (1.12433,1.23979) (1.12907,1.25064) (1.13385,1.26110) (1.13868,1.27118) (1.14354,1.28088) (1.14843,1.29018) (1.15336,1.29909) (1.15833,1.30761) (1.16332,1.31573) (1.16835,1.32344) (1.17341,1.33075) (1.17849,1.33766) (1.18359,1.34416) (1.18873,1.35025) (1.19388,1.35593) (1.19906,1.36120) (1.20425,1.36605) (1.20946,1.37048) (1.21469,1.37450) (1.21993,1.37810) (1.22519,1.38128) (1.23045,1.38404) (1.23573,1.38638) (1.24101,1.38829) (1.24630,1.38978) (1.25160,1.39085) (1.25690,1.39150) (1.26220,1.39172) (1.26750,1.39152) (1.27280,1.39089) (1.27810,1.38984)};
\node[courseaccent,inner sep=0pt] at (1.26231,0.00000) {$\times$};
\draw[->,courseaccent,line width=.7pt] plot coordinates {(0.00000,0.00000) (1.32544,1.36133)};
\filldraw[fill=courseink,draw=coursepaper,line width=.4pt] (1.32544,1.36133) circle[radius=2.3pt];
\draw[->,courseaccent,line width=.7pt] plot coordinates {(0.00000,0.00000) (1.32544,-1.36133)};
\filldraw[fill=courseink,draw=coursepaper,line width=.4pt] (1.32544,-1.36133) circle[radius=2.3pt];
\node[below left,inner sep=2pt] at (0,0) {$0$};
\node[align=center,text width=3.45cm,anchor=north] at (0,-3.15) {б) Невырожденное\\(NDC)\\Смещённый эллипс};
\end{scope}
\begin{scope}[shift={(12.10,0)}]
\fill[coursetint!45!white] (0,0) circle[radius=1.9cm];
\filldraw[fill=courseaccent,fill opacity=.09,draw=courseaccent!30] (2.06598,1.20229)--(-0.52317,2.50634)--(-2.06598,-1.20229)--(0.52317,-2.50634)--cycle;
\draw[coursemuted!40,thin,densely dashed] plot coordinates {(-1.45548,-0.51614) (-1.41198,-0.53729) (-1.36675,-0.55779) (-1.31985,-0.57761) (-1.27135,-0.59673) (-1.22130,-0.61511) (-1.16976,-0.63275) (-1.11679,-0.64962) (-1.06247,-0.66570) (-1.00685,-0.68096) (-0.95000,-0.69540) (-0.89200,-0.70898) (-0.83291,-0.72171) (-0.77280,-0.73355) (-0.71175,-0.74451) (-0.64984,-0.75455) (-0.58713,-0.76367) (-0.52371,-0.77187) (-0.45965,-0.77912) (-0.39503,-0.78543) (-0.32993,-0.79078) (-0.26443,-0.79516) (-0.19860,-0.79858) (-0.13254,-0.80102) (-0.06631,-0.80249) (-0.00000,-0.80297) (0.06631,-0.80249) (0.13254,-0.80102) (0.19860,-0.79858) (0.26443,-0.79516) (0.32993,-0.79078) (0.39503,-0.78543) (0.45965,-0.77912) (0.52371,-0.77187) (0.58713,-0.76367) (0.64984,-0.75455) (0.71175,-0.74451) (0.77280,-0.73355) (0.83291,-0.72171) (0.89200,-0.70898) (0.95000,-0.69540) (1.00685,-0.68096) (1.06247,-0.66570) (1.11679,-0.64962) (1.16976,-0.63275) (1.22130,-0.61511) (1.27135,-0.59673) (1.31985,-0.57761) (1.36675,-0.55779) (1.41198,-0.53729) (1.45548,-0.51614) (1.49722,-0.49436) (1.53713,-0.47198) (1.57517,-0.44902) (1.61129,-0.42551) (1.64545,-0.40149) (1.67760,-0.37697) (1.70771,-0.35200) (1.73574,-0.32660) (1.76165,-0.30080) (1.78542,-0.27463) (1.80701,-0.24813) (1.82640,-0.22133) (1.84356,-0.19426) (1.85848,-0.16695) (1.87113,-0.13944) (1.88151,-0.11175) (1.88959,-0.08393) (1.89537,-0.05601) (1.89884,-0.02802) (1.90000,-0.00000) (1.89884,0.02802) (1.89537,0.05601) (1.88959,0.08393) (1.88151,0.11175) (1.87113,0.13944) (1.85848,0.16695) (1.84356,0.19426) (1.82640,0.22133) (1.80701,0.24813) (1.78542,0.27463) (1.76165,0.30080) (1.73574,0.32660) (1.70771,0.35200) (1.67760,0.37697) (1.64545,0.40149) (1.61129,0.42551) (1.57517,0.44902) (1.53713,0.47198) (1.49722,0.49436) (1.45548,0.51614) (1.41198,0.53729) (1.36675,0.55779) (1.31985,0.57761) (1.27135,0.59673) (1.22130,0.61511) (1.16976,0.63275) (1.11679,0.64962) (1.06247,0.66570) (1.00685,0.68096) (0.95000,0.69540) (0.89200,0.70898) (0.83291,0.72171) (0.77280,0.73355) (0.71175,0.74451) (0.64984,0.75455) (0.58713,0.76367) (0.52371,0.77187) (0.45965,0.77912) (0.39503,0.78543) (0.32993,0.79078) (0.26443,0.79516) (0.19860,0.79858) (0.13254,0.80102) (0.06631,0.80249) (0.00000,0.80297) (-0.06631,0.80249) (-0.13254,0.80102) (-0.19860,0.79858) (-0.26443,0.79516) (-0.32993,0.79078) (-0.39503,0.78543) (-0.45965,0.77912) (-0.52371,0.77187) (-0.58713,0.76367) (-0.64984,0.75455) (-0.71175,0.74451) (-0.77280,0.73355) (-0.83291,0.72171) (-0.89200,0.70898) (-0.95000,0.69540) (-1.00685,0.68096) (-1.06247,0.66570) (-1.11679,0.64962) (-1.16976,0.63275) (-1.22130,0.61511) (-1.27135,0.59673) (-1.31985,0.57761) (-1.36675,0.55779) (-1.41198,0.53729) (-1.45548,0.51614) (-1.49722,0.49436) (-1.53713,0.47198) (-1.57517,0.44902) (-1.61129,0.42551) (-1.64545,0.40149) (-1.67760,0.37697) (-1.70771,0.35200) (-1.73574,0.32660) (-1.76165,0.30080) (-1.78542,0.27463) (-1.80701,0.24813) (-1.82640,0.22133) (-1.84356,0.19426) (-1.85848,0.16695) (-1.87113,0.13944) (-1.88151,0.11175) (-1.88959,0.08393) (-1.89537,0.05601) (-1.89884,0.02802) (-1.90000,0.00000) (-1.89884,-0.02802) (-1.89537,-0.05601) (-1.88959,-0.08393) (-1.88151,-0.11175) (-1.87113,-0.13944) (-1.85848,-0.16695) (-1.84356,-0.19426) (-1.82640,-0.22133) (-1.80701,-0.24813) (-1.78542,-0.27463) (-1.76165,-0.30080) (-1.73574,-0.32660) (-1.70771,-0.35200) (-1.67760,-0.37697) (-1.64545,-0.40149) (-1.61129,-0.42551) (-1.57517,-0.44902) (-1.53713,-0.47198) (-1.49722,-0.49436) (-1.45548,-0.51614)};
\draw[coursemuted!55,thin] plot coordinates {(-1.90000,-0.00000) (-1.89884,-0.02802) (-1.89537,-0.05601) (-1.88959,-0.08393) (-1.88151,-0.11175) (-1.87113,-0.13944) (-1.85848,-0.16695) (-1.84356,-0.19426) (-1.82640,-0.22133) (-1.80701,-0.24813) (-1.78542,-0.27463) (-1.76165,-0.30080) (-1.73574,-0.32660) (-1.70771,-0.35200) (-1.67760,-0.37697) (-1.64545,-0.40149) (-1.61129,-0.42551) (-1.57517,-0.44902) (-1.53713,-0.47198) (-1.49722,-0.49436) (-1.45548,-0.51614) (-1.41198,-0.53729) (-1.36675,-0.55779) (-1.31985,-0.57761) (-1.27135,-0.59673) (-1.22130,-0.61511) (-1.16976,-0.63275) (-1.11679,-0.64962) (-1.06247,-0.66570) (-1.00685,-0.68096) (-0.95000,-0.69540) (-0.89200,-0.70898) (-0.83291,-0.72171) (-0.77280,-0.73355) (-0.71175,-0.74451) (-0.64984,-0.75455) (-0.58713,-0.76367) (-0.52371,-0.77187) (-0.45965,-0.77912) (-0.39503,-0.78543) (-0.32993,-0.79078) (-0.26443,-0.79516) (-0.19860,-0.79858) (-0.13254,-0.80102) (-0.06631,-0.80249) (-0.00000,-0.80297) (0.06631,-0.80249) (0.13254,-0.80102) (0.19860,-0.79858) (0.26443,-0.79516) (0.32993,-0.79078) (0.39503,-0.78543) (0.45965,-0.77912) (0.52371,-0.77187) (0.58713,-0.76367) (0.64984,-0.75455) (0.71175,-0.74451) (0.77280,-0.73355) (0.83291,-0.72171) (0.89200,-0.70898) (0.95000,-0.69540) (1.00685,-0.68096) (1.06247,-0.66570) (1.11679,-0.64962) (1.16976,-0.63275) (1.22130,-0.61511) (1.27135,-0.59673) (1.31985,-0.57761) (1.36675,-0.55779) (1.41198,-0.53729) (1.45548,-0.51614) (1.49722,-0.49436) (1.53713,-0.47198) (1.57517,-0.44902) (1.61129,-0.42551) (1.64545,-0.40149) (1.67760,-0.37697) (1.70771,-0.35200) (1.73574,-0.32660) (1.76165,-0.30080) (1.78542,-0.27463) (1.80701,-0.24813) (1.82640,-0.22133) (1.84356,-0.19426) (1.85848,-0.16695) (1.87113,-0.13944) (1.88151,-0.11175) (1.88959,-0.08393) (1.89537,-0.05601) (1.89884,-0.02802) (1.90000,-0.00000)};
\draw[coursemuted!40,thin,densely dashed] plot coordinates {(-1.45548,-0.51614) (-1.45460,-0.45573) (-1.45194,-0.39477) (-1.44751,-0.33332) (-1.44132,-0.27147) (-1.43337,-0.20928) (-1.42368,-0.14684) (-1.41225,-0.08422) (-1.39910,-0.02150) (-1.38425,0.04124) (-1.36771,0.10394) (-1.34950,0.16651) (-1.32965,0.22887) (-1.30818,0.29096) (-1.28512,0.35270) (-1.26049,0.41400) (-1.23432,0.47480) (-1.20665,0.53502) (-1.17751,0.59459) (-1.14694,0.65343) (-1.11497,0.71148) (-1.08164,0.76866) (-1.04699,0.82491) (-1.01106,0.88015) (-0.97391,0.93432) (-0.93557,0.98735) (-0.89609,1.03917) (-0.85551,1.08973) (-0.81390,1.13897) (-0.77129,1.18681) (-0.72774,1.23321) (-0.68331,1.27811) (-0.63804,1.32145) (-0.59200,1.36318) (-0.54523,1.40325) (-0.49780,1.44161) (-0.44977,1.47821) (-0.40119,1.51301) (-0.35211,1.54597) (-0.30261,1.57704) (-0.25274,1.60620) (-0.20256,1.63339) (-0.15214,1.65860) (-0.10153,1.68179) (-0.05080,1.70292) (-0.00000,1.72198) (0.05080,1.73895) (0.10153,1.75379) (0.15214,1.76650) (0.20256,1.77706) (0.25274,1.78545) (0.30261,1.79167) (0.35211,1.79570) (0.40119,1.79755) (0.44977,1.79720) (0.49780,1.79467) (0.54523,1.78995) (0.59200,1.78305) (0.63804,1.77397) (0.68331,1.76274) (0.72774,1.74935) (0.77129,1.73384) (0.81390,1.71621) (0.85551,1.69650) (0.89609,1.67471) (0.93557,1.65089) (0.97391,1.62505) (1.01106,1.59723) (1.04699,1.56747) (1.08164,1.53580) (1.11497,1.50226) (1.14694,1.46689) (1.17751,1.42973) (1.20665,1.39082) (1.23432,1.35023) (1.26049,1.30798) (1.28512,1.26415) (1.30818,1.21877) (1.32965,1.17191) (1.34950,1.12363) (1.36771,1.07397) (1.38425,1.02300) (1.39910,0.97079) (1.41225,0.91740) (1.42368,0.86288) (1.43337,0.80732) (1.44132,0.75077) (1.44751,0.69331) (1.45194,0.63500) (1.45460,0.57592) (1.45548,0.51614) (1.45460,0.45573) (1.45194,0.39477) (1.44751,0.33332) (1.44132,0.27147) (1.43337,0.20928) (1.42368,0.14684) (1.41225,0.08422) (1.39910,0.02150) (1.38425,-0.04124) (1.36771,-0.10394) (1.34950,-0.16651) (1.32965,-0.22887) (1.30818,-0.29096) (1.28512,-0.35270) (1.26049,-0.41400) (1.23432,-0.47480) (1.20665,-0.53502) (1.17751,-0.59459) (1.14694,-0.65343) (1.11497,-0.71148) (1.08164,-0.76866) (1.04699,-0.82491) (1.01106,-0.88015) (0.97391,-0.93432) (0.93557,-0.98735) (0.89609,-1.03917) (0.85551,-1.08973) (0.81390,-1.13897) (0.77129,-1.18681) (0.72774,-1.23321) (0.68331,-1.27811) (0.63804,-1.32145) (0.59200,-1.36318) (0.54523,-1.40325) (0.49780,-1.44161) (0.44977,-1.47821) (0.40119,-1.51301) (0.35211,-1.54597) (0.30261,-1.57704) (0.25274,-1.60620) (0.20256,-1.63339) (0.15214,-1.65860) (0.10153,-1.68179) (0.05080,-1.70292) (0.00000,-1.72198) (-0.05080,-1.73895) (-0.10153,-1.75379) (-0.15214,-1.76650) (-0.20256,-1.77706) (-0.25274,-1.78545) (-0.30261,-1.79167) (-0.35211,-1.79570) (-0.40119,-1.79755) (-0.44977,-1.79720) (-0.49780,-1.79467) (-0.54523,-1.78995) (-0.59200,-1.78305) (-0.63804,-1.77397) (-0.68331,-1.76274) (-0.72774,-1.74935) (-0.77129,-1.73384) (-0.81390,-1.71621) (-0.85551,-1.69650) (-0.89609,-1.67471) (-0.93557,-1.65089) (-0.97391,-1.62505) (-1.01106,-1.59723) (-1.04699,-1.56747) (-1.08164,-1.53580) (-1.11497,-1.50226) (-1.14694,-1.46689) (-1.17751,-1.42973) (-1.20665,-1.39082) (-1.23432,-1.35023) (-1.26049,-1.30798) (-1.28512,-1.26415) (-1.30818,-1.21877) (-1.32965,-1.17191) (-1.34950,-1.12363) (-1.36771,-1.07397) (-1.38425,-1.02300) (-1.39910,-0.97079) (-1.41225,-0.91740) (-1.42368,-0.86288) (-1.43337,-0.80732) (-1.44132,-0.75077) (-1.44751,-0.69331) (-1.45194,-0.63500) (-1.45460,-0.57592) (-1.45548,-0.51614)};
\draw[coursemuted!55,thin] plot coordinates {(-0.85467,-1.69692) (-0.89526,-1.67518) (-0.93477,-1.65140) (-0.97313,-1.62560) (-1.01031,-1.59783) (-1.04626,-1.56810) (-1.08094,-1.53647) (-1.11429,-1.50297) (-1.14629,-1.46763) (-1.17690,-1.43051) (-1.20607,-1.39164) (-1.23377,-1.35108) (-1.25996,-1.30887) (-1.28463,-1.26507) (-1.30772,-1.21972) (-1.32923,-1.17289) (-1.34911,-1.12463) (-1.36735,-1.07500) (-1.38392,-1.02406) (-1.39881,-0.97188) (-1.41200,-0.91851) (-1.42346,-0.86402) (-1.43319,-0.80847) (-1.44117,-0.75195) (-1.44740,-0.69450) (-1.45187,-0.63621) (-1.45456,-0.57715) (-1.45548,-0.51738) (-1.45463,-0.45698) (-1.45201,-0.39602) (-1.44762,-0.33459) (-1.44146,-0.27274) (-1.43355,-0.21056) (-1.42390,-0.14813) (-1.41250,-0.08551) (-1.39939,-0.02280) (-1.38457,0.03995) (-1.36807,0.10265) (-1.34989,0.16522) (-1.33008,0.22759) (-1.30864,0.28969) (-1.28561,0.35143) (-1.26101,0.41274) (-1.23487,0.47355) (-1.20724,0.53379) (-1.17813,0.59337) (-1.14758,0.65223) (-1.11564,0.71030) (-1.08233,0.76750) (-1.04771,0.82376) (-1.01182,0.87902) (-0.97469,0.93321) (-0.93637,0.98627) (-0.89691,1.03812) (-0.85636,1.08871) (-0.81476,1.13797) (-0.77218,1.18584) (-0.72865,1.23227) (-0.68423,1.27720) (-0.63898,1.32057) (-0.59295,1.36234) (-0.54620,1.40244) (-0.49879,1.44083) (-0.45076,1.47747) (-0.40219,1.51231) (-0.35313,1.54531) (-0.30363,1.57642) (-0.25377,1.60562) (-0.20360,1.63285) (-0.15318,1.65810) (-0.10257,1.68133) (-0.05184,1.70251) (-0.00105,1.72161) (0.04975,1.73862) (0.10049,1.75351) (0.15110,1.76626) (0.20153,1.77686) (0.25171,1.78530) (0.30159,1.79156) (0.35110,1.79564) (0.40018,1.79753) (0.44878,1.79723) (0.49682,1.79474) (0.54426,1.79007) (0.59104,1.78321) (0.63710,1.77418) (0.68239,1.76299) (0.72684,1.74965) (0.77040,1.73418) (0.81303,1.71660) (0.85467,1.69692)};
\draw[coursemuted!40,thin,densely dashed] plot coordinates {(1.22130,-0.61511) (1.22055,-0.55464) (1.21832,-0.49350) (1.21461,-0.43175) (1.20941,-0.36947) (1.20274,-0.30675) (1.19461,-0.24365) (1.18502,-0.18026) (1.17399,-0.11664) (1.16152,-0.05289) (1.14764,0.01094) (1.13237,0.07474) (1.11571,0.13846) (1.09769,0.20201) (1.07834,0.26531) (1.05767,0.32829) (1.03572,0.39087) (1.01250,0.45297) (0.98805,0.51452) (0.96239,0.57544) (0.93557,0.63567) (0.90760,0.69511) (0.87853,0.75371) (0.84838,0.81140) (0.81721,0.86809) (0.78503,0.92373) (0.75191,0.97824) (0.71786,1.03156) (0.68294,1.08362) (0.64719,1.13437) (0.61065,1.18373) (0.57336,1.23164) (0.53538,1.27806) (0.49675,1.32292) (0.45751,1.36617) (0.41771,1.40775) (0.37740,1.44762) (0.33663,1.48573) (0.29546,1.52202) (0.25392,1.55647) (0.21208,1.58901) (0.16997,1.61962) (0.12766,1.64825) (0.08519,1.67488) (0.04262,1.69947) (0.00000,1.72198) (-0.04262,1.74240) (-0.08519,1.76070) (-0.12766,1.77685) (-0.16997,1.79083) (-0.21208,1.80264) (-0.25392,1.81224) (-0.29546,1.81964) (-0.33663,1.82483) (-0.37740,1.82779) (-0.41771,1.82852) (-0.45751,1.82702) (-0.49675,1.82330) (-0.53538,1.81736) (-0.57336,1.80920) (-0.61065,1.79884) (-0.64719,1.78629) (-0.68294,1.77156) (-0.71786,1.75467) (-0.75191,1.73564) (-0.78503,1.71450) (-0.81721,1.69128) (-0.84838,1.66599) (-0.87853,1.63867) (-0.90760,1.60935) (-0.93557,1.57808) (-0.96239,1.54488) (-0.98805,1.50980) (-1.01250,1.47287) (-1.03572,1.43416) (-1.05767,1.39370) (-1.07834,1.35154) (-1.09769,1.30773) (-1.11571,1.26233) (-1.13237,1.21539) (-1.14764,1.16697) (-1.16152,1.11713) (-1.17399,1.06593) (-1.18502,1.01343) (-1.19461,0.95969) (-1.20274,0.90479) (-1.20941,0.84878) (-1.21461,0.79174) (-1.21832,0.73374) (-1.22055,0.67484) (-1.22130,0.61511) (-1.22055,0.55464) (-1.21832,0.49350) (-1.21461,0.43175) (-1.20941,0.36947) (-1.20274,0.30675) (-1.19461,0.24365) (-1.18502,0.18026) (-1.17399,0.11664) (-1.16152,0.05289) (-1.14764,-0.01094) (-1.13237,-0.07474) (-1.11571,-0.13846) (-1.09769,-0.20201) (-1.07834,-0.26531) (-1.05767,-0.32829) (-1.03572,-0.39087) (-1.01250,-0.45297) (-0.98805,-0.51452) (-0.96239,-0.57544) (-0.93557,-0.63567) (-0.90760,-0.69511) (-0.87853,-0.75371) (-0.84838,-0.81140) (-0.81721,-0.86809) (-0.78503,-0.92373) (-0.75191,-0.97824) (-0.71786,-1.03156) (-0.68294,-1.08362) (-0.64719,-1.13437) (-0.61065,-1.18373) (-0.57336,-1.23164) (-0.53538,-1.27806) (-0.49675,-1.32292) (-0.45751,-1.36617) (-0.41771,-1.40775) (-0.37740,-1.44762) (-0.33663,-1.48573) (-0.29546,-1.52202) (-0.25392,-1.55647) (-0.21208,-1.58901) (-0.16997,-1.61962) (-0.12766,-1.64825) (-0.08519,-1.67488) (-0.04262,-1.69947) (-0.00000,-1.72198) (0.04262,-1.74240) (0.08519,-1.76070) (0.12766,-1.77685) (0.16997,-1.79083) (0.21208,-1.80264) (0.25392,-1.81224) (0.29546,-1.81964) (0.33663,-1.82483) (0.37740,-1.82779) (0.41771,-1.82852) (0.45751,-1.82702) (0.49675,-1.82330) (0.53538,-1.81736) (0.57336,-1.80920) (0.61065,-1.79884) (0.64719,-1.78629) (0.68294,-1.77156) (0.71786,-1.75467) (0.75191,-1.73564) (0.78503,-1.71450) (0.81721,-1.69128) (0.84838,-1.66599) (0.87853,-1.63867) (0.90760,-1.60935) (0.93557,-1.57808) (0.96239,-1.54488) (0.98805,-1.50980) (1.01250,-1.47287) (1.03572,-1.43416) (1.05767,-1.39370) (1.07834,-1.35154) (1.09769,-1.30773) (1.11571,-1.26233) (1.13237,-1.21539) (1.14764,-1.16697) (1.16152,-1.11713) (1.17399,-1.06593) (1.18502,-1.01343) (1.19461,-0.95969) (1.20274,-0.90479) (1.20941,-0.84878) (1.21461,-0.79174) (1.21832,-0.73374) (1.22055,-0.67484) (1.22130,-0.61511)};
\draw[coursemuted!55,thin] plot coordinates {(0.63503,-1.79074) (0.67105,-1.77674) (0.70625,-1.76057) (0.74060,-1.74226) (0.77404,-1.72182) (0.80653,-1.69929) (0.83805,-1.67469) (0.86854,-1.64805) (0.89798,-1.61940) (0.92632,-1.58877) (0.95353,-1.55621) (0.97958,-1.52176) (1.00444,-1.48545) (1.02808,-1.44733) (1.05046,-1.40745) (1.07156,-1.36585) (1.09136,-1.32259) (1.10982,-1.27772) (1.12694,-1.23129) (1.14268,-1.18336) (1.15703,-1.13399) (1.16997,-1.08324) (1.18148,-1.03116) (1.19156,-0.97783) (1.20018,-0.92331) (1.20734,-0.86767) (1.21303,-0.81097) (1.21724,-0.75328) (1.21997,-0.69467) (1.22121,-0.63522) (1.22097,-0.57499) (1.21923,-0.51406) (1.21602,-0.45250) (1.21132,-0.39040) (1.20514,-0.32782) (1.19750,-0.26483) (1.18839,-0.20153) (1.17784,-0.13798) (1.16586,-0.07426) (1.15245,-0.01045) (1.13764,0.05337) (1.12144,0.11712) (1.10388,0.18074) (1.08497,0.24413) (1.06474,0.30722) (1.04322,0.36994) (1.02042,0.43222) (0.99638,0.49396) (0.97112,0.55510) (0.94469,0.61557) (0.91710,0.67528) (0.88839,0.73418) (0.85860,0.79217) (0.82777,0.84921) (0.79592,0.90521) (0.76311,0.96010) (0.72937,1.01383) (0.69474,1.06632) (0.65926,1.11751) (0.62298,1.16734) (0.58594,1.21575) (0.54818,1.26268) (0.50976,1.30807) (0.47072,1.35186) (0.43110,1.39401) (0.39096,1.43446) (0.35034,1.47316) (0.30930,1.51007) (0.26788,1.54513) (0.22613,1.57832) (0.18411,1.60958) (0.14186,1.63888) (0.09944,1.66618) (0.05690,1.69146) (0.01429,1.71467) (-0.02834,1.73580) (-0.07094,1.75481) (-0.11344,1.77168) (-0.15581,1.78639) (-0.19799,1.79893) (-0.23993,1.80927) (-0.28158,1.81741) (-0.32288,1.82334) (-0.36379,1.82704) (-0.40425,1.82852) (-0.44423,1.82777) (-0.48366,1.82480) (-0.52250,1.81960) (-0.56071,1.81218) (-0.59823,1.80256) (-0.63503,1.79074)};
\draw[coursemuted!50,thin,densely dashed] plot coordinates {(1.67381,0.59356) (0.00000,0.00000)};
\draw[->,courseink!75,thin] plot coordinates {(0.00000,0.00000) (-1.96490,-0.69679)};
\node[above left,inner sep=2pt] at (-1.96490,-0.69679) {$s_r$};
\draw[coursemuted!50,thin,densely dashed] plot coordinates {(-1.40449,0.70738) (0.00000,0.00000)};
\draw[->,courseink!75,thin] plot coordinates {(0.00000,0.00000) (1.64875,-0.83040)};
\node[below right,inner sep=2pt] at (1.64875,-0.83040) {$s_t$};
\draw[coursemuted!50,thin,densely dashed] plot coordinates {(0.00000,-1.98028) (0.00000,0.00000)};
\draw[->,courseink!75,thin] plot coordinates {(0.00000,0.00000) (0.00000,2.32468)};
\node[above,inner sep=2pt] at (0.00000,2.32468) {$s_z$};
\draw[courseaccent,line width=1.2pt] plot coordinates {(0.72774,1.74935) (0.74895,1.73835) (0.76992,1.72682) (0.79066,1.71476) (0.81116,1.70218) (0.83142,1.68909) (0.85142,1.67547) (0.87116,1.66135) (0.89063,1.64672) (0.90984,1.63159) (0.92876,1.61596) (0.94741,1.59984) (0.96576,1.58324) (0.98382,1.56615) (1.00158,1.54858) (1.01904,1.53054) (1.03619,1.51204) (1.05302,1.49307) (1.06953,1.47365) (1.08571,1.45378) (1.10156,1.43347) (1.11708,1.41273) (1.13226,1.39155) (1.14709,1.36994) (1.16157,1.34792) (1.17570,1.32549) (1.18947,1.30266) (1.20288,1.27943) (1.21592,1.25581) (1.22859,1.23181) (1.24089,1.20743) (1.25281,1.18268) (1.26435,1.15758) (1.27550,1.13212) (1.28627,1.10631) (1.29664,1.08017) (1.30662,1.05370) (1.31620,1.02691) (1.32537,0.99981) (1.33415,0.97240) (1.34252,0.94469) (1.35048,0.91670) (1.35802,0.88843) (1.36516,0.85989) (1.37188,0.83109) (1.37818,0.80203) (1.38406,0.77273) (1.38952,0.74319) (1.39455,0.71343) (1.39917,0.68345) (1.40335,0.65326) (1.40711,0.62287) (1.41044,0.59229) (1.41334,0.56154) (1.41581,0.53061) (1.41784,0.49952) (1.41945,0.46827) (1.42062,0.43689) (1.42136,0.40537) (1.42167,0.37373) (1.42154,0.34197) (1.42099,0.31011) (1.42000,0.27816) (1.41857,0.24612) (1.41672,0.21401) (1.41443,0.18183) (1.41171,0.14959) (1.40856,0.11731) (1.40498,0.08500) (1.40098,0.05265) (1.39655,0.02030) (1.39169,-0.01207) (1.38641,-0.04443) (1.38070,-0.07678) (1.37458,-0.10910) (1.36804,-0.14139) (1.36108,-0.17364) (1.35370,-0.20583) (1.34591,-0.23796) (1.33772,-0.27002) (1.32911,-0.30200) (1.32010,-0.33388) (1.31069,-0.36567) (1.30088,-0.39734) (1.29068,-0.42889) (1.28008,-0.46031) (1.26909,-0.49159) (1.25771,-0.52272) (1.24595,-0.55369) (1.23381,-0.58449) (1.22130,-0.61511) (1.20841,-0.64555) (1.19515,-0.67579) (1.18154,-0.70583) (1.16756,-0.73564) (1.15322,-0.76524) (1.13854,-0.79460) (1.12350,-0.82372) (1.10813,-0.85259) (1.09242,-0.88120) (1.07637,-0.90954) (1.06000,-0.93761) (1.04330,-0.96538) (1.02629,-0.99287) (1.00896,-1.02005) (0.99133,-1.04692) (0.97339,-1.07347) (0.95516,-1.09970) (0.93664,-1.12559) (0.91783,-1.15114) (0.89874,-1.17633) (0.87938,-1.20117) (0.85975,-1.22564) (0.83986,-1.24974) (0.81971,-1.27346) (0.79931,-1.29679) (0.77867,-1.31973) (0.75780,-1.34226) (0.73669,-1.36439) (0.71535,-1.38609) (0.69380,-1.40738) (0.67204,-1.42824) (0.65007,-1.44866) (0.62791,-1.46864) (0.60555,-1.48818) (0.58301,-1.50726) (0.56029,-1.52588) (0.53740,-1.54404) (0.51435,-1.56173) (0.49114,-1.57894) (0.46778,-1.59567) (0.44428,-1.61191) (0.42065,-1.62767) (0.39688,-1.64292) (0.37299,-1.65768) (0.34900,-1.67193) (0.32489,-1.68567) (0.30069,-1.69890) (0.27639,-1.71161) (0.25201,-1.72381) (0.22755,-1.73547) (0.20302,-1.74661) (0.17844,-1.75721) (0.15379,-1.76728) (0.12910,-1.77681) (0.10438,-1.78580) (0.07961,-1.79425) (0.05483,-1.80215) (0.03003,-1.80950) (0.00522,-1.81630) (-0.01960,-1.82254) (-0.04440,-1.82823) (-0.06919,-1.83337) (-0.09397,-1.83794) (-0.11871,-1.84196) (-0.14342,-1.84541) (-0.16808,-1.84830) (-0.19269,-1.85063) (-0.21725,-1.85240) (-0.24173,-1.85360) (-0.26615,-1.85424) (-0.29048,-1.85431) (-0.31472,-1.85381) (-0.33887,-1.85276) (-0.36292,-1.85114) (-0.38685,-1.84895) (-0.41067,-1.84620) (-0.43436,-1.84289) (-0.45792,-1.83902) (-0.48134,-1.83458) (-0.50461,-1.82959) (-0.52773,-1.82404) (-0.55069,-1.81794) (-0.57348,-1.81128) (-0.59610,-1.80407) (-0.61853,-1.79631) (-0.64078,-1.78800) (-0.66283,-1.77915) (-0.68468,-1.76976) (-0.70632,-1.75982) (-0.72774,-1.74935) (-0.74895,-1.73835) (-0.76992,-1.72682) (-0.79066,-1.71476) (-0.81116,-1.70218) (-0.83142,-1.68909) (-0.85142,-1.67547) (-0.87116,-1.66135) (-0.89063,-1.64672) (-0.90984,-1.63159) (-0.92876,-1.61596) (-0.94741,-1.59984) (-0.96576,-1.58324) (-0.98382,-1.56615) (-1.00158,-1.54858) (-1.01904,-1.53054) (-1.03619,-1.51204) (-1.05302,-1.49307) (-1.06953,-1.47365) (-1.08571,-1.45378) (-1.10156,-1.43347) (-1.11708,-1.41273) (-1.13226,-1.39155) (-1.14709,-1.36994) (-1.16157,-1.34792) (-1.17570,-1.32549) (-1.18947,-1.30266) (-1.20288,-1.27943) (-1.21592,-1.25581) (-1.22859,-1.23181) (-1.24089,-1.20743) (-1.25281,-1.18268) (-1.26435,-1.15758) (-1.27550,-1.13212) (-1.28627,-1.10631) (-1.29664,-1.08017) (-1.30662,-1.05370) (-1.31620,-1.02691) (-1.32537,-0.99981) (-1.33415,-0.97240) (-1.34252,-0.94469) (-1.35048,-0.91670) (-1.35802,-0.88843) (-1.36516,-0.85989) (-1.37188,-0.83109) (-1.37818,-0.80203) (-1.38406,-0.77273) (-1.38952,-0.74319) (-1.39455,-0.71343) (-1.39917,-0.68345) (-1.40335,-0.65326) (-1.40711,-0.62287) (-1.41044,-0.59229) (-1.41334,-0.56154) (-1.41581,-0.53061) (-1.41784,-0.49952) (-1.41945,-0.46827) (-1.42062,-0.43689) (-1.42136,-0.40537) (-1.42167,-0.37373) (-1.42154,-0.34197) (-1.42099,-0.31011) (-1.42000,-0.27816) (-1.41857,-0.24612) (-1.41672,-0.21401) (-1.41443,-0.18183) (-1.41171,-0.14959) (-1.40856,-0.11731) (-1.40498,-0.08500) (-1.40098,-0.05265) (-1.39655,-0.02030) (-1.39169,0.01207) (-1.38641,0.04443) (-1.38070,0.07678) (-1.37458,0.10910) (-1.36804,0.14139) (-1.36108,0.17364) (-1.35370,0.20583) (-1.34591,0.23796) (-1.33772,0.27002) (-1.32911,0.30200) (-1.32010,0.33388) (-1.31069,0.36567) (-1.30088,0.39734) (-1.29068,0.42889) (-1.28008,0.46031) (-1.26909,0.49159) (-1.25771,0.52272) (-1.24595,0.55369) (-1.23381,0.58449) (-1.22130,0.61511) (-1.20841,0.64555) (-1.19515,0.67579) (-1.18154,0.70583) (-1.16756,0.73564) (-1.15322,0.76524) (-1.13854,0.79460) (-1.12350,0.82372) (-1.10813,0.85259) (-1.09242,0.88120) (-1.07637,0.90954) (-1.06000,0.93761) (-1.04330,0.96538) (-1.02629,0.99287) (-1.00896,1.02005) (-0.99133,1.04692) (-0.97339,1.07347) (-0.95516,1.09970) (-0.93664,1.12559) (-0.91783,1.15114) (-0.89874,1.17633) (-0.87938,1.20117) (-0.85975,1.22564) (-0.83986,1.24974) (-0.81971,1.27346) (-0.79931,1.29679) (-0.77867,1.31973) (-0.75780,1.34226) (-0.73669,1.36439) (-0.71535,1.38609) (-0.69380,1.40738) (-0.67204,1.42824) (-0.65007,1.44866) (-0.62791,1.46864) (-0.60555,1.48818) (-0.58301,1.50726) (-0.56029,1.52588) (-0.53740,1.54404) (-0.51435,1.56173) (-0.49114,1.57894) (-0.46778,1.59567) (-0.44428,1.61191) (-0.42065,1.62767) (-0.39688,1.64292) (-0.37299,1.65768) (-0.34900,1.67193) (-0.32489,1.68567) (-0.30069,1.69890) (-0.27639,1.71161) (-0.25201,1.72381) (-0.22755,1.73547) (-0.20302,1.74661) (-0.17844,1.75721) (-0.15379,1.76728) (-0.12910,1.77681) (-0.10438,1.78580) (-0.07961,1.79425) (-0.05483,1.80215) (-0.03003,1.80950) (-0.00522,1.81630) (0.01960,1.82254) (0.04440,1.82823) (0.06919,1.83337) (0.09397,1.83794) (0.11871,1.84196) (0.14342,1.84541) (0.16808,1.84830) (0.19269,1.85063) (0.21725,1.85240) (0.24173,1.85360) (0.26615,1.85424) (0.29048,1.85431) (0.31472,1.85381) (0.33887,1.85276) (0.36292,1.85114) (0.38685,1.84895) (0.41067,1.84620) (0.43436,1.84289) (0.45792,1.83902) (0.48134,1.83458) (0.50461,1.82959) (0.52773,1.82404) (0.55069,1.81794) (0.57348,1.81128) (0.59610,1.80407) (0.61853,1.79631) (0.64078,1.78800) (0.66283,1.77915) (0.68468,1.76976) (0.70632,1.75982) (0.72774,1.74935)};
\draw[coursemuted!65,thin] (0,0) circle[radius=1.9cm];
\node[below left,inner sep=2pt] at (0,0) {$0$};
\node[align=center,text width=3.45cm,anchor=north] at (0,-3.15) {в) Особое (SC)\\Окружность на сфере};
\end{scope}
\begin{scope}[shift={(18.15,0)}]
\filldraw[fill=coursetint!45!white,draw=coursemuted!65,line width=.65pt] (0.00000,0.00000) circle[radius=1.9cm];
\draw[->,courseink!75,thin] plot coordinates {(0.00000,0.00000) (-2.06206,-1.42758)};
\node[below left,inner sep=2pt] at (-2.06206,-1.42758) {$s_r$};
\draw[->,courseink!75,thin] plot coordinates {(0.00000,0.00000) (2.50800,0.00000)};
\node[right,inner sep=2pt] at (2.50800,0.00000) {$s_t$};
\draw[->,courseink!75,thin] plot coordinates {(0.00000,0.00000) (0.00000,2.50800)};
\node[above,inner sep=2pt] at (0.00000,2.50800) {$s_z$};
\draw[courseaccent,line width=1.2pt] plot coordinates {(1.34350,1.34350) (1.34350,1.34330) (1.34350,1.34268) (1.34350,1.34166) (1.34350,1.34023) (1.34350,1.33839) (1.34350,1.33614) (1.34350,1.33349) (1.34350,1.33043) (1.34350,1.32696) (1.34350,1.32309) (1.34350,1.31882) (1.34350,1.31414) (1.34350,1.30907) (1.34350,1.30360) (1.34350,1.29772) (1.34350,1.29146) (1.34350,1.28480) (1.34350,1.27775) (1.34350,1.27031) (1.34350,1.26248) (1.34350,1.25427) (1.34350,1.24567) (1.34350,1.23670) (1.34350,1.22735) (1.34350,1.21763) (1.34350,1.20753) (1.34350,1.19707) (1.34350,1.18624) (1.34350,1.17505) (1.34350,1.16351) (1.34350,1.15161) (1.34350,1.13936) (1.34350,1.12676) (1.34350,1.11381) (1.34350,1.10053) (1.34350,1.08692) (1.34350,1.07297) (1.34350,1.05869) (1.34350,1.04410) (1.34350,1.02918) (1.34350,1.01395) (1.34350,0.99842) (1.34350,0.98258) (1.34350,0.96644) (1.34350,0.95000) (1.34350,0.93328) (1.34350,0.91627) (1.34350,0.89898) (1.34350,0.88142) (1.34350,0.86359) (1.34350,0.84549) (1.34350,0.82714) (1.34350,0.80854) (1.34350,0.78969) (1.34350,0.77060) (1.34350,0.75128) (1.34350,0.73172) (1.34350,0.71195) (1.34350,0.69196) (1.34350,0.67175) (1.34350,0.65134) (1.34350,0.63074) (1.34350,0.60994) (1.34350,0.58895) (1.34350,0.56779) (1.34350,0.54645) (1.34350,0.52495) (1.34350,0.50329) (1.34350,0.48147) (1.34350,0.45951) (1.34350,0.43740) (1.34350,0.41517) (1.34350,0.39280) (1.34350,0.37032) (1.34350,0.34772) (1.34350,0.32502) (1.34350,0.30222) (1.34350,0.27933) (1.34350,0.25635) (1.34350,0.23330) (1.34350,0.21017) (1.34350,0.18698) (1.34350,0.16373) (1.34350,0.14043) (1.34350,0.11709) (1.34350,0.09372) (1.34350,0.07031) (1.34350,0.04689) (1.34350,0.02345) (1.34350,0.00000) (1.34350,-0.02345) (1.34350,-0.04689) (1.34350,-0.07031) (1.34350,-0.09372) (1.34350,-0.11709) (1.34350,-0.14043) (1.34350,-0.16373) (1.34350,-0.18698) (1.34350,-0.21017) (1.34350,-0.23330) (1.34350,-0.25635) (1.34350,-0.27933) (1.34350,-0.30222) (1.34350,-0.32502) (1.34350,-0.34772) (1.34350,-0.37032) (1.34350,-0.39280) (1.34350,-0.41517) (1.34350,-0.43740) (1.34350,-0.45951) (1.34350,-0.48147) (1.34350,-0.50329) (1.34350,-0.52495) (1.34350,-0.54645) (1.34350,-0.56779) (1.34350,-0.58895) (1.34350,-0.60994) (1.34350,-0.63074) (1.34350,-0.65134) (1.34350,-0.67175) (1.34350,-0.69196) (1.34350,-0.71195) (1.34350,-0.73172) (1.34350,-0.75128) (1.34350,-0.77060) (1.34350,-0.78969) (1.34350,-0.80854) (1.34350,-0.82714) (1.34350,-0.84549) (1.34350,-0.86359) (1.34350,-0.88142) (1.34350,-0.89898) (1.34350,-0.91627) (1.34350,-0.93328) (1.34350,-0.95000) (1.34350,-0.96644) (1.34350,-0.98258) (1.34350,-0.99842) (1.34350,-1.01395) (1.34350,-1.02918) (1.34350,-1.04410) (1.34350,-1.05869) (1.34350,-1.07297) (1.34350,-1.08692) (1.34350,-1.10053) (1.34350,-1.11381) (1.34350,-1.12676) (1.34350,-1.13936) (1.34350,-1.15161) (1.34350,-1.16351) (1.34350,-1.17505) (1.34350,-1.18624) (1.34350,-1.19707) (1.34350,-1.20753) (1.34350,-1.21763) (1.34350,-1.22735) (1.34350,-1.23670) (1.34350,-1.24567) (1.34350,-1.25427) (1.34350,-1.26248) (1.34350,-1.27031) (1.34350,-1.27775) (1.34350,-1.28480) (1.34350,-1.29146) (1.34350,-1.29772) (1.34350,-1.30360) (1.34350,-1.30907) (1.34350,-1.31414) (1.34350,-1.31882) (1.34350,-1.32309) (1.34350,-1.32696) (1.34350,-1.33043) (1.34350,-1.33349) (1.34350,-1.33614) (1.34350,-1.33839) (1.34350,-1.34023) (1.34350,-1.34166) (1.34350,-1.34268) (1.34350,-1.34330) (1.34350,-1.34350) (1.34350,-1.34330) (1.34350,-1.34268) (1.34350,-1.34166) (1.34350,-1.34023) (1.34350,-1.33839) (1.34350,-1.33614) (1.34350,-1.33349) (1.34350,-1.33043) (1.34350,-1.32696) (1.34350,-1.32309) (1.34350,-1.31882) (1.34350,-1.31414) (1.34350,-1.30907) (1.34350,-1.30360) (1.34350,-1.29772) (1.34350,-1.29146) (1.34350,-1.28480) (1.34350,-1.27775) (1.34350,-1.27031) (1.34350,-1.26248) (1.34350,-1.25427) (1.34350,-1.24567) (1.34350,-1.23670) (1.34350,-1.22735) (1.34350,-1.21763) (1.34350,-1.20753) (1.34350,-1.19707) (1.34350,-1.18624) (1.34350,-1.17505) (1.34350,-1.16351) (1.34350,-1.15161) (1.34350,-1.13936) (1.34350,-1.12676) (1.34350,-1.11381) (1.34350,-1.10053) (1.34350,-1.08692) (1.34350,-1.07297) (1.34350,-1.05869) (1.34350,-1.04410) (1.34350,-1.02918) (1.34350,-1.01395) (1.34350,-0.99842) (1.34350,-0.98258) (1.34350,-0.96644) (1.34350,-0.95000) (1.34350,-0.93328) (1.34350,-0.91627) (1.34350,-0.89898) (1.34350,-0.88142) (1.34350,-0.86359) (1.34350,-0.84549) (1.34350,-0.82714) (1.34350,-0.80854) (1.34350,-0.78969) (1.34350,-0.77060) (1.34350,-0.75128) (1.34350,-0.73172) (1.34350,-0.71195) (1.34350,-0.69196) (1.34350,-0.67175) (1.34350,-0.65134) (1.34350,-0.63074) (1.34350,-0.60994) (1.34350,-0.58895) (1.34350,-0.56779) (1.34350,-0.54645) (1.34350,-0.52495) (1.34350,-0.50329) (1.34350,-0.48147) (1.34350,-0.45951) (1.34350,-0.43740) (1.34350,-0.41517) (1.34350,-0.39280) (1.34350,-0.37032) (1.34350,-0.34772) (1.34350,-0.32502) (1.34350,-0.30222) (1.34350,-0.27933) (1.34350,-0.25635) (1.34350,-0.23330) (1.34350,-0.21017) (1.34350,-0.18698) (1.34350,-0.16373) (1.34350,-0.14043) (1.34350,-0.11709) (1.34350,-0.09372) (1.34350,-0.07031) (1.34350,-0.04689) (1.34350,-0.02345) (1.34350,-0.00000) (1.34350,0.02345) (1.34350,0.04689) (1.34350,0.07031) (1.34350,0.09372) (1.34350,0.11709) (1.34350,0.14043) (1.34350,0.16373) (1.34350,0.18698) (1.34350,0.21017) (1.34350,0.23330) (1.34350,0.25635) (1.34350,0.27933) (1.34350,0.30222) (1.34350,0.32502) (1.34350,0.34772) (1.34350,0.37032) (1.34350,0.39280) (1.34350,0.41517) (1.34350,0.43740) (1.34350,0.45951) (1.34350,0.48147) (1.34350,0.50329) (1.34350,0.52495) (1.34350,0.54645) (1.34350,0.56779) (1.34350,0.58895) (1.34350,0.60994) (1.34350,0.63074) (1.34350,0.65134) (1.34350,0.67175) (1.34350,0.69196) (1.34350,0.71195) (1.34350,0.73172) (1.34350,0.75128) (1.34350,0.77060) (1.34350,0.78969) (1.34350,0.80854) (1.34350,0.82714) (1.34350,0.84549) (1.34350,0.86359) (1.34350,0.88142) (1.34350,0.89898) (1.34350,0.91627) (1.34350,0.93328) (1.34350,0.95000) (1.34350,0.96644) (1.34350,0.98258) (1.34350,0.99842) (1.34350,1.01395) (1.34350,1.02918) (1.34350,1.04410) (1.34350,1.05869) (1.34350,1.07297) (1.34350,1.08692) (1.34350,1.10053) (1.34350,1.11381) (1.34350,1.12676) (1.34350,1.13936) (1.34350,1.15161) (1.34350,1.16351) (1.34350,1.17505) (1.34350,1.18624) (1.34350,1.19707) (1.34350,1.20753) (1.34350,1.21763) (1.34350,1.22735) (1.34350,1.23670) (1.34350,1.24567) (1.34350,1.25427) (1.34350,1.26248) (1.34350,1.27031) (1.34350,1.27775) (1.34350,1.28480) (1.34350,1.29146) (1.34350,1.29772) (1.34350,1.30360) (1.34350,1.30907) (1.34350,1.31414) (1.34350,1.31882) (1.34350,1.32309) (1.34350,1.32696) (1.34350,1.33043) (1.34350,1.33349) (1.34350,1.33614) (1.34350,1.33839) (1.34350,1.34023) (1.34350,1.34166) (1.34350,1.34268) (1.34350,1.34330) (1.34350,1.34350)};
\node[courseaccent,inner sep=0pt] at (1.34350,0.00000) {$\times$};
\draw[->,courseaccent,line width=.7pt] plot coordinates {(0.00000,0.00000) (1.34350,1.34350)};
\filldraw[fill=courseink,draw=coursepaper,line width=.4pt] (1.34350,1.34350) circle[radius=2.3pt];
\draw[->,courseaccent,line width=.7pt] plot coordinates {(0.00000,0.00000) (1.34350,-1.34350)};
\filldraw[fill=courseink,draw=coursepaper,line width=.4pt] (1.34350,-1.34350) circle[radius=2.3pt];
\node[below left,inner sep=2pt] at (0,0) {$0$};
\node[align=center,text width=3.45cm,anchor=north] at (0,-3.15) {г) Предельный NDC\\Отрезок};
\end{scope}
\end{tikzpicture}

\end{frame}

% Источник: eq:plane-change-boundaries, eq:spatial-hodograph.
\begin{frame}{Годограф в пространстве базис-вектора}
\relax
\derivationmathspacing
\[
\begin{pmatrix}s_r\\s_t\\s_z\end{pmatrix}
=\begin{pmatrix}0\\2\lambda_1\\0\end{pmatrix}
+\begin{pmatrix}-\lambda_2\\2\lambda_3\\\lambda_5\end{pmatrix}\cos\varphi
+\begin{pmatrix}\lambda_3\\2\lambda_2\\\lambda_4\end{pmatrix}\sin\varphi.
\]
\justifying
При $\Delta i>0$ ось $X$ направлена по линии пересечения плоскостей;
$\Delta i$ --- их взаимный угол.
\par\smallskip
\[
\sum_j\Delta\widehat V_{zj}\sin\varphi_j=0,\qquad
\sum_j\Delta\widehat V_{zj}\cos\varphi_j=\Delta i.
\]
\justifying
Постоянное слагаемое задаёт центр, гармонические --- плоскость
и размеры эллипса. Точки касания единичной сферы задают места импульсов.
\end{frame}

% Источник: sec:NC, строки 502--521 методички.
\begin{frame}{Узловое решение: геометрия}
\relax
\derivationmathspacing
При $\lambda_1=0$ центр эллипса находится в начале координат.
\[
\varphi_2=\varphi_1+\pi,\qquad
\boldsymbol s(\varphi_2)=-\boldsymbol s(\varphi_1).
\]
\justifying
Два противоположных места приложения нормальных импульсов лежат
на одной прямой, проходящей через центр. Оба импульса меняют плоскость
орбиты вокруг этой прямой. Чтобы получить заданную конечную плоскость,
эта прямая должна совпасть с линией пересечения плоскостей.
\par
\begin{block}{Места импульсов}
\[
\varphi_1=0,\qquad \varphi_2=\pi.
\]
\end{block}
\justifying
Направления противоположны в локальных компонентах;
модули импульсов определяются граничными условиями.
\end{frame}

\begin{frame}{Узловое решение: годограф}
\relax
\centering
\begin{tikzpicture}[course diagram,scale=.72]
\node at (2.65,6.85) {Годограф};
\node at (11.6,6.85) {Модули импульсов};
\node at (11.6,4.1) {Модуль базис-вектора};
\fill[coursetint!45!coursepaper] (2.65,3.6) circle[radius=1.7cm];
\filldraw[fill=courseaccent,fill opacity=.09,draw=courseaccent!30] (3.85370,3.65614)--(3.05981,6.09937)--(1.44630,3.54386)--(2.24019,1.10063)--cycle;
\draw[coursemuted!40,thin,densely dashed] plot coordinates {(1.34772,3.13819) (1.38665,3.11926) (1.42712,3.10092) (1.46908,3.08319) (1.51248,3.06609) (1.55726,3.04963) (1.60338,3.03385) (1.65077,3.01876) (1.69937,3.00438) (1.74914,2.99072) (1.80000,2.97780) (1.85190,2.96565) (1.90477,2.95426) (1.95855,2.94366) (2.01317,2.93386) (2.06857,2.92488) (2.12467,2.91671) (2.18142,2.90938) (2.23873,2.90289) (2.29655,2.89725) (2.35480,2.89246) (2.41341,2.88854) (2.47230,2.88548) (2.53141,2.88330) (2.59067,2.88199) (2.65000,2.88155) (2.70933,2.88199) (2.76859,2.88330) (2.82770,2.88548) (2.88659,2.88854) (2.94520,2.89246) (3.00345,2.89725) (3.06127,2.90289) (3.11858,2.90938) (3.17533,2.91671) (3.23143,2.92488) (3.28683,2.93386) (3.34145,2.94366) (3.39523,2.95426) (3.44810,2.96565) (3.50000,2.97780) (3.55086,2.99072) (3.60063,3.00438) (3.64923,3.01876) (3.69662,3.03385) (3.74274,3.04963) (3.78752,3.06609) (3.83092,3.08319) (3.87288,3.10092) (3.91335,3.11926) (3.95228,3.13819) (3.98962,3.15768) (4.02533,3.17771) (4.05936,3.19825) (4.09168,3.21928) (4.12224,3.24077) (4.15101,3.26271) (4.17795,3.28505) (4.20303,3.30778) (4.22621,3.33086) (4.24748,3.35428) (4.26680,3.37799) (4.28414,3.40197) (4.29950,3.42619) (4.31285,3.45063) (4.32417,3.47524) (4.33346,3.50001) (4.34069,3.52490) (4.34586,3.54988) (4.34896,3.57493) (4.35000,3.60000) (4.34896,3.62507) (4.34586,3.65012) (4.34069,3.67510) (4.33346,3.69999) (4.32417,3.72476) (4.31285,3.74937) (4.29950,3.77381) (4.28414,3.79803) (4.26680,3.82201) (4.24748,3.84572) (4.22621,3.86914) (4.20303,3.89222) (4.17795,3.91495) (4.15101,3.93729) (4.12224,3.95923) (4.09168,3.98072) (4.05936,4.00175) (4.02533,4.02229) (3.98962,4.04232) (3.95228,4.06181) (3.91335,4.08074) (3.87288,4.09908) (3.83092,4.11681) (3.78752,4.13391) (3.74274,4.15037) (3.69662,4.16615) (3.64923,4.18124) (3.60063,4.19562) (3.55086,4.20928) (3.50000,4.22220) (3.44810,4.23435) (3.39523,4.24574) (3.34145,4.25634) (3.28683,4.26614) (3.23143,4.27512) (3.17533,4.28329) (3.11858,4.29062) (3.06127,4.29711) (3.00345,4.30275) (2.94520,4.30754) (2.88659,4.31146) (2.82770,4.31452) (2.76859,4.31670) (2.70933,4.31801) (2.65000,4.31845) (2.59067,4.31801) (2.53141,4.31670) (2.47230,4.31452) (2.41341,4.31146) (2.35480,4.30754) (2.29655,4.30275) (2.23873,4.29711) (2.18142,4.29062) (2.12467,4.28329) (2.06857,4.27512) (2.01317,4.26614) (1.95855,4.25634) (1.90477,4.24574) (1.85190,4.23435) (1.80000,4.22220) (1.74914,4.20928) (1.69937,4.19562) (1.65077,4.18124) (1.60338,4.16615) (1.55726,4.15037) (1.51248,4.13391) (1.46908,4.11681) (1.42712,4.09908) (1.38665,4.08074) (1.34772,4.06181) (1.31038,4.04232) (1.27467,4.02229) (1.24064,4.00175) (1.20832,3.98072) (1.17776,3.95923) (1.14899,3.93729) (1.12205,3.91495) (1.09697,3.89222) (1.07379,3.86914) (1.05252,3.84572) (1.03320,3.82201) (1.01586,3.79803) (1.00050,3.77381) (0.98715,3.74937) (0.97583,3.72476) (0.96654,3.69999) (0.95931,3.67510) (0.95414,3.65012) (0.95104,3.62507) (0.95000,3.60000) (0.95104,3.57493) (0.95414,3.54988) (0.95931,3.52490) (0.96654,3.50001) (0.97583,3.47524) (0.98715,3.45063) (1.00050,3.42619) (1.01586,3.40197) (1.03320,3.37799) (1.05252,3.35428) (1.07379,3.33086) (1.09697,3.30778) (1.12205,3.28505) (1.14899,3.26271) (1.17776,3.24077) (1.20832,3.21928) (1.24064,3.19825) (1.27467,3.17771) (1.31038,3.15768) (1.34772,3.13819)};
\draw[coursemuted!55,thin] plot coordinates {(0.95000,3.60000) (0.95104,3.57493) (0.95414,3.54988) (0.95931,3.52490) (0.96654,3.50001) (0.97583,3.47524) (0.98715,3.45063) (1.00050,3.42619) (1.01586,3.40197) (1.03320,3.37799) (1.05252,3.35428) (1.07379,3.33086) (1.09697,3.30778) (1.12205,3.28505) (1.14899,3.26271) (1.17776,3.24077) (1.20832,3.21928) (1.24064,3.19825) (1.27467,3.17771) (1.31038,3.15768) (1.34772,3.13819) (1.38665,3.11926) (1.42712,3.10092) (1.46908,3.08319) (1.51248,3.06609) (1.55726,3.04963) (1.60338,3.03385) (1.65077,3.01876) (1.69937,3.00438) (1.74914,2.99072) (1.80000,2.97780) (1.85190,2.96565) (1.90477,2.95426) (1.95855,2.94366) (2.01317,2.93386) (2.06857,2.92488) (2.12467,2.91671) (2.18142,2.90938) (2.23873,2.90289) (2.29655,2.89725) (2.35480,2.89246) (2.41341,2.88854) (2.47230,2.88548) (2.53141,2.88330) (2.59067,2.88199) (2.65000,2.88155) (2.70933,2.88199) (2.76859,2.88330) (2.82770,2.88548) (2.88659,2.88854) (2.94520,2.89246) (3.00345,2.89725) (3.06127,2.90289) (3.11858,2.90938) (3.17533,2.91671) (3.23143,2.92488) (3.28683,2.93386) (3.34145,2.94366) (3.39523,2.95426) (3.44810,2.96565) (3.50000,2.97780) (3.55086,2.99072) (3.60063,3.00438) (3.64923,3.01876) (3.69662,3.03385) (3.74274,3.04963) (3.78752,3.06609) (3.83092,3.08319) (3.87288,3.10092) (3.91335,3.11926) (3.95228,3.13819) (3.98962,3.15768) (4.02533,3.17771) (4.05936,3.19825) (4.09168,3.21928) (4.12224,3.24077) (4.15101,3.26271) (4.17795,3.28505) (4.20303,3.30778) (4.22621,3.33086) (4.24748,3.35428) (4.26680,3.37799) (4.28414,3.40197) (4.29950,3.42619) (4.31285,3.45063) (4.32417,3.47524) (4.33346,3.50001) (4.34069,3.52490) (4.34586,3.54988) (4.34896,3.57493) (4.35000,3.60000)};
\draw[coursemuted!40,thin,densely dashed] plot coordinates {(1.34772,3.13819) (1.34852,3.19224) (1.35090,3.24679) (1.35486,3.30177) (1.36040,3.35711) (1.36751,3.41275) (1.37618,3.46861) (1.38641,3.52464) (1.39817,3.58076) (1.41146,3.63690) (1.42626,3.69300) (1.44255,3.74898) (1.46031,3.80478) (1.47952,3.86034) (1.50016,3.91557) (1.52220,3.97042) (1.54561,4.02482) (1.57036,4.07870) (1.59644,4.13200) (1.62379,4.18465) (1.65240,4.23659) (1.68222,4.28775) (1.71322,4.33808) (1.74536,4.38750) (1.77861,4.43597) (1.81291,4.48342) (1.84824,4.52979) (1.88454,4.57503) (1.92178,4.61908) (1.95990,4.66188) (1.99886,4.70340) (2.03862,4.74357) (2.07912,4.78235) (2.12032,4.81969) (2.16216,4.85554) (2.20460,4.88986) (2.24757,4.92261) (2.29104,4.95375) (2.33495,4.98323) (2.37924,5.01104) (2.42386,5.03712) (2.46876,5.06146) (2.51388,5.08401) (2.55916,5.10476) (2.60455,5.12367) (2.65000,5.14072) (2.69545,5.15590) (2.74084,5.16918) (2.78612,5.18056) (2.83124,5.19000) (2.87614,5.19751) (2.92076,5.20307) (2.96505,5.20668) (3.00896,5.20833) (3.05243,5.20802) (3.09540,5.20576) (3.13784,5.20153) (3.17968,5.19536) (3.22088,5.18724) (3.26138,5.17719) (3.30114,5.16521) (3.34010,5.15133) (3.37822,5.13556) (3.41546,5.11792) (3.45176,5.09843) (3.48709,5.07711) (3.52139,5.05399) (3.55464,5.02910) (3.58678,5.00248) (3.61778,4.97414) (3.64760,4.94413) (3.67621,4.91248) (3.70356,4.87923) (3.72964,4.84442) (3.75439,4.80810) (3.77780,4.77030) (3.79984,4.73108) (3.82048,4.69048) (3.83969,4.64855) (3.85745,4.60535) (3.87374,4.56092) (3.88854,4.51532) (3.90183,4.46860) (3.91359,4.42083) (3.92382,4.37205) (3.93249,4.32234) (3.93960,4.27174) (3.94514,4.22033) (3.94910,4.16816) (3.95148,4.11530) (3.95228,4.06181) (3.95148,4.00776) (3.94910,3.95321) (3.94514,3.89823) (3.93960,3.84289) (3.93249,3.78725) (3.92382,3.73139) (3.91359,3.67536) (3.90183,3.61924) (3.88854,3.56310) (3.87374,3.50700) (3.85745,3.45102) (3.83969,3.39522) (3.82048,3.33966) (3.79984,3.28443) (3.77780,3.22958) (3.75439,3.17518) (3.72964,3.12130) (3.70356,3.06800) (3.67621,3.01535) (3.64760,2.96341) (3.61778,2.91225) (3.58678,2.86192) (3.55464,2.81250) (3.52139,2.76403) (3.48709,2.71658) (3.45176,2.67021) (3.41546,2.62497) (3.37822,2.58092) (3.34010,2.53812) (3.30114,2.49660) (3.26138,2.45643) (3.22088,2.41765) (3.17968,2.38031) (3.13784,2.34446) (3.09540,2.31014) (3.05243,2.27739) (3.00896,2.24625) (2.96505,2.21677) (2.92076,2.18896) (2.87614,2.16288) (2.83124,2.13854) (2.78612,2.11599) (2.74084,2.09524) (2.69545,2.07633) (2.65000,2.05928) (2.60455,2.04410) (2.55916,2.03082) (2.51388,2.01944) (2.46876,2.01000) (2.42386,2.00249) (2.37924,1.99693) (2.33495,1.99332) (2.29104,1.99167) (2.24757,1.99198) (2.20460,1.99424) (2.16216,1.99847) (2.12032,2.00464) (2.07912,2.01276) (2.03862,2.02281) (1.99886,2.03479) (1.95990,2.04867) (1.92178,2.06444) (1.88454,2.08208) (1.84824,2.10157) (1.81291,2.12289) (1.77861,2.14601) (1.74536,2.17090) (1.71322,2.19752) (1.68222,2.22586) (1.65240,2.25587) (1.62379,2.28752) (1.59644,2.32077) (1.57036,2.35558) (1.54561,2.39190) (1.52220,2.42970) (1.50016,2.46892) (1.47952,2.50952) (1.46031,2.55145) (1.44255,2.59465) (1.42626,2.63908) (1.41146,2.68468) (1.39817,2.73140) (1.38641,2.77917) (1.37618,2.82795) (1.36751,2.87766) (1.36040,2.92826) (1.35486,2.97967) (1.35090,3.03184) (1.34852,3.08470) (1.34772,3.13819)};
\draw[coursemuted!55,thin] plot coordinates {(1.88530,2.08170) (1.84898,2.10115) (1.81363,2.12243) (1.77930,2.14551) (1.74604,2.17037) (1.71387,2.19696) (1.68285,2.22526) (1.65300,2.25524) (1.62437,2.28686) (1.59699,2.32007) (1.57089,2.35485) (1.54610,2.39114) (1.52266,2.42891) (1.50060,2.46810) (1.47993,2.50867) (1.46069,2.55057) (1.44290,2.59375) (1.42658,2.63815) (1.41175,2.68373) (1.39843,2.73043) (1.38663,2.77818) (1.37638,2.82693) (1.36767,2.87663) (1.36053,2.92721) (1.35496,2.97860) (1.35096,3.03076) (1.34855,3.08361) (1.34772,3.13708) (1.34849,3.19112) (1.35083,3.24566) (1.35476,3.30063) (1.36027,3.35597) (1.36735,3.41160) (1.37599,3.46746) (1.38618,3.52349) (1.39791,3.57960) (1.41117,3.63575) (1.42594,3.69184) (1.44220,3.74783) (1.45993,3.80364) (1.47911,3.85920) (1.49972,3.91444) (1.52173,3.96930) (1.54511,4.02371) (1.56984,4.07760) (1.59589,4.13091) (1.62322,4.18358) (1.65180,4.23553) (1.68160,4.28671) (1.71257,4.33705) (1.74469,4.38650) (1.77791,4.43498) (1.81220,4.48245) (1.84750,4.52884) (1.88379,4.57411) (1.92100,4.61818) (1.95911,4.66102) (1.99805,4.70256) (2.03779,4.74276) (2.07828,4.78156) (2.11946,4.81893) (2.16129,4.85481) (2.20372,4.88917) (2.24669,4.92195) (2.29015,4.95312) (2.33404,4.98265) (2.37833,5.01048) (2.42294,5.03660) (2.46783,5.06097) (2.51295,5.08356) (2.55822,5.10435) (2.60362,5.12330) (2.64906,5.14039) (2.69451,5.15561) (2.73991,5.16893) (2.78519,5.18034) (2.83032,5.18983) (2.87522,5.19737) (2.91984,5.20298) (2.96414,5.20662) (3.00806,5.20832) (3.05154,5.20805) (3.09453,5.20582) (3.13697,5.20164) (3.17883,5.19550) (3.22004,5.18742) (3.26056,5.17741) (3.30033,5.16548) (3.33931,5.15163) (3.37745,5.13590) (3.41470,5.11830)};
\draw[coursemuted!40,thin,densely dashed] plot coordinates {(3.74274,3.04963) (3.74207,3.10374) (3.74008,3.15845) (3.73675,3.21370) (3.73210,3.26942) (3.72614,3.32554) (3.71886,3.38200) (3.71028,3.43872) (3.70041,3.49564) (3.68926,3.55268) (3.67684,3.60978) (3.66317,3.66688) (3.64827,3.72388) (3.63215,3.78074) (3.61483,3.83738) (3.59634,3.89373) (3.57670,3.94972) (3.55592,4.00529) (3.53404,4.06036) (3.51109,4.11487) (3.48709,4.16875) (3.46206,4.22194) (3.43605,4.27438) (3.40908,4.32599) (3.38119,4.37671) (3.35240,4.42649) (3.32276,4.47527) (3.29230,4.52297) (3.26105,4.56956) (3.22906,4.61496) (3.19637,4.65912) (3.16301,4.70200) (3.12903,4.74353) (3.09446,4.78367) (3.05935,4.82236) (3.02374,4.85957) (2.98767,4.89524) (2.95120,4.92934) (2.91436,4.96181) (2.87719,4.99263) (2.83975,5.02175) (2.80208,5.04913) (2.76422,5.07475) (2.72623,5.09858) (2.68814,5.12058) (2.65000,5.14072) (2.61186,5.15899) (2.57377,5.17536) (2.53578,5.18981) (2.49792,5.20233) (2.46025,5.21289) (2.42281,5.22148) (2.38564,5.22810) (2.34880,5.23274) (2.31233,5.23539) (2.27626,5.23604) (2.24065,5.23470) (2.20554,5.23137) (2.17097,5.22606) (2.13699,5.21876) (2.10363,5.20949) (2.07094,5.19826) (2.03895,5.18508) (2.00770,5.16997) (1.97724,5.15295) (1.94760,5.13403) (1.91881,5.11325) (1.89092,5.09062) (1.86395,5.06618) (1.83794,5.03995) (1.81291,5.01196) (1.78891,4.98226) (1.76596,4.95087) (1.74408,4.91784) (1.72330,4.88320) (1.70366,4.84699) (1.68517,4.80927) (1.66785,4.77007) (1.65173,4.72945) (1.63683,4.68746) (1.62316,4.64413) (1.61074,4.59954) (1.59959,4.55373) (1.58972,4.50675) (1.58114,4.45867) (1.57386,4.40955) (1.56790,4.35944) (1.56325,4.30840) (1.55992,4.25650) (1.55793,4.20380) (1.55726,4.15037) (1.55793,4.09626) (1.55992,4.04155) (1.56325,3.98630) (1.56790,3.93058) (1.57386,3.87446) (1.58114,3.81800) (1.58972,3.76128) (1.59959,3.70436) (1.61074,3.64732) (1.62316,3.59022) (1.63683,3.53312) (1.65173,3.47612) (1.66785,3.41926) (1.68517,3.36262) (1.70366,3.30627) (1.72330,3.25028) (1.74408,3.19471) (1.76596,3.13964) (1.78891,3.08513) (1.81291,3.03125) (1.83794,2.97806) (1.86395,2.92562) (1.89092,2.87401) (1.91881,2.82329) (1.94760,2.77351) (1.97724,2.72473) (2.00770,2.67703) (2.03895,2.63044) (2.07094,2.58504) (2.10363,2.54088) (2.13699,2.49800) (2.17097,2.45647) (2.20554,2.41633) (2.24065,2.37764) (2.27626,2.34043) (2.31233,2.30476) (2.34880,2.27066) (2.38564,2.23819) (2.42281,2.20737) (2.46025,2.17825) (2.49792,2.15087) (2.53578,2.12525) (2.57377,2.10142) (2.61186,2.07942) (2.65000,2.05928) (2.68814,2.04101) (2.72623,2.02464) (2.76422,2.01019) (2.80208,1.99767) (2.83975,1.98711) (2.87719,1.97852) (2.91436,1.97190) (2.95120,1.96726) (2.98767,1.96461) (3.02374,1.96396) (3.05935,1.96530) (3.09446,1.96863) (3.12903,1.97394) (3.16301,1.98124) (3.19637,1.99051) (3.22906,2.00174) (3.26105,2.01492) (3.29230,2.03003) (3.32276,2.04705) (3.35240,2.06597) (3.38119,2.08675) (3.40908,2.10938) (3.43605,2.13382) (3.46206,2.16005) (3.48709,2.18804) (3.51109,2.21774) (3.53404,2.24913) (3.55592,2.28216) (3.57670,2.31680) (3.59634,2.35301) (3.61483,2.39073) (3.63215,2.42993) (3.64827,2.47055) (3.66317,2.51254) (3.67684,2.55587) (3.68926,2.60046) (3.70041,2.64627) (3.71028,2.69325) (3.71886,2.74133) (3.72614,2.79045) (3.73210,2.84056) (3.73675,2.89160) (3.74008,2.94350) (3.74207,2.99620) (3.74274,3.04963)};
\draw[coursemuted!55,thin] plot coordinates {(3.21818,1.99776) (3.25041,2.01029) (3.28191,2.02475) (3.31264,2.04114) (3.34256,2.05942) (3.37164,2.07958) (3.39983,2.10159) (3.42712,2.12543) (3.45345,2.15107) (3.47881,2.17847) (3.50316,2.20760) (3.52647,2.23843) (3.54871,2.27091) (3.56986,2.30502) (3.58988,2.34070) (3.60876,2.37792) (3.62648,2.41663) (3.64300,2.45678) (3.65831,2.49832) (3.67240,2.54120) (3.68524,2.58538) (3.69681,2.63079) (3.70712,2.67738) (3.71613,2.72510) (3.72384,2.77388) (3.73025,2.82366) (3.73534,2.87440) (3.73911,2.92602) (3.74155,2.97845) (3.74266,3.03165) (3.74244,3.08554) (3.74089,3.14005) (3.73801,3.19513) (3.73381,3.25070) (3.72828,3.30669) (3.72144,3.36304) (3.71330,3.41968) (3.70386,3.47654) (3.69314,3.53355) (3.68114,3.59065) (3.66789,3.64775) (3.65340,3.70479) (3.63768,3.76171) (3.62077,3.81843) (3.60267,3.87488) (3.58341,3.93100) (3.56301,3.98672) (3.54150,4.04196) (3.51890,4.09667) (3.49525,4.15077) (3.47056,4.20420) (3.44488,4.25689) (3.41822,4.30879) (3.39063,4.35982) (3.36214,4.40992) (3.33278,4.45904) (3.30259,4.50711) (3.27161,4.55408) (3.23986,4.59988) (3.20740,4.64446) (3.17426,4.68778) (3.14048,4.72976) (3.10610,4.77037) (3.07117,4.80956) (3.03572,4.84727) (2.99981,4.88346) (2.96347,4.91809) (2.92674,4.95111) (2.88968,4.98249) (2.85233,5.01218) (2.81473,5.04015) (2.77693,5.06637) (2.73897,5.09080) (2.70091,5.11341) (2.66278,5.13418) (2.62464,5.15308) (2.58653,5.17009) (2.54850,5.18518) (2.51059,5.19835) (2.47285,5.20957) (2.43533,5.21882) (2.39806,5.22610) (2.36111,5.23141) (2.32450,5.23472) (2.28830,5.23604) (2.25253,5.23537) (2.21725,5.23271) (2.18250,5.22806) (2.14831,5.22142) (2.11474,5.21281) (2.08182,5.20224)};
\draw[coursemuted!50,thin,densely dashed] plot coordinates {(4.14762,4.13108) (2.65000,3.60000)};
\draw[->,courseink!75,thin] plot coordinates {(2.65000,3.60000) (0.82681,2.95346)};
\node[above right,inner sep=1pt] at (0.82681,2.95346) {$s_r$};
\draw[coursemuted!50,thin,densely dashed] plot coordinates {(1.39335,4.23292) (2.65000,3.60000)};
\draw[->,courseink!75,thin] plot coordinates {(2.65000,3.60000) (4.17983,2.82949)};
\node[below right,inner sep=1pt] at (4.17983,2.82949) {$s_t$};
\draw[coursemuted!50,thin,densely dashed] plot coordinates {(2.65000,1.82817) (2.65000,3.60000)};
\draw[->,courseink!75,thin] plot coordinates {(2.65000,3.60000) (2.65000,5.75701)};
\node[above,inner sep=1pt] at (2.65000,5.75701) {$s_z$};
\draw[courseaccent,line width=1.15pt] plot coordinates {(3.41109,4.80543) (3.41434,4.79489) (3.41735,4.78400) (3.42013,4.77274) (3.42268,4.76112) (3.42499,4.74915) (3.42707,4.73683) (3.42891,4.72416) (3.43051,4.71115) (3.43187,4.69780) (3.43299,4.68412) (3.43388,4.67011) (3.43453,4.65577) (3.43494,4.64111) (3.43511,4.62613) (3.43504,4.61084) (3.43473,4.59524) (3.43418,4.57934) (3.43340,4.56314) (3.43237,4.54665) (3.43111,4.52987) (3.42961,4.51281) (3.42787,4.49546) (3.42589,4.47785) (3.42368,4.45997) (3.42123,4.44182) (3.41855,4.42342) (3.41563,4.40477) (3.41248,4.38587) (3.40910,4.36674) (3.40549,4.34737) (3.40165,4.32777) (3.39757,4.30795) (3.39327,4.28791) (3.38875,4.26767) (3.38399,4.24722) (3.37902,4.22658) (3.37382,4.20574) (3.36840,4.18472) (3.36277,4.16352) (3.35691,4.14215) (3.35084,4.12062) (3.34456,4.09892) (3.33807,4.07708) (3.33136,4.05509) (3.32445,4.03296) (3.31734,4.01070) (3.31002,3.98831) (3.30250,3.96581) (3.29478,3.94319) (3.28686,3.92047) (3.27875,3.89765) (3.27045,3.87474) (3.26196,3.85175) (3.25328,3.82868) (3.24442,3.80554) (3.23538,3.78234) (3.22616,3.75908) (3.21676,3.73577) (3.20719,3.71243) (3.19745,3.68905) (3.18755,3.66564) (3.17748,3.64221) (3.16725,3.61877) (3.15687,3.59532) (3.14632,3.57187) (3.13563,3.54844) (3.12479,3.52502) (3.11381,3.50162) (3.10268,3.47825) (3.09142,3.45492) (3.08002,3.43163) (3.06849,3.40839) (3.05683,3.38522) (3.04505,3.36210) (3.03315,3.33906) (3.02113,3.31610) (3.00900,3.29323) (2.99676,3.27045) (2.98441,3.24777) (2.97196,3.22520) (2.95942,3.20274) (2.94678,3.18040) (2.93405,3.15819) (2.92123,3.13612) (2.90833,3.11418) (2.89535,3.09240) (2.88230,3.07077) (2.86917,3.04930) (2.85598,3.02799) (2.84273,3.00687) (2.82942,2.98592) (2.81605,2.96516) (2.80263,2.94459) (2.78917,2.92422) (2.77566,2.90406) (2.76212,2.88411) (2.74854,2.86438) (2.73493,2.84487) (2.72130,2.82560) (2.70764,2.80656) (2.69397,2.78776) (2.68028,2.76920) (2.66658,2.75090) (2.65288,2.73286) (2.63918,2.71509) (2.62548,2.69758) (2.61179,2.68035) (2.59811,2.66340) (2.58444,2.64673) (2.57080,2.63035) (2.55718,2.61427) (2.54359,2.59849) (2.53003,2.58302) (2.51650,2.56785) (2.50302,2.55300) (2.48959,2.53847) (2.47620,2.52426) (2.46286,2.51038) (2.44958,2.49683) (2.43636,2.48362) (2.42321,2.47074) (2.41013,2.45821) (2.39712,2.44603) (2.38418,2.43420) (2.37133,2.42273) (2.35856,2.41161) (2.34589,2.40086) (2.33330,2.39047) (2.32081,2.38045) (2.30842,2.37080) (2.29614,2.36152) (2.28396,2.35263) (2.27189,2.34411) (2.25994,2.33597) (2.24811,2.32822) (2.23640,2.32086) (2.22482,2.31389) (2.21336,2.30731) (2.20204,2.30112) (2.19086,2.29533) (2.17981,2.28993) (2.16891,2.28494) (2.15815,2.28034) (2.14755,2.27615) (2.13710,2.27236) (2.12680,2.26898) (2.11666,2.26600) (2.10669,2.26342) (2.09688,2.26126) (2.08724,2.25950) (2.07777,2.25815) (2.06848,2.25721) (2.05936,2.25668) (2.05042,2.25655) (2.04167,2.25684) (2.03310,2.25754) (2.02472,2.25864) (2.01653,2.26015) (2.00853,2.26207) (2.00073,2.26440) (1.99312,2.26714) (1.98572,2.27028) (1.97851,2.27383) (1.97152,2.27778) (1.96473,2.28213) (1.95814,2.28688) (1.95177,2.29204) (1.94561,2.29759) (1.93967,2.30354) (1.93394,2.30989) (1.92843,2.31662) (1.92314,2.32375) (1.91807,2.33127) (1.91322,2.33917) (1.90860,2.34746) (1.90421,2.35613) (1.90004,2.36518) (1.89610,2.37460) (1.89239,2.38440) (1.88891,2.39457) (1.88566,2.40511) (1.88265,2.41600) (1.87987,2.42726) (1.87732,2.43888) (1.87501,2.45085) (1.87293,2.46317) (1.87109,2.47584) (1.86949,2.48885) (1.86813,2.50220) (1.86701,2.51588) (1.86612,2.52989) (1.86547,2.54423) (1.86506,2.55889) (1.86489,2.57387) (1.86496,2.58916) (1.86527,2.60476) (1.86582,2.62066) (1.86660,2.63686) (1.86763,2.65335) (1.86889,2.67013) (1.87039,2.68719) (1.87213,2.70454) (1.87411,2.72215) (1.87632,2.74003) (1.87877,2.75818) (1.88145,2.77658) (1.88437,2.79523) (1.88752,2.81413) (1.89090,2.83326) (1.89451,2.85263) (1.89835,2.87223) (1.90243,2.89205) (1.90673,2.91209) (1.91125,2.93233) (1.91601,2.95278) (1.92098,2.97342) (1.92618,2.99426) (1.93160,3.01528) (1.93723,3.03648) (1.94309,3.05785) (1.94916,3.07938) (1.95544,3.10108) (1.96193,3.12292) (1.96864,3.14491) (1.97555,3.16704) (1.98266,3.18930) (1.98998,3.21169) (1.99750,3.23419) (2.00522,3.25681) (2.01314,3.27953) (2.02125,3.30235) (2.02955,3.32526) (2.03804,3.34825) (2.04672,3.37132) (2.05558,3.39446) (2.06462,3.41766) (2.07384,3.44092) (2.08324,3.46423) (2.09281,3.48757) (2.10255,3.51095) (2.11245,3.53436) (2.12252,3.55779) (2.13275,3.58123) (2.14313,3.60468) (2.15368,3.62813) (2.16437,3.65156) (2.17521,3.67498) (2.18619,3.69838) (2.19732,3.72175) (2.20858,3.74508) (2.21998,3.76837) (2.23151,3.79161) (2.24317,3.81478) (2.25495,3.83790) (2.26685,3.86094) (2.27887,3.88390) (2.29100,3.90677) (2.30324,3.92955) (2.31559,3.95223) (2.32804,3.97480) (2.34058,3.99726) (2.35322,4.01960) (2.36595,4.04181) (2.37877,4.06388) (2.39167,4.08582) (2.40465,4.10760) (2.41770,4.12923) (2.43083,4.15070) (2.44402,4.17201) (2.45727,4.19313) (2.47058,4.21408) (2.48395,4.23484) (2.49737,4.25541) (2.51083,4.27578) (2.52434,4.29594) (2.53788,4.31589) (2.55146,4.33562) (2.56507,4.35513) (2.57870,4.37440) (2.59236,4.39344) (2.60603,4.41224) (2.61972,4.43080) (2.63342,4.44910) (2.64712,4.46714) (2.66082,4.48491) (2.67452,4.50242) (2.68821,4.51965) (2.70189,4.53660) (2.71556,4.55327) (2.72920,4.56965) (2.74282,4.58573) (2.75641,4.60151) (2.76997,4.61698) (2.78350,4.63215) (2.79698,4.64700) (2.81041,4.66153) (2.82380,4.67574) (2.83714,4.68962) (2.85042,4.70317) (2.86364,4.71638) (2.87679,4.72926) (2.88987,4.74179) (2.90288,4.75397) (2.91582,4.76580) (2.92867,4.77727) (2.94144,4.78839) (2.95411,4.79914) (2.96670,4.80953) (2.97919,4.81955) (2.99158,4.82920) (3.00386,4.83848) (3.01604,4.84737) (3.02811,4.85589) (3.04006,4.86403) (3.05189,4.87178) (3.06360,4.87914) (3.07518,4.88611) (3.08664,4.89269) (3.09796,4.89888) (3.10914,4.90467) (3.12019,4.91007) (3.13109,4.91506) (3.14185,4.91966) (3.15245,4.92385) (3.16290,4.92764) (3.17320,4.93102) (3.18334,4.93400) (3.19331,4.93658) (3.20312,4.93874) (3.21276,4.94050) (3.22223,4.94185) (3.23152,4.94279) (3.24064,4.94332) (3.24958,4.94345) (3.25833,4.94316) (3.26690,4.94246) (3.27528,4.94136) (3.28347,4.93985) (3.29147,4.93793) (3.29927,4.93560) (3.30688,4.93286) (3.31428,4.92972) (3.32149,4.92617) (3.32848,4.92222) (3.33527,4.91787) (3.34186,4.91312) (3.34823,4.90796) (3.35439,4.90241) (3.36033,4.89646) (3.36606,4.89011) (3.37157,4.88338) (3.37686,4.87625) (3.38193,4.86873) (3.38678,4.86083) (3.39140,4.85254) (3.39579,4.84387) (3.39996,4.83482) (3.40390,4.82540) (3.40761,4.81560) (3.41109,4.80543)};
\draw[->,courseaccent,line width=.55pt] plot coordinates {(2.65000,3.60000) (3.41109,4.80543)};
\filldraw[fill=courseink,draw=coursepaper,line width=.4pt] (3.41109,4.80543) circle[radius=2pt];
\node[fill=coursepaper,fill opacity=.85,text opacity=1,inner sep=.6pt] at (3.54109,4.98543) {1};
\draw[->,courseaccent,line width=.55pt] plot coordinates {(2.65000,3.60000) (1.88891,2.39457)};
\filldraw[fill=courseink,draw=coursepaper,line width=.4pt] (1.88891,2.39457) circle[radius=2pt];
\node[fill=coursepaper,fill opacity=.85,text opacity=1,inner sep=.6pt] at (1.75891,2.21457) {2};
\draw[coursemuted!65,thin] (2.65,3.6) circle[radius=1.7cm];
\node[below left,inner sep=2pt] at (2.65,3.6) {$0$};
\draw[coursemuted!25,thin] plot coordinates {(7.10000,0.85000) (16.50000,0.85000)};
\node[left,inner sep=3pt] at (7.1,0.85) {$0$};
\draw[coursemuted!25,thin] plot coordinates {(7.10000,2.15000) (16.50000,2.15000)};
\node[left,inner sep=3pt] at (7.1,2.15) {$0{,}5$};
\draw[courseink!65,dashed] plot coordinates {(7.10000,3.45000) (16.50000,3.45000)};
\node[left,inner sep=3pt] at (7.1,3.45) {$1$};
\draw[->,coursemuted] plot coordinates {(7.10000,0.85000) (16.60000,0.85000)};
\draw[->,coursemuted] plot coordinates {(7.10000,0.85000) (7.10000,3.84000)};
\node[rotate=90] at (5.9,2.2) {$\|\boldsymbol s(\varphi)\|$};
\draw[courseaccent,line width=1.15pt] plot coordinates {(7.10000,3.45000) (7.11306,3.44993) (7.12611,3.44971) (7.13917,3.44935) (7.15222,3.44884) (7.16528,3.44818) (7.17833,3.44738) (7.19139,3.44644) (7.20444,3.44535) (7.21750,3.44411) (7.23056,3.44273) (7.24361,3.44121) (7.25667,3.43954) (7.26972,3.43772) (7.28278,3.43577) (7.29583,3.43367) (7.30889,3.43142) (7.32194,3.42904) (7.33500,3.42651) (7.34806,3.42384) (7.36111,3.42102) (7.37417,3.41807) (7.38722,3.41497) (7.40028,3.41173) (7.41333,3.40836) (7.42639,3.40484) (7.43944,3.40118) (7.45250,3.39739) (7.46556,3.39345) (7.47861,3.38938) (7.49167,3.38518) (7.50472,3.38083) (7.51778,3.37635) (7.53083,3.37173) (7.54389,3.36698) (7.55694,3.36210) (7.57000,3.35708) (7.58306,3.35193) (7.59611,3.34665) (7.60917,3.34124) (7.62222,3.33570) (7.63528,3.33003) (7.64833,3.32423) (7.66139,3.31830) (7.67444,3.31224) (7.68750,3.30606) (7.70056,3.29976) (7.71361,3.29333) (7.72667,3.28678) (7.73972,3.28010) (7.75278,3.27331) (7.76583,3.26640) (7.77889,3.25936) (7.79194,3.25221) (7.80500,3.24495) (7.81806,3.23756) (7.83111,3.23007) (7.84417,3.22246) (7.85722,3.21474) (7.87028,3.20691) (7.88333,3.19897) (7.89639,3.19092) (7.90944,3.18277) (7.92250,3.17451) (7.93556,3.16614) (7.94861,3.15768) (7.96167,3.14911) (7.97472,3.14045) (7.98778,3.13168) (8.00083,3.12283) (8.01389,3.11387) (8.02694,3.10483) (8.04000,3.09569) (8.05306,3.08646) (8.06611,3.07714) (8.07917,3.06774) (8.09222,3.05826) (8.10528,3.04869) (8.11833,3.03904) (8.13139,3.02931) (8.14444,3.01951) (8.15750,3.00963) (8.17056,2.99967) (8.18361,2.98965) (8.19667,2.97956) (8.20972,2.96940) (8.22278,2.95918) (8.23583,2.94889) (8.24889,2.93854) (8.26194,2.92814) (8.27500,2.91768) (8.28806,2.90717) (8.30111,2.89660) (8.31417,2.88599) (8.32722,2.87533) (8.34028,2.86463) (8.35333,2.85389) (8.36639,2.84311) (8.37944,2.83230) (8.39250,2.82145) (8.40556,2.81057) (8.41861,2.79967) (8.43167,2.78874) (8.44472,2.77779) (8.45778,2.76682) (8.47083,2.75584) (8.48389,2.74484) (8.49694,2.73384) (8.51000,2.72283) (8.52306,2.71182) (8.53611,2.70081) (8.54917,2.68980) (8.56222,2.67880) (8.57528,2.66781) (8.58833,2.65684) (8.60139,2.64589) (8.61444,2.63496) (8.62750,2.62405) (8.64056,2.61318) (8.65361,2.60234) (8.66667,2.59154) (8.67972,2.58077) (8.69278,2.57006) (8.70583,2.55940) (8.71889,2.54878) (8.73194,2.53823) (8.74500,2.52774) (8.75806,2.51732) (8.77111,2.50697) (8.78417,2.49670) (8.79722,2.48651) (8.81028,2.47640) (8.82333,2.46639) (8.83639,2.45647) (8.84944,2.44665) (8.86250,2.43693) (8.87556,2.42732) (8.88861,2.41783) (8.90167,2.40846) (8.91472,2.39922) (8.92778,2.39010) (8.94083,2.38112) (8.95389,2.37228) (8.96694,2.36358) (8.98000,2.35504) (8.99306,2.34665) (9.00611,2.33842) (9.01917,2.33035) (9.03222,2.32246) (9.04528,2.31475) (9.05833,2.30722) (9.07139,2.29987) (9.08444,2.29272) (9.09750,2.28576) (9.11056,2.27901) (9.12361,2.27246) (9.13667,2.26612) (9.14972,2.26001) (9.16278,2.25411) (9.17583,2.24844) (9.18889,2.24300) (9.20194,2.23779) (9.21500,2.23283) (9.22806,2.22810) (9.24111,2.22363) (9.25417,2.21940) (9.26722,2.21543) (9.28028,2.21172) (9.29333,2.20827) (9.30639,2.20508) (9.31944,2.20216) (9.33250,2.19951) (9.34556,2.19714) (9.35861,2.19503) (9.37167,2.19321) (9.38472,2.19166) (9.39778,2.19039) (9.41083,2.18940) (9.42389,2.18870) (9.43694,2.18827) (9.45000,2.18813) (9.46306,2.18827) (9.47611,2.18870) (9.48917,2.18940) (9.50222,2.19039) (9.51528,2.19166) (9.52833,2.19321) (9.54139,2.19503) (9.55444,2.19714) (9.56750,2.19951) (9.58056,2.20216) (9.59361,2.20508) (9.60667,2.20827) (9.61972,2.21172) (9.63278,2.21543) (9.64583,2.21940) (9.65889,2.22363) (9.67194,2.22810) (9.68500,2.23283) (9.69806,2.23779) (9.71111,2.24300) (9.72417,2.24844) (9.73722,2.25411) (9.75028,2.26001) (9.76333,2.26612) (9.77639,2.27246) (9.78944,2.27901) (9.80250,2.28576) (9.81556,2.29272) (9.82861,2.29987) (9.84167,2.30722) (9.85472,2.31475) (9.86778,2.32246) (9.88083,2.33035) (9.89389,2.33842) (9.90694,2.34665) (9.92000,2.35504) (9.93306,2.36358) (9.94611,2.37228) (9.95917,2.38112) (9.97222,2.39010) (9.98528,2.39922) (9.99833,2.40846) (10.01139,2.41783) (10.02444,2.42732) (10.03750,2.43693) (10.05056,2.44665) (10.06361,2.45647) (10.07667,2.46639) (10.08972,2.47640) (10.10278,2.48651) (10.11583,2.49670) (10.12889,2.50697) (10.14194,2.51732) (10.15500,2.52774) (10.16806,2.53823) (10.18111,2.54878) (10.19417,2.55940) (10.20722,2.57006) (10.22028,2.58077) (10.23333,2.59154) (10.24639,2.60234) (10.25944,2.61318) (10.27250,2.62405) (10.28556,2.63496) (10.29861,2.64589) (10.31167,2.65684) (10.32472,2.66781) (10.33778,2.67880) (10.35083,2.68980) (10.36389,2.70081) (10.37694,2.71182) (10.39000,2.72283) (10.40306,2.73384) (10.41611,2.74484) (10.42917,2.75584) (10.44222,2.76682) (10.45528,2.77779) (10.46833,2.78874) (10.48139,2.79967) (10.49444,2.81057) (10.50750,2.82145) (10.52056,2.83230) (10.53361,2.84311) (10.54667,2.85389) (10.55972,2.86463) (10.57278,2.87533) (10.58583,2.88599) (10.59889,2.89660) (10.61194,2.90717) (10.62500,2.91768) (10.63806,2.92814) (10.65111,2.93854) (10.66417,2.94889) (10.67722,2.95918) (10.69028,2.96940) (10.70333,2.97956) (10.71639,2.98965) (10.72944,2.99967) (10.74250,3.00963) (10.75556,3.01951) (10.76861,3.02931) (10.78167,3.03904) (10.79472,3.04869) (10.80778,3.05826) (10.82083,3.06774) (10.83389,3.07714) (10.84694,3.08646) (10.86000,3.09569) (10.87306,3.10483) (10.88611,3.11387) (10.89917,3.12283) (10.91222,3.13168) (10.92528,3.14045) (10.93833,3.14911) (10.95139,3.15768) (10.96444,3.16614) (10.97750,3.17451) (10.99056,3.18277) (11.00361,3.19092) (11.01667,3.19897) (11.02972,3.20691) (11.04278,3.21474) (11.05583,3.22246) (11.06889,3.23007) (11.08194,3.23756) (11.09500,3.24495) (11.10806,3.25221) (11.12111,3.25936) (11.13417,3.26640) (11.14722,3.27331) (11.16028,3.28010) (11.17333,3.28678) (11.18639,3.29333) (11.19944,3.29976) (11.21250,3.30606) (11.22556,3.31224) (11.23861,3.31830) (11.25167,3.32423) (11.26472,3.33003) (11.27778,3.33570) (11.29083,3.34124) (11.30389,3.34665) (11.31694,3.35193) (11.33000,3.35708) (11.34306,3.36210) (11.35611,3.36698) (11.36917,3.37173) (11.38222,3.37635) (11.39528,3.38083) (11.40833,3.38518) (11.42139,3.38938) (11.43444,3.39345) (11.44750,3.39739) (11.46056,3.40118) (11.47361,3.40484) (11.48667,3.40836) (11.49972,3.41173) (11.51278,3.41497) (11.52583,3.41807) (11.53889,3.42102) (11.55194,3.42384) (11.56500,3.42651) (11.57806,3.42904) (11.59111,3.43142) (11.60417,3.43367) (11.61722,3.43577) (11.63028,3.43772) (11.64333,3.43954) (11.65639,3.44121) (11.66944,3.44273) (11.68250,3.44411) (11.69556,3.44535) (11.70861,3.44644) (11.72167,3.44738) (11.73472,3.44818) (11.74778,3.44884) (11.76083,3.44935) (11.77389,3.44971) (11.78694,3.44993) (11.80000,3.45000) (11.81306,3.44993) (11.82611,3.44971) (11.83917,3.44935) (11.85222,3.44884) (11.86528,3.44818) (11.87833,3.44738) (11.89139,3.44644) (11.90444,3.44535) (11.91750,3.44411) (11.93056,3.44273) (11.94361,3.44121) (11.95667,3.43954) (11.96972,3.43772) (11.98278,3.43577) (11.99583,3.43367) (12.00889,3.43142) (12.02194,3.42904) (12.03500,3.42651) (12.04806,3.42384) (12.06111,3.42102) (12.07417,3.41807) (12.08722,3.41497) (12.10028,3.41173) (12.11333,3.40836) (12.12639,3.40484) (12.13944,3.40118) (12.15250,3.39739) (12.16556,3.39345) (12.17861,3.38938) (12.19167,3.38518) (12.20472,3.38083) (12.21778,3.37635) (12.23083,3.37173) (12.24389,3.36698) (12.25694,3.36210) (12.27000,3.35708) (12.28306,3.35193) (12.29611,3.34665) (12.30917,3.34124) (12.32222,3.33570) (12.33528,3.33003) (12.34833,3.32423) (12.36139,3.31830) (12.37444,3.31224) (12.38750,3.30606) (12.40056,3.29976) (12.41361,3.29333) (12.42667,3.28678) (12.43972,3.28010) (12.45278,3.27331) (12.46583,3.26640) (12.47889,3.25936) (12.49194,3.25221) (12.50500,3.24495) (12.51806,3.23756) (12.53111,3.23007) (12.54417,3.22246) (12.55722,3.21474) (12.57028,3.20691) (12.58333,3.19897) (12.59639,3.19092) (12.60944,3.18277) (12.62250,3.17451) (12.63556,3.16614) (12.64861,3.15768) (12.66167,3.14911) (12.67472,3.14045) (12.68778,3.13168) (12.70083,3.12283) (12.71389,3.11387) (12.72694,3.10483) (12.74000,3.09569) (12.75306,3.08646) (12.76611,3.07714) (12.77917,3.06774) (12.79222,3.05826) (12.80528,3.04869) (12.81833,3.03904) (12.83139,3.02931) (12.84444,3.01951) (12.85750,3.00963) (12.87056,2.99967) (12.88361,2.98965) (12.89667,2.97956) (12.90972,2.96940) (12.92278,2.95918) (12.93583,2.94889) (12.94889,2.93854) (12.96194,2.92814) (12.97500,2.91768) (12.98806,2.90717) (13.00111,2.89660) (13.01417,2.88599) (13.02722,2.87533) (13.04028,2.86463) (13.05333,2.85389) (13.06639,2.84311) (13.07944,2.83230) (13.09250,2.82145) (13.10556,2.81057) (13.11861,2.79967) (13.13167,2.78874) (13.14472,2.77779) (13.15778,2.76682) (13.17083,2.75584) (13.18389,2.74484) (13.19694,2.73384) (13.21000,2.72283) (13.22306,2.71182) (13.23611,2.70081) (13.24917,2.68980) (13.26222,2.67880) (13.27528,2.66781) (13.28833,2.65684) (13.30139,2.64589) (13.31444,2.63496) (13.32750,2.62405) (13.34056,2.61318) (13.35361,2.60234) (13.36667,2.59154) (13.37972,2.58077) (13.39278,2.57006) (13.40583,2.55940) (13.41889,2.54878) (13.43194,2.53823) (13.44500,2.52774) (13.45806,2.51732) (13.47111,2.50697) (13.48417,2.49670) (13.49722,2.48651) (13.51028,2.47640) (13.52333,2.46639) (13.53639,2.45647) (13.54944,2.44665) (13.56250,2.43693) (13.57556,2.42732) (13.58861,2.41783) (13.60167,2.40846) (13.61472,2.39922) (13.62778,2.39010) (13.64083,2.38112) (13.65389,2.37228) (13.66694,2.36358) (13.68000,2.35504) (13.69306,2.34665) (13.70611,2.33842) (13.71917,2.33035) (13.73222,2.32246) (13.74528,2.31475) (13.75833,2.30722) (13.77139,2.29987) (13.78444,2.29272) (13.79750,2.28576) (13.81056,2.27901) (13.82361,2.27246) (13.83667,2.26612) (13.84972,2.26001) (13.86278,2.25411) (13.87583,2.24844) (13.88889,2.24300) (13.90194,2.23779) (13.91500,2.23283) (13.92806,2.22810) (13.94111,2.22363) (13.95417,2.21940) (13.96722,2.21543) (13.98028,2.21172) (13.99333,2.20827) (14.00639,2.20508) (14.01944,2.20216) (14.03250,2.19951) (14.04556,2.19714) (14.05861,2.19503) (14.07167,2.19321) (14.08472,2.19166) (14.09778,2.19039) (14.11083,2.18940) (14.12389,2.18870) (14.13694,2.18827) (14.15000,2.18813) (14.16306,2.18827) (14.17611,2.18870) (14.18917,2.18940) (14.20222,2.19039) (14.21528,2.19166) (14.22833,2.19321) (14.24139,2.19503) (14.25444,2.19714) (14.26750,2.19951) (14.28056,2.20216) (14.29361,2.20508) (14.30667,2.20827) (14.31972,2.21172) (14.33278,2.21543) (14.34583,2.21940) (14.35889,2.22363) (14.37194,2.22810) (14.38500,2.23283) (14.39806,2.23779) (14.41111,2.24300) (14.42417,2.24844) (14.43722,2.25411) (14.45028,2.26001) (14.46333,2.26612) (14.47639,2.27246) (14.48944,2.27901) (14.50250,2.28576) (14.51556,2.29272) (14.52861,2.29987) (14.54167,2.30722) (14.55472,2.31475) (14.56778,2.32246) (14.58083,2.33035) (14.59389,2.33842) (14.60694,2.34665) (14.62000,2.35504) (14.63306,2.36358) (14.64611,2.37228) (14.65917,2.38112) (14.67222,2.39010) (14.68528,2.39922) (14.69833,2.40846) (14.71139,2.41783) (14.72444,2.42732) (14.73750,2.43693) (14.75056,2.44665) (14.76361,2.45647) (14.77667,2.46639) (14.78972,2.47640) (14.80278,2.48651) (14.81583,2.49670) (14.82889,2.50697) (14.84194,2.51732) (14.85500,2.52774) (14.86806,2.53823) (14.88111,2.54878) (14.89417,2.55940) (14.90722,2.57006) (14.92028,2.58077) (14.93333,2.59154) (14.94639,2.60234) (14.95944,2.61318) (14.97250,2.62405) (14.98556,2.63496) (14.99861,2.64589) (15.01167,2.65684) (15.02472,2.66781) (15.03778,2.67880) (15.05083,2.68980) (15.06389,2.70081) (15.07694,2.71182) (15.09000,2.72283) (15.10306,2.73384) (15.11611,2.74484) (15.12917,2.75584) (15.14222,2.76682) (15.15528,2.77779) (15.16833,2.78874) (15.18139,2.79967) (15.19444,2.81057) (15.20750,2.82145) (15.22056,2.83230) (15.23361,2.84311) (15.24667,2.85389) (15.25972,2.86463) (15.27278,2.87533) (15.28583,2.88599) (15.29889,2.89660) (15.31194,2.90717) (15.32500,2.91768) (15.33806,2.92814) (15.35111,2.93854) (15.36417,2.94889) (15.37722,2.95918) (15.39028,2.96940) (15.40333,2.97956) (15.41639,2.98965) (15.42944,2.99967) (15.44250,3.00963) (15.45556,3.01951) (15.46861,3.02931) (15.48167,3.03904) (15.49472,3.04869) (15.50778,3.05826) (15.52083,3.06774) (15.53389,3.07714) (15.54694,3.08646) (15.56000,3.09569) (15.57306,3.10483) (15.58611,3.11387) (15.59917,3.12283) (15.61222,3.13168) (15.62528,3.14045) (15.63833,3.14911) (15.65139,3.15768) (15.66444,3.16614) (15.67750,3.17451) (15.69056,3.18277) (15.70361,3.19092) (15.71667,3.19897) (15.72972,3.20691) (15.74278,3.21474) (15.75583,3.22246) (15.76889,3.23007) (15.78194,3.23756) (15.79500,3.24495) (15.80806,3.25221) (15.82111,3.25936) (15.83417,3.26640) (15.84722,3.27331) (15.86028,3.28010) (15.87333,3.28678) (15.88639,3.29333) (15.89944,3.29976) (15.91250,3.30606) (15.92556,3.31224) (15.93861,3.31830) (15.95167,3.32423) (15.96472,3.33003) (15.97778,3.33570) (15.99083,3.34124) (16.00389,3.34665) (16.01694,3.35193) (16.03000,3.35708) (16.04306,3.36210) (16.05611,3.36698) (16.06917,3.37173) (16.08222,3.37635) (16.09528,3.38083) (16.10833,3.38518) (16.12139,3.38938) (16.13444,3.39345) (16.14750,3.39739) (16.16056,3.40118) (16.17361,3.40484) (16.18667,3.40836) (16.19972,3.41173) (16.21278,3.41497) (16.22583,3.41807) (16.23889,3.42102) (16.25194,3.42384) (16.26500,3.42651) (16.27806,3.42904) (16.29111,3.43142) (16.30417,3.43367) (16.31722,3.43577) (16.33028,3.43772) (16.34333,3.43954) (16.35639,3.44121) (16.36944,3.44273) (16.38250,3.44411) (16.39556,3.44535) (16.40861,3.44644) (16.42167,3.44738) (16.43472,3.44818) (16.44778,3.44884) (16.46083,3.44935) (16.47389,3.44971) (16.48694,3.44993) (16.50000,3.45000)};
\draw[->,coursemuted] plot coordinates {(7.10000,4.65000) (16.60000,4.65000)};
\draw[->,coursemuted] plot coordinates {(7.10000,4.65000) (7.10000,6.50000)};
\draw[coursemuted!25,thin] plot coordinates {(7.10000,4.65000) (16.50000,4.65000)};
\node[left,inner sep=3pt] at (7.1,4.65000) {$0$};
\draw[coursemuted!25,thin] plot coordinates {(7.10000,5.45000) (16.50000,5.45000)};
\node[left,inner sep=3pt] at (7.1,5.45000) {$8$};
\draw[coursemuted!25,thin] plot coordinates {(7.10000,6.25000) (16.50000,6.25000)};
\node[left,inner sep=3pt] at (7.1,6.25000) {$16$};
\node[rotate=90] at (5.9,5.55) {$\Delta V_j$, м/с};
\draw[coursemuted!16,thin] plot coordinates {(7.10000,0.80000) (7.10000,6.25000)};
\node[below,inner sep=3pt] at (7.10000,0.85) {$0$};
\draw[coursemuted!16,thin] plot coordinates {(9.45000,0.80000) (9.45000,6.25000)};
\node[below,inner sep=3pt] at (9.45000,0.85) {$\pi/2$};
\draw[coursemuted!16,thin] plot coordinates {(11.80000,0.80000) (11.80000,6.25000)};
\node[below,inner sep=3pt] at (11.80000,0.85) {$\pi$};
\draw[coursemuted!16,thin] plot coordinates {(14.15000,0.80000) (14.15000,6.25000)};
\node[below,inner sep=3pt] at (14.15000,0.85) {$3\pi/2$};
\draw[coursemuted!16,thin] plot coordinates {(16.50000,0.80000) (16.50000,6.25000)};
\node[below,inner sep=3pt] at (16.50000,0.85) {$2\pi$};
\node at (11.8,.15) {Угол $\varphi$, рад};
\draw[courseink!70,densely dashed,thin] plot coordinates {(7.10000,0.85000) (7.10000,3.45000)};
\draw[courseink,line width=1pt] plot coordinates {(7.10000,4.65000) (7.10000,5.94676)};
\fill[courseink] (7.10000,5.94676) circle[radius=1.8pt];
\fill[courseink] (7.10000,3.45000) circle[radius=1.8pt];
\node[above right,inner sep=2pt] at (7.10000,5.94676) {1};
\draw[courseink!70,densely dashed,thin] plot coordinates {(11.80000,0.85000) (11.80000,3.45000)};
\draw[courseink,line width=1pt] plot coordinates {(11.80000,4.65000) (11.80000,5.08225)};
\fill[courseink] (11.80000,5.08225) circle[radius=1.8pt];
\fill[courseink] (11.80000,3.45000) circle[radius=1.8pt];
\node[above,inner sep=2pt] at (11.80000,5.08225) {2};
\end{tikzpicture}

\par
\figcaption{$\Delta a=0{,}001$, $\Delta e_x=0{,}002$, $\Delta e_y=0{,}0005$, $\Delta i=0{,}002$ рад.\\
Номера точек годографа соответствуют номерам импульсов на графиках.}
\end{frame}

% Источник: eq:nc-balances, eq:nc-cost.
% Фрагмент для замены двух расчётных слайдов NC после рисунка.
% Источник: manual/chapter-3-algorithms.tex, sec:NC, eq:nc-balances--eq:nc-region.
% Фрагмент для замены двух расчётных слайдов NC после рисунка.
% Источник: manual/chapter-3-algorithms.tex, sec:NC, eq:nc-balances--eq:nc-region.
\begin{frame}{Узловое решение}
\relax
\derivationmathspacing
Пусть $\boldsymbol d=(d_r,d_t,d_z)$ --- единичное направление
первого импульса. Второй направлен вдоль $-\boldsymbol d$:
\[
\varphi_1=0,\qquad \varphi_2=\pi,\qquad
d_r^2+d_t^2+d_z^2=1.
\]
\[
\begin{aligned}
\Delta\widehat{\boldsymbol V}_1
&=\Delta\widehat V_1(d_r,d_t,d_z),\\
\Delta\widehat{\boldsymbol V}_2
&=\Delta\widehat V_2(-d_r,-d_t,-d_z).
\end{aligned}
\]
\begin{block}{Модули и суммарные затраты}
\[
\Delta\widehat V_1\geq0,\qquad
\Delta\widehat V_2\geq0,\qquad
J=\Delta\widehat V_1+\Delta\widehat V_2.
\]
\end{block}
{\justifying
Противоположность направлений ещё не означает равенства модулей.
Подставим компоненты обоих импульсов в пять граничных условий.
\par}
\end{frame}

\begin{frame}{Узловое решение}
\relax
\derivationmathspacing
\[
\begin{aligned}
\Delta a
&=2(d_t\Delta\widehat V_1-d_t\Delta\widehat V_2)
 =2d_t(\Delta\widehat V_1-\Delta\widehat V_2),\\[3pt]
\Delta e_x
&=d_r\Delta\widehat V_1\sin0
  +2d_t\Delta\widehat V_1\cos0\\
&\quad-d_r\Delta\widehat V_2\sin\pi
  -2d_t\Delta\widehat V_2\cos\pi\\
&=2d_t(\Delta\widehat V_1+\Delta\widehat V_2),\\[3pt]
\Delta e_y
&=-d_r\Delta\widehat V_1\cos0
  +2d_t\Delta\widehat V_1\sin0\\
&\quad+d_r\Delta\widehat V_2\cos\pi
  -2d_t\Delta\widehat V_2\sin\pi\\
&=-d_r(\Delta\widehat V_1+\Delta\widehat V_2).
\end{aligned}
\]
{\justifying
В условиях на эксцентриситет смена направления второго импульса
сочетается со сменой знаков синуса и косинуса при добавлении $\pi$.
\par}
\end{frame}

\begin{frame}{Узловое решение}
\relax
\derivationmathspacing
\[
\begin{aligned}
0
&=d_z\Delta\widehat V_1\sin0
  -d_z\Delta\widehat V_2\sin\pi=0,\\[4pt]
\Delta i
&=d_z\Delta\widehat V_1\cos0
  -d_z\Delta\widehat V_2\cos\pi\\
&=d_z(\Delta\widehat V_1+\Delta\widehat V_2)=Jd_z.
\end{aligned}
\]
{\justifying
Условие с синусом выполнено автоматически. Остались четыре баланса:
\par}
\[
\begin{aligned}
\Delta a&=2d_t(\Delta\widehat V_1-\Delta\widehat V_2),\\
\Delta e_x&=2Jd_t,\qquad
\Delta e_y=-Jd_r,\qquad \Delta i=Jd_z.
\end{aligned}
\]
{\justifying
Большая полуось задаёт разность модулей; три остальных условия
связывают их сумму с направлением импульсов.
\par}
\end{frame}

\begin{frame}{Узловое решение}
\relax
\derivationmathspacing
Три баланса задают компоненты направления через сумму модулей:
\[
Jd_r=-\Delta e_y,\qquad
Jd_t=\frac{\Delta e_x}{2},\qquad
Jd_z=\Delta i.
\]
Возводим их в квадрат и складываем, используя нормировку:
\[
J^2\underbrace{(d_r^2+d_t^2+d_z^2)}_{=1}
=(\Delta e_y)^2+\frac{(\Delta e_x)^2}{4}+(\Delta i)^2.
\]
\[
J_{\rm NC}=\sqrt{\frac{(\Delta e_x)^2}{4}
                         +(\Delta e_y)^2+(\Delta i)^2},
\]
\[
\boldsymbol d=\frac{1}{J_{\rm NC}}
       \bigl(-\Delta e_y,\ \Delta e_x/2,\ \Delta i\bigr).
\]
\end{frame}

\begin{frame}{Узловое решение}
\relax
\derivationmathspacing
\[
\Delta\widehat V_1+\Delta\widehat V_2=J_{\rm NC}.
\]
В условие по большой полуоси подставляем $d_t=\Delta e_x/(2J_{\rm NC})$:
\[
\Delta a=\frac{\Delta e_x}{J_{\rm NC}}
          (\Delta\widehat V_1-\Delta\widehat V_2).
\]
При $\Delta e_x\ne0$ получаем разность:
\[
\Delta\widehat V_1-\Delta\widehat V_2
 =J_{\rm NC}\frac{\Delta a}{\Delta e_x}.
\]
Сложение и вычитание суммы и разности дают
\[
2\Delta\widehat V_1=J_{\rm NC}
  \left(1+\frac{\Delta a}{\Delta e_x}\right),\qquad
2\Delta\widehat V_2=J_{\rm NC}
  \left(1-\frac{\Delta a}{\Delta e_x}\right).
\]
\end{frame}

\begin{frame}{Узловое решение}
\relax
\derivationmathspacing
Делим найденные равенства на два:
\[
\begin{aligned}
\Delta\widehat V_1&=\frac{J_{\rm NC}}2
 \left(1+\frac{\Delta a}{\Delta e_x}\right),\\
\Delta\widehat V_2&=\frac{J_{\rm NC}}2
 \left(1-\frac{\Delta a}{\Delta e_x}\right).
\end{aligned}
\]
Векторы импульсов восстанавливаются по их направлениям:
\[
\Delta\widehat{\boldsymbol V}_1=\Delta\widehat V_1\boldsymbol d,
\qquad
\Delta\widehat{\boldsymbol V}_2=-\Delta\widehat V_2\boldsymbol d.
\]
Модули должны быть неотрицательными, поэтому
\[
1+\frac{\Delta a}{\Delta e_x}\geq0,\qquad
1-\frac{\Delta a}{\Delta e_x}\geq0
\quad\Longleftrightarrow\quad
|\Delta a|\leq|\Delta e_x|.
\]
{\justifying
Таким образом, условие по большой полуоси определяет распределение
уже найденных суммарных затрат между двумя импульсами.
\par}
\end{frame}

% Развёрнутая подстановка, описанная после eq:nc-magnitudes:
% eq:primer-components и eq:primer-stationarity.
\begin{frame}{Узловое решение}
\relax
\derivationmathspacing
При $\lambda_1=0$ и $J=J_{\rm NC}$ контакт в первом узле даёт
\[
\boldsymbol s(0)=(-\lambda_2,2\lambda_3,\lambda_5)
 =\boldsymbol d=\frac1J
 \bigl(-\Delta e_y,\ \Delta e_x/2,\ \Delta i\bigr).
\]
Поэтому три множителя уже известны:
\[
\lambda_2=\frac{\Delta e_y}{J},\qquad
\lambda_3=\frac{\Delta e_x}{4J},\qquad
\lambda_5=\frac{\Delta i}{J}.
\]
Оставшийся множитель находим из стационарности модуля в узле:
\[
\boldsymbol s'(0)=(\lambda_3,2\lambda_2,\lambda_4),\qquad
\boldsymbol s(0)\mathbin\cdot\boldsymbol s'(0)=0,
\]
\[
3\lambda_2\lambda_3+\lambda_4\lambda_5=0
\quad\Longrightarrow\quad
\lambda_4=-\frac{3\Delta e_x\Delta e_y}{4J\Delta i},
\qquad \Delta i>0.
\]
\end{frame}

\begin{frame}{Узловое решение}
\relax
\derivationmathspacing
{\justifying
При $\lambda_1=0$ свободного члена нет. Коэффициенты при косинусе
равны $\boldsymbol s(0)=\boldsymbol d$, при синусе ---
$\boldsymbol s'(0)$, поэтому
\par}
\[
\boldsymbol s(\varphi)
 =\boldsymbol d\cos\varphi+\boldsymbol s'(0)\sin\varphi.
\]
Из условия стационарности и нормировки направления следует
\[
\boldsymbol d\mathbin\cdot\boldsymbol s'(0)=0,
\qquad \|\boldsymbol d\|=1,
\]
\[
\|\boldsymbol s(\varphi)\|^2
 =\cos^2\varphi+\|\boldsymbol s'(0)\|^2\sin^2\varphi.
\]
Для нахождения всего годографа внутри сферы необходимо и достаточно
\[
\|\boldsymbol s(\varphi)\|\leq1\quad\text{при всех }\varphi
\quad\Longleftrightarrow\quad
\|\boldsymbol s'(0)\|^2\leq1.
\]
\end{frame}

\begin{frame}{Узловое решение}
\relax
\derivationmathspacing
\[
\|\boldsymbol s'(0)\|^2
 =\frac{(\Delta e_x)^2}{16J^2}
  +\frac{4(\Delta e_y)^2}{J^2}
  +\frac{9(\Delta e_x)^2(\Delta e_y)^2}{16J^2(\Delta i)^2}.
\]
При $J^2=(\Delta e_x)^2/4+(\Delta e_y)^2+(\Delta i)^2$:
\[
\begin{aligned}
J^2\bigl(1-\|\boldsymbol s'(0)\|^2\bigr)
 &=\frac{3(\Delta e_x)^2}{16}-3(\Delta e_y)^2+(\Delta i)^2\\
 &\quad-\frac{9(\Delta e_x)^2(\Delta e_y)^2}{16(\Delta i)^2}\\
 &=\bigl((\Delta i)^2-3(\Delta e_y)^2\bigr)
   \left(1+\frac{3(\Delta e_x)^2}{16(\Delta i)^2}\right).
\end{aligned}
\]
Второй множитель положителен при $\Delta i>0$, поэтому
\[
\|\boldsymbol s'(0)\|^2\leq1
\quad\Longleftrightarrow\quad
\Delta i\geq\sqrt3\,|\Delta e_y|.
\]
\end{frame}

\begin{frame}{Узловое решение}
\relax
\derivationmathspacing
{\justifying
Объединяем неотрицательность модулей с условием формы годографа.
Касающаяся сферы ось должна быть большой осью эллипса,
а вторая полуось не должна выходить за сферу.
\par}
\begin{block}{Область применимости}
\[
|\Delta a|\leq|\Delta e_x|,\qquad
\Delta i\geq\sqrt3\,|\Delta e_y|.
\]
\end{block}
{\justifying
Первое ограничение относится к модулям импульсов, второе ---
к форме годографа. При равенстве во втором условии эллипс
превращается в окружность на сфере.
\par}
\end{frame}

\begin{frame}{Узловое решение}
\relax
\derivationmathspacing
{\justifying
Если $\Delta e_x=0$, первое условие требует $\Delta a=0$.
Формулы с делением на $\Delta e_x$ не применяем: достаточно
одного импульса в узле:
\par}
\[
\Delta\widehat{\boldsymbol V}_1=(-\Delta e_y,0,\Delta i),
\qquad \varphi_1=0.
\]
Пять граничных условий выполняются, а модуль импульса равен
\[
\Delta\widehat V_1
 =\sqrt{(\Delta e_y)^2+(\Delta i)^2}=J_{\rm NC}.
\]
\end{frame}



% Источник: eq:ndc-region, eq:ndc-cost. Только область и прямой алгоритм.
\begin{frame}{Невырожденное решение}
\relax
\derivationmathspacing
Смещённый эллипс касается сферы в двух симметричных точках.
\[
\begin{gathered}
\Delta e_y\ne0,\qquad |\Delta a|>|\Delta e_x|,\\
(\Delta a)^2>\Delta e^2+
\frac{2\Delta i|\Delta e_y|}{\sqrt3}-(\Delta i)^2.
\end{gathered}
\]
\[
\begin{aligned}
J_{\rm NDC}^2={}&\frac{(\Delta i)^2+\Delta e^2}{2}
-\frac{(\Delta a)^2}{4}\\
&+\frac12\sqrt{
\bigl[(\Delta a)^2+(\Delta i)^2-\Delta e^2\bigr]^2
+4(\Delta i)^2(\Delta e_y)^2}.
\end{aligned}
\]
\medskip
Берём положительный корень; далее $J=J_{\rm NDC}$.
На границах области используются узловое или особое решение.
\end{frame}

\begin{frame}{Невырожденное решение: годограф}
\relax
\centering
\begin{tikzpicture}[course diagram,scale=.72]
\node at (2.65,6.85) {Годограф};
\node at (11.6,6.85) {Модули импульсов};
\node at (11.6,4.1) {Модуль базис-вектора};
\filldraw[fill=coursetint!45!coursepaper,draw=coursemuted!65,line width=.65pt] (2.65000,3.60000) circle[radius=1.7cm];
\draw[->,courseink!75,thin] plot coordinates {(2.65000,3.60000) (0.94963,2.13566)};
\node[below left,inner sep=2pt] at (0.94963,2.13566) {$s_r$};
\draw[->,courseink!75,thin] plot coordinates {(2.65000,3.60000) (4.89400,3.60000)};
\node[right,inner sep=2pt] at (4.89400,3.60000) {$s_t$};
\draw[->,courseink!75,thin] plot coordinates {(2.65000,3.60000) (2.48204,5.83771)};
\node[above,inner sep=2pt] at (2.48204,5.83771) {$s_z$};
\draw[coursemuted!60,thin,densely dashed] plot coordinates {(3.83592,4.81803) (3.83592,2.38197)};
\draw[courseaccent,line width=1.2pt] plot coordinates {(3.79356,4.84354) (3.79829,4.84222) (3.80302,4.84052) (3.80774,4.83845) (3.81245,4.83600) (3.81716,4.83317) (3.82185,4.82996) (3.82653,4.82639) (3.83119,4.82243) (3.83584,4.81811) (3.84047,4.81341) (3.84508,4.80835) (3.84967,4.80291) (3.85424,4.79711) (3.85879,4.79095) (3.86332,4.78442) (3.86781,4.77753) (3.87229,4.77029) (3.87673,4.76268) (3.88114,4.75473) (3.88553,4.74642) (3.88988,4.73776) (3.89419,4.72875) (3.89847,4.71940) (3.90272,4.70971) (3.90693,4.69969) (3.91110,4.68932) (3.91522,4.67863) (3.91931,4.66760) (3.92336,4.65626) (3.92736,4.64459) (3.93131,4.63260) (3.93522,4.62029) (3.93909,4.60768) (3.94290,4.59476) (3.94666,4.58154) (3.95038,4.56801) (3.95404,4.55420) (3.95764,4.54009) (3.96120,4.52569) (3.96470,4.51102) (3.96814,4.49606) (3.97152,4.48084) (3.97485,4.46534) (3.97811,4.44958) (3.98132,4.43356) (3.98446,4.41729) (3.98754,4.40077) (3.99056,4.38401) (3.99352,4.36700) (3.99641,4.34977) (3.99923,4.33230) (4.00198,4.31461) (4.00467,4.29671) (4.00729,4.27859) (4.00984,4.26026) (4.01232,4.24174) (4.01473,4.22302) (4.01707,4.20411) (4.01933,4.18501) (4.02152,4.16574) (4.02364,4.14629) (4.02568,4.12668) (4.02765,4.10691) (4.02955,4.08698) (4.03136,4.06690) (4.03310,4.04669) (4.03476,4.02633) (4.03635,4.00585) (4.03786,3.98524) (4.03928,3.96452) (4.04063,3.94368) (4.04190,3.92274) (4.04309,3.90170) (4.04420,3.88057) (4.04523,3.85935) (4.04618,3.83806) (4.04704,3.81669) (4.04783,3.79526) (4.04853,3.77377) (4.04915,3.75222) (4.04969,3.73063) (4.05014,3.70900) (4.05052,3.68733) (4.05081,3.66564) (4.05102,3.64393) (4.05114,3.62220) (4.05119,3.60047) (4.05115,3.57874) (4.05103,3.55701) (4.05082,3.53530) (4.05053,3.51361) (4.05016,3.49194) (4.04971,3.47031) (4.04917,3.44871) (4.04856,3.42716) (4.04786,3.40567) (4.04708,3.38423) (4.04621,3.36286) (4.04527,3.34156) (4.04425,3.32034) (4.04314,3.29921) (4.04195,3.27817) (4.04069,3.25722) (4.03934,3.23638) (4.03792,3.21565) (4.03642,3.19504) (4.03483,3.17455) (4.03317,3.15419) (4.03144,3.13397) (4.02962,3.11388) (4.02774,3.09395) (4.02577,3.07417) (4.02373,3.05455) (4.02162,3.03510) (4.01943,3.01582) (4.01717,2.99671) (4.01483,2.97780) (4.01243,2.95907) (4.00995,2.94053) (4.00740,2.92220) (4.00479,2.90407) (4.00210,2.88615) (3.99935,2.86846) (3.99653,2.85098) (3.99364,2.83373) (3.99069,2.81672) (3.98768,2.79995) (3.98460,2.78341) (3.98146,2.76713) (3.97825,2.75110) (3.97499,2.73533) (3.97167,2.71983) (3.96828,2.70459) (3.96485,2.68962) (3.96135,2.67493) (3.95780,2.66053) (3.95419,2.64641) (3.95054,2.63258) (3.94682,2.61904) (3.94306,2.60580) (3.93925,2.59287) (3.93539,2.58024) (3.93148,2.56793) (3.92753,2.55592) (3.92353,2.54424) (3.91949,2.53288) (3.91540,2.52184) (3.91127,2.51113) (3.90711,2.50075) (3.90290,2.49071) (3.89866,2.48101) (3.89438,2.47164) (3.89006,2.46262) (3.88571,2.45395) (3.88133,2.44563) (3.87692,2.43765) (3.87248,2.43003) (3.86801,2.42277) (3.86351,2.41587) (3.85899,2.40933) (3.85444,2.40315) (3.84987,2.39733) (3.84528,2.39188) (3.84067,2.38680) (3.83604,2.38209) (3.83139,2.37775) (3.82673,2.37378) (3.82205,2.37018) (3.81736,2.36696) (3.81266,2.36412) (3.80795,2.36165) (3.80322,2.35956) (3.79850,2.35784) (3.79376,2.35651) (3.78902,2.35555) (3.78428,2.35498) (3.77954,2.35478) (3.77480,2.35496) (3.77006,2.35552) (3.76532,2.35646) (3.76058,2.35778) (3.75586,2.35948) (3.75113,2.36155) (3.74642,2.36400) (3.74172,2.36683) (3.73703,2.37004) (3.73235,2.37361) (3.72769,2.37757) (3.72304,2.38189) (3.71841,2.38659) (3.71380,2.39165) (3.70920,2.39709) (3.70463,2.40289) (3.70008,2.40905) (3.69556,2.41558) (3.69106,2.42247) (3.68659,2.42971) (3.68215,2.43732) (3.67773,2.44527) (3.67335,2.45358) (3.66900,2.46224) (3.66469,2.47125) (3.66040,2.48060) (3.65616,2.49029) (3.65195,2.50031) (3.64778,2.51068) (3.64365,2.52137) (3.63957,2.53240) (3.63552,2.54374) (3.63152,2.55541) (3.62756,2.56740) (3.62365,2.57971) (3.61979,2.59232) (3.61598,2.60524) (3.61221,2.61846) (3.60850,2.63199) (3.60484,2.64580) (3.60123,2.65991) (3.59768,2.67431) (3.59418,2.68898) (3.59074,2.70394) (3.58736,2.71916) (3.58403,2.73466) (3.58076,2.75042) (3.57756,2.76644) (3.57441,2.78271) (3.57133,2.79923) (3.56831,2.81599) (3.56536,2.83300) (3.56247,2.85023) (3.55965,2.86770) (3.55689,2.88539) (3.55421,2.90329) (3.55159,2.92141) (3.54904,2.93974) (3.54656,2.95826) (3.54415,2.97698) (3.54181,2.99589) (3.53955,3.01499) (3.53735,3.03426) (3.53524,3.05371) (3.53319,3.07332) (3.53123,3.09309) (3.52933,3.11302) (3.52752,3.13310) (3.52578,3.15331) (3.52411,3.17367) (3.52253,3.19415) (3.52102,3.21476) (3.51959,3.23548) (3.51824,3.25632) (3.51698,3.27726) (3.51579,3.29830) (3.51468,3.31943) (3.51365,3.34065) (3.51270,3.36194) (3.51184,3.38331) (3.51105,3.40474) (3.51035,3.42623) (3.50973,3.44778) (3.50919,3.46937) (3.50873,3.49100) (3.50836,3.51267) (3.50807,3.53436) (3.50786,3.55607) (3.50773,3.57780) (3.50769,3.59953) (3.50773,3.62126) (3.50785,3.64299) (3.50806,3.66470) (3.50834,3.68639) (3.50871,3.70806) (3.50917,3.72969) (3.50970,3.75129) (3.51032,3.77284) (3.51102,3.79433) (3.51180,3.81577) (3.51266,3.83714) (3.51361,3.85844) (3.51463,3.87966) (3.51574,3.90079) (3.51692,3.92183) (3.51819,3.94278) (3.51953,3.96362) (3.52096,3.98435) (3.52246,4.00496) (3.52404,4.02545) (3.52570,4.04581) (3.52744,4.06603) (3.52925,4.08612) (3.53114,4.10605) (3.53311,4.12583) (3.53515,4.14545) (3.53726,4.16490) (3.53945,4.18418) (3.54171,4.20329) (3.54405,4.22220) (3.54645,4.24093) (3.54893,4.25947) (3.55147,4.27780) (3.55409,4.29593) (3.55678,4.31385) (3.55953,4.33154) (3.56235,4.34902) (3.56523,4.36627) (3.56819,4.38328) (3.57120,4.40005) (3.57428,4.41659) (3.57742,4.43287) (3.58062,4.44890) (3.58389,4.46467) (3.58721,4.48017) (3.59059,4.49541) (3.59403,4.51038) (3.59753,4.52507) (3.60108,4.53947) (3.60468,4.55359) (3.60834,4.56742) (3.61205,4.58096) (3.61581,4.59420) (3.61963,4.60713) (3.62349,4.61976) (3.62739,4.63207) (3.63135,4.64408) (3.63535,4.65576) (3.63939,4.66712) (3.64348,4.67816) (3.64760,4.68887) (3.65177,4.69925) (3.65598,4.70929) (3.66022,4.71899) (3.66450,4.72836) (3.66881,4.73738) (3.67316,4.74605) (3.67754,4.75437) (3.68196,4.76235) (3.68640,4.76997) (3.69087,4.77723) (3.69537,4.78413) (3.69989,4.79067) (3.70444,4.79685) (3.70901,4.80267) (3.71360,4.80812) (3.71821,4.81320) (3.72284,4.81791) (3.72749,4.82225) (3.73215,4.82622) (3.73683,4.82982) (3.74152,4.83304) (3.74622,4.83588) (3.75093,4.83835) (3.75565,4.84044) (3.76038,4.84216) (3.76511,4.84349) (3.76985,4.84445) (3.77459,4.84502) (3.77934,4.84522) (3.78408,4.84504) (3.78882,4.84448) (3.79356,4.84354)};
\node[courseaccent,inner sep=0pt] at (3.77944,3.60000) {$\times$};
\draw[->,courseaccent,line width=.7pt] plot coordinates {(2.65000,3.60000) (3.83592,4.81803)};
\filldraw[fill=courseink,draw=coursepaper,line width=.4pt] (3.83592,4.81803) circle[radius=2.3pt];
\node[above right,inner sep=3pt] at (3.83592,4.81803) {1};
\draw[->,courseaccent,line width=.7pt] plot coordinates {(2.65000,3.60000) (3.83592,2.38197)};
\filldraw[fill=courseink,draw=coursepaper,line width=.4pt] (3.83592,2.38197) circle[radius=2.3pt];
\node[below right,inner sep=3pt] at (3.83592,2.38197) {2};
\node[below left,inner sep=2pt] at (2.65,3.6) {$0$};
\draw[coursemuted!25,thin] plot coordinates {(7.10000,0.85000) (16.50000,0.85000)};
\node[left,inner sep=3pt] at (7.1,0.85) {$0$};
\draw[coursemuted!25,thin] plot coordinates {(7.10000,2.15000) (16.50000,2.15000)};
\node[left,inner sep=3pt] at (7.1,2.15) {$0{,}5$};
\draw[courseink!65,dashed] plot coordinates {(7.10000,3.45000) (16.50000,3.45000)};
\node[left,inner sep=3pt] at (7.1,3.45) {$1$};
\draw[->,coursemuted] plot coordinates {(7.10000,0.85000) (16.60000,0.85000)};
\draw[->,coursemuted] plot coordinates {(7.10000,0.85000) (7.10000,3.84000)};
\node[rotate=90] at (5.9,2.2) {$\|\boldsymbol s(\varphi)\|$};
\draw[courseaccent,line width=1.15pt] plot coordinates {(7.10000,3.43381) (7.11306,3.43557) (7.12611,3.43723) (7.13917,3.43879) (7.15222,3.44025) (7.16528,3.44160) (7.17833,3.44285) (7.19139,3.44400) (7.20444,3.44504) (7.21750,3.44599) (7.23056,3.44683) (7.24361,3.44758) (7.25667,3.44822) (7.26972,3.44876) (7.28278,3.44921) (7.29583,3.44955) (7.30889,3.44980) (7.32194,3.44995) (7.33500,3.45000) (7.34806,3.44995) (7.36111,3.44981) (7.37417,3.44958) (7.38722,3.44924) (7.40028,3.44882) (7.41333,3.44830) (7.42639,3.44768) (7.43944,3.44698) (7.45250,3.44618) (7.46556,3.44529) (7.47861,3.44431) (7.49167,3.44324) (7.50472,3.44208) (7.51778,3.44083) (7.53083,3.43950) (7.54389,3.43808) (7.55694,3.43657) (7.57000,3.43498) (7.58306,3.43330) (7.59611,3.43154) (7.60917,3.42969) (7.62222,3.42777) (7.63528,3.42576) (7.64833,3.42368) (7.66139,3.42152) (7.67444,3.41927) (7.68750,3.41695) (7.70056,3.41456) (7.71361,3.41209) (7.72667,3.40955) (7.73972,3.40693) (7.75278,3.40424) (7.76583,3.40148) (7.77889,3.39866) (7.79194,3.39576) (7.80500,3.39280) (7.81806,3.38977) (7.83111,3.38667) (7.84417,3.38351) (7.85722,3.38029) (7.87028,3.37701) (7.88333,3.37367) (7.89639,3.37027) (7.90944,3.36681) (7.92250,3.36330) (7.93556,3.35973) (7.94861,3.35610) (7.96167,3.35243) (7.97472,3.34870) (7.98778,3.34493) (8.00083,3.34110) (8.01389,3.33723) (8.02694,3.33332) (8.04000,3.32936) (8.05306,3.32536) (8.06611,3.32132) (8.07917,3.31723) (8.09222,3.31311) (8.10528,3.30896) (8.11833,3.30477) (8.13139,3.30054) (8.14444,3.29628) (8.15750,3.29200) (8.17056,3.28768) (8.18361,3.28334) (8.19667,3.27897) (8.20972,3.27458) (8.22278,3.27016) (8.23583,3.26573) (8.24889,3.26127) (8.26194,3.25680) (8.27500,3.25232) (8.28806,3.24782) (8.30111,3.24330) (8.31417,3.23878) (8.32722,3.23425) (8.34028,3.22971) (8.35333,3.22516) (8.36639,3.22061) (8.37944,3.21606) (8.39250,3.21151) (8.40556,3.20696) (8.41861,3.20242) (8.43167,3.19788) (8.44472,3.19334) (8.45778,3.18882) (8.47083,3.18430) (8.48389,3.17980) (8.49694,3.17531) (8.51000,3.17084) (8.52306,3.16639) (8.53611,3.16195) (8.54917,3.15754) (8.56222,3.15315) (8.57528,3.14878) (8.58833,3.14444) (8.60139,3.14013) (8.61444,3.13585) (8.62750,3.13160) (8.64056,3.12739) (8.65361,3.12321) (8.66667,3.11907) (8.67972,3.11497) (8.69278,3.11090) (8.70583,3.10688) (8.71889,3.10291) (8.73194,3.09898) (8.74500,3.09510) (8.75806,3.09127) (8.77111,3.08749) (8.78417,3.08376) (8.79722,3.08009) (8.81028,3.07647) (8.82333,3.07291) (8.83639,3.06941) (8.84944,3.06597) (8.86250,3.06260) (8.87556,3.05929) (8.88861,3.05604) (8.90167,3.05286) (8.91472,3.04975) (8.92778,3.04671) (8.94083,3.04374) (8.95389,3.04084) (8.96694,3.03802) (8.98000,3.03527) (8.99306,3.03260) (9.00611,3.03000) (9.01917,3.02749) (9.03222,3.02505) (9.04528,3.02270) (9.05833,3.02043) (9.07139,3.01825) (9.08444,3.01615) (9.09750,3.01413) (9.11056,3.01220) (9.12361,3.01036) (9.13667,3.00861) (9.14972,3.00695) (9.16278,3.00537) (9.17583,3.00389) (9.18889,3.00250) (9.20194,3.00121) (9.21500,3.00000) (9.22806,2.99889) (9.24111,2.99788) (9.25417,2.99696) (9.26722,2.99613) (9.28028,2.99540) (9.29333,2.99477) (9.30639,2.99423) (9.31944,2.99379) (9.33250,2.99344) (9.34556,2.99320) (9.35861,2.99305) (9.37167,2.99299) (9.38472,2.99304) (9.39778,2.99318) (9.41083,2.99342) (9.42389,2.99375) (9.43694,2.99419) (9.45000,2.99472) (9.46306,2.99534) (9.47611,2.99606) (9.48917,2.99688) (9.50222,2.99779) (9.51528,2.99880) (9.52833,2.99990) (9.54139,3.00110) (9.55444,3.00239) (9.56750,3.00377) (9.58056,3.00524) (9.59361,3.00681) (9.60667,3.00846) (9.61972,3.01021) (9.63278,3.01204) (9.64583,3.01396) (9.65889,3.01597) (9.67194,3.01806) (9.68500,3.02024) (9.69806,3.02250) (9.71111,3.02485) (9.72417,3.02728) (9.73722,3.02978) (9.75028,3.03237) (9.76333,3.03503) (9.77639,3.03778) (9.78944,3.04059) (9.80250,3.04348) (9.81556,3.04645) (9.82861,3.04948) (9.84167,3.05259) (9.85472,3.05576) (9.86778,3.05900) (9.88083,3.06231) (9.89389,3.06568) (9.90694,3.06911) (9.92000,3.07261) (9.93306,3.07616) (9.94611,3.07977) (9.95917,3.08344) (9.97222,3.08717) (9.98528,3.09094) (9.99833,3.09477) (10.01139,3.09864) (10.02444,3.10257) (10.03750,3.10654) (10.05056,3.11056) (10.06361,3.11461) (10.07667,3.11871) (10.08972,3.12285) (10.10278,3.12703) (10.11583,3.13124) (10.12889,3.13548) (10.14194,3.13976) (10.15500,3.14407) (10.16806,3.14841) (10.18111,3.15277) (10.19417,3.15716) (10.20722,3.16157) (10.22028,3.16600) (10.23333,3.17046) (10.24639,3.17493) (10.25944,3.17941) (10.27250,3.18391) (10.28556,3.18843) (10.29861,3.19295) (10.31167,3.19748) (10.32472,3.20202) (10.33778,3.20657) (10.35083,3.21112) (10.36389,3.21567) (10.37694,3.22022) (10.39000,3.22477) (10.40306,3.22931) (10.41611,3.23385) (10.42917,3.23839) (10.44222,3.24291) (10.45528,3.24743) (10.46833,3.25193) (10.48139,3.25642) (10.49444,3.26089) (10.50750,3.26534) (10.52056,3.26978) (10.53361,3.27420) (10.54667,3.27859) (10.55972,3.28296) (10.57278,3.28731) (10.58583,3.29163) (10.59889,3.29592) (10.61194,3.30017) (10.62500,3.30440) (10.63806,3.30860) (10.65111,3.31276) (10.66417,3.31688) (10.67722,3.32096) (10.69028,3.32501) (10.70333,3.32902) (10.71639,3.33298) (10.72944,3.33690) (10.74250,3.34077) (10.75556,3.34460) (10.76861,3.34838) (10.78167,3.35211) (10.79472,3.35579) (10.80778,3.35942) (10.82083,3.36299) (10.83389,3.36651) (10.84694,3.36997) (10.86000,3.37338) (10.87306,3.37673) (10.88611,3.38001) (10.89917,3.38324) (10.91222,3.38640) (10.92528,3.38950) (10.93833,3.39254) (10.95139,3.39551) (10.96444,3.39841) (10.97750,3.40124) (10.99056,3.40401) (11.00361,3.40670) (11.01667,3.40932) (11.02972,3.41187) (11.04278,3.41435) (11.05583,3.41675) (11.06889,3.41908) (11.08194,3.42132) (11.09500,3.42350) (11.10806,3.42559) (11.12111,3.42760) (11.13417,3.42953) (11.14722,3.43138) (11.16028,3.43315) (11.17333,3.43483) (11.18639,3.43643) (11.19944,3.43795) (11.21250,3.43938) (11.22556,3.44072) (11.23861,3.44198) (11.25167,3.44314) (11.26472,3.44422) (11.27778,3.44521) (11.29083,3.44611) (11.30389,3.44691) (11.31694,3.44763) (11.33000,3.44825) (11.34306,3.44878) (11.35611,3.44921) (11.36917,3.44955) (11.38222,3.44980) (11.39528,3.44995) (11.40833,3.45000) (11.42139,3.44996) (11.43444,3.44982) (11.44750,3.44958) (11.46056,3.44924) (11.47361,3.44880) (11.48667,3.44827) (11.49972,3.44763) (11.51278,3.44690) (11.52583,3.44606) (11.53889,3.44513) (11.55194,3.44409) (11.56500,3.44295) (11.57806,3.44171) (11.59111,3.44037) (11.60417,3.43892) (11.61722,3.43737) (11.63028,3.43572) (11.64333,3.43396) (11.65639,3.43211) (11.66944,3.43014) (11.68250,3.42808) (11.69556,3.42591) (11.70861,3.42363) (11.72167,3.42125) (11.73472,3.41877) (11.74778,3.41618) (11.76083,3.41348) (11.77389,3.41068) (11.78694,3.40778) (11.80000,3.40477) (11.81306,3.40166) (11.82611,3.39844) (11.83917,3.39511) (11.85222,3.39169) (11.86528,3.38815) (11.87833,3.38451) (11.89139,3.38077) (11.90444,3.37692) (11.91750,3.37297) (11.93056,3.36892) (11.94361,3.36476) (11.95667,3.36049) (11.96972,3.35613) (11.98278,3.35165) (11.99583,3.34708) (12.00889,3.34240) (12.02194,3.33762) (12.03500,3.33274) (12.04806,3.32776) (12.06111,3.32267) (12.07417,3.31748) (12.08722,3.31219) (12.10028,3.30680) (12.11333,3.30131) (12.12639,3.29573) (12.13944,3.29004) (12.15250,3.28425) (12.16556,3.27836) (12.17861,3.27238) (12.19167,3.26630) (12.20472,3.26012) (12.21778,3.25385) (12.23083,3.24748) (12.24389,3.24102) (12.25694,3.23446) (12.27000,3.22781) (12.28306,3.22107) (12.29611,3.21423) (12.30917,3.20730) (12.32222,3.20028) (12.33528,3.19317) (12.34833,3.18598) (12.36139,3.17869) (12.37444,3.17132) (12.38750,3.16386) (12.40056,3.15631) (12.41361,3.14868) (12.42667,3.14096) (12.43972,3.13317) (12.45278,3.12529) (12.46583,3.11733) (12.47889,3.10928) (12.49194,3.10116) (12.50500,3.09297) (12.51806,3.08469) (12.53111,3.07634) (12.54417,3.06792) (12.55722,3.05942) (12.57028,3.05085) (12.58333,3.04221) (12.59639,3.03350) (12.60944,3.02472) (12.62250,3.01587) (12.63556,3.00696) (12.64861,2.99798) (12.66167,2.98894) (12.67472,2.97984) (12.68778,2.97068) (12.70083,2.96146) (12.71389,2.95219) (12.72694,2.94286) (12.74000,2.93347) (12.75306,2.92403) (12.76611,2.91454) (12.77917,2.90500) (12.79222,2.89542) (12.80528,2.88579) (12.81833,2.87611) (12.83139,2.86640) (12.84444,2.85664) (12.85750,2.84684) (12.87056,2.83701) (12.88361,2.82714) (12.89667,2.81725) (12.90972,2.80732) (12.92278,2.79736) (12.93583,2.78738) (12.94889,2.77737) (12.96194,2.76734) (12.97500,2.75729) (12.98806,2.74722) (13.00111,2.73714) (13.01417,2.72705) (13.02722,2.71694) (13.04028,2.70683) (13.05333,2.69671) (13.06639,2.68659) (13.07944,2.67647) (13.09250,2.66635) (13.10556,2.65623) (13.11861,2.64612) (13.13167,2.63603) (13.14472,2.62594) (13.15778,2.61588) (13.17083,2.60583) (13.18389,2.59580) (13.19694,2.58580) (13.21000,2.57583) (13.22306,2.56588) (13.23611,2.55597) (13.24917,2.54610) (13.26222,2.53627) (13.27528,2.52648) (13.28833,2.51674) (13.30139,2.50705) (13.31444,2.49742) (13.32750,2.48784) (13.34056,2.47832) (13.35361,2.46887) (13.36667,2.45949) (13.37972,2.45018) (13.39278,2.44094) (13.40583,2.43178) (13.41889,2.42271) (13.43194,2.41372) (13.44500,2.40483) (13.45806,2.39603) (13.47111,2.38732) (13.48417,2.37872) (13.49722,2.37023) (13.51028,2.36185) (13.52333,2.35358) (13.53639,2.34544) (13.54944,2.33741) (13.56250,2.32952) (13.57556,2.32175) (13.58861,2.31412) (13.60167,2.30663) (13.61472,2.29928) (13.62778,2.29209) (13.64083,2.28504) (13.65389,2.27815) (13.66694,2.27142) (13.68000,2.26486) (13.69306,2.25846) (13.70611,2.25224) (13.71917,2.24619) (13.73222,2.24032) (13.74528,2.23464) (13.75833,2.22915) (13.77139,2.22385) (13.78444,2.21874) (13.79750,2.21383) (13.81056,2.20913) (13.82361,2.20463) (13.83667,2.20034) (13.84972,2.19626) (13.86278,2.19240) (13.87583,2.18875) (13.88889,2.18533) (13.90194,2.18213) (13.91500,2.17915) (13.92806,2.17641) (13.94111,2.17389) (13.95417,2.17161) (13.96722,2.16956) (13.98028,2.16775) (13.99333,2.16617) (14.00639,2.16484) (14.01944,2.16374) (14.03250,2.16288) (14.04556,2.16227) (14.05861,2.16189) (14.07167,2.16176) (14.08472,2.16187) (14.09778,2.16222) (14.11083,2.16282) (14.12389,2.16365) (14.13694,2.16473) (14.15000,2.16605) (14.16306,2.16760) (14.17611,2.16939) (14.18917,2.17142) (14.20222,2.17369) (14.21528,2.17618) (14.22833,2.17891) (14.24139,2.18186) (14.25444,2.18504) (14.26750,2.18845) (14.28056,2.19207) (14.29361,2.19592) (14.30667,2.19998) (14.31972,2.20425) (14.33278,2.20873) (14.34583,2.21342) (14.35889,2.21831) (14.37194,2.22340) (14.38500,2.22868) (14.39806,2.23416) (14.41111,2.23983) (14.42417,2.24568) (14.43722,2.25171) (14.45028,2.25792) (14.46333,2.26430) (14.47639,2.27085) (14.48944,2.27756) (14.50250,2.28444) (14.51556,2.29147) (14.52861,2.29866) (14.54167,2.30599) (14.55472,2.31347) (14.56778,2.32109) (14.58083,2.32884) (14.59389,2.33673) (14.60694,2.34474) (14.62000,2.35288) (14.63306,2.36113) (14.64611,2.36950) (14.65917,2.37799) (14.67222,2.38658) (14.68528,2.39527) (14.69833,2.40406) (14.71139,2.41295) (14.72444,2.42193) (14.73750,2.43100) (14.75056,2.44015) (14.76361,2.44938) (14.77667,2.45868) (14.78972,2.46806) (14.80278,2.47751) (14.81583,2.48702) (14.82889,2.49659) (14.84194,2.50622) (14.85500,2.51591) (14.86806,2.52564) (14.88111,2.53542) (14.89417,2.54525) (14.90722,2.55512) (14.92028,2.56503) (14.93333,2.57497) (14.94639,2.58494) (14.95944,2.59494) (14.97250,2.60496) (14.98556,2.61501) (14.99861,2.62508) (15.01167,2.63516) (15.02472,2.64525) (15.03778,2.65536) (15.05083,2.66547) (15.06389,2.67559) (15.07694,2.68571) (15.09000,2.69584) (15.10306,2.70596) (15.11611,2.71607) (15.12917,2.72617) (15.14222,2.73627) (15.15528,2.74635) (15.16833,2.75642) (15.18139,2.76647) (15.19444,2.77650) (15.20750,2.78651) (15.22056,2.79650) (15.23361,2.80646) (15.24667,2.81639) (15.25972,2.82629) (15.27278,2.83616) (15.28583,2.84600) (15.29889,2.85579) (15.31194,2.86556) (15.32500,2.87528) (15.33806,2.88495) (15.35111,2.89459) (15.36417,2.90418) (15.37722,2.91372) (15.39028,2.92321) (15.40333,2.93266) (15.41639,2.94205) (15.42944,2.95138) (15.44250,2.96067) (15.45556,2.96989) (15.46861,2.97906) (15.48167,2.98816) (15.49472,2.99721) (15.50778,3.00619) (15.52083,3.01511) (15.53389,3.02396) (15.54694,3.03274) (15.56000,3.04146) (15.57306,3.05011) (15.58611,3.05868) (15.59917,3.06719) (15.61222,3.07562) (15.62528,3.08397) (15.63833,3.09226) (15.65139,3.10046) (15.66444,3.10859) (15.67750,3.11663) (15.69056,3.12460) (15.70361,3.13249) (15.71667,3.14029) (15.72972,3.14802) (15.74278,3.15565) (15.75583,3.16321) (15.76889,3.17068) (15.78194,3.17806) (15.79500,3.18535) (15.80806,3.19256) (15.82111,3.19967) (15.83417,3.20670) (15.84722,3.21364) (15.86028,3.22048) (15.87333,3.22723) (15.88639,3.23389) (15.89944,3.24046) (15.91250,3.24693) (15.92556,3.25331) (15.93861,3.25959) (15.95167,3.26577) (15.96472,3.27186) (15.97778,3.27785) (15.99083,3.28375) (16.00389,3.28954) (16.01694,3.29524) (16.03000,3.30084) (16.04306,3.30633) (16.05611,3.31173) (16.06917,3.31703) (16.08222,3.32223) (16.09528,3.32732) (16.10833,3.33231) (16.12139,3.33720) (16.13444,3.34199) (16.14750,3.34668) (16.16056,3.35126) (16.17361,3.35574) (16.18667,3.36012) (16.19972,3.36439) (16.21278,3.36856) (16.22583,3.37263) (16.23889,3.37659) (16.25194,3.38044) (16.26500,3.38420) (16.27806,3.38784) (16.29111,3.39139) (16.30417,3.39482) (16.31722,3.39816) (16.33028,3.40138) (16.34333,3.40451) (16.35639,3.40752) (16.36944,3.41044) (16.38250,3.41324) (16.39556,3.41595) (16.40861,3.41855) (16.42167,3.42104) (16.43472,3.42343) (16.44778,3.42571) (16.46083,3.42789) (16.47389,3.42997) (16.48694,3.43194) (16.50000,3.43381)};
\draw[->,coursemuted] plot coordinates {(7.10000,4.65000) (16.60000,4.65000)};
\draw[->,coursemuted] plot coordinates {(7.10000,4.65000) (7.10000,6.50000)};
\draw[coursemuted!25,thin] plot coordinates {(7.10000,4.65000) (16.50000,4.65000)};
\node[left,inner sep=3pt] at (7.1,4.65000) {$0$};
\draw[coursemuted!25,thin] plot coordinates {(7.10000,5.45000) (16.50000,5.45000)};
\node[left,inner sep=3pt] at (7.1,5.45000) {$8{,}5$};
\draw[coursemuted!25,thin] plot coordinates {(7.10000,6.25000) (16.50000,6.25000)};
\node[left,inner sep=3pt] at (7.1,6.25000) {$17$};
\node[rotate=90] at (5.9,5.55) {$\Delta V_j$, м/с};
\draw[coursemuted!16,thin] plot coordinates {(7.10000,0.80000) (7.10000,6.25000)};
\node[below,inner sep=3pt] at (7.10000,0.85) {$0$};
\draw[coursemuted!16,thin] plot coordinates {(9.45000,0.80000) (9.45000,6.25000)};
\node[below,inner sep=3pt] at (9.45000,0.85) {$\pi/2$};
\draw[coursemuted!16,thin] plot coordinates {(11.80000,0.80000) (11.80000,6.25000)};
\node[below,inner sep=3pt] at (11.80000,0.85) {$\pi$};
\draw[coursemuted!16,thin] plot coordinates {(14.15000,0.80000) (14.15000,6.25000)};
\node[below,inner sep=3pt] at (14.15000,0.85) {$3\pi/2$};
\draw[coursemuted!16,thin] plot coordinates {(16.50000,0.80000) (16.50000,6.25000)};
\node[below,inner sep=3pt] at (16.50000,0.85) {$2\pi$};
\node at (11.8,.15) {Угол $\varphi$, рад};
\draw[courseink!70,densely dashed,thin] plot coordinates {(7.33545,0.85000) (7.33545,3.45000)};
\draw[courseink,line width=1pt] plot coordinates {(7.33545,4.65000) (7.33545,5.91739)};
\fill[courseink] (7.33545,5.91739) circle[radius=1.8pt];
\fill[courseink] (7.33545,3.45000) circle[radius=1.8pt];
\node[above,inner sep=2pt] at (7.33545,5.91739) {1};
\draw[courseink!70,densely dashed,thin] plot coordinates {(11.40901,0.85000) (11.40901,3.45000)};
\draw[courseink,line width=1pt] plot coordinates {(11.40901,4.65000) (11.40901,5.41879)};
\fill[courseink] (11.40901,5.41879) circle[radius=1.8pt];
\fill[courseink] (11.40901,3.45000) circle[radius=1.8pt];
\node[above,inner sep=2pt] at (11.40901,5.41879) {2};
\end{tikzpicture}

\par
\figcaption{$\Delta a=0{,}004$, $\Delta e_x=0{,}001$, $\Delta e_y=0{,}001$, $\Delta i=0{,}002$ рад.\\
Годограф показан в своей плоскости; оси --- в косоугольной проекции.}
\end{frame}

% Источник: прямой расчёт направлений после eq:ndc-cost, строки 1032--1045.
\begin{frame}{Невырожденное решение}
\relax
\derivationmathspacing
Сначала вычисляем отношение компонент направления:
\[
\frac{d_z}{d_r}=-\frac{\Delta i\,\Delta e_y}
{J^2-(\Delta a)^2/4-(\Delta i)^2}.
\]
\[
\begin{aligned}
d_t&=\frac{\Delta a}{2J},\\
d_r&=\sqrt{\frac{1-d_t^2}{1+(d_z/d_r)^2}}>0,\\
d_z&=-\frac{\Delta i\,\Delta e_y\,d_r}
{J^2-(\Delta a)^2/4-(\Delta i)^2}.
\end{aligned}
\]
\medskip
Направления двух импульсов:
$(d_r,d_t,d_z)$ и $(-d_r,d_t,-d_z)$.
\end{frame}

% Источник: eq:ndc-angles-via-directions, eq:ndc-symmetry.
\begin{frame}{Невырожденное решение}
\relax
\derivationmathspacing
После вычисления $d_z/d_r$ находим углы:
\[
\begin{aligned}
\cos\alpha&=\frac{\Delta e_x}
{\sqrt{(\Delta e_x)^2+(\Delta e_y-\Delta i\,d_z/d_r)^2}},\\
\sin\alpha&=\frac{\Delta e_y-\Delta i\,d_z/d_r}
{\sqrt{(\Delta e_x)^2+(\Delta e_y-\Delta i\,d_z/d_r)^2}},\\
\cos\beta&=\frac{\Delta a}
{\sqrt{(\Delta e_x)^2+(\Delta e_y-\Delta i\,d_z/d_r)^2}}.
\end{aligned}
\]
\[
0<\beta<\pi,\qquad
\varphi_+=\alpha+\beta,\qquad \varphi_-=\alpha-\beta.
\]
\medskip
Угол $\alpha$ определяется совместно по синусу и косинусу;
углы $\varphi_\pm$ приводятся по модулю $2\pi$.
\end{frame}

% Источник: eq:ndc-impulses; замечания после eq:ndc-directions.
\begin{frame}{Невырожденное решение}
\relax
\derivationmathspacing
\[
\begin{aligned}
\Delta\widehat V_+&=\frac J2\left[
1-\frac{\cos\alpha\sin\beta}{\sin\alpha\cos\beta}\right],\\
\Delta\widehat V_-&=\frac J2\left[
1+\frac{\cos\alpha\sin\beta}{\sin\alpha\cos\beta}\right].
\end{aligned}
\]
\[
\begin{aligned}
\Delta\widehat{\boldsymbol V}_+
&=\Delta\widehat V_+(d_r,d_t,d_z),\\
\Delta\widehat{\boldsymbol V}_-
&=\Delta\widehat V_-(-d_r,d_t,-d_z).
\end{aligned}
\]
\medskip
\justifying
В указанной области оба модуля положительны.
При $\Delta e_x=0$ они равны. Порядок реализации выбирается после расчёта.
\end{frame}

% Источник: sec:line, eq:line-primer, eq:line-impulses.
\begin{frame}{Нулевые радиальные компоненты}
\relax
\derivationmathspacing
Пусть $\Delta e_y=0$ и $|\Delta a|>|\Delta e_x|$.
\[
\lambda_2=\lambda_3=0.
\]
{\justifying
Радиальная компонента базис-вектора исчезает, трансверсальная
становится постоянной, а нормальная остаётся гармонической.
Эллипс вырождается в отрезок, параллельный оси $s_z$.
Его концы лежат на единичной сфере.
\par\medskip
В концах отрезка трансверсальные компоненты одинаковы,
нормальные противоположны. Концы достигаются через пол-оборота;
импульсы прикладываются в узлах:
\par}
\[
\varphi_1=0,\qquad\varphi_2=\pi.
\]
\end{frame}

\begin{frame}{Нулевые радиальные компоненты: годограф}
\relax
\centering
\begin{tikzpicture}[course diagram,scale=.72]
\node at (2.65,6.85) {Годограф};
\node at (11.6,6.85) {Модули импульсов};
\node at (11.6,4.1) {Модуль базис-вектора};
\filldraw[fill=coursetint!45!coursepaper,draw=coursemuted!65,line width=.65pt] (2.65000,3.60000) circle[radius=1.7cm];
\draw[->,courseink!75,thin] plot coordinates {(2.65000,3.60000) (0.80500,2.32269)};
\node[below left,inner sep=2pt] at (0.80500,2.32269) {$s_r$};
\draw[->,courseink!75,thin] plot coordinates {(2.65000,3.60000) (4.89400,3.60000)};
\node[right,inner sep=2pt] at (4.89400,3.60000) {$s_t$};
\draw[->,courseink!75,thin] plot coordinates {(2.65000,3.60000) (2.65000,5.84400)};
\node[above,inner sep=2pt] at (2.65000,5.84400) {$s_z$};
\draw[courseaccent,line width=1.2pt] plot coordinates {(3.85208,4.80208) (3.85208,4.80190) (3.85208,4.80135) (3.85208,4.80043) (3.85208,4.79915) (3.85208,4.79751) (3.85208,4.79550) (3.85208,4.79312) (3.85208,4.79038) (3.85208,4.78728) (3.85208,4.78382) (3.85208,4.78000) (3.85208,4.77581) (3.85208,4.77127) (3.85208,4.76637) (3.85208,4.76112) (3.85208,4.75551) (3.85208,4.74956) (3.85208,4.74325) (3.85208,4.73659) (3.85208,4.72959) (3.85208,4.72224) (3.85208,4.71455) (3.85208,4.70652) (3.85208,4.69816) (3.85208,4.68946) (3.85208,4.68042) (3.85208,4.67106) (3.85208,4.66137) (3.85208,4.65136) (3.85208,4.64103) (3.85208,4.63038) (3.85208,4.61942) (3.85208,4.60815) (3.85208,4.59657) (3.85208,4.58469) (3.85208,4.57250) (3.85208,4.56002) (3.85208,4.54725) (3.85208,4.53419) (3.85208,4.52085) (3.85208,4.50722) (3.85208,4.49332) (3.85208,4.47915) (3.85208,4.46471) (3.85208,4.45000) (3.85208,4.43504) (3.85208,4.41982) (3.85208,4.40435) (3.85208,4.38864) (3.85208,4.37268) (3.85208,4.35649) (3.85208,4.34008) (3.85208,4.32343) (3.85208,4.30657) (3.85208,4.28949) (3.85208,4.27220) (3.85208,4.25470) (3.85208,4.23701) (3.85208,4.21912) (3.85208,4.20104) (3.85208,4.18278) (3.85208,4.16434) (3.85208,4.14573) (3.85208,4.12696) (3.85208,4.10802) (3.85208,4.08893) (3.85208,4.06969) (3.85208,4.05031) (3.85208,4.03079) (3.85208,4.01114) (3.85208,3.99136) (3.85208,3.97146) (3.85208,3.95145) (3.85208,3.93134) (3.85208,3.91112) (3.85208,3.89081) (3.85208,3.87041) (3.85208,3.84993) (3.85208,3.82937) (3.85208,3.80874) (3.85208,3.78805) (3.85208,3.76730) (3.85208,3.74650) (3.85208,3.72565) (3.85208,3.70477) (3.85208,3.68385) (3.85208,3.66291) (3.85208,3.64195) (3.85208,3.62098) (3.85208,3.60000) (3.85208,3.57902) (3.85208,3.55805) (3.85208,3.53709) (3.85208,3.51615) (3.85208,3.49523) (3.85208,3.47435) (3.85208,3.45350) (3.85208,3.43270) (3.85208,3.41195) (3.85208,3.39126) (3.85208,3.37063) (3.85208,3.35007) (3.85208,3.32959) (3.85208,3.30919) (3.85208,3.28888) (3.85208,3.26866) (3.85208,3.24855) (3.85208,3.22854) (3.85208,3.20864) (3.85208,3.18886) (3.85208,3.16921) (3.85208,3.14969) (3.85208,3.13031) (3.85208,3.11107) (3.85208,3.09198) (3.85208,3.07304) (3.85208,3.05427) (3.85208,3.03566) (3.85208,3.01722) (3.85208,2.99896) (3.85208,2.98088) (3.85208,2.96299) (3.85208,2.94530) (3.85208,2.92780) (3.85208,2.91051) (3.85208,2.89343) (3.85208,2.87657) (3.85208,2.85992) (3.85208,2.84351) (3.85208,2.82732) (3.85208,2.81136) (3.85208,2.79565) (3.85208,2.78018) (3.85208,2.76496) (3.85208,2.75000) (3.85208,2.73529) (3.85208,2.72085) (3.85208,2.70668) (3.85208,2.69278) (3.85208,2.67915) (3.85208,2.66581) (3.85208,2.65275) (3.85208,2.63998) (3.85208,2.62750) (3.85208,2.61531) (3.85208,2.60343) (3.85208,2.59185) (3.85208,2.58058) (3.85208,2.56962) (3.85208,2.55897) (3.85208,2.54864) (3.85208,2.53863) (3.85208,2.52894) (3.85208,2.51958) (3.85208,2.51054) (3.85208,2.50184) (3.85208,2.49348) (3.85208,2.48545) (3.85208,2.47776) (3.85208,2.47041) (3.85208,2.46341) (3.85208,2.45675) (3.85208,2.45044) (3.85208,2.44449) (3.85208,2.43888) (3.85208,2.43363) (3.85208,2.42873) (3.85208,2.42419) (3.85208,2.42000) (3.85208,2.41618) (3.85208,2.41272) (3.85208,2.40962) (3.85208,2.40688) (3.85208,2.40450) (3.85208,2.40249) (3.85208,2.40085) (3.85208,2.39957) (3.85208,2.39865) (3.85208,2.39810) (3.85208,2.39792) (3.85208,2.39810) (3.85208,2.39865) (3.85208,2.39957) (3.85208,2.40085) (3.85208,2.40249) (3.85208,2.40450) (3.85208,2.40688) (3.85208,2.40962) (3.85208,2.41272) (3.85208,2.41618) (3.85208,2.42000) (3.85208,2.42419) (3.85208,2.42873) (3.85208,2.43363) (3.85208,2.43888) (3.85208,2.44449) (3.85208,2.45044) (3.85208,2.45675) (3.85208,2.46341) (3.85208,2.47041) (3.85208,2.47776) (3.85208,2.48545) (3.85208,2.49348) (3.85208,2.50184) (3.85208,2.51054) (3.85208,2.51958) (3.85208,2.52894) (3.85208,2.53863) (3.85208,2.54864) (3.85208,2.55897) (3.85208,2.56962) (3.85208,2.58058) (3.85208,2.59185) (3.85208,2.60343) (3.85208,2.61531) (3.85208,2.62750) (3.85208,2.63998) (3.85208,2.65275) (3.85208,2.66581) (3.85208,2.67915) (3.85208,2.69278) (3.85208,2.70668) (3.85208,2.72085) (3.85208,2.73529) (3.85208,2.75000) (3.85208,2.76496) (3.85208,2.78018) (3.85208,2.79565) (3.85208,2.81136) (3.85208,2.82732) (3.85208,2.84351) (3.85208,2.85992) (3.85208,2.87657) (3.85208,2.89343) (3.85208,2.91051) (3.85208,2.92780) (3.85208,2.94530) (3.85208,2.96299) (3.85208,2.98088) (3.85208,2.99896) (3.85208,3.01722) (3.85208,3.03566) (3.85208,3.05427) (3.85208,3.07304) (3.85208,3.09198) (3.85208,3.11107) (3.85208,3.13031) (3.85208,3.14969) (3.85208,3.16921) (3.85208,3.18886) (3.85208,3.20864) (3.85208,3.22854) (3.85208,3.24855) (3.85208,3.26866) (3.85208,3.28888) (3.85208,3.30919) (3.85208,3.32959) (3.85208,3.35007) (3.85208,3.37063) (3.85208,3.39126) (3.85208,3.41195) (3.85208,3.43270) (3.85208,3.45350) (3.85208,3.47435) (3.85208,3.49523) (3.85208,3.51615) (3.85208,3.53709) (3.85208,3.55805) (3.85208,3.57902) (3.85208,3.60000) (3.85208,3.62098) (3.85208,3.64195) (3.85208,3.66291) (3.85208,3.68385) (3.85208,3.70477) (3.85208,3.72565) (3.85208,3.74650) (3.85208,3.76730) (3.85208,3.78805) (3.85208,3.80874) (3.85208,3.82937) (3.85208,3.84993) (3.85208,3.87041) (3.85208,3.89081) (3.85208,3.91112) (3.85208,3.93134) (3.85208,3.95145) (3.85208,3.97146) (3.85208,3.99136) (3.85208,4.01114) (3.85208,4.03079) (3.85208,4.05031) (3.85208,4.06969) (3.85208,4.08893) (3.85208,4.10802) (3.85208,4.12696) (3.85208,4.14573) (3.85208,4.16434) (3.85208,4.18278) (3.85208,4.20104) (3.85208,4.21912) (3.85208,4.23701) (3.85208,4.25470) (3.85208,4.27220) (3.85208,4.28949) (3.85208,4.30657) (3.85208,4.32343) (3.85208,4.34008) (3.85208,4.35649) (3.85208,4.37268) (3.85208,4.38864) (3.85208,4.40435) (3.85208,4.41982) (3.85208,4.43504) (3.85208,4.45000) (3.85208,4.46471) (3.85208,4.47915) (3.85208,4.49332) (3.85208,4.50722) (3.85208,4.52085) (3.85208,4.53419) (3.85208,4.54725) (3.85208,4.56002) (3.85208,4.57250) (3.85208,4.58469) (3.85208,4.59657) (3.85208,4.60815) (3.85208,4.61942) (3.85208,4.63038) (3.85208,4.64103) (3.85208,4.65136) (3.85208,4.66137) (3.85208,4.67106) (3.85208,4.68042) (3.85208,4.68946) (3.85208,4.69816) (3.85208,4.70652) (3.85208,4.71455) (3.85208,4.72224) (3.85208,4.72959) (3.85208,4.73659) (3.85208,4.74325) (3.85208,4.74956) (3.85208,4.75551) (3.85208,4.76112) (3.85208,4.76637) (3.85208,4.77127) (3.85208,4.77581) (3.85208,4.78000) (3.85208,4.78382) (3.85208,4.78728) (3.85208,4.79038) (3.85208,4.79312) (3.85208,4.79550) (3.85208,4.79751) (3.85208,4.79915) (3.85208,4.80043) (3.85208,4.80135) (3.85208,4.80190) (3.85208,4.80208)};
\node[courseaccent,inner sep=0pt] at (3.85208,3.60000) {$\times$};
\draw[->,courseaccent,line width=.7pt] plot coordinates {(2.65000,3.60000) (3.85208,4.80208)};
\filldraw[fill=courseink,draw=coursepaper,line width=.4pt] (3.85208,4.80208) circle[radius=2.3pt];
\node[above right,inner sep=3pt] at (3.85208,4.80208) {1};
\draw[->,courseaccent,line width=.7pt] plot coordinates {(2.65000,3.60000) (3.85208,2.39792)};
\filldraw[fill=courseink,draw=coursepaper,line width=.4pt] (3.85208,2.39792) circle[radius=2.3pt];
\node[below right,inner sep=3pt] at (3.85208,2.39792) {2};
\node[below left,inner sep=2pt] at (2.65,3.6) {$0$};
\draw[coursemuted!25,thin] plot coordinates {(7.10000,0.85000) (16.50000,0.85000)};
\node[left,inner sep=3pt] at (7.1,0.85) {$0$};
\draw[coursemuted!25,thin] plot coordinates {(7.10000,2.15000) (16.50000,2.15000)};
\node[left,inner sep=3pt] at (7.1,2.15) {$0{,}5$};
\draw[courseink!65,dashed] plot coordinates {(7.10000,3.45000) (16.50000,3.45000)};
\node[left,inner sep=3pt] at (7.1,3.45) {$1$};
\draw[->,coursemuted] plot coordinates {(7.10000,0.85000) (16.60000,0.85000)};
\draw[->,coursemuted] plot coordinates {(7.10000,0.85000) (7.10000,3.84000)};
\node[rotate=90] at (5.9,2.2) {$\|\boldsymbol s(\varphi)\|$};
\draw[courseaccent,line width=1.15pt] plot coordinates {(7.10000,3.45000) (7.11306,3.44995) (7.12611,3.44980) (7.13917,3.44955) (7.15222,3.44921) (7.16528,3.44876) (7.17833,3.44822) (7.19139,3.44758) (7.20444,3.44684) (7.21750,3.44600) (7.23056,3.44506) (7.24361,3.44402) (7.25667,3.44289) (7.26972,3.44166) (7.28278,3.44033) (7.29583,3.43890) (7.30889,3.43738) (7.32194,3.43576) (7.33500,3.43404) (7.34806,3.43223) (7.36111,3.43033) (7.37417,3.42832) (7.38722,3.42623) (7.40028,3.42403) (7.41333,3.42175) (7.42639,3.41937) (7.43944,3.41690) (7.45250,3.41433) (7.46556,3.41168) (7.47861,3.40893) (7.49167,3.40609) (7.50472,3.40316) (7.51778,3.40014) (7.53083,3.39703) (7.54389,3.39383) (7.55694,3.39054) (7.57000,3.38717) (7.58306,3.38371) (7.59611,3.38017) (7.60917,3.37653) (7.62222,3.37282) (7.63528,3.36902) (7.64833,3.36514) (7.66139,3.36117) (7.67444,3.35713) (7.68750,3.35300) (7.70056,3.34879) (7.71361,3.34451) (7.72667,3.34015) (7.73972,3.33571) (7.75278,3.33119) (7.76583,3.32660) (7.77889,3.32194) (7.79194,3.31720) (7.80500,3.31239) (7.81806,3.30751) (7.83111,3.30256) (7.84417,3.29754) (7.85722,3.29245) (7.87028,3.28730) (7.88333,3.28208) (7.89639,3.27679) (7.90944,3.27145) (7.92250,3.26604) (7.93556,3.26057) (7.94861,3.25504) (7.96167,3.24946) (7.97472,3.24381) (7.98778,3.23811) (8.00083,3.23236) (8.01389,3.22656) (8.02694,3.22070) (8.04000,3.21479) (8.05306,3.20884) (8.06611,3.20283) (8.07917,3.19679) (8.09222,3.19069) (8.10528,3.18456) (8.11833,3.17838) (8.13139,3.17217) (8.14444,3.16592) (8.15750,3.15963) (8.17056,3.15330) (8.18361,3.14695) (8.19667,3.14056) (8.20972,3.13414) (8.22278,3.12769) (8.23583,3.12122) (8.24889,3.11473) (8.26194,3.10821) (8.27500,3.10167) (8.28806,3.09511) (8.30111,3.08853) (8.31417,3.08194) (8.32722,3.07533) (8.34028,3.06872) (8.35333,3.06209) (8.36639,3.05546) (8.37944,3.04882) (8.39250,3.04217) (8.40556,3.03553) (8.41861,3.02888) (8.43167,3.02224) (8.44472,3.01560) (8.45778,3.00897) (8.47083,3.00235) (8.48389,2.99573) (8.49694,2.98913) (8.51000,2.98255) (8.52306,2.97598) (8.53611,2.96943) (8.54917,2.96290) (8.56222,2.95640) (8.57528,2.94992) (8.58833,2.94347) (8.60139,2.93705) (8.61444,2.93066) (8.62750,2.92431) (8.64056,2.91799) (8.65361,2.91172) (8.66667,2.90548) (8.67972,2.89929) (8.69278,2.89314) (8.70583,2.88705) (8.71889,2.88100) (8.73194,2.87501) (8.74500,2.86907) (8.75806,2.86319) (8.77111,2.85737) (8.78417,2.85161) (8.79722,2.84592) (8.81028,2.84029) (8.82333,2.83473) (8.83639,2.82925) (8.84944,2.82384) (8.86250,2.81850) (8.87556,2.81324) (8.88861,2.80806) (8.90167,2.80297) (8.91472,2.79796) (8.92778,2.79303) (8.94083,2.78820) (8.95389,2.78346) (8.96694,2.77881) (8.98000,2.77426) (8.99306,2.76980) (9.00611,2.76544) (9.01917,2.76119) (9.03222,2.75704) (9.04528,2.75299) (9.05833,2.74906) (9.07139,2.74523) (9.08444,2.74151) (9.09750,2.73791) (9.11056,2.73442) (9.12361,2.73105) (9.13667,2.72779) (9.14972,2.72466) (9.16278,2.72165) (9.17583,2.71876) (9.18889,2.71599) (9.20194,2.71335) (9.21500,2.71084) (9.22806,2.70845) (9.24111,2.70620) (9.25417,2.70407) (9.26722,2.70208) (9.28028,2.70022) (9.29333,2.69849) (9.30639,2.69690) (9.31944,2.69545) (9.33250,2.69413) (9.34556,2.69295) (9.35861,2.69190) (9.37167,2.69099) (9.38472,2.69023) (9.39778,2.68960) (9.41083,2.68911) (9.42389,2.68876) (9.43694,2.68855) (9.45000,2.68848) (9.46306,2.68855) (9.47611,2.68876) (9.48917,2.68911) (9.50222,2.68960) (9.51528,2.69023) (9.52833,2.69099) (9.54139,2.69190) (9.55444,2.69295) (9.56750,2.69413) (9.58056,2.69545) (9.59361,2.69690) (9.60667,2.69849) (9.61972,2.70022) (9.63278,2.70208) (9.64583,2.70407) (9.65889,2.70620) (9.67194,2.70845) (9.68500,2.71084) (9.69806,2.71335) (9.71111,2.71599) (9.72417,2.71876) (9.73722,2.72165) (9.75028,2.72466) (9.76333,2.72779) (9.77639,2.73105) (9.78944,2.73442) (9.80250,2.73791) (9.81556,2.74151) (9.82861,2.74523) (9.84167,2.74906) (9.85472,2.75299) (9.86778,2.75704) (9.88083,2.76119) (9.89389,2.76544) (9.90694,2.76980) (9.92000,2.77426) (9.93306,2.77881) (9.94611,2.78346) (9.95917,2.78820) (9.97222,2.79303) (9.98528,2.79796) (9.99833,2.80297) (10.01139,2.80806) (10.02444,2.81324) (10.03750,2.81850) (10.05056,2.82384) (10.06361,2.82925) (10.07667,2.83473) (10.08972,2.84029) (10.10278,2.84592) (10.11583,2.85161) (10.12889,2.85737) (10.14194,2.86319) (10.15500,2.86907) (10.16806,2.87501) (10.18111,2.88100) (10.19417,2.88705) (10.20722,2.89314) (10.22028,2.89929) (10.23333,2.90548) (10.24639,2.91172) (10.25944,2.91799) (10.27250,2.92431) (10.28556,2.93066) (10.29861,2.93705) (10.31167,2.94347) (10.32472,2.94992) (10.33778,2.95640) (10.35083,2.96290) (10.36389,2.96943) (10.37694,2.97598) (10.39000,2.98255) (10.40306,2.98913) (10.41611,2.99573) (10.42917,3.00235) (10.44222,3.00897) (10.45528,3.01560) (10.46833,3.02224) (10.48139,3.02888) (10.49444,3.03553) (10.50750,3.04217) (10.52056,3.04882) (10.53361,3.05546) (10.54667,3.06209) (10.55972,3.06872) (10.57278,3.07533) (10.58583,3.08194) (10.59889,3.08853) (10.61194,3.09511) (10.62500,3.10167) (10.63806,3.10821) (10.65111,3.11473) (10.66417,3.12122) (10.67722,3.12769) (10.69028,3.13414) (10.70333,3.14056) (10.71639,3.14695) (10.72944,3.15330) (10.74250,3.15963) (10.75556,3.16592) (10.76861,3.17217) (10.78167,3.17838) (10.79472,3.18456) (10.80778,3.19069) (10.82083,3.19679) (10.83389,3.20283) (10.84694,3.20884) (10.86000,3.21479) (10.87306,3.22070) (10.88611,3.22656) (10.89917,3.23236) (10.91222,3.23811) (10.92528,3.24381) (10.93833,3.24946) (10.95139,3.25504) (10.96444,3.26057) (10.97750,3.26604) (10.99056,3.27145) (11.00361,3.27679) (11.01667,3.28208) (11.02972,3.28730) (11.04278,3.29245) (11.05583,3.29754) (11.06889,3.30256) (11.08194,3.30751) (11.09500,3.31239) (11.10806,3.31720) (11.12111,3.32194) (11.13417,3.32660) (11.14722,3.33119) (11.16028,3.33571) (11.17333,3.34015) (11.18639,3.34451) (11.19944,3.34879) (11.21250,3.35300) (11.22556,3.35713) (11.23861,3.36117) (11.25167,3.36514) (11.26472,3.36902) (11.27778,3.37282) (11.29083,3.37653) (11.30389,3.38017) (11.31694,3.38371) (11.33000,3.38717) (11.34306,3.39054) (11.35611,3.39383) (11.36917,3.39703) (11.38222,3.40014) (11.39528,3.40316) (11.40833,3.40609) (11.42139,3.40893) (11.43444,3.41168) (11.44750,3.41433) (11.46056,3.41690) (11.47361,3.41937) (11.48667,3.42175) (11.49972,3.42403) (11.51278,3.42623) (11.52583,3.42832) (11.53889,3.43033) (11.55194,3.43223) (11.56500,3.43404) (11.57806,3.43576) (11.59111,3.43738) (11.60417,3.43890) (11.61722,3.44033) (11.63028,3.44166) (11.64333,3.44289) (11.65639,3.44402) (11.66944,3.44506) (11.68250,3.44600) (11.69556,3.44684) (11.70861,3.44758) (11.72167,3.44822) (11.73472,3.44876) (11.74778,3.44921) (11.76083,3.44955) (11.77389,3.44980) (11.78694,3.44995) (11.80000,3.45000) (11.81306,3.44995) (11.82611,3.44980) (11.83917,3.44955) (11.85222,3.44921) (11.86528,3.44876) (11.87833,3.44822) (11.89139,3.44758) (11.90444,3.44684) (11.91750,3.44600) (11.93056,3.44506) (11.94361,3.44402) (11.95667,3.44289) (11.96972,3.44166) (11.98278,3.44033) (11.99583,3.43890) (12.00889,3.43738) (12.02194,3.43576) (12.03500,3.43404) (12.04806,3.43223) (12.06111,3.43033) (12.07417,3.42832) (12.08722,3.42623) (12.10028,3.42403) (12.11333,3.42175) (12.12639,3.41937) (12.13944,3.41690) (12.15250,3.41433) (12.16556,3.41168) (12.17861,3.40893) (12.19167,3.40609) (12.20472,3.40316) (12.21778,3.40014) (12.23083,3.39703) (12.24389,3.39383) (12.25694,3.39054) (12.27000,3.38717) (12.28306,3.38371) (12.29611,3.38017) (12.30917,3.37653) (12.32222,3.37282) (12.33528,3.36902) (12.34833,3.36514) (12.36139,3.36117) (12.37444,3.35713) (12.38750,3.35300) (12.40056,3.34879) (12.41361,3.34451) (12.42667,3.34015) (12.43972,3.33571) (12.45278,3.33119) (12.46583,3.32660) (12.47889,3.32194) (12.49194,3.31720) (12.50500,3.31239) (12.51806,3.30751) (12.53111,3.30256) (12.54417,3.29754) (12.55722,3.29245) (12.57028,3.28730) (12.58333,3.28208) (12.59639,3.27679) (12.60944,3.27145) (12.62250,3.26604) (12.63556,3.26057) (12.64861,3.25504) (12.66167,3.24946) (12.67472,3.24381) (12.68778,3.23811) (12.70083,3.23236) (12.71389,3.22656) (12.72694,3.22070) (12.74000,3.21479) (12.75306,3.20884) (12.76611,3.20283) (12.77917,3.19679) (12.79222,3.19069) (12.80528,3.18456) (12.81833,3.17838) (12.83139,3.17217) (12.84444,3.16592) (12.85750,3.15963) (12.87056,3.15330) (12.88361,3.14695) (12.89667,3.14056) (12.90972,3.13414) (12.92278,3.12769) (12.93583,3.12122) (12.94889,3.11473) (12.96194,3.10821) (12.97500,3.10167) (12.98806,3.09511) (13.00111,3.08853) (13.01417,3.08194) (13.02722,3.07533) (13.04028,3.06872) (13.05333,3.06209) (13.06639,3.05546) (13.07944,3.04882) (13.09250,3.04217) (13.10556,3.03553) (13.11861,3.02888) (13.13167,3.02224) (13.14472,3.01560) (13.15778,3.00897) (13.17083,3.00235) (13.18389,2.99573) (13.19694,2.98913) (13.21000,2.98255) (13.22306,2.97598) (13.23611,2.96943) (13.24917,2.96290) (13.26222,2.95640) (13.27528,2.94992) (13.28833,2.94347) (13.30139,2.93705) (13.31444,2.93066) (13.32750,2.92431) (13.34056,2.91799) (13.35361,2.91172) (13.36667,2.90548) (13.37972,2.89929) (13.39278,2.89314) (13.40583,2.88705) (13.41889,2.88100) (13.43194,2.87501) (13.44500,2.86907) (13.45806,2.86319) (13.47111,2.85737) (13.48417,2.85161) (13.49722,2.84592) (13.51028,2.84029) (13.52333,2.83473) (13.53639,2.82925) (13.54944,2.82384) (13.56250,2.81850) (13.57556,2.81324) (13.58861,2.80806) (13.60167,2.80297) (13.61472,2.79796) (13.62778,2.79303) (13.64083,2.78820) (13.65389,2.78346) (13.66694,2.77881) (13.68000,2.77426) (13.69306,2.76980) (13.70611,2.76544) (13.71917,2.76119) (13.73222,2.75704) (13.74528,2.75299) (13.75833,2.74906) (13.77139,2.74523) (13.78444,2.74151) (13.79750,2.73791) (13.81056,2.73442) (13.82361,2.73105) (13.83667,2.72779) (13.84972,2.72466) (13.86278,2.72165) (13.87583,2.71876) (13.88889,2.71599) (13.90194,2.71335) (13.91500,2.71084) (13.92806,2.70845) (13.94111,2.70620) (13.95417,2.70407) (13.96722,2.70208) (13.98028,2.70022) (13.99333,2.69849) (14.00639,2.69690) (14.01944,2.69545) (14.03250,2.69413) (14.04556,2.69295) (14.05861,2.69190) (14.07167,2.69099) (14.08472,2.69023) (14.09778,2.68960) (14.11083,2.68911) (14.12389,2.68876) (14.13694,2.68855) (14.15000,2.68848) (14.16306,2.68855) (14.17611,2.68876) (14.18917,2.68911) (14.20222,2.68960) (14.21528,2.69023) (14.22833,2.69099) (14.24139,2.69190) (14.25444,2.69295) (14.26750,2.69413) (14.28056,2.69545) (14.29361,2.69690) (14.30667,2.69849) (14.31972,2.70022) (14.33278,2.70208) (14.34583,2.70407) (14.35889,2.70620) (14.37194,2.70845) (14.38500,2.71084) (14.39806,2.71335) (14.41111,2.71599) (14.42417,2.71876) (14.43722,2.72165) (14.45028,2.72466) (14.46333,2.72779) (14.47639,2.73105) (14.48944,2.73442) (14.50250,2.73791) (14.51556,2.74151) (14.52861,2.74523) (14.54167,2.74906) (14.55472,2.75299) (14.56778,2.75704) (14.58083,2.76119) (14.59389,2.76544) (14.60694,2.76980) (14.62000,2.77426) (14.63306,2.77881) (14.64611,2.78346) (14.65917,2.78820) (14.67222,2.79303) (14.68528,2.79796) (14.69833,2.80297) (14.71139,2.80806) (14.72444,2.81324) (14.73750,2.81850) (14.75056,2.82384) (14.76361,2.82925) (14.77667,2.83473) (14.78972,2.84029) (14.80278,2.84592) (14.81583,2.85161) (14.82889,2.85737) (14.84194,2.86319) (14.85500,2.86907) (14.86806,2.87501) (14.88111,2.88100) (14.89417,2.88705) (14.90722,2.89314) (14.92028,2.89929) (14.93333,2.90548) (14.94639,2.91172) (14.95944,2.91799) (14.97250,2.92431) (14.98556,2.93066) (14.99861,2.93705) (15.01167,2.94347) (15.02472,2.94992) (15.03778,2.95640) (15.05083,2.96290) (15.06389,2.96943) (15.07694,2.97598) (15.09000,2.98255) (15.10306,2.98913) (15.11611,2.99573) (15.12917,3.00235) (15.14222,3.00897) (15.15528,3.01560) (15.16833,3.02224) (15.18139,3.02888) (15.19444,3.03553) (15.20750,3.04217) (15.22056,3.04882) (15.23361,3.05546) (15.24667,3.06209) (15.25972,3.06872) (15.27278,3.07533) (15.28583,3.08194) (15.29889,3.08853) (15.31194,3.09511) (15.32500,3.10167) (15.33806,3.10821) (15.35111,3.11473) (15.36417,3.12122) (15.37722,3.12769) (15.39028,3.13414) (15.40333,3.14056) (15.41639,3.14695) (15.42944,3.15330) (15.44250,3.15963) (15.45556,3.16592) (15.46861,3.17217) (15.48167,3.17838) (15.49472,3.18456) (15.50778,3.19069) (15.52083,3.19679) (15.53389,3.20283) (15.54694,3.20884) (15.56000,3.21479) (15.57306,3.22070) (15.58611,3.22656) (15.59917,3.23236) (15.61222,3.23811) (15.62528,3.24381) (15.63833,3.24946) (15.65139,3.25504) (15.66444,3.26057) (15.67750,3.26604) (15.69056,3.27145) (15.70361,3.27679) (15.71667,3.28208) (15.72972,3.28730) (15.74278,3.29245) (15.75583,3.29754) (15.76889,3.30256) (15.78194,3.30751) (15.79500,3.31239) (15.80806,3.31720) (15.82111,3.32194) (15.83417,3.32660) (15.84722,3.33119) (15.86028,3.33571) (15.87333,3.34015) (15.88639,3.34451) (15.89944,3.34879) (15.91250,3.35300) (15.92556,3.35713) (15.93861,3.36117) (15.95167,3.36514) (15.96472,3.36902) (15.97778,3.37282) (15.99083,3.37653) (16.00389,3.38017) (16.01694,3.38371) (16.03000,3.38717) (16.04306,3.39054) (16.05611,3.39383) (16.06917,3.39703) (16.08222,3.40014) (16.09528,3.40316) (16.10833,3.40609) (16.12139,3.40893) (16.13444,3.41168) (16.14750,3.41433) (16.16056,3.41690) (16.17361,3.41937) (16.18667,3.42175) (16.19972,3.42403) (16.21278,3.42623) (16.22583,3.42832) (16.23889,3.43033) (16.25194,3.43223) (16.26500,3.43404) (16.27806,3.43576) (16.29111,3.43738) (16.30417,3.43890) (16.31722,3.44033) (16.33028,3.44166) (16.34333,3.44289) (16.35639,3.44402) (16.36944,3.44506) (16.38250,3.44600) (16.39556,3.44684) (16.40861,3.44758) (16.42167,3.44822) (16.43472,3.44876) (16.44778,3.44921) (16.46083,3.44955) (16.47389,3.44980) (16.48694,3.44995) (16.50000,3.45000)};
\draw[->,coursemuted] plot coordinates {(7.10000,4.65000) (16.60000,4.65000)};
\draw[->,coursemuted] plot coordinates {(7.10000,4.65000) (7.10000,6.50000)};
\draw[coursemuted!25,thin] plot coordinates {(7.10000,4.65000) (16.50000,4.65000)};
\node[left,inner sep=3pt] at (7.1,4.65000) {$0$};
\draw[coursemuted!25,thin] plot coordinates {(7.10000,5.45000) (16.50000,5.45000)};
\node[left,inner sep=3pt] at (7.1,5.45000) {$8{,}5$};
\draw[coursemuted!25,thin] plot coordinates {(7.10000,6.25000) (16.50000,6.25000)};
\node[left,inner sep=3pt] at (7.1,6.25000) {$17$};
\node[rotate=90] at (5.9,5.55) {$\Delta V_j$, м/с};
\draw[coursemuted!16,thin] plot coordinates {(7.10000,0.80000) (7.10000,6.25000)};
\node[below,inner sep=3pt] at (7.10000,0.85) {$0$};
\draw[coursemuted!16,thin] plot coordinates {(9.45000,0.80000) (9.45000,6.25000)};
\node[below,inner sep=3pt] at (9.45000,0.85) {$\pi/2$};
\draw[coursemuted!16,thin] plot coordinates {(11.80000,0.80000) (11.80000,6.25000)};
\node[below,inner sep=3pt] at (11.80000,0.85) {$\pi$};
\draw[coursemuted!16,thin] plot coordinates {(14.15000,0.80000) (14.15000,6.25000)};
\node[below,inner sep=3pt] at (14.15000,0.85) {$3\pi/2$};
\draw[coursemuted!16,thin] plot coordinates {(16.50000,0.80000) (16.50000,6.25000)};
\node[below,inner sep=3pt] at (16.50000,0.85) {$2\pi$};
\node at (11.8,.15) {Угол $\varphi$, рад};
\draw[courseink!70,densely dashed,thin] plot coordinates {(7.10000,0.85000) (7.10000,3.45000)};
\draw[courseink,line width=1pt] plot coordinates {(7.10000,4.65000) (7.10000,5.90550)};
\fill[courseink] (7.10000,5.90550) circle[radius=1.8pt];
\fill[courseink] (7.10000,3.45000) circle[radius=1.8pt];
\node[above right,inner sep=2pt] at (7.10000,5.90550) {1};
\draw[courseink!70,densely dashed,thin] plot coordinates {(11.80000,0.85000) (11.80000,3.45000)};
\draw[courseink,line width=1pt] plot coordinates {(11.80000,4.65000) (11.80000,5.40330)};
\fill[courseink] (11.80000,5.40330) circle[radius=1.8pt];
\fill[courseink] (11.80000,3.45000) circle[radius=1.8pt];
\node[above,inner sep=2pt] at (11.80000,5.40330) {2};
\end{tikzpicture}

\par
\figcaption{$\Delta a=0{,}004$, $\Delta e_x=0{,}001$, $\Delta e_y=0$, $\Delta i=0{,}002$ рад.\\
Отрезок показан в плоскости $(s_t,s_z)$; оси --- в косоугольной проекции.}
\end{frame}

\begin{frame}{Нулевые радиальные компоненты}
\relax
\derivationmathspacing
Для двух импульсов в узлах $\varphi_1=0$, $\varphi_2=\pi$:
\[
J=\sqrt{\frac{(\Delta a)^2}{4}+(\Delta i)^2},\qquad
\boldsymbol s(\varphi)=\frac1J(0,\Delta a/2,\Delta i\cos\varphi).
\]
\[
\begin{aligned}
\Delta\widehat V_1&=\frac J2\left(1+\frac{\Delta e_x}{\Delta a}\right),&
\Delta\widehat V_2&=\frac J2\left(1-\frac{\Delta e_x}{\Delta a}\right),\\
\Delta\widehat{\boldsymbol V}_1&=
\frac{\Delta\widehat V_1}{J}(0,\Delta a/2,\Delta i),\\
\Delta\widehat{\boldsymbol V}_2&=
\frac{\Delta\widehat V_2}{J}(0,\Delta a/2,-\Delta i).
\end{aligned}
\]
\medskip
\justifying
Трансверсальные компоненты имеют один знак, нормальные ---
противоположные. При $\Delta e_x=0$ модули равны.
\end{frame}

% Источник: eq:sc-region, eq:sc-cost. Без вывода стоимости и области.
\begin{frame}{Особое решение}
\relax
\derivationmathspacing
\[
\begin{gathered}
\Delta i\leq\sqrt3\,|\Delta e_y|,\\
(\Delta a)^2\leq\Delta e^2+
\frac{2\Delta i|\Delta e_y|}{\sqrt3}-(\Delta i)^2.
\end{gathered}
\]
\[
J_s=\frac12\sqrt{(\Delta e_x)^2+
\bigl(|\Delta e_y|+\sqrt3\Delta i\bigr)^2}.
\]
\[
\cos\alpha=\frac{\Delta e_x}{2J_s},\qquad
\sin\alpha=\frac{\operatorname{sgn}(\Delta e_y)
\bigl(|\Delta e_y|+\sqrt3\Delta i\bigr)}{2J_s}.
\]
\medskip
\justifying
При $\Delta i>0$ в этой области $\Delta e_y\ne0$.
Изменение $\Delta a$ обеспечивается перераспределением модулей
и ограничивает область существования решения.
\end{frame}

% Источник: eq:sc-primer, eq:sc-angle.
\begin{frame}{Особое решение: геометрия}
\relax
\derivationmathspacing
\[
\boldsymbol s(\varphi)=
\left(\tfrac12\sin\psi,\ \cos\psi,\
-\operatorname{sgn}(\Delta e_y)\tfrac{\sqrt3}{2}\sin\psi\right),
\quad \psi=\varphi-\alpha.
\]
\medskip
{\justifying
Годограф --- окружность на единичной сфере.
Модуль базис-вектора равен единице при любом угле:
\par}
\[
\|\boldsymbol s(\varphi)\|=1\qquad\text{при всех }\varphi.
\]
\medskip
{\justifying
Отдельные точки касания не определяют места импульсов.
Пять граничных условий по-прежнему должны выполняться.
\par}
\end{frame}

\begin{frame}{Особое решение: годограф}
\relax
\centering
\begin{tikzpicture}[course diagram,scale=.72]
\node at (2.65,6.85) {Годограф};
\node at (11.6,6.85) {Модули импульсов};
\node at (11.6,4.1) {Модуль базис-вектора};
\fill[coursetint!45!coursepaper] (2.65,3.6) circle[radius=1.7cm];
\filldraw[fill=courseaccent,fill opacity=.09,draw=courseaccent!30] (4.49851,4.67574)--(2.18190,5.84251)--(0.80149,2.52426)--(3.11810,1.35749)--cycle;
\draw[coursemuted!40,thin,densely dashed] plot coordinates {(1.34772,3.13819) (1.38665,3.11926) (1.42712,3.10092) (1.46908,3.08319) (1.51248,3.06609) (1.55726,3.04963) (1.60338,3.03385) (1.65077,3.01876) (1.69937,3.00438) (1.74914,2.99072) (1.80000,2.97780) (1.85190,2.96565) (1.90477,2.95426) (1.95855,2.94366) (2.01317,2.93386) (2.06857,2.92488) (2.12467,2.91671) (2.18142,2.90938) (2.23873,2.90289) (2.29655,2.89725) (2.35480,2.89246) (2.41341,2.88854) (2.47230,2.88548) (2.53141,2.88330) (2.59067,2.88199) (2.65000,2.88155) (2.70933,2.88199) (2.76859,2.88330) (2.82770,2.88548) (2.88659,2.88854) (2.94520,2.89246) (3.00345,2.89725) (3.06127,2.90289) (3.11858,2.90938) (3.17533,2.91671) (3.23143,2.92488) (3.28683,2.93386) (3.34145,2.94366) (3.39523,2.95426) (3.44810,2.96565) (3.50000,2.97780) (3.55086,2.99072) (3.60063,3.00438) (3.64923,3.01876) (3.69662,3.03385) (3.74274,3.04963) (3.78752,3.06609) (3.83092,3.08319) (3.87288,3.10092) (3.91335,3.11926) (3.95228,3.13819) (3.98962,3.15768) (4.02533,3.17771) (4.05936,3.19825) (4.09168,3.21928) (4.12224,3.24077) (4.15101,3.26271) (4.17795,3.28505) (4.20303,3.30778) (4.22621,3.33086) (4.24748,3.35428) (4.26680,3.37799) (4.28414,3.40197) (4.29950,3.42619) (4.31285,3.45063) (4.32417,3.47524) (4.33346,3.50001) (4.34069,3.52490) (4.34586,3.54988) (4.34896,3.57493) (4.35000,3.60000) (4.34896,3.62507) (4.34586,3.65012) (4.34069,3.67510) (4.33346,3.69999) (4.32417,3.72476) (4.31285,3.74937) (4.29950,3.77381) (4.28414,3.79803) (4.26680,3.82201) (4.24748,3.84572) (4.22621,3.86914) (4.20303,3.89222) (4.17795,3.91495) (4.15101,3.93729) (4.12224,3.95923) (4.09168,3.98072) (4.05936,4.00175) (4.02533,4.02229) (3.98962,4.04232) (3.95228,4.06181) (3.91335,4.08074) (3.87288,4.09908) (3.83092,4.11681) (3.78752,4.13391) (3.74274,4.15037) (3.69662,4.16615) (3.64923,4.18124) (3.60063,4.19562) (3.55086,4.20928) (3.50000,4.22220) (3.44810,4.23435) (3.39523,4.24574) (3.34145,4.25634) (3.28683,4.26614) (3.23143,4.27512) (3.17533,4.28329) (3.11858,4.29062) (3.06127,4.29711) (3.00345,4.30275) (2.94520,4.30754) (2.88659,4.31146) (2.82770,4.31452) (2.76859,4.31670) (2.70933,4.31801) (2.65000,4.31845) (2.59067,4.31801) (2.53141,4.31670) (2.47230,4.31452) (2.41341,4.31146) (2.35480,4.30754) (2.29655,4.30275) (2.23873,4.29711) (2.18142,4.29062) (2.12467,4.28329) (2.06857,4.27512) (2.01317,4.26614) (1.95855,4.25634) (1.90477,4.24574) (1.85190,4.23435) (1.80000,4.22220) (1.74914,4.20928) (1.69937,4.19562) (1.65077,4.18124) (1.60338,4.16615) (1.55726,4.15037) (1.51248,4.13391) (1.46908,4.11681) (1.42712,4.09908) (1.38665,4.08074) (1.34772,4.06181) (1.31038,4.04232) (1.27467,4.02229) (1.24064,4.00175) (1.20832,3.98072) (1.17776,3.95923) (1.14899,3.93729) (1.12205,3.91495) (1.09697,3.89222) (1.07379,3.86914) (1.05252,3.84572) (1.03320,3.82201) (1.01586,3.79803) (1.00050,3.77381) (0.98715,3.74937) (0.97583,3.72476) (0.96654,3.69999) (0.95931,3.67510) (0.95414,3.65012) (0.95104,3.62507) (0.95000,3.60000) (0.95104,3.57493) (0.95414,3.54988) (0.95931,3.52490) (0.96654,3.50001) (0.97583,3.47524) (0.98715,3.45063) (1.00050,3.42619) (1.01586,3.40197) (1.03320,3.37799) (1.05252,3.35428) (1.07379,3.33086) (1.09697,3.30778) (1.12205,3.28505) (1.14899,3.26271) (1.17776,3.24077) (1.20832,3.21928) (1.24064,3.19825) (1.27467,3.17771) (1.31038,3.15768) (1.34772,3.13819)};
\draw[coursemuted!55,thin] plot coordinates {(0.95000,3.60000) (0.95104,3.57493) (0.95414,3.54988) (0.95931,3.52490) (0.96654,3.50001) (0.97583,3.47524) (0.98715,3.45063) (1.00050,3.42619) (1.01586,3.40197) (1.03320,3.37799) (1.05252,3.35428) (1.07379,3.33086) (1.09697,3.30778) (1.12205,3.28505) (1.14899,3.26271) (1.17776,3.24077) (1.20832,3.21928) (1.24064,3.19825) (1.27467,3.17771) (1.31038,3.15768) (1.34772,3.13819) (1.38665,3.11926) (1.42712,3.10092) (1.46908,3.08319) (1.51248,3.06609) (1.55726,3.04963) (1.60338,3.03385) (1.65077,3.01876) (1.69937,3.00438) (1.74914,2.99072) (1.80000,2.97780) (1.85190,2.96565) (1.90477,2.95426) (1.95855,2.94366) (2.01317,2.93386) (2.06857,2.92488) (2.12467,2.91671) (2.18142,2.90938) (2.23873,2.90289) (2.29655,2.89725) (2.35480,2.89246) (2.41341,2.88854) (2.47230,2.88548) (2.53141,2.88330) (2.59067,2.88199) (2.65000,2.88155) (2.70933,2.88199) (2.76859,2.88330) (2.82770,2.88548) (2.88659,2.88854) (2.94520,2.89246) (3.00345,2.89725) (3.06127,2.90289) (3.11858,2.90938) (3.17533,2.91671) (3.23143,2.92488) (3.28683,2.93386) (3.34145,2.94366) (3.39523,2.95426) (3.44810,2.96565) (3.50000,2.97780) (3.55086,2.99072) (3.60063,3.00438) (3.64923,3.01876) (3.69662,3.03385) (3.74274,3.04963) (3.78752,3.06609) (3.83092,3.08319) (3.87288,3.10092) (3.91335,3.11926) (3.95228,3.13819) (3.98962,3.15768) (4.02533,3.17771) (4.05936,3.19825) (4.09168,3.21928) (4.12224,3.24077) (4.15101,3.26271) (4.17795,3.28505) (4.20303,3.30778) (4.22621,3.33086) (4.24748,3.35428) (4.26680,3.37799) (4.28414,3.40197) (4.29950,3.42619) (4.31285,3.45063) (4.32417,3.47524) (4.33346,3.50001) (4.34069,3.52490) (4.34586,3.54988) (4.34896,3.57493) (4.35000,3.60000)};
\draw[coursemuted!40,thin,densely dashed] plot coordinates {(1.34772,3.13819) (1.34852,3.19224) (1.35090,3.24679) (1.35486,3.30177) (1.36040,3.35711) (1.36751,3.41275) (1.37618,3.46861) (1.38641,3.52464) (1.39817,3.58076) (1.41146,3.63690) (1.42626,3.69300) (1.44255,3.74898) (1.46031,3.80478) (1.47952,3.86034) (1.50016,3.91557) (1.52220,3.97042) (1.54561,4.02482) (1.57036,4.07870) (1.59644,4.13200) (1.62379,4.18465) (1.65240,4.23659) (1.68222,4.28775) (1.71322,4.33808) (1.74536,4.38750) (1.77861,4.43597) (1.81291,4.48342) (1.84824,4.52979) (1.88454,4.57503) (1.92178,4.61908) (1.95990,4.66188) (1.99886,4.70340) (2.03862,4.74357) (2.07912,4.78235) (2.12032,4.81969) (2.16216,4.85554) (2.20460,4.88986) (2.24757,4.92261) (2.29104,4.95375) (2.33495,4.98323) (2.37924,5.01104) (2.42386,5.03712) (2.46876,5.06146) (2.51388,5.08401) (2.55916,5.10476) (2.60455,5.12367) (2.65000,5.14072) (2.69545,5.15590) (2.74084,5.16918) (2.78612,5.18056) (2.83124,5.19000) (2.87614,5.19751) (2.92076,5.20307) (2.96505,5.20668) (3.00896,5.20833) (3.05243,5.20802) (3.09540,5.20576) (3.13784,5.20153) (3.17968,5.19536) (3.22088,5.18724) (3.26138,5.17719) (3.30114,5.16521) (3.34010,5.15133) (3.37822,5.13556) (3.41546,5.11792) (3.45176,5.09843) (3.48709,5.07711) (3.52139,5.05399) (3.55464,5.02910) (3.58678,5.00248) (3.61778,4.97414) (3.64760,4.94413) (3.67621,4.91248) (3.70356,4.87923) (3.72964,4.84442) (3.75439,4.80810) (3.77780,4.77030) (3.79984,4.73108) (3.82048,4.69048) (3.83969,4.64855) (3.85745,4.60535) (3.87374,4.56092) (3.88854,4.51532) (3.90183,4.46860) (3.91359,4.42083) (3.92382,4.37205) (3.93249,4.32234) (3.93960,4.27174) (3.94514,4.22033) (3.94910,4.16816) (3.95148,4.11530) (3.95228,4.06181) (3.95148,4.00776) (3.94910,3.95321) (3.94514,3.89823) (3.93960,3.84289) (3.93249,3.78725) (3.92382,3.73139) (3.91359,3.67536) (3.90183,3.61924) (3.88854,3.56310) (3.87374,3.50700) (3.85745,3.45102) (3.83969,3.39522) (3.82048,3.33966) (3.79984,3.28443) (3.77780,3.22958) (3.75439,3.17518) (3.72964,3.12130) (3.70356,3.06800) (3.67621,3.01535) (3.64760,2.96341) (3.61778,2.91225) (3.58678,2.86192) (3.55464,2.81250) (3.52139,2.76403) (3.48709,2.71658) (3.45176,2.67021) (3.41546,2.62497) (3.37822,2.58092) (3.34010,2.53812) (3.30114,2.49660) (3.26138,2.45643) (3.22088,2.41765) (3.17968,2.38031) (3.13784,2.34446) (3.09540,2.31014) (3.05243,2.27739) (3.00896,2.24625) (2.96505,2.21677) (2.92076,2.18896) (2.87614,2.16288) (2.83124,2.13854) (2.78612,2.11599) (2.74084,2.09524) (2.69545,2.07633) (2.65000,2.05928) (2.60455,2.04410) (2.55916,2.03082) (2.51388,2.01944) (2.46876,2.01000) (2.42386,2.00249) (2.37924,1.99693) (2.33495,1.99332) (2.29104,1.99167) (2.24757,1.99198) (2.20460,1.99424) (2.16216,1.99847) (2.12032,2.00464) (2.07912,2.01276) (2.03862,2.02281) (1.99886,2.03479) (1.95990,2.04867) (1.92178,2.06444) (1.88454,2.08208) (1.84824,2.10157) (1.81291,2.12289) (1.77861,2.14601) (1.74536,2.17090) (1.71322,2.19752) (1.68222,2.22586) (1.65240,2.25587) (1.62379,2.28752) (1.59644,2.32077) (1.57036,2.35558) (1.54561,2.39190) (1.52220,2.42970) (1.50016,2.46892) (1.47952,2.50952) (1.46031,2.55145) (1.44255,2.59465) (1.42626,2.63908) (1.41146,2.68468) (1.39817,2.73140) (1.38641,2.77917) (1.37618,2.82795) (1.36751,2.87766) (1.36040,2.92826) (1.35486,2.97967) (1.35090,3.03184) (1.34852,3.08470) (1.34772,3.13819)};
\draw[coursemuted!55,thin] plot coordinates {(1.88530,2.08170) (1.84898,2.10115) (1.81363,2.12243) (1.77930,2.14551) (1.74604,2.17037) (1.71387,2.19696) (1.68285,2.22526) (1.65300,2.25524) (1.62437,2.28686) (1.59699,2.32007) (1.57089,2.35485) (1.54610,2.39114) (1.52266,2.42891) (1.50060,2.46810) (1.47993,2.50867) (1.46069,2.55057) (1.44290,2.59375) (1.42658,2.63815) (1.41175,2.68373) (1.39843,2.73043) (1.38663,2.77818) (1.37638,2.82693) (1.36767,2.87663) (1.36053,2.92721) (1.35496,2.97860) (1.35096,3.03076) (1.34855,3.08361) (1.34772,3.13708) (1.34849,3.19112) (1.35083,3.24566) (1.35476,3.30063) (1.36027,3.35597) (1.36735,3.41160) (1.37599,3.46746) (1.38618,3.52349) (1.39791,3.57960) (1.41117,3.63575) (1.42594,3.69184) (1.44220,3.74783) (1.45993,3.80364) (1.47911,3.85920) (1.49972,3.91444) (1.52173,3.96930) (1.54511,4.02371) (1.56984,4.07760) (1.59589,4.13091) (1.62322,4.18358) (1.65180,4.23553) (1.68160,4.28671) (1.71257,4.33705) (1.74469,4.38650) (1.77791,4.43498) (1.81220,4.48245) (1.84750,4.52884) (1.88379,4.57411) (1.92100,4.61818) (1.95911,4.66102) (1.99805,4.70256) (2.03779,4.74276) (2.07828,4.78156) (2.11946,4.81893) (2.16129,4.85481) (2.20372,4.88917) (2.24669,4.92195) (2.29015,4.95312) (2.33404,4.98265) (2.37833,5.01048) (2.42294,5.03660) (2.46783,5.06097) (2.51295,5.08356) (2.55822,5.10435) (2.60362,5.12330) (2.64906,5.14039) (2.69451,5.15561) (2.73991,5.16893) (2.78519,5.18034) (2.83032,5.18983) (2.87522,5.19737) (2.91984,5.20298) (2.96414,5.20662) (3.00806,5.20832) (3.05154,5.20805) (3.09453,5.20582) (3.13697,5.20164) (3.17883,5.19550) (3.22004,5.18742) (3.26056,5.17741) (3.30033,5.16548) (3.33931,5.15163) (3.37745,5.13590) (3.41470,5.11830)};
\draw[coursemuted!40,thin,densely dashed] plot coordinates {(3.74274,3.04963) (3.74207,3.10374) (3.74008,3.15845) (3.73675,3.21370) (3.73210,3.26942) (3.72614,3.32554) (3.71886,3.38200) (3.71028,3.43872) (3.70041,3.49564) (3.68926,3.55268) (3.67684,3.60978) (3.66317,3.66688) (3.64827,3.72388) (3.63215,3.78074) (3.61483,3.83738) (3.59634,3.89373) (3.57670,3.94972) (3.55592,4.00529) (3.53404,4.06036) (3.51109,4.11487) (3.48709,4.16875) (3.46206,4.22194) (3.43605,4.27438) (3.40908,4.32599) (3.38119,4.37671) (3.35240,4.42649) (3.32276,4.47527) (3.29230,4.52297) (3.26105,4.56956) (3.22906,4.61496) (3.19637,4.65912) (3.16301,4.70200) (3.12903,4.74353) (3.09446,4.78367) (3.05935,4.82236) (3.02374,4.85957) (2.98767,4.89524) (2.95120,4.92934) (2.91436,4.96181) (2.87719,4.99263) (2.83975,5.02175) (2.80208,5.04913) (2.76422,5.07475) (2.72623,5.09858) (2.68814,5.12058) (2.65000,5.14072) (2.61186,5.15899) (2.57377,5.17536) (2.53578,5.18981) (2.49792,5.20233) (2.46025,5.21289) (2.42281,5.22148) (2.38564,5.22810) (2.34880,5.23274) (2.31233,5.23539) (2.27626,5.23604) (2.24065,5.23470) (2.20554,5.23137) (2.17097,5.22606) (2.13699,5.21876) (2.10363,5.20949) (2.07094,5.19826) (2.03895,5.18508) (2.00770,5.16997) (1.97724,5.15295) (1.94760,5.13403) (1.91881,5.11325) (1.89092,5.09062) (1.86395,5.06618) (1.83794,5.03995) (1.81291,5.01196) (1.78891,4.98226) (1.76596,4.95087) (1.74408,4.91784) (1.72330,4.88320) (1.70366,4.84699) (1.68517,4.80927) (1.66785,4.77007) (1.65173,4.72945) (1.63683,4.68746) (1.62316,4.64413) (1.61074,4.59954) (1.59959,4.55373) (1.58972,4.50675) (1.58114,4.45867) (1.57386,4.40955) (1.56790,4.35944) (1.56325,4.30840) (1.55992,4.25650) (1.55793,4.20380) (1.55726,4.15037) (1.55793,4.09626) (1.55992,4.04155) (1.56325,3.98630) (1.56790,3.93058) (1.57386,3.87446) (1.58114,3.81800) (1.58972,3.76128) (1.59959,3.70436) (1.61074,3.64732) (1.62316,3.59022) (1.63683,3.53312) (1.65173,3.47612) (1.66785,3.41926) (1.68517,3.36262) (1.70366,3.30627) (1.72330,3.25028) (1.74408,3.19471) (1.76596,3.13964) (1.78891,3.08513) (1.81291,3.03125) (1.83794,2.97806) (1.86395,2.92562) (1.89092,2.87401) (1.91881,2.82329) (1.94760,2.77351) (1.97724,2.72473) (2.00770,2.67703) (2.03895,2.63044) (2.07094,2.58504) (2.10363,2.54088) (2.13699,2.49800) (2.17097,2.45647) (2.20554,2.41633) (2.24065,2.37764) (2.27626,2.34043) (2.31233,2.30476) (2.34880,2.27066) (2.38564,2.23819) (2.42281,2.20737) (2.46025,2.17825) (2.49792,2.15087) (2.53578,2.12525) (2.57377,2.10142) (2.61186,2.07942) (2.65000,2.05928) (2.68814,2.04101) (2.72623,2.02464) (2.76422,2.01019) (2.80208,1.99767) (2.83975,1.98711) (2.87719,1.97852) (2.91436,1.97190) (2.95120,1.96726) (2.98767,1.96461) (3.02374,1.96396) (3.05935,1.96530) (3.09446,1.96863) (3.12903,1.97394) (3.16301,1.98124) (3.19637,1.99051) (3.22906,2.00174) (3.26105,2.01492) (3.29230,2.03003) (3.32276,2.04705) (3.35240,2.06597) (3.38119,2.08675) (3.40908,2.10938) (3.43605,2.13382) (3.46206,2.16005) (3.48709,2.18804) (3.51109,2.21774) (3.53404,2.24913) (3.55592,2.28216) (3.57670,2.31680) (3.59634,2.35301) (3.61483,2.39073) (3.63215,2.42993) (3.64827,2.47055) (3.66317,2.51254) (3.67684,2.55587) (3.68926,2.60046) (3.70041,2.64627) (3.71028,2.69325) (3.71886,2.74133) (3.72614,2.79045) (3.73210,2.84056) (3.73675,2.89160) (3.74008,2.94350) (3.74207,2.99620) (3.74274,3.04963)};
\draw[coursemuted!55,thin] plot coordinates {(3.21818,1.99776) (3.25041,2.01029) (3.28191,2.02475) (3.31264,2.04114) (3.34256,2.05942) (3.37164,2.07958) (3.39983,2.10159) (3.42712,2.12543) (3.45345,2.15107) (3.47881,2.17847) (3.50316,2.20760) (3.52647,2.23843) (3.54871,2.27091) (3.56986,2.30502) (3.58988,2.34070) (3.60876,2.37792) (3.62648,2.41663) (3.64300,2.45678) (3.65831,2.49832) (3.67240,2.54120) (3.68524,2.58538) (3.69681,2.63079) (3.70712,2.67738) (3.71613,2.72510) (3.72384,2.77388) (3.73025,2.82366) (3.73534,2.87440) (3.73911,2.92602) (3.74155,2.97845) (3.74266,3.03165) (3.74244,3.08554) (3.74089,3.14005) (3.73801,3.19513) (3.73381,3.25070) (3.72828,3.30669) (3.72144,3.36304) (3.71330,3.41968) (3.70386,3.47654) (3.69314,3.53355) (3.68114,3.59065) (3.66789,3.64775) (3.65340,3.70479) (3.63768,3.76171) (3.62077,3.81843) (3.60267,3.87488) (3.58341,3.93100) (3.56301,3.98672) (3.54150,4.04196) (3.51890,4.09667) (3.49525,4.15077) (3.47056,4.20420) (3.44488,4.25689) (3.41822,4.30879) (3.39063,4.35982) (3.36214,4.40992) (3.33278,4.45904) (3.30259,4.50711) (3.27161,4.55408) (3.23986,4.59988) (3.20740,4.64446) (3.17426,4.68778) (3.14048,4.72976) (3.10610,4.77037) (3.07117,4.80956) (3.03572,4.84727) (2.99981,4.88346) (2.96347,4.91809) (2.92674,4.95111) (2.88968,4.98249) (2.85233,5.01218) (2.81473,5.04015) (2.77693,5.06637) (2.73897,5.09080) (2.70091,5.11341) (2.66278,5.13418) (2.62464,5.15308) (2.58653,5.17009) (2.54850,5.18518) (2.51059,5.19835) (2.47285,5.20957) (2.43533,5.21882) (2.39806,5.22610) (2.36111,5.23141) (2.32450,5.23472) (2.28830,5.23604) (2.25253,5.23537) (2.21725,5.23271) (2.18250,5.22806) (2.14831,5.22142) (2.11474,5.21281) (2.08182,5.20224)};
\draw[coursemuted!50,thin,densely dashed] plot coordinates {(4.14762,4.13108) (2.65000,3.60000)};
\draw[->,courseink!75,thin] plot coordinates {(2.65000,3.60000) (0.82681,2.95346)};
\node[above right,inner sep=1pt] at (0.82681,2.95346) {$s_r$};
\draw[coursemuted!50,thin,densely dashed] plot coordinates {(1.39335,4.23292) (2.65000,3.60000)};
\draw[->,courseink!75,thin] plot coordinates {(2.65000,3.60000) (4.17983,2.82949)};
\node[below right,inner sep=1pt] at (4.17983,2.82949) {$s_t$};
\draw[coursemuted!50,thin,densely dashed] plot coordinates {(2.65000,1.82817) (2.65000,3.60000)};
\draw[->,courseink!75,thin] plot coordinates {(2.65000,3.60000) (2.65000,5.75701)};
\node[above,inner sep=1pt] at (2.65000,5.75701) {$s_z$};
\draw[courseaccent,line width=1.15pt] plot coordinates {(3.30114,5.16521) (3.32011,5.15537) (3.33888,5.14505) (3.35743,5.13426) (3.37578,5.12301) (3.39390,5.11129) (3.41179,5.09911) (3.42946,5.08647) (3.44688,5.07338) (3.46406,5.05984) (3.48100,5.04586) (3.49768,5.03144) (3.51410,5.01658) (3.53026,5.00129) (3.54615,4.98557) (3.56177,4.96943) (3.57711,4.95288) (3.59217,4.93591) (3.60694,4.91853) (3.62142,4.90075) (3.63561,4.88258) (3.64949,4.86402) (3.66307,4.84507) (3.67634,4.82574) (3.68930,4.80604) (3.70194,4.78597) (3.71426,4.76554) (3.72626,4.74475) (3.73793,4.72362) (3.74927,4.70214) (3.76027,4.68033) (3.77094,4.65819) (3.78126,4.63573) (3.79124,4.61295) (3.80087,4.58986) (3.81015,4.56647) (3.81908,4.54279) (3.82765,4.51882) (3.83586,4.49456) (3.84371,4.47004) (3.85120,4.44525) (3.85832,4.42021) (3.86507,4.39491) (3.87146,4.36937) (3.87747,4.34360) (3.88311,4.31760) (3.88837,4.29139) (3.89325,4.26496) (3.89776,4.23833) (3.90189,4.21150) (3.90563,4.18449) (3.90899,4.15731) (3.91197,4.12995) (3.91456,4.10243) (3.91677,4.07475) (3.91860,4.04694) (3.92003,4.01898) (3.92108,3.99090) (3.92175,3.96270) (3.92202,3.93439) (3.92191,3.90598) (3.92141,3.87747) (3.92052,3.84888) (3.91925,3.82021) (3.91759,3.79148) (3.91554,3.76269) (3.91311,3.73384) (3.91029,3.70496) (3.90709,3.67605) (3.90351,3.64711) (3.89954,3.61816) (3.89519,3.58920) (3.89047,3.56025) (3.88537,3.53131) (3.87989,3.50239) (3.87403,3.47349) (3.86780,3.44464) (3.86121,3.41584) (3.85424,3.38709) (3.84691,3.35840) (3.83921,3.32979) (3.83115,3.30126) (3.82273,3.27283) (3.81395,3.24449) (3.80482,3.21626) (3.79533,3.18815) (3.78550,3.16016) (3.77532,3.13231) (3.76480,3.10459) (3.75394,3.07704) (3.74274,3.04963) (3.73121,3.02240) (3.71935,2.99534) (3.70716,2.96847) (3.69466,2.94179) (3.68183,2.91531) (3.66869,2.88904) (3.65524,2.86299) (3.64148,2.83716) (3.62743,2.81156) (3.61307,2.78620) (3.59842,2.76109) (3.58348,2.73624) (3.56826,2.71164) (3.55276,2.68732) (3.53698,2.66328) (3.52093,2.63952) (3.50462,2.61606) (3.48804,2.59289) (3.47122,2.57004) (3.45414,2.54749) (3.43681,2.52527) (3.41925,2.50337) (3.40145,2.48181) (3.38343,2.46059) (3.36518,2.43971) (3.34671,2.41919) (3.32803,2.39903) (3.30914,2.37923) (3.29005,2.35981) (3.27077,2.34076) (3.25130,2.32210) (3.23164,2.30383) (3.21181,2.28595) (3.19181,2.26847) (3.17164,2.25140) (3.15132,2.23474) (3.13084,2.21849) (3.11021,2.20267) (3.08944,2.18727) (3.06854,2.17230) (3.04752,2.15776) (3.02637,2.14367) (3.00510,2.13002) (2.98373,2.11681) (2.96226,2.10406) (2.94069,2.09177) (2.91903,2.07993) (2.89730,2.06856) (2.87548,2.05765) (2.85360,2.04721) (2.83165,2.03725) (2.80965,2.02776) (2.78760,2.01875) (2.76551,2.01022) (2.74339,2.00218) (2.72123,1.99462) (2.69906,1.98755) (2.67687,1.98098) (2.65467,1.97489) (2.63247,1.96930) (2.61027,1.96421) (2.58809,1.95962) (2.56593,1.95553) (2.54379,1.95193) (2.52168,1.94884) (2.49961,1.94625) (2.47759,1.94417) (2.45562,1.94259) (2.43371,1.94152) (2.41187,1.94095) (2.39010,1.94088) (2.36841,1.94132) (2.34680,1.94227) (2.32529,1.94372) (2.30387,1.94568) (2.28256,1.94814) (2.26136,1.95110) (2.24028,1.95457) (2.21933,1.95853) (2.19851,1.96300) (2.17782,1.96796) (2.15728,1.97343) (2.13689,1.97938) (2.11665,1.98583) (2.09658,1.99278) (2.07667,2.00021) (2.05694,2.00813) (2.03739,2.01653) (2.01803,2.02542) (1.99886,2.03479) (1.97989,2.04463) (1.96112,2.05495) (1.94257,2.06574) (1.92422,2.07699) (1.90610,2.08871) (1.88821,2.10089) (1.87054,2.11353) (1.85312,2.12662) (1.83594,2.14016) (1.81900,2.15414) (1.80232,2.16856) (1.78590,2.18342) (1.76974,2.19871) (1.75385,2.21443) (1.73823,2.23057) (1.72289,2.24712) (1.70783,2.26409) (1.69306,2.28147) (1.67858,2.29925) (1.66439,2.31742) (1.65051,2.33598) (1.63693,2.35493) (1.62366,2.37426) (1.61070,2.39396) (1.59806,2.41403) (1.58574,2.43446) (1.57374,2.45525) (1.56207,2.47638) (1.55073,2.49786) (1.53973,2.51967) (1.52906,2.54181) (1.51874,2.56427) (1.50876,2.58705) (1.49913,2.61014) (1.48985,2.63353) (1.48092,2.65721) (1.47235,2.68118) (1.46414,2.70544) (1.45629,2.72996) (1.44880,2.75475) (1.44168,2.77979) (1.43493,2.80509) (1.42854,2.83063) (1.42253,2.85640) (1.41689,2.88240) (1.41163,2.90861) (1.40675,2.93504) (1.40224,2.96167) (1.39811,2.98850) (1.39437,3.01551) (1.39101,3.04269) (1.38803,3.07005) (1.38544,3.09757) (1.38323,3.12525) (1.38140,3.15306) (1.37997,3.18102) (1.37892,3.20910) (1.37825,3.23730) (1.37798,3.26561) (1.37809,3.29402) (1.37859,3.32253) (1.37948,3.35112) (1.38075,3.37979) (1.38241,3.40852) (1.38446,3.43731) (1.38689,3.46616) (1.38971,3.49504) (1.39291,3.52395) (1.39649,3.55289) (1.40046,3.58184) (1.40481,3.61080) (1.40953,3.63975) (1.41463,3.66869) (1.42011,3.69761) (1.42597,3.72651) (1.43220,3.75536) (1.43879,3.78416) (1.44576,3.81291) (1.45309,3.84160) (1.46079,3.87021) (1.46885,3.89874) (1.47727,3.92717) (1.48605,3.95551) (1.49518,3.98374) (1.50467,4.01185) (1.51450,4.03984) (1.52468,4.06769) (1.53520,4.09541) (1.54606,4.12296) (1.55726,4.15037) (1.56879,4.17760) (1.58065,4.20466) (1.59284,4.23153) (1.60534,4.25821) (1.61817,4.28469) (1.63131,4.31096) (1.64476,4.33701) (1.65852,4.36284) (1.67257,4.38844) (1.68693,4.41380) (1.70158,4.43891) (1.71652,4.46376) (1.73174,4.48836) (1.74724,4.51268) (1.76302,4.53672) (1.77907,4.56048) (1.79538,4.58394) (1.81196,4.60711) (1.82878,4.62996) (1.84586,4.65251) (1.86319,4.67473) (1.88075,4.69663) (1.89855,4.71819) (1.91657,4.73941) (1.93482,4.76029) (1.95329,4.78081) (1.97197,4.80097) (1.99086,4.82077) (2.00995,4.84019) (2.02923,4.85924) (2.04870,4.87790) (2.06836,4.89617) (2.08819,4.91405) (2.10819,4.93153) (2.12836,4.94860) (2.14868,4.96526) (2.16916,4.98151) (2.18979,4.99733) (2.21056,5.01273) (2.23146,5.02770) (2.25248,5.04224) (2.27363,5.05633) (2.29490,5.06998) (2.31627,5.08319) (2.33774,5.09594) (2.35931,5.10823) (2.38097,5.12007) (2.40270,5.13144) (2.42452,5.14235) (2.44640,5.15279) (2.46835,5.16275) (2.49035,5.17224) (2.51240,5.18125) (2.53449,5.18978) (2.55661,5.19782) (2.57877,5.20538) (2.60094,5.21245) (2.62313,5.21902) (2.64533,5.22511) (2.66753,5.23070) (2.68973,5.23579) (2.71191,5.24038) (2.73407,5.24447) (2.75621,5.24807) (2.77832,5.25116) (2.80039,5.25375) (2.82241,5.25583) (2.84438,5.25741) (2.86629,5.25848) (2.88813,5.25905) (2.90990,5.25912) (2.93159,5.25868) (2.95320,5.25773) (2.97471,5.25628) (2.99613,5.25432) (3.01744,5.25186) (3.03864,5.24890) (3.05972,5.24543) (3.08067,5.24147) (3.10149,5.23700) (3.12218,5.23204) (3.14272,5.22657) (3.16311,5.22062) (3.18335,5.21417) (3.20342,5.20722) (3.22333,5.19979) (3.24306,5.19187) (3.26261,5.18347) (3.28197,5.17458) (3.30114,5.16521)};
\draw[->,courseaccent,line width=.55pt] plot coordinates {(2.65000,3.60000) (3.91707,4.07066)};
\filldraw[fill=courseink,draw=coursepaper,line width=.4pt] (3.91707,4.07066) circle[radius=2pt];
\node[fill=coursepaper,fill opacity=.85,text opacity=1,inner sep=.6pt] at (4.04707,4.25066) {1};
\draw[->,courseaccent,line width=.55pt] plot coordinates {(2.65000,3.60000) (3.15432,2.23717)};
\filldraw[fill=courseink,draw=coursepaper,line width=.4pt] (3.15432,2.23717) circle[radius=2pt];
\node[fill=coursepaper,fill opacity=.85,text opacity=1,inner sep=.6pt] at (3.28432,2.05717) {2};
\draw[->,courseaccent,line width=.55pt] plot coordinates {(2.65000,3.60000) (1.38293,3.12934)};
\filldraw[fill=courseink,draw=coursepaper,line width=.4pt] (1.38293,3.12934) circle[radius=2pt];
\node[fill=coursepaper,fill opacity=.85,text opacity=1,inner sep=.6pt] at (1.25293,2.94934) {3};
\draw[coursemuted!65,thin] (2.65,3.6) circle[radius=1.7cm];
\node[below left,inner sep=2pt] at (2.65,3.6) {$0$};
\draw[coursemuted!25,thin] plot coordinates {(7.10000,0.85000) (16.50000,0.85000)};
\node[left,inner sep=3pt] at (7.1,0.85) {$0$};
\draw[coursemuted!25,thin] plot coordinates {(7.10000,2.15000) (16.50000,2.15000)};
\node[left,inner sep=3pt] at (7.1,2.15) {$0{,}5$};
\draw[courseink!65,dashed] plot coordinates {(7.10000,3.45000) (16.50000,3.45000)};
\node[left,inner sep=3pt] at (7.1,3.45) {$1$};
\draw[->,coursemuted] plot coordinates {(7.10000,0.85000) (16.60000,0.85000)};
\draw[->,coursemuted] plot coordinates {(7.10000,0.85000) (7.10000,3.84000)};
\node[rotate=90] at (5.9,2.2) {$\|\boldsymbol s(\varphi)\|$};
\draw[courseaccent,line width=1.15pt] plot coordinates {(7.10000,3.45000) (7.11306,3.45000) (7.12611,3.45000) (7.13917,3.45000) (7.15222,3.45000) (7.16528,3.45000) (7.17833,3.45000) (7.19139,3.45000) (7.20444,3.45000) (7.21750,3.45000) (7.23056,3.45000) (7.24361,3.45000) (7.25667,3.45000) (7.26972,3.45000) (7.28278,3.45000) (7.29583,3.45000) (7.30889,3.45000) (7.32194,3.45000) (7.33500,3.45000) (7.34806,3.45000) (7.36111,3.45000) (7.37417,3.45000) (7.38722,3.45000) (7.40028,3.45000) (7.41333,3.45000) (7.42639,3.45000) (7.43944,3.45000) (7.45250,3.45000) (7.46556,3.45000) (7.47861,3.45000) (7.49167,3.45000) (7.50472,3.45000) (7.51778,3.45000) (7.53083,3.45000) (7.54389,3.45000) (7.55694,3.45000) (7.57000,3.45000) (7.58306,3.45000) (7.59611,3.45000) (7.60917,3.45000) (7.62222,3.45000) (7.63528,3.45000) (7.64833,3.45000) (7.66139,3.45000) (7.67444,3.45000) (7.68750,3.45000) (7.70056,3.45000) (7.71361,3.45000) (7.72667,3.45000) (7.73972,3.45000) (7.75278,3.45000) (7.76583,3.45000) (7.77889,3.45000) (7.79194,3.45000) (7.80500,3.45000) (7.81806,3.45000) (7.83111,3.45000) (7.84417,3.45000) (7.85722,3.45000) (7.87028,3.45000) (7.88333,3.45000) (7.89639,3.45000) (7.90944,3.45000) (7.92250,3.45000) (7.93556,3.45000) (7.94861,3.45000) (7.96167,3.45000) (7.97472,3.45000) (7.98778,3.45000) (8.00083,3.45000) (8.01389,3.45000) (8.02694,3.45000) (8.04000,3.45000) (8.05306,3.45000) (8.06611,3.45000) (8.07917,3.45000) (8.09222,3.45000) (8.10528,3.45000) (8.11833,3.45000) (8.13139,3.45000) (8.14444,3.45000) (8.15750,3.45000) (8.17056,3.45000) (8.18361,3.45000) (8.19667,3.45000) (8.20972,3.45000) (8.22278,3.45000) (8.23583,3.45000) (8.24889,3.45000) (8.26194,3.45000) (8.27500,3.45000) (8.28806,3.45000) (8.30111,3.45000) (8.31417,3.45000) (8.32722,3.45000) (8.34028,3.45000) (8.35333,3.45000) (8.36639,3.45000) (8.37944,3.45000) (8.39250,3.45000) (8.40556,3.45000) (8.41861,3.45000) (8.43167,3.45000) (8.44472,3.45000) (8.45778,3.45000) (8.47083,3.45000) (8.48389,3.45000) (8.49694,3.45000) (8.51000,3.45000) (8.52306,3.45000) (8.53611,3.45000) (8.54917,3.45000) (8.56222,3.45000) (8.57528,3.45000) (8.58833,3.45000) (8.60139,3.45000) (8.61444,3.45000) (8.62750,3.45000) (8.64056,3.45000) (8.65361,3.45000) (8.66667,3.45000) (8.67972,3.45000) (8.69278,3.45000) (8.70583,3.45000) (8.71889,3.45000) (8.73194,3.45000) (8.74500,3.45000) (8.75806,3.45000) (8.77111,3.45000) (8.78417,3.45000) (8.79722,3.45000) (8.81028,3.45000) (8.82333,3.45000) (8.83639,3.45000) (8.84944,3.45000) (8.86250,3.45000) (8.87556,3.45000) (8.88861,3.45000) (8.90167,3.45000) (8.91472,3.45000) (8.92778,3.45000) (8.94083,3.45000) (8.95389,3.45000) (8.96694,3.45000) (8.98000,3.45000) (8.99306,3.45000) (9.00611,3.45000) (9.01917,3.45000) (9.03222,3.45000) (9.04528,3.45000) (9.05833,3.45000) (9.07139,3.45000) (9.08444,3.45000) (9.09750,3.45000) (9.11056,3.45000) (9.12361,3.45000) (9.13667,3.45000) (9.14972,3.45000) (9.16278,3.45000) (9.17583,3.45000) (9.18889,3.45000) (9.20194,3.45000) (9.21500,3.45000) (9.22806,3.45000) (9.24111,3.45000) (9.25417,3.45000) (9.26722,3.45000) (9.28028,3.45000) (9.29333,3.45000) (9.30639,3.45000) (9.31944,3.45000) (9.33250,3.45000) (9.34556,3.45000) (9.35861,3.45000) (9.37167,3.45000) (9.38472,3.45000) (9.39778,3.45000) (9.41083,3.45000) (9.42389,3.45000) (9.43694,3.45000) (9.45000,3.45000) (9.46306,3.45000) (9.47611,3.45000) (9.48917,3.45000) (9.50222,3.45000) (9.51528,3.45000) (9.52833,3.45000) (9.54139,3.45000) (9.55444,3.45000) (9.56750,3.45000) (9.58056,3.45000) (9.59361,3.45000) (9.60667,3.45000) (9.61972,3.45000) (9.63278,3.45000) (9.64583,3.45000) (9.65889,3.45000) (9.67194,3.45000) (9.68500,3.45000) (9.69806,3.45000) (9.71111,3.45000) (9.72417,3.45000) (9.73722,3.45000) (9.75028,3.45000) (9.76333,3.45000) (9.77639,3.45000) (9.78944,3.45000) (9.80250,3.45000) (9.81556,3.45000) (9.82861,3.45000) (9.84167,3.45000) (9.85472,3.45000) (9.86778,3.45000) (9.88083,3.45000) (9.89389,3.45000) (9.90694,3.45000) (9.92000,3.45000) (9.93306,3.45000) (9.94611,3.45000) (9.95917,3.45000) (9.97222,3.45000) (9.98528,3.45000) (9.99833,3.45000) (10.01139,3.45000) (10.02444,3.45000) (10.03750,3.45000) (10.05056,3.45000) (10.06361,3.45000) (10.07667,3.45000) (10.08972,3.45000) (10.10278,3.45000) (10.11583,3.45000) (10.12889,3.45000) (10.14194,3.45000) (10.15500,3.45000) (10.16806,3.45000) (10.18111,3.45000) (10.19417,3.45000) (10.20722,3.45000) (10.22028,3.45000) (10.23333,3.45000) (10.24639,3.45000) (10.25944,3.45000) (10.27250,3.45000) (10.28556,3.45000) (10.29861,3.45000) (10.31167,3.45000) (10.32472,3.45000) (10.33778,3.45000) (10.35083,3.45000) (10.36389,3.45000) (10.37694,3.45000) (10.39000,3.45000) (10.40306,3.45000) (10.41611,3.45000) (10.42917,3.45000) (10.44222,3.45000) (10.45528,3.45000) (10.46833,3.45000) (10.48139,3.45000) (10.49444,3.45000) (10.50750,3.45000) (10.52056,3.45000) (10.53361,3.45000) (10.54667,3.45000) (10.55972,3.45000) (10.57278,3.45000) (10.58583,3.45000) (10.59889,3.45000) (10.61194,3.45000) (10.62500,3.45000) (10.63806,3.45000) (10.65111,3.45000) (10.66417,3.45000) (10.67722,3.45000) (10.69028,3.45000) (10.70333,3.45000) (10.71639,3.45000) (10.72944,3.45000) (10.74250,3.45000) (10.75556,3.45000) (10.76861,3.45000) (10.78167,3.45000) (10.79472,3.45000) (10.80778,3.45000) (10.82083,3.45000) (10.83389,3.45000) (10.84694,3.45000) (10.86000,3.45000) (10.87306,3.45000) (10.88611,3.45000) (10.89917,3.45000) (10.91222,3.45000) (10.92528,3.45000) (10.93833,3.45000) (10.95139,3.45000) (10.96444,3.45000) (10.97750,3.45000) (10.99056,3.45000) (11.00361,3.45000) (11.01667,3.45000) (11.02972,3.45000) (11.04278,3.45000) (11.05583,3.45000) (11.06889,3.45000) (11.08194,3.45000) (11.09500,3.45000) (11.10806,3.45000) (11.12111,3.45000) (11.13417,3.45000) (11.14722,3.45000) (11.16028,3.45000) (11.17333,3.45000) (11.18639,3.45000) (11.19944,3.45000) (11.21250,3.45000) (11.22556,3.45000) (11.23861,3.45000) (11.25167,3.45000) (11.26472,3.45000) (11.27778,3.45000) (11.29083,3.45000) (11.30389,3.45000) (11.31694,3.45000) (11.33000,3.45000) (11.34306,3.45000) (11.35611,3.45000) (11.36917,3.45000) (11.38222,3.45000) (11.39528,3.45000) (11.40833,3.45000) (11.42139,3.45000) (11.43444,3.45000) (11.44750,3.45000) (11.46056,3.45000) (11.47361,3.45000) (11.48667,3.45000) (11.49972,3.45000) (11.51278,3.45000) (11.52583,3.45000) (11.53889,3.45000) (11.55194,3.45000) (11.56500,3.45000) (11.57806,3.45000) (11.59111,3.45000) (11.60417,3.45000) (11.61722,3.45000) (11.63028,3.45000) (11.64333,3.45000) (11.65639,3.45000) (11.66944,3.45000) (11.68250,3.45000) (11.69556,3.45000) (11.70861,3.45000) (11.72167,3.45000) (11.73472,3.45000) (11.74778,3.45000) (11.76083,3.45000) (11.77389,3.45000) (11.78694,3.45000) (11.80000,3.45000) (11.81306,3.45000) (11.82611,3.45000) (11.83917,3.45000) (11.85222,3.45000) (11.86528,3.45000) (11.87833,3.45000) (11.89139,3.45000) (11.90444,3.45000) (11.91750,3.45000) (11.93056,3.45000) (11.94361,3.45000) (11.95667,3.45000) (11.96972,3.45000) (11.98278,3.45000) (11.99583,3.45000) (12.00889,3.45000) (12.02194,3.45000) (12.03500,3.45000) (12.04806,3.45000) (12.06111,3.45000) (12.07417,3.45000) (12.08722,3.45000) (12.10028,3.45000) (12.11333,3.45000) (12.12639,3.45000) (12.13944,3.45000) (12.15250,3.45000) (12.16556,3.45000) (12.17861,3.45000) (12.19167,3.45000) (12.20472,3.45000) (12.21778,3.45000) (12.23083,3.45000) (12.24389,3.45000) (12.25694,3.45000) (12.27000,3.45000) (12.28306,3.45000) (12.29611,3.45000) (12.30917,3.45000) (12.32222,3.45000) (12.33528,3.45000) (12.34833,3.45000) (12.36139,3.45000) (12.37444,3.45000) (12.38750,3.45000) (12.40056,3.45000) (12.41361,3.45000) (12.42667,3.45000) (12.43972,3.45000) (12.45278,3.45000) (12.46583,3.45000) (12.47889,3.45000) (12.49194,3.45000) (12.50500,3.45000) (12.51806,3.45000) (12.53111,3.45000) (12.54417,3.45000) (12.55722,3.45000) (12.57028,3.45000) (12.58333,3.45000) (12.59639,3.45000) (12.60944,3.45000) (12.62250,3.45000) (12.63556,3.45000) (12.64861,3.45000) (12.66167,3.45000) (12.67472,3.45000) (12.68778,3.45000) (12.70083,3.45000) (12.71389,3.45000) (12.72694,3.45000) (12.74000,3.45000) (12.75306,3.45000) (12.76611,3.45000) (12.77917,3.45000) (12.79222,3.45000) (12.80528,3.45000) (12.81833,3.45000) (12.83139,3.45000) (12.84444,3.45000) (12.85750,3.45000) (12.87056,3.45000) (12.88361,3.45000) (12.89667,3.45000) (12.90972,3.45000) (12.92278,3.45000) (12.93583,3.45000) (12.94889,3.45000) (12.96194,3.45000) (12.97500,3.45000) (12.98806,3.45000) (13.00111,3.45000) (13.01417,3.45000) (13.02722,3.45000) (13.04028,3.45000) (13.05333,3.45000) (13.06639,3.45000) (13.07944,3.45000) (13.09250,3.45000) (13.10556,3.45000) (13.11861,3.45000) (13.13167,3.45000) (13.14472,3.45000) (13.15778,3.45000) (13.17083,3.45000) (13.18389,3.45000) (13.19694,3.45000) (13.21000,3.45000) (13.22306,3.45000) (13.23611,3.45000) (13.24917,3.45000) (13.26222,3.45000) (13.27528,3.45000) (13.28833,3.45000) (13.30139,3.45000) (13.31444,3.45000) (13.32750,3.45000) (13.34056,3.45000) (13.35361,3.45000) (13.36667,3.45000) (13.37972,3.45000) (13.39278,3.45000) (13.40583,3.45000) (13.41889,3.45000) (13.43194,3.45000) (13.44500,3.45000) (13.45806,3.45000) (13.47111,3.45000) (13.48417,3.45000) (13.49722,3.45000) (13.51028,3.45000) (13.52333,3.45000) (13.53639,3.45000) (13.54944,3.45000) (13.56250,3.45000) (13.57556,3.45000) (13.58861,3.45000) (13.60167,3.45000) (13.61472,3.45000) (13.62778,3.45000) (13.64083,3.45000) (13.65389,3.45000) (13.66694,3.45000) (13.68000,3.45000) (13.69306,3.45000) (13.70611,3.45000) (13.71917,3.45000) (13.73222,3.45000) (13.74528,3.45000) (13.75833,3.45000) (13.77139,3.45000) (13.78444,3.45000) (13.79750,3.45000) (13.81056,3.45000) (13.82361,3.45000) (13.83667,3.45000) (13.84972,3.45000) (13.86278,3.45000) (13.87583,3.45000) (13.88889,3.45000) (13.90194,3.45000) (13.91500,3.45000) (13.92806,3.45000) (13.94111,3.45000) (13.95417,3.45000) (13.96722,3.45000) (13.98028,3.45000) (13.99333,3.45000) (14.00639,3.45000) (14.01944,3.45000) (14.03250,3.45000) (14.04556,3.45000) (14.05861,3.45000) (14.07167,3.45000) (14.08472,3.45000) (14.09778,3.45000) (14.11083,3.45000) (14.12389,3.45000) (14.13694,3.45000) (14.15000,3.45000) (14.16306,3.45000) (14.17611,3.45000) (14.18917,3.45000) (14.20222,3.45000) (14.21528,3.45000) (14.22833,3.45000) (14.24139,3.45000) (14.25444,3.45000) (14.26750,3.45000) (14.28056,3.45000) (14.29361,3.45000) (14.30667,3.45000) (14.31972,3.45000) (14.33278,3.45000) (14.34583,3.45000) (14.35889,3.45000) (14.37194,3.45000) (14.38500,3.45000) (14.39806,3.45000) (14.41111,3.45000) (14.42417,3.45000) (14.43722,3.45000) (14.45028,3.45000) (14.46333,3.45000) (14.47639,3.45000) (14.48944,3.45000) (14.50250,3.45000) (14.51556,3.45000) (14.52861,3.45000) (14.54167,3.45000) (14.55472,3.45000) (14.56778,3.45000) (14.58083,3.45000) (14.59389,3.45000) (14.60694,3.45000) (14.62000,3.45000) (14.63306,3.45000) (14.64611,3.45000) (14.65917,3.45000) (14.67222,3.45000) (14.68528,3.45000) (14.69833,3.45000) (14.71139,3.45000) (14.72444,3.45000) (14.73750,3.45000) (14.75056,3.45000) (14.76361,3.45000) (14.77667,3.45000) (14.78972,3.45000) (14.80278,3.45000) (14.81583,3.45000) (14.82889,3.45000) (14.84194,3.45000) (14.85500,3.45000) (14.86806,3.45000) (14.88111,3.45000) (14.89417,3.45000) (14.90722,3.45000) (14.92028,3.45000) (14.93333,3.45000) (14.94639,3.45000) (14.95944,3.45000) (14.97250,3.45000) (14.98556,3.45000) (14.99861,3.45000) (15.01167,3.45000) (15.02472,3.45000) (15.03778,3.45000) (15.05083,3.45000) (15.06389,3.45000) (15.07694,3.45000) (15.09000,3.45000) (15.10306,3.45000) (15.11611,3.45000) (15.12917,3.45000) (15.14222,3.45000) (15.15528,3.45000) (15.16833,3.45000) (15.18139,3.45000) (15.19444,3.45000) (15.20750,3.45000) (15.22056,3.45000) (15.23361,3.45000) (15.24667,3.45000) (15.25972,3.45000) (15.27278,3.45000) (15.28583,3.45000) (15.29889,3.45000) (15.31194,3.45000) (15.32500,3.45000) (15.33806,3.45000) (15.35111,3.45000) (15.36417,3.45000) (15.37722,3.45000) (15.39028,3.45000) (15.40333,3.45000) (15.41639,3.45000) (15.42944,3.45000) (15.44250,3.45000) (15.45556,3.45000) (15.46861,3.45000) (15.48167,3.45000) (15.49472,3.45000) (15.50778,3.45000) (15.52083,3.45000) (15.53389,3.45000) (15.54694,3.45000) (15.56000,3.45000) (15.57306,3.45000) (15.58611,3.45000) (15.59917,3.45000) (15.61222,3.45000) (15.62528,3.45000) (15.63833,3.45000) (15.65139,3.45000) (15.66444,3.45000) (15.67750,3.45000) (15.69056,3.45000) (15.70361,3.45000) (15.71667,3.45000) (15.72972,3.45000) (15.74278,3.45000) (15.75583,3.45000) (15.76889,3.45000) (15.78194,3.45000) (15.79500,3.45000) (15.80806,3.45000) (15.82111,3.45000) (15.83417,3.45000) (15.84722,3.45000) (15.86028,3.45000) (15.87333,3.45000) (15.88639,3.45000) (15.89944,3.45000) (15.91250,3.45000) (15.92556,3.45000) (15.93861,3.45000) (15.95167,3.45000) (15.96472,3.45000) (15.97778,3.45000) (15.99083,3.45000) (16.00389,3.45000) (16.01694,3.45000) (16.03000,3.45000) (16.04306,3.45000) (16.05611,3.45000) (16.06917,3.45000) (16.08222,3.45000) (16.09528,3.45000) (16.10833,3.45000) (16.12139,3.45000) (16.13444,3.45000) (16.14750,3.45000) (16.16056,3.45000) (16.17361,3.45000) (16.18667,3.45000) (16.19972,3.45000) (16.21278,3.45000) (16.22583,3.45000) (16.23889,3.45000) (16.25194,3.45000) (16.26500,3.45000) (16.27806,3.45000) (16.29111,3.45000) (16.30417,3.45000) (16.31722,3.45000) (16.33028,3.45000) (16.34333,3.45000) (16.35639,3.45000) (16.36944,3.45000) (16.38250,3.45000) (16.39556,3.45000) (16.40861,3.45000) (16.42167,3.45000) (16.43472,3.45000) (16.44778,3.45000) (16.46083,3.45000) (16.47389,3.45000) (16.48694,3.45000) (16.50000,3.45000)};
\draw[->,coursemuted] plot coordinates {(7.10000,4.65000) (16.60000,4.65000)};
\draw[->,coursemuted] plot coordinates {(7.10000,4.65000) (7.10000,6.50000)};
\draw[coursemuted!25,thin] plot coordinates {(7.10000,4.65000) (16.50000,4.65000)};
\node[left,inner sep=3pt] at (7.1,4.65000) {$0$};
\draw[coursemuted!25,thin] plot coordinates {(7.10000,5.45000) (16.50000,5.45000)};
\node[left,inner sep=3pt] at (7.1,5.45000) {$2$};
\draw[coursemuted!25,thin] plot coordinates {(7.10000,6.25000) (16.50000,6.25000)};
\node[left,inner sep=3pt] at (7.1,6.25000) {$4$};
\node[rotate=90] at (5.9,5.55) {$\Delta V_j$, м/с};
\draw[coursemuted!16,thin] plot coordinates {(7.10000,0.80000) (7.10000,6.25000)};
\node[below,inner sep=3pt] at (7.10000,0.85) {$0$};
\draw[coursemuted!16,thin] plot coordinates {(9.45000,0.80000) (9.45000,6.25000)};
\node[below,inner sep=3pt] at (9.45000,0.85) {$\pi/2$};
\draw[coursemuted!16,thin] plot coordinates {(11.80000,0.80000) (11.80000,6.25000)};
\node[below,inner sep=3pt] at (11.80000,0.85) {$\pi$};
\draw[coursemuted!16,thin] plot coordinates {(14.15000,0.80000) (14.15000,6.25000)};
\node[below,inner sep=3pt] at (14.15000,0.85) {$3\pi/2$};
\draw[coursemuted!16,thin] plot coordinates {(16.50000,0.80000) (16.50000,6.25000)};
\node[below,inner sep=3pt] at (16.50000,0.85) {$2\pi$};
\node at (11.8,.15) {Угол $\varphi$, рад};
\draw[courseink!70,densely dashed,thin] plot coordinates {(8.51385,0.85000) (8.51385,3.45000)};
\draw[courseink,line width=1pt] plot coordinates {(8.51385,4.65000) (8.51385,5.58101)};
\fill[courseink] (8.51385,5.58101) circle[radius=1.8pt];
\fill[courseink] (8.51385,3.45000) circle[radius=1.8pt];
\node[above,inner sep=2pt] at (8.51385,5.58101) {1};
\draw[courseink!70,densely dashed,thin] plot coordinates {(10.38615,0.85000) (10.38615,3.45000)};
\draw[courseink,line width=1pt] plot coordinates {(10.38615,4.65000) (10.38615,5.66601)};
\fill[courseink] (10.38615,5.66601) circle[radius=1.8pt];
\fill[courseink] (10.38615,3.45000) circle[radius=1.8pt];
\node[above,inner sep=2pt] at (10.38615,5.66601) {2};
\draw[courseink!70,densely dashed,thin] plot coordinates {(13.21385,0.85000) (13.21385,3.45000)};
\draw[courseink,line width=1pt] plot coordinates {(13.21385,4.65000) (13.21385,4.73500)};
\fill[courseink] (13.21385,4.73500) circle[radius=1.8pt];
\fill[courseink] (13.21385,3.45000) circle[radius=1.8pt];
\node[above,inner sep=2pt] at (13.21385,4.73500) {3};
\end{tikzpicture}

\par
\figcaption{$\Delta a=0{,}001$, $\Delta e_x=0$, $\Delta e_y=0{,}001$, $\Delta i=0{,}0002$ рад.\\
Весь годограф лежит на сфере; модуль базис-вектора равен единице на всём обороте.}
\end{frame}

\begin{frame}{Особое решение}
\relax
\derivationmathspacing
Для явного построения выбираем одинаковые $\sin^2\psi_j$:
\[
\begin{aligned}
\sin^2\beta&=\frac{\Delta i(|\Delta e_y|+\sqrt3\Delta i)}{\sqrt3J_s^2},\\
\beta&=\arcsin\sqrt{\frac{\Delta i(|\Delta e_y|+\sqrt3\Delta i)}
{\sqrt3J_s^2}},
\qquad 0\leq\beta\leq\frac\pi2.
\end{aligned}
\]
\medskip
При $0<\beta<\pi/2$ используем четыре угла:
\[
\psi_1=\pi-\beta,\quad \psi_2=-\beta,\quad
\psi_3=\beta,\quad \psi_4=\pi+\beta.
\]
\medskip
\justifying
При $\cos\beta=0$ используется узловое решение с одним импульсом;
$\sin\beta=0$ соответствует компланарному случаю.
\end{frame}

% Источник: eq:sc-first-three.
\begin{frame}{Особое решение}
\relax
\derivationmathspacing
Выбираем $\Delta\widehat V_4=0$, если
\[
-\frac{2\operatorname{sgn}(\Delta e_y)\Delta i\Delta e_x}{\sqrt3J_s}
\leq\Delta a\sin\beta.
\]
Тогда модули остальных импульсов равны
\[
\begin{aligned}
\Delta\widehat V_1&=\frac{J_s}{2}-\frac{\Delta a}{4\cos\beta},\\
\Delta\widehat V_2&=\frac{\Delta a}{4\cos\beta}
+\frac{\operatorname{sgn}(\Delta e_y)\Delta i\Delta e_x}
{2\sqrt3J_s\sin\beta\cos\beta},\\
\Delta\widehat V_3&=\frac{J_s}{2}
-\frac{\operatorname{sgn}(\Delta e_y)\Delta i\Delta e_x}
{2\sqrt3J_s\sin\beta\cos\beta}.
\end{aligned}
\]
\medskip
\justifying
Не более трёх ненулевых импульсов. Особое решение не единственно.
\end{frame}

% Источник: eq:sc-last-three, eq:sc-final-impulses;
% предельные случаи после eq:sc-region.
\begin{frame}{Особое решение}
\relax
\derivationmathspacing
Выбираем $\Delta\widehat V_2=0$, если
\[
-\frac{2\operatorname{sgn}(\Delta e_y)\Delta i\Delta e_x}{\sqrt3J_s}
>\Delta a\sin\beta.
\]
\[
\begin{aligned}
\Delta\widehat V_1&=\frac{J_s}{2}
+\frac{\operatorname{sgn}(\Delta e_y)\Delta i\Delta e_x}
{2\sqrt3J_s\sin\beta\cos\beta},\\
\Delta\widehat V_3&=\frac{J_s}{2}+\frac{\Delta a}{4\cos\beta},\\
\Delta\widehat V_4&=-\frac{\operatorname{sgn}(\Delta e_y)\Delta i\Delta e_x}
{2\sqrt3J_s\sin\beta\cos\beta}
-\frac{\Delta a}{4\cos\beta}.
\end{aligned}
\]
\[
\varphi_j=\alpha+\psi_j\pmod{2\pi},\qquad
\Delta\widehat{\boldsymbol V}_j=
\Delta\widehat V_j\boldsymbol s(\varphi_j).
\]
\medskip
\justifying
Нулевые импульсы исключаем.
\end{frame}

% Источник: sec:transfer-check, порядок из методички сохранён.
\begin{frame}{Порядок расчёта перехода}
\relax
\derivationmathspacing
После вычисления $\Delta a,\Delta e_x,\Delta e_y,\Delta i$:
\begin{enumerate}
\item При $\Delta i=0$ используем компланарное решение.
Если также $\Delta a=\Delta e=0$, импульсы не нужны.
\item Если выполнены условия узлового решения, используем его.
При $\Delta e_x=0$ применяем одноимпульсную формулу.
\item Если узловое решение не подходит и $\Delta e_y=0$,
используем решение с нулевыми радиальными компонентами.
\item Если выполнены условия особого решения, используем его.
\item В остальных случаях используем невырожденное решение.
\end{enumerate}
\end{frame}

% Итоги околокругового приближения.
\begin{frame}[plain]
\begin{tikzpicture}[remember picture,overlay]
\node[text=courseink,font=\usebeamerfont{title}]
  at (current page.center) {Итоги};
\end{tikzpicture}
\end{frame}

% Полная финальная система: manual/chapter-2-optimality.tex,
% eq:opt-boundaries-full; обозначения и приближения главы 1.
\begin{frame}{Система Баранова}
\relax
\derivationmathspacing
\begin{columns}[T,onlytextwidth]
\begin{column}{.64\textwidth}
{\eqcompact\derivationmathspacing
\centering\(\displaystyle
\begin{aligned}[t]
 &2\sum\nolimits_j\Delta\widehat V_{tj}=\Delta a,\\
 &\sum\nolimits_j(\Delta\widehat V_{rj}\sin\varphi_j+
              2\Delta\widehat V_{tj}\cos\varphi_j)=\Delta e_x,\\
 &\sum\nolimits_j(-\Delta\widehat V_{rj}\cos\varphi_j+
              2\Delta\widehat V_{tj}\sin\varphi_j)=\Delta e_y,\\
 &-\sum\nolimits_j\Delta\widehat V_{zj}\sin\varphi_j=0,\\
 &\sum\nolimits_j\Delta\widehat V_{zj}\cos\varphi_j=\Delta i,\\
 &\sum\nolimits_j\bigl[2\Delta\widehat V_{rj}(1-\cos\varphi_j)
       +\Delta\widehat V_{tj}(-3\varphi_j+4\sin\varphi_j)\bigr]=\\[2pt]
 &\quad\Delta t-\frac32\varphi_f\Delta a
       +2\Delta e_x\sin\varphi_f+2\Delta e_y(1-\cos\varphi_f).
\end{aligned}
\)\par}
\smallskip
\figcaption{$\Delta t=n_0(t_f-t_0)$ --- запаздывание прихода.\\
$t_0$ --- момент пассивного прохождения конечного луча.}
\end{column}
\begin{column}{.33\textwidth}
\begin{block}{Принятые приближения}
Центральное ньютоновское поле. Один опорный круг $r_0$.\par\smallskip
Малы $e$, $i$, $|a-r_0|/r_0$ и $\Delta\widehat V_j$.\par\smallskip
На всём интервале малы нормированные отклонения состояния,
включая $|u-\varphi|$.\par\smallskip
Первый порядок малости.
\end{block}
\end{column}
\end{columns}
\end{frame}

% Общий случай, включая вырожденный: eq:primer-phase-shift.
\begin{frame}{Базис-вектор в общем случае}
\relax
\derivationmathspacing
При $\lambda_2^2+\lambda_3^2>0$ вводим вспомогательный угол:
\[
\cos\alpha=\frac{\lambda_3}{\sqrt{\lambda_2^2+\lambda_3^2}},\qquad
\sin\alpha=\frac{\lambda_2}{\sqrt{\lambda_2^2+\lambda_3^2}}.
\]
\[
\boxed{\begin{aligned}
s_r&=\sqrt{\lambda_2^2+\lambda_3^2}\sin(\varphi-\alpha)+2\lambda_6,\\
s_t&=2\lambda_1+2\sqrt{\lambda_2^2+\lambda_3^2}
              \cos(\varphi-\alpha)-3\lambda_6\varphi,\\
s_z&=(\lambda_4\cos\alpha-\lambda_5\sin\alpha)\sin(\varphi-\alpha)\\
   &\quad+(\lambda_4\sin\alpha+\lambda_5\cos\alpha)\cos(\varphi-\alpha).
\end{aligned}}
\]
При $\lambda_2=\lambda_3=0$ угол $\alpha$ не определён; используем
\[
s_r=2\lambda_6,\quad
s_t=2\lambda_1-3\lambda_6\varphi,\quad
s_z=\lambda_4\sin\varphi+\lambda_5\cos\varphi.
\]
\end{frame}

% Достаточность: sec:optimality-check-full в главе 2.
\begin{frame}{Достаточные условия оптимальности}
\relax
\derivationmathspacing
\begin{enumerate}
\setlength{\itemsep}{2pt}
\item Выполнены все граничные условия задачи, фазы импульсов допустимы:
$\varphi_j\in[\varphi_i,\varphi_f]$.
\item Существует один набор постоянных $\lambda_1,\ldots,\lambda_6$,
для которого каждый ненулевой импульс удовлетворяет
\[
\boldsymbol s(\varphi_j)
 =\frac{\Delta\widehat{\boldsymbol V}_j}{\Delta\widehat V_j},\qquad
\|\boldsymbol s(\varphi_j)\|=1.
\]
\item Для того же базис-вектора на всём допустимом интервале
\[
\|\boldsymbol s(\varphi)\|\leq1
\qquad\text{при всех }\varphi\in[\varphi_i,\varphi_f].
\]
\end{enumerate}
\begin{block}{Глобальный минимум в околокруговом приближении}
$J=\sum_j\Delta\widehat V_j$ минимально среди допустимых импульсных манёвров.
\end{block}
\end{frame}

% sec:planar, eq:planar-intersection; условия без новых обозначений.
\begin{frame}{Классификация компланарных переходов}
\relax
\derivationmathspacing
$\Delta i=0$, $\Delta e=\sqrt{(\Delta e_x)^2+(\Delta e_y)^2}$.
\par\smallskip
\centering
% Классификация по manual/chapter-3-algorithms.tex, sec:planar.
\begin{tikzpicture}[course diagram,x=1mm,y=1mm,
  condition/.style={draw=courserule,fill=coursetint,line width=.6pt,
    text width=43mm,minimum height=10mm,align=center,inner sep=1.5mm,
    font=\normalsize},
  result/.style={draw=courserule,fill=coursetint,line width=.6pt,
    text width=76mm,minimum height=12mm,align=center,inner sep=1.5mm,
    font=\normalsize},
  branch/.style={draw=courseaccent,line width=.7pt,->}]
  \path[use as bounding box] (0,0) rectangle (148,42);
  \node[condition] (tangent-test) at (26,35) {$|\Delta a|=\Delta e>0$};
  \node[condition] (disjoint-test) at (26,21) {$|\Delta a|>\Delta e$};
  \node[condition] (crossing-test) at (26,7) {$|\Delta a|<\Delta e$};
  \node[result] (tangent) at (108,35)
    {Касающиеся орбиты\\
     {\fontsize{9.5}{11.5}\selectfont Один трансверсальный импульс}};
  \node[result] (disjoint) at (108,21)
    {Непересекающиеся орбиты\\
     {\fontsize{9.5}{11.5}\selectfont Два импульса одного знака}};
  \node[result] (crossing) at (108,7)
    {Пересекающиеся орбиты\\
     {\fontsize{9.5}{11.5}\selectfont Два импульса противоположных знаков}};
  \draw[branch] (tangent-test.east)--(tangent.west);
  \draw[branch] (disjoint-test.east)--(disjoint.west);
  \draw[branch] (crossing-test.east)--(crossing.west);
\end{tikzpicture}

\par\smallskip
При $\Delta a=\Delta e=0$ орбиты совпадают; манёвр не требуется.
\end{frame}

% sec:transfer-check; границы проверяются в порядке методички.
\begin{frame}{Классификация некомпланарных переходов}
\relax
\derivationmathspacing
При $\Delta i>0$ проверяем условия сверху вниз:
\par
\centering
% Порядок проверок: manual/chapter-3-algorithms.tex, sec:transfer-check.
% Последовательность отдаёт общие границы сначала NC, затем SC.
\begin{tikzpicture}[course diagram,x=1mm,y=1mm,
  condition/.style={draw=courserule,fill=coursetint,line width=.6pt,
    text width=85mm,align=center,inner sep=1.5mm},
  result/.style={draw=courserule,fill=coursetint,line width=.6pt,
    text width=42mm,minimum height=11mm,align=center,inner sep=1.5mm},
  branch/.style={draw=courseaccent,line width=.7pt,->}]
  \path[use as bounding box] (0,-7) rectangle (148,46.5);
  \node[condition,minimum height=11mm] (nc-test) at (44,40)
    {$\begin{gathered}
       |\Delta a|\leq|\Delta e_x|,\\
       \Delta i\geq\sqrt3\,|\Delta e_y|
     \end{gathered}$};
  \node[result] (nc) at (124.5,40)
    {Узловое решение\\
     {\fontsize{9.5}{11.5}\selectfont $\Delta e_x=0$: один импульс}};
  \node[condition,minimum height=10mm] (line-test) at (44,26)
    {$\Delta e_y=0$};
  \node[result] (line-case) at (124.5,26)
    {Нулевые радиальные\\компоненты};
  \node[condition,minimum height=16mm] (sc-test) at (44,10)
    {$\displaystyle\begin{gathered}
       \Delta i\leq\sqrt3\,|\Delta e_y|,\\
       (\Delta a)^2\leq\Delta e^2+
         \frac{2\Delta i|\Delta e_y|}{\sqrt3}-(\Delta i)^2
     \end{gathered}$};
  \node[result] (sc) at (124.5,14) {Особое решение};
  \node[result] (ndc) at (124.5,-1) {Невырожденное\\решение};
  \draw[branch] (nc-test.east)--node[above] {да}(nc.west);
  \draw[branch] (nc-test.south)--node[right] {нет}(line-test.north);
  \draw[branch] (line-test.east)--node[above] {да}(line-case.west);
  \draw[branch] (line-test.south)--node[right] {нет}(sc-test.north);
  \draw[branch] (sc-test.east)--node[above] {да}(sc.west);
  \draw[branch] (sc-test.south)--(44,-1)--
    node[below=1pt,inner sep=.5pt] {нет}(ndc.west);
\end{tikzpicture}

\end{frame}

\begin{frame}{Источники}
\relax
\setbeamertemplate{bibliography item}{\insertbiblabel}
\begin{thebibliography}{9}
\bibitem{Baranov}
Баранов А.~А. \textit{Маневрирование космических аппаратов
в окрестности круговой орбиты}.
М.: Спутник+, 2016. 511~с. Гл.~1--2.
\bibitem{Lawden}
Lawden D.~F. \textit{Optimal Trajectories for Space Navigation}.
London: Butterworths, 1963. Гл.~3, \S~5.1, гл.~6.
\bibitem{Edelbaum}
Edelbaum T.~N.
\href{https://ntrs.nasa.gov/citations/19670020044}{A General Solution for Minimum Impulse Transfers
in the Near Vicinity of a Circular Orbit} //
\textit{First Compilation of Papers on Trajectory Analysis and Guidance Theory}.
NASA SP-141, 1967. P.~37--71.
Невырожденное решение: с.~43--44, 48--53;
особое решение: с.~44--45, 54--56.
\bibitem{Gobetz}
Gobetz F.~W., Washington M., Edelbaum T.~N.
\href{https://ntrs.nasa.gov/citations/19660030604}{\textit{Minimum-Impulse Time-Free Transfer Between Elliptic Orbits}}.
NASA CR-636, 1966. Околокруговое приближение: с.~25--42;
приложение~C: с.~45--49.
\end{thebibliography}
\end{frame}

\end{document}
